% ch10 — Spectral Methods
\chapter{Spectral Methods}
\label{ch:spectral-methods}

The finite difference stencils of \cref{ch:finite-differences} approximate
derivatives by combining a handful of neighbouring grid values --- five points
for a fourth-order first derivative, seven for the dissipation operator.
Adding more points to the stencil raises the order of accuracy, but the
improvement is only algebraic: doubling the stencil width roughly doubles the
convergence rate.  Spectral methods take the idea to its logical conclusion
by using \emph{every} point in the element to approximate the derivative at
each node.  The resulting differentiation matrix is dense rather than banded,
and for smooth functions the convergence is no longer algebraic but
\emph{exponential} --- each additional node can buy an extra digit of
accuracy rather than a fixed fraction of one.  This is the accuracy advantage
that makes high-order discontinuous Galerkin (DG) methods competitive with
finite differences at a fraction of the grid resolution.

The practical question is which polynomials and which nodes.  The chapter
opens with the interpolation theory that governs the choice, showing how
equidistant nodes fail catastrophically at high order (the Runge phenomenon)
and how clustered nodes cure the instability.  Two polynomial families
provide the clustering: Chebyshev polynomials, whose zeros and extrema
deliver near-optimal interpolation points for collocation methods, and
Legendre polynomials, whose orthogonality under the unit weight function
makes them the natural basis for the DG weak form.  The Legendre--Gauss--Lobatto
(LGL) nodes that the codebase uses sit at the intersection of these two
threads --- they inherit the endpoint clustering that tames the Runge
phenomenon and the quadrature exactness that makes mass lumping diagonal.

The chapter assembles the LGL toolbox piece by piece: interpolation and
the Lebesgue constant (\cref{sec:spec-interpolation}), Chebyshev polynomial
theory (\cref{sec:spec-chebyshev}), spectral differentiation matrices
(\cref{sec:spec-diffmat}), Legendre polynomials and Gaussian quadrature
(\cref{sec:spec-legendre}), convergence rates tied to function smoothness
(\cref{sec:spec-convergence}), the LGL differentiation and mass matrices
used in the DG weak form (\cref{sec:spec-lgl}), and the Vandermonde
transform that converts between nodal and modal representations for spectral
filtering (\cref{sec:spec-vandermonde}).

% ---- Drafted sections ----
% sec:spec-interpolation — Polynomial Interpolation
\section{Polynomial Interpolation}
\label{sec:spec-interpolation}

The discontinuous Galerkin solver at the heart of the numerical
relativity engine represents the solution within each element as a
polynomial.  Every operation the solver performs --- differentiating,
integrating, exchanging data across element faces --- reduces to
manipulating polynomials defined by their values at a set of nodes.
The quality of this representation depends entirely on where those
nodes are placed.  This section develops the interpolation theory that
governs the choice: Lagrange interpolation on general node sets, the
error bound that ties accuracy to node placement, and the spectacular
failure that results from the most natural choice of all ---
equidistant points.

\paragraph{Lagrange interpolation.}

Suppose we are given $N + 1$ distinct
nodes on an interval $[a, b]$ and the corresponding function values.
There exists a unique polynomial $p_N$ of degree at most $N$ that
passes through every data point $(x_j, f(x_j))$.  The idea is to construct $N + 1$ ``selector'' polynomials,
each of which equals one at exactly one node and vanishes at all the
others; the interpolant is then a weighted combination that picks out
the correct function value at every node.  The explicit construction
uses the \emph{Lagrange basis polynomials}
%
\begin{equation}
  \ell_j(x)
    = \prod_{\substack{k=0 \\ k \neq j}}^{N}
      \frac{x - x_k}{x_j - x_k} \,,
  \qquad j = 0, 1, \ldots, N \,.
  \label{eq:lagrange-basis}
\end{equation}
%
Each basis polynomial $\ell_j$ equals one at $x_j$ and vanishes at
every other node: $\ell_j(x_k) = \delta_{jk}$.  The interpolant is
then
%
\begin{equation}
  p_N(x)
    = \sum_{j=0}^{N} f(x_j)\,\ell_j(x) \,.
  \label{eq:lagrange-interp}
\end{equation}
%
This is not an approximation --- $p_N$ agrees with $f$ exactly at the
nodes.  The question is how well it approximates $f$ between them.

\paragraph{The interpolation error.}

For a function $f \in C^{N+1}[a,b]$, the pointwise error satisfies
%
\begin{equation}
  f(x) - p_N(x)
    = \frac{f^{(N+1)}(\xi)}{(N+1)!}\;\omega_{N+1}(x) \,,
  \label{eq:interp-error}
\end{equation}
%
where $\xi \in [a,b]$ depends on $x$ and $\omega_{N+1}$ is the
\emph{node polynomial}
%
\begin{equation}
  \omega_{N+1}(x)
    = \prod_{j=0}^{N} (x - x_j) \,.
  \label{eq:node-polynomial}
\end{equation}
%
\Cref{eq:interp-error} separates the error into two factors: the
smoothness of $f$ (through its $(N+1)$-th derivative) and the
geometry of the node set (through $\omega_{N+1}$).  The first factor
is fixed by the function being interpolated; the second is under our
control.  Choosing nodes to minimise $\max_{x \in [a,b]}
\abs{\omega_{N+1}(x)}$ is the central problem of interpolation
theory, and the subject of the next section.
\physnote{The error formula \cref{eq:interp-error} has the same
structure as the remainder term in Taylor's theorem --- both express
the approximation error through a higher derivative evaluated at an
unknown intermediate point.  The difference is that Taylor's theorem
uses a single expansion point, while Lagrange interpolation
distributes the information across $N + 1$ nodes.}

\paragraph{The Lebesgue constant.}

The error formula \cref{eq:interp-error} requires $(N+1)$-th
derivative information that is rarely available in practice.  A more
operational measure of interpolation quality is the \emph{Lebesgue
constant} $\Lambda_N$, defined as
%
\begin{equation}
  \Lambda_N
    = \max_{x \in [a,b]}
      \sum_{j=0}^{N} \abs{\ell_j(x)} \,.
  \label{eq:lebesgue-constant}
\end{equation}
%
Informally, $\Lambda_N$ measures the worst-case amplification: if the
function values $f(x_j)$ are perturbed by a small amount $\epsilon$,
the interpolant can change by as much as $\Lambda_N\,\epsilon$.  A
Lebesgue constant near one means the interpolant is stable; a
Lebesgue constant of $10^4$ means tiny perturbations at the nodes
produce enormous swings between them.  The Lebesgue constant also
controls how close the interpolant is to the best possible polynomial
approximation: if $p_N^*$ denotes the best approximation of degree $N$
in the supremum norm, then
%
\begin{equation}
  \norm{f - p_N}_\infty
    \;\leq\; (1 + \Lambda_N)\,\norm{f - p_N^*}_\infty \,.
  \label{eq:lebesgue-bound}
\end{equation}
%
A small Lebesgue constant means the interpolant is near-optimal;
a large one means the node set amplifies errors unnecessarily.  For
equidistant nodes, $\Lambda_N$ grows exponentially with $N$ --- a fact
that underlies the Runge phenomenon described below.
\warnnote{The Lebesgue constant is a property of the node set alone,
not of any particular function.  A large $\Lambda_N$ means the node
set is a poor choice for \emph{every} smooth function, not just
pathological ones.}

\paragraph{The Runge phenomenon.}

The most natural choice of nodes --- equally spaced on $[a, b]$ ---
turns out to be catastrophically bad at high polynomial degree.  Carl
Runge demonstrated this in 1901 using the innocent-looking function
%
\begin{equation}
  f(x) = \frac{1}{1 + 25\,x^2}
  \label{eq:runge-function}
\end{equation}
%
on $[-1, 1]$.  This function is smooth, bounded, and gently peaked at
the origin.  Yet polynomial interpolation on equidistant nodes
diverges as $N$ increases: the interpolant develops wild oscillations
near the endpoints that grow without bound.  At degree 15 the
endpoint error already exceeds the function's maximum value, and by
degree 25 the oscillations dwarf the function itself --- adding more
data makes the approximation worse, not better.
\histnote{Runge published this counterexample in 1901, nearly half a
century after Chebyshev had identified the optimal node distribution
in the 1850s.  The phenomenon surprised many who assumed that more
data points would always improve polynomial approximation.}

The node polynomial of \cref{eq:node-polynomial} reveals the
mechanism.  For equidistant nodes, $\omega_{N+1}$
is much larger near the endpoints than near the centre of the
interval.  The ratio between its endpoint and midpoint values grows
exponentially with $N$, so the error formula \cref{eq:interp-error}
produces large errors near $x = \pm 1$ even when the derivative
factor $f^{(N+1)}$ remains moderate.  The Lebesgue constant for
equidistant nodes grows as
$\Lambda_N \sim 2^N / (e\,N\ln N)$ \cite{Trefethen2000}, confirming
that the instability is not specific to Runge's function but is an
intrinsic property of the node distribution.

\begin{figure}[t]
\centering
%% Creator: Matplotlib, PGF backend
%%
%% To include the figure in your LaTeX document, write
%%   \input{<filename>.pgf}
%%
%% Make sure the required packages are loaded in your preamble
%%   \usepackage{pgf}
%%
%% Also ensure that all the required font packages are loaded; for instance,
%% the lmodern package is sometimes necessary when using math font.
%%   \usepackage{lmodern}
%%
%% Figures using additional raster images can only be included by \input if
%% they are in the same directory as the main LaTeX file. For loading figures
%% from other directories you can use the `import` package
%%   \usepackage{import}
%%
%% and then include the figures with
%%   \import{<path to file>}{<filename>.pgf}
%%
%% Matplotlib used the following preamble
%%   \def\mathdefault#1{#1}
%%   \everymath=\expandafter{\the\everymath\displaystyle}
%%   \IfFileExists{scrextend.sty}{
%%     \usepackage[fontsize=10.000000pt]{scrextend}
%%   }{
%%     \renewcommand{\normalsize}{\fontsize{10.000000}{12.000000}\selectfont}
%%     \normalsize
%%   }
%%   \usepackage{fontspec}
%%   \usepackage{unicode-math}
%%   \setmainfont{STIX Two Text}[Numbers=OldStyle, Ligatures=TeX]
%%   \setmathfont{STIX Two Math}
%%   \setmonofont{Source Code Pro}[Scale=MatchLowercase]
%%   \ifdefined\pdftexversion\else  % non-pdftex case.
%%     \usepackage{fontspec}
%%     \setmainfont{DejaVuSerif.ttf}[Path=\detokenize{/Users/gvn/.pyenv/versions/3.12.11/lib/python3.12/site-packages/matplotlib/mpl-data/fonts/ttf/}]
%%     \setsansfont{DejaVuSans.ttf}[Path=\detokenize{/Users/gvn/.pyenv/versions/3.12.11/lib/python3.12/site-packages/matplotlib/mpl-data/fonts/ttf/}]
%%     \setmonofont{DejaVuSansMono.ttf}[Path=\detokenize{/Users/gvn/.pyenv/versions/3.12.11/lib/python3.12/site-packages/matplotlib/mpl-data/fonts/ttf/}]
%%   \fi
%%   \makeatletter\@ifpackageloaded{underscore}{}{\usepackage[strings]{underscore}}\makeatother
%%
\begingroup%
\makeatletter%
\begin{pgfpicture}%
\pgfpathrectangle{\pgfpointorigin}{\pgfqpoint{3.788821in}{2.311421in}}%
\pgfusepath{use as bounding box, clip}%
\begin{pgfscope}%
\pgfsetbuttcap%
\pgfsetmiterjoin%
\definecolor{currentfill}{rgb}{1.000000,1.000000,1.000000}%
\pgfsetfillcolor{currentfill}%
\pgfsetlinewidth{0.000000pt}%
\definecolor{currentstroke}{rgb}{1.000000,1.000000,1.000000}%
\pgfsetstrokecolor{currentstroke}%
\pgfsetdash{}{0pt}%
\pgfpathmoveto{\pgfqpoint{0.000000in}{0.000000in}}%
\pgfpathlineto{\pgfqpoint{3.788821in}{0.000000in}}%
\pgfpathlineto{\pgfqpoint{3.788821in}{2.311421in}}%
\pgfpathlineto{\pgfqpoint{0.000000in}{2.311421in}}%
\pgfpathlineto{\pgfqpoint{0.000000in}{0.000000in}}%
\pgfpathclose%
\pgfusepath{fill}%
\end{pgfscope}%
\begin{pgfscope}%
\pgfsetbuttcap%
\pgfsetmiterjoin%
\definecolor{currentfill}{rgb}{1.000000,1.000000,1.000000}%
\pgfsetfillcolor{currentfill}%
\pgfsetlinewidth{0.000000pt}%
\definecolor{currentstroke}{rgb}{0.000000,0.000000,0.000000}%
\pgfsetstrokecolor{currentstroke}%
\pgfsetstrokeopacity{0.000000}%
\pgfsetdash{}{0pt}%
\pgfpathmoveto{\pgfqpoint{0.507620in}{0.352669in}}%
\pgfpathlineto{\pgfqpoint{1.857769in}{0.352669in}}%
\pgfpathlineto{\pgfqpoint{1.857769in}{2.123669in}}%
\pgfpathlineto{\pgfqpoint{0.507620in}{2.123669in}}%
\pgfpathlineto{\pgfqpoint{0.507620in}{0.352669in}}%
\pgfpathclose%
\pgfusepath{fill}%
\end{pgfscope}%
\begin{pgfscope}%
\pgfsetbuttcap%
\pgfsetroundjoin%
\definecolor{currentfill}{rgb}{0.000000,0.000000,0.000000}%
\pgfsetfillcolor{currentfill}%
\pgfsetlinewidth{0.501875pt}%
\definecolor{currentstroke}{rgb}{0.000000,0.000000,0.000000}%
\pgfsetstrokecolor{currentstroke}%
\pgfsetdash{}{0pt}%
\pgfsys@defobject{currentmarker}{\pgfqpoint{0.000000in}{0.000000in}}{\pgfqpoint{0.000000in}{0.041667in}}{%
\pgfpathmoveto{\pgfqpoint{0.000000in}{0.000000in}}%
\pgfpathlineto{\pgfqpoint{0.000000in}{0.041667in}}%
\pgfusepath{stroke,fill}%
}%
\begin{pgfscope}%
\pgfsys@transformshift{0.507620in}{0.352669in}%
\pgfsys@useobject{currentmarker}{}%
\end{pgfscope}%
\end{pgfscope}%
\begin{pgfscope}%
\definecolor{textcolor}{rgb}{0.000000,0.000000,0.000000}%
\pgfsetstrokecolor{textcolor}%
\pgfsetfillcolor{textcolor}%
\pgftext[x=0.507620in,y=0.304058in,,top]{\color{textcolor}{\sffamily\fontsize{8.000000}{9.600000}\selectfont\catcode`\^=\active\def^{\ifmmode\sp\else\^{}\fi}\catcode`\%=\active\def%{\%}\ensuremath{-}1.0}}%
\end{pgfscope}%
\begin{pgfscope}%
\pgfsetbuttcap%
\pgfsetroundjoin%
\definecolor{currentfill}{rgb}{0.000000,0.000000,0.000000}%
\pgfsetfillcolor{currentfill}%
\pgfsetlinewidth{0.501875pt}%
\definecolor{currentstroke}{rgb}{0.000000,0.000000,0.000000}%
\pgfsetstrokecolor{currentstroke}%
\pgfsetdash{}{0pt}%
\pgfsys@defobject{currentmarker}{\pgfqpoint{0.000000in}{0.000000in}}{\pgfqpoint{0.000000in}{0.041667in}}{%
\pgfpathmoveto{\pgfqpoint{0.000000in}{0.000000in}}%
\pgfpathlineto{\pgfqpoint{0.000000in}{0.041667in}}%
\pgfusepath{stroke,fill}%
}%
\begin{pgfscope}%
\pgfsys@transformshift{0.845157in}{0.352669in}%
\pgfsys@useobject{currentmarker}{}%
\end{pgfscope}%
\end{pgfscope}%
\begin{pgfscope}%
\definecolor{textcolor}{rgb}{0.000000,0.000000,0.000000}%
\pgfsetstrokecolor{textcolor}%
\pgfsetfillcolor{textcolor}%
\pgftext[x=0.845157in,y=0.304058in,,top]{\color{textcolor}{\sffamily\fontsize{8.000000}{9.600000}\selectfont\catcode`\^=\active\def^{\ifmmode\sp\else\^{}\fi}\catcode`\%=\active\def%{\%}\ensuremath{-}0.5}}%
\end{pgfscope}%
\begin{pgfscope}%
\pgfsetbuttcap%
\pgfsetroundjoin%
\definecolor{currentfill}{rgb}{0.000000,0.000000,0.000000}%
\pgfsetfillcolor{currentfill}%
\pgfsetlinewidth{0.501875pt}%
\definecolor{currentstroke}{rgb}{0.000000,0.000000,0.000000}%
\pgfsetstrokecolor{currentstroke}%
\pgfsetdash{}{0pt}%
\pgfsys@defobject{currentmarker}{\pgfqpoint{0.000000in}{0.000000in}}{\pgfqpoint{0.000000in}{0.041667in}}{%
\pgfpathmoveto{\pgfqpoint{0.000000in}{0.000000in}}%
\pgfpathlineto{\pgfqpoint{0.000000in}{0.041667in}}%
\pgfusepath{stroke,fill}%
}%
\begin{pgfscope}%
\pgfsys@transformshift{1.182694in}{0.352669in}%
\pgfsys@useobject{currentmarker}{}%
\end{pgfscope}%
\end{pgfscope}%
\begin{pgfscope}%
\definecolor{textcolor}{rgb}{0.000000,0.000000,0.000000}%
\pgfsetstrokecolor{textcolor}%
\pgfsetfillcolor{textcolor}%
\pgftext[x=1.182694in,y=0.304058in,,top]{\color{textcolor}{\sffamily\fontsize{8.000000}{9.600000}\selectfont\catcode`\^=\active\def^{\ifmmode\sp\else\^{}\fi}\catcode`\%=\active\def%{\%}0.0}}%
\end{pgfscope}%
\begin{pgfscope}%
\pgfsetbuttcap%
\pgfsetroundjoin%
\definecolor{currentfill}{rgb}{0.000000,0.000000,0.000000}%
\pgfsetfillcolor{currentfill}%
\pgfsetlinewidth{0.501875pt}%
\definecolor{currentstroke}{rgb}{0.000000,0.000000,0.000000}%
\pgfsetstrokecolor{currentstroke}%
\pgfsetdash{}{0pt}%
\pgfsys@defobject{currentmarker}{\pgfqpoint{0.000000in}{0.000000in}}{\pgfqpoint{0.000000in}{0.041667in}}{%
\pgfpathmoveto{\pgfqpoint{0.000000in}{0.000000in}}%
\pgfpathlineto{\pgfqpoint{0.000000in}{0.041667in}}%
\pgfusepath{stroke,fill}%
}%
\begin{pgfscope}%
\pgfsys@transformshift{1.520231in}{0.352669in}%
\pgfsys@useobject{currentmarker}{}%
\end{pgfscope}%
\end{pgfscope}%
\begin{pgfscope}%
\definecolor{textcolor}{rgb}{0.000000,0.000000,0.000000}%
\pgfsetstrokecolor{textcolor}%
\pgfsetfillcolor{textcolor}%
\pgftext[x=1.520231in,y=0.304058in,,top]{\color{textcolor}{\sffamily\fontsize{8.000000}{9.600000}\selectfont\catcode`\^=\active\def^{\ifmmode\sp\else\^{}\fi}\catcode`\%=\active\def%{\%}0.5}}%
\end{pgfscope}%
\begin{pgfscope}%
\pgfsetbuttcap%
\pgfsetroundjoin%
\definecolor{currentfill}{rgb}{0.000000,0.000000,0.000000}%
\pgfsetfillcolor{currentfill}%
\pgfsetlinewidth{0.501875pt}%
\definecolor{currentstroke}{rgb}{0.000000,0.000000,0.000000}%
\pgfsetstrokecolor{currentstroke}%
\pgfsetdash{}{0pt}%
\pgfsys@defobject{currentmarker}{\pgfqpoint{0.000000in}{0.000000in}}{\pgfqpoint{0.000000in}{0.041667in}}{%
\pgfpathmoveto{\pgfqpoint{0.000000in}{0.000000in}}%
\pgfpathlineto{\pgfqpoint{0.000000in}{0.041667in}}%
\pgfusepath{stroke,fill}%
}%
\begin{pgfscope}%
\pgfsys@transformshift{1.857769in}{0.352669in}%
\pgfsys@useobject{currentmarker}{}%
\end{pgfscope}%
\end{pgfscope}%
\begin{pgfscope}%
\definecolor{textcolor}{rgb}{0.000000,0.000000,0.000000}%
\pgfsetstrokecolor{textcolor}%
\pgfsetfillcolor{textcolor}%
\pgftext[x=1.857769in,y=0.304058in,,top]{\color{textcolor}{\sffamily\fontsize{8.000000}{9.600000}\selectfont\catcode`\^=\active\def^{\ifmmode\sp\else\^{}\fi}\catcode`\%=\active\def%{\%}1.0}}%
\end{pgfscope}%
\begin{pgfscope}%
\pgfsetbuttcap%
\pgfsetroundjoin%
\definecolor{currentfill}{rgb}{0.000000,0.000000,0.000000}%
\pgfsetfillcolor{currentfill}%
\pgfsetlinewidth{0.301125pt}%
\definecolor{currentstroke}{rgb}{0.000000,0.000000,0.000000}%
\pgfsetstrokecolor{currentstroke}%
\pgfsetdash{}{0pt}%
\pgfsys@defobject{currentmarker}{\pgfqpoint{0.000000in}{0.000000in}}{\pgfqpoint{0.000000in}{0.027778in}}{%
\pgfpathmoveto{\pgfqpoint{0.000000in}{0.000000in}}%
\pgfpathlineto{\pgfqpoint{0.000000in}{0.027778in}}%
\pgfusepath{stroke,fill}%
}%
\begin{pgfscope}%
\pgfsys@transformshift{0.575127in}{0.352669in}%
\pgfsys@useobject{currentmarker}{}%
\end{pgfscope}%
\end{pgfscope}%
\begin{pgfscope}%
\pgfsetbuttcap%
\pgfsetroundjoin%
\definecolor{currentfill}{rgb}{0.000000,0.000000,0.000000}%
\pgfsetfillcolor{currentfill}%
\pgfsetlinewidth{0.301125pt}%
\definecolor{currentstroke}{rgb}{0.000000,0.000000,0.000000}%
\pgfsetstrokecolor{currentstroke}%
\pgfsetdash{}{0pt}%
\pgfsys@defobject{currentmarker}{\pgfqpoint{0.000000in}{0.000000in}}{\pgfqpoint{0.000000in}{0.027778in}}{%
\pgfpathmoveto{\pgfqpoint{0.000000in}{0.000000in}}%
\pgfpathlineto{\pgfqpoint{0.000000in}{0.027778in}}%
\pgfusepath{stroke,fill}%
}%
\begin{pgfscope}%
\pgfsys@transformshift{0.642634in}{0.352669in}%
\pgfsys@useobject{currentmarker}{}%
\end{pgfscope}%
\end{pgfscope}%
\begin{pgfscope}%
\pgfsetbuttcap%
\pgfsetroundjoin%
\definecolor{currentfill}{rgb}{0.000000,0.000000,0.000000}%
\pgfsetfillcolor{currentfill}%
\pgfsetlinewidth{0.301125pt}%
\definecolor{currentstroke}{rgb}{0.000000,0.000000,0.000000}%
\pgfsetstrokecolor{currentstroke}%
\pgfsetdash{}{0pt}%
\pgfsys@defobject{currentmarker}{\pgfqpoint{0.000000in}{0.000000in}}{\pgfqpoint{0.000000in}{0.027778in}}{%
\pgfpathmoveto{\pgfqpoint{0.000000in}{0.000000in}}%
\pgfpathlineto{\pgfqpoint{0.000000in}{0.027778in}}%
\pgfusepath{stroke,fill}%
}%
\begin{pgfscope}%
\pgfsys@transformshift{0.710142in}{0.352669in}%
\pgfsys@useobject{currentmarker}{}%
\end{pgfscope}%
\end{pgfscope}%
\begin{pgfscope}%
\pgfsetbuttcap%
\pgfsetroundjoin%
\definecolor{currentfill}{rgb}{0.000000,0.000000,0.000000}%
\pgfsetfillcolor{currentfill}%
\pgfsetlinewidth{0.301125pt}%
\definecolor{currentstroke}{rgb}{0.000000,0.000000,0.000000}%
\pgfsetstrokecolor{currentstroke}%
\pgfsetdash{}{0pt}%
\pgfsys@defobject{currentmarker}{\pgfqpoint{0.000000in}{0.000000in}}{\pgfqpoint{0.000000in}{0.027778in}}{%
\pgfpathmoveto{\pgfqpoint{0.000000in}{0.000000in}}%
\pgfpathlineto{\pgfqpoint{0.000000in}{0.027778in}}%
\pgfusepath{stroke,fill}%
}%
\begin{pgfscope}%
\pgfsys@transformshift{0.777649in}{0.352669in}%
\pgfsys@useobject{currentmarker}{}%
\end{pgfscope}%
\end{pgfscope}%
\begin{pgfscope}%
\pgfsetbuttcap%
\pgfsetroundjoin%
\definecolor{currentfill}{rgb}{0.000000,0.000000,0.000000}%
\pgfsetfillcolor{currentfill}%
\pgfsetlinewidth{0.301125pt}%
\definecolor{currentstroke}{rgb}{0.000000,0.000000,0.000000}%
\pgfsetstrokecolor{currentstroke}%
\pgfsetdash{}{0pt}%
\pgfsys@defobject{currentmarker}{\pgfqpoint{0.000000in}{0.000000in}}{\pgfqpoint{0.000000in}{0.027778in}}{%
\pgfpathmoveto{\pgfqpoint{0.000000in}{0.000000in}}%
\pgfpathlineto{\pgfqpoint{0.000000in}{0.027778in}}%
\pgfusepath{stroke,fill}%
}%
\begin{pgfscope}%
\pgfsys@transformshift{0.912664in}{0.352669in}%
\pgfsys@useobject{currentmarker}{}%
\end{pgfscope}%
\end{pgfscope}%
\begin{pgfscope}%
\pgfsetbuttcap%
\pgfsetroundjoin%
\definecolor{currentfill}{rgb}{0.000000,0.000000,0.000000}%
\pgfsetfillcolor{currentfill}%
\pgfsetlinewidth{0.301125pt}%
\definecolor{currentstroke}{rgb}{0.000000,0.000000,0.000000}%
\pgfsetstrokecolor{currentstroke}%
\pgfsetdash{}{0pt}%
\pgfsys@defobject{currentmarker}{\pgfqpoint{0.000000in}{0.000000in}}{\pgfqpoint{0.000000in}{0.027778in}}{%
\pgfpathmoveto{\pgfqpoint{0.000000in}{0.000000in}}%
\pgfpathlineto{\pgfqpoint{0.000000in}{0.027778in}}%
\pgfusepath{stroke,fill}%
}%
\begin{pgfscope}%
\pgfsys@transformshift{0.980172in}{0.352669in}%
\pgfsys@useobject{currentmarker}{}%
\end{pgfscope}%
\end{pgfscope}%
\begin{pgfscope}%
\pgfsetbuttcap%
\pgfsetroundjoin%
\definecolor{currentfill}{rgb}{0.000000,0.000000,0.000000}%
\pgfsetfillcolor{currentfill}%
\pgfsetlinewidth{0.301125pt}%
\definecolor{currentstroke}{rgb}{0.000000,0.000000,0.000000}%
\pgfsetstrokecolor{currentstroke}%
\pgfsetdash{}{0pt}%
\pgfsys@defobject{currentmarker}{\pgfqpoint{0.000000in}{0.000000in}}{\pgfqpoint{0.000000in}{0.027778in}}{%
\pgfpathmoveto{\pgfqpoint{0.000000in}{0.000000in}}%
\pgfpathlineto{\pgfqpoint{0.000000in}{0.027778in}}%
\pgfusepath{stroke,fill}%
}%
\begin{pgfscope}%
\pgfsys@transformshift{1.047679in}{0.352669in}%
\pgfsys@useobject{currentmarker}{}%
\end{pgfscope}%
\end{pgfscope}%
\begin{pgfscope}%
\pgfsetbuttcap%
\pgfsetroundjoin%
\definecolor{currentfill}{rgb}{0.000000,0.000000,0.000000}%
\pgfsetfillcolor{currentfill}%
\pgfsetlinewidth{0.301125pt}%
\definecolor{currentstroke}{rgb}{0.000000,0.000000,0.000000}%
\pgfsetstrokecolor{currentstroke}%
\pgfsetdash{}{0pt}%
\pgfsys@defobject{currentmarker}{\pgfqpoint{0.000000in}{0.000000in}}{\pgfqpoint{0.000000in}{0.027778in}}{%
\pgfpathmoveto{\pgfqpoint{0.000000in}{0.000000in}}%
\pgfpathlineto{\pgfqpoint{0.000000in}{0.027778in}}%
\pgfusepath{stroke,fill}%
}%
\begin{pgfscope}%
\pgfsys@transformshift{1.115187in}{0.352669in}%
\pgfsys@useobject{currentmarker}{}%
\end{pgfscope}%
\end{pgfscope}%
\begin{pgfscope}%
\pgfsetbuttcap%
\pgfsetroundjoin%
\definecolor{currentfill}{rgb}{0.000000,0.000000,0.000000}%
\pgfsetfillcolor{currentfill}%
\pgfsetlinewidth{0.301125pt}%
\definecolor{currentstroke}{rgb}{0.000000,0.000000,0.000000}%
\pgfsetstrokecolor{currentstroke}%
\pgfsetdash{}{0pt}%
\pgfsys@defobject{currentmarker}{\pgfqpoint{0.000000in}{0.000000in}}{\pgfqpoint{0.000000in}{0.027778in}}{%
\pgfpathmoveto{\pgfqpoint{0.000000in}{0.000000in}}%
\pgfpathlineto{\pgfqpoint{0.000000in}{0.027778in}}%
\pgfusepath{stroke,fill}%
}%
\begin{pgfscope}%
\pgfsys@transformshift{1.250201in}{0.352669in}%
\pgfsys@useobject{currentmarker}{}%
\end{pgfscope}%
\end{pgfscope}%
\begin{pgfscope}%
\pgfsetbuttcap%
\pgfsetroundjoin%
\definecolor{currentfill}{rgb}{0.000000,0.000000,0.000000}%
\pgfsetfillcolor{currentfill}%
\pgfsetlinewidth{0.301125pt}%
\definecolor{currentstroke}{rgb}{0.000000,0.000000,0.000000}%
\pgfsetstrokecolor{currentstroke}%
\pgfsetdash{}{0pt}%
\pgfsys@defobject{currentmarker}{\pgfqpoint{0.000000in}{0.000000in}}{\pgfqpoint{0.000000in}{0.027778in}}{%
\pgfpathmoveto{\pgfqpoint{0.000000in}{0.000000in}}%
\pgfpathlineto{\pgfqpoint{0.000000in}{0.027778in}}%
\pgfusepath{stroke,fill}%
}%
\begin{pgfscope}%
\pgfsys@transformshift{1.317709in}{0.352669in}%
\pgfsys@useobject{currentmarker}{}%
\end{pgfscope}%
\end{pgfscope}%
\begin{pgfscope}%
\pgfsetbuttcap%
\pgfsetroundjoin%
\definecolor{currentfill}{rgb}{0.000000,0.000000,0.000000}%
\pgfsetfillcolor{currentfill}%
\pgfsetlinewidth{0.301125pt}%
\definecolor{currentstroke}{rgb}{0.000000,0.000000,0.000000}%
\pgfsetstrokecolor{currentstroke}%
\pgfsetdash{}{0pt}%
\pgfsys@defobject{currentmarker}{\pgfqpoint{0.000000in}{0.000000in}}{\pgfqpoint{0.000000in}{0.027778in}}{%
\pgfpathmoveto{\pgfqpoint{0.000000in}{0.000000in}}%
\pgfpathlineto{\pgfqpoint{0.000000in}{0.027778in}}%
\pgfusepath{stroke,fill}%
}%
\begin{pgfscope}%
\pgfsys@transformshift{1.385216in}{0.352669in}%
\pgfsys@useobject{currentmarker}{}%
\end{pgfscope}%
\end{pgfscope}%
\begin{pgfscope}%
\pgfsetbuttcap%
\pgfsetroundjoin%
\definecolor{currentfill}{rgb}{0.000000,0.000000,0.000000}%
\pgfsetfillcolor{currentfill}%
\pgfsetlinewidth{0.301125pt}%
\definecolor{currentstroke}{rgb}{0.000000,0.000000,0.000000}%
\pgfsetstrokecolor{currentstroke}%
\pgfsetdash{}{0pt}%
\pgfsys@defobject{currentmarker}{\pgfqpoint{0.000000in}{0.000000in}}{\pgfqpoint{0.000000in}{0.027778in}}{%
\pgfpathmoveto{\pgfqpoint{0.000000in}{0.000000in}}%
\pgfpathlineto{\pgfqpoint{0.000000in}{0.027778in}}%
\pgfusepath{stroke,fill}%
}%
\begin{pgfscope}%
\pgfsys@transformshift{1.452724in}{0.352669in}%
\pgfsys@useobject{currentmarker}{}%
\end{pgfscope}%
\end{pgfscope}%
\begin{pgfscope}%
\pgfsetbuttcap%
\pgfsetroundjoin%
\definecolor{currentfill}{rgb}{0.000000,0.000000,0.000000}%
\pgfsetfillcolor{currentfill}%
\pgfsetlinewidth{0.301125pt}%
\definecolor{currentstroke}{rgb}{0.000000,0.000000,0.000000}%
\pgfsetstrokecolor{currentstroke}%
\pgfsetdash{}{0pt}%
\pgfsys@defobject{currentmarker}{\pgfqpoint{0.000000in}{0.000000in}}{\pgfqpoint{0.000000in}{0.027778in}}{%
\pgfpathmoveto{\pgfqpoint{0.000000in}{0.000000in}}%
\pgfpathlineto{\pgfqpoint{0.000000in}{0.027778in}}%
\pgfusepath{stroke,fill}%
}%
\begin{pgfscope}%
\pgfsys@transformshift{1.587739in}{0.352669in}%
\pgfsys@useobject{currentmarker}{}%
\end{pgfscope}%
\end{pgfscope}%
\begin{pgfscope}%
\pgfsetbuttcap%
\pgfsetroundjoin%
\definecolor{currentfill}{rgb}{0.000000,0.000000,0.000000}%
\pgfsetfillcolor{currentfill}%
\pgfsetlinewidth{0.301125pt}%
\definecolor{currentstroke}{rgb}{0.000000,0.000000,0.000000}%
\pgfsetstrokecolor{currentstroke}%
\pgfsetdash{}{0pt}%
\pgfsys@defobject{currentmarker}{\pgfqpoint{0.000000in}{0.000000in}}{\pgfqpoint{0.000000in}{0.027778in}}{%
\pgfpathmoveto{\pgfqpoint{0.000000in}{0.000000in}}%
\pgfpathlineto{\pgfqpoint{0.000000in}{0.027778in}}%
\pgfusepath{stroke,fill}%
}%
\begin{pgfscope}%
\pgfsys@transformshift{1.655246in}{0.352669in}%
\pgfsys@useobject{currentmarker}{}%
\end{pgfscope}%
\end{pgfscope}%
\begin{pgfscope}%
\pgfsetbuttcap%
\pgfsetroundjoin%
\definecolor{currentfill}{rgb}{0.000000,0.000000,0.000000}%
\pgfsetfillcolor{currentfill}%
\pgfsetlinewidth{0.301125pt}%
\definecolor{currentstroke}{rgb}{0.000000,0.000000,0.000000}%
\pgfsetstrokecolor{currentstroke}%
\pgfsetdash{}{0pt}%
\pgfsys@defobject{currentmarker}{\pgfqpoint{0.000000in}{0.000000in}}{\pgfqpoint{0.000000in}{0.027778in}}{%
\pgfpathmoveto{\pgfqpoint{0.000000in}{0.000000in}}%
\pgfpathlineto{\pgfqpoint{0.000000in}{0.027778in}}%
\pgfusepath{stroke,fill}%
}%
\begin{pgfscope}%
\pgfsys@transformshift{1.722754in}{0.352669in}%
\pgfsys@useobject{currentmarker}{}%
\end{pgfscope}%
\end{pgfscope}%
\begin{pgfscope}%
\pgfsetbuttcap%
\pgfsetroundjoin%
\definecolor{currentfill}{rgb}{0.000000,0.000000,0.000000}%
\pgfsetfillcolor{currentfill}%
\pgfsetlinewidth{0.301125pt}%
\definecolor{currentstroke}{rgb}{0.000000,0.000000,0.000000}%
\pgfsetstrokecolor{currentstroke}%
\pgfsetdash{}{0pt}%
\pgfsys@defobject{currentmarker}{\pgfqpoint{0.000000in}{0.000000in}}{\pgfqpoint{0.000000in}{0.027778in}}{%
\pgfpathmoveto{\pgfqpoint{0.000000in}{0.000000in}}%
\pgfpathlineto{\pgfqpoint{0.000000in}{0.027778in}}%
\pgfusepath{stroke,fill}%
}%
\begin{pgfscope}%
\pgfsys@transformshift{1.790261in}{0.352669in}%
\pgfsys@useobject{currentmarker}{}%
\end{pgfscope}%
\end{pgfscope}%
\begin{pgfscope}%
\definecolor{textcolor}{rgb}{0.000000,0.000000,0.000000}%
\pgfsetstrokecolor{textcolor}%
\pgfsetfillcolor{textcolor}%
\pgftext[x=1.182694in,y=0.140972in,,top]{\color{textcolor}{\sffamily\fontsize{9.000000}{10.800000}\selectfont\catcode`\^=\active\def^{\ifmmode\sp\else\^{}\fi}\catcode`\%=\active\def%{\%}$x$}}%
\end{pgfscope}%
\begin{pgfscope}%
\pgfsetbuttcap%
\pgfsetroundjoin%
\definecolor{currentfill}{rgb}{0.000000,0.000000,0.000000}%
\pgfsetfillcolor{currentfill}%
\pgfsetlinewidth{0.501875pt}%
\definecolor{currentstroke}{rgb}{0.000000,0.000000,0.000000}%
\pgfsetstrokecolor{currentstroke}%
\pgfsetdash{}{0pt}%
\pgfsys@defobject{currentmarker}{\pgfqpoint{0.000000in}{0.000000in}}{\pgfqpoint{0.041667in}{0.000000in}}{%
\pgfpathmoveto{\pgfqpoint{0.000000in}{0.000000in}}%
\pgfpathlineto{\pgfqpoint{0.041667in}{0.000000in}}%
\pgfusepath{stroke,fill}%
}%
\begin{pgfscope}%
\pgfsys@transformshift{0.507620in}{0.352669in}%
\pgfsys@useobject{currentmarker}{}%
\end{pgfscope}%
\end{pgfscope}%
\begin{pgfscope}%
\definecolor{textcolor}{rgb}{0.000000,0.000000,0.000000}%
\pgfsetstrokecolor{textcolor}%
\pgfsetfillcolor{textcolor}%
\pgftext[x=0.196527in, y=0.310460in, left, base]{\color{textcolor}{\sffamily\fontsize{8.000000}{9.600000}\selectfont\catcode`\^=\active\def^{\ifmmode\sp\else\^{}\fi}\catcode`\%=\active\def%{\%}\ensuremath{-}1.5}}%
\end{pgfscope}%
\begin{pgfscope}%
\pgfsetbuttcap%
\pgfsetroundjoin%
\definecolor{currentfill}{rgb}{0.000000,0.000000,0.000000}%
\pgfsetfillcolor{currentfill}%
\pgfsetlinewidth{0.501875pt}%
\definecolor{currentstroke}{rgb}{0.000000,0.000000,0.000000}%
\pgfsetstrokecolor{currentstroke}%
\pgfsetdash{}{0pt}%
\pgfsys@defobject{currentmarker}{\pgfqpoint{0.000000in}{0.000000in}}{\pgfqpoint{0.041667in}{0.000000in}}{%
\pgfpathmoveto{\pgfqpoint{0.000000in}{0.000000in}}%
\pgfpathlineto{\pgfqpoint{0.041667in}{0.000000in}}%
\pgfusepath{stroke,fill}%
}%
\begin{pgfscope}%
\pgfsys@transformshift{0.507620in}{0.605669in}%
\pgfsys@useobject{currentmarker}{}%
\end{pgfscope}%
\end{pgfscope}%
\begin{pgfscope}%
\definecolor{textcolor}{rgb}{0.000000,0.000000,0.000000}%
\pgfsetstrokecolor{textcolor}%
\pgfsetfillcolor{textcolor}%
\pgftext[x=0.196527in, y=0.563460in, left, base]{\color{textcolor}{\sffamily\fontsize{8.000000}{9.600000}\selectfont\catcode`\^=\active\def^{\ifmmode\sp\else\^{}\fi}\catcode`\%=\active\def%{\%}\ensuremath{-}1.0}}%
\end{pgfscope}%
\begin{pgfscope}%
\pgfsetbuttcap%
\pgfsetroundjoin%
\definecolor{currentfill}{rgb}{0.000000,0.000000,0.000000}%
\pgfsetfillcolor{currentfill}%
\pgfsetlinewidth{0.501875pt}%
\definecolor{currentstroke}{rgb}{0.000000,0.000000,0.000000}%
\pgfsetstrokecolor{currentstroke}%
\pgfsetdash{}{0pt}%
\pgfsys@defobject{currentmarker}{\pgfqpoint{0.000000in}{0.000000in}}{\pgfqpoint{0.041667in}{0.000000in}}{%
\pgfpathmoveto{\pgfqpoint{0.000000in}{0.000000in}}%
\pgfpathlineto{\pgfqpoint{0.041667in}{0.000000in}}%
\pgfusepath{stroke,fill}%
}%
\begin{pgfscope}%
\pgfsys@transformshift{0.507620in}{0.858669in}%
\pgfsys@useobject{currentmarker}{}%
\end{pgfscope}%
\end{pgfscope}%
\begin{pgfscope}%
\definecolor{textcolor}{rgb}{0.000000,0.000000,0.000000}%
\pgfsetstrokecolor{textcolor}%
\pgfsetfillcolor{textcolor}%
\pgftext[x=0.196527in, y=0.816460in, left, base]{\color{textcolor}{\sffamily\fontsize{8.000000}{9.600000}\selectfont\catcode`\^=\active\def^{\ifmmode\sp\else\^{}\fi}\catcode`\%=\active\def%{\%}\ensuremath{-}0.5}}%
\end{pgfscope}%
\begin{pgfscope}%
\pgfsetbuttcap%
\pgfsetroundjoin%
\definecolor{currentfill}{rgb}{0.000000,0.000000,0.000000}%
\pgfsetfillcolor{currentfill}%
\pgfsetlinewidth{0.501875pt}%
\definecolor{currentstroke}{rgb}{0.000000,0.000000,0.000000}%
\pgfsetstrokecolor{currentstroke}%
\pgfsetdash{}{0pt}%
\pgfsys@defobject{currentmarker}{\pgfqpoint{0.000000in}{0.000000in}}{\pgfqpoint{0.041667in}{0.000000in}}{%
\pgfpathmoveto{\pgfqpoint{0.000000in}{0.000000in}}%
\pgfpathlineto{\pgfqpoint{0.041667in}{0.000000in}}%
\pgfusepath{stroke,fill}%
}%
\begin{pgfscope}%
\pgfsys@transformshift{0.507620in}{1.111669in}%
\pgfsys@useobject{currentmarker}{}%
\end{pgfscope}%
\end{pgfscope}%
\begin{pgfscope}%
\definecolor{textcolor}{rgb}{0.000000,0.000000,0.000000}%
\pgfsetstrokecolor{textcolor}%
\pgfsetfillcolor{textcolor}%
\pgftext[x=0.282305in, y=1.069460in, left, base]{\color{textcolor}{\sffamily\fontsize{8.000000}{9.600000}\selectfont\catcode`\^=\active\def^{\ifmmode\sp\else\^{}\fi}\catcode`\%=\active\def%{\%}0.0}}%
\end{pgfscope}%
\begin{pgfscope}%
\pgfsetbuttcap%
\pgfsetroundjoin%
\definecolor{currentfill}{rgb}{0.000000,0.000000,0.000000}%
\pgfsetfillcolor{currentfill}%
\pgfsetlinewidth{0.501875pt}%
\definecolor{currentstroke}{rgb}{0.000000,0.000000,0.000000}%
\pgfsetstrokecolor{currentstroke}%
\pgfsetdash{}{0pt}%
\pgfsys@defobject{currentmarker}{\pgfqpoint{0.000000in}{0.000000in}}{\pgfqpoint{0.041667in}{0.000000in}}{%
\pgfpathmoveto{\pgfqpoint{0.000000in}{0.000000in}}%
\pgfpathlineto{\pgfqpoint{0.041667in}{0.000000in}}%
\pgfusepath{stroke,fill}%
}%
\begin{pgfscope}%
\pgfsys@transformshift{0.507620in}{1.364669in}%
\pgfsys@useobject{currentmarker}{}%
\end{pgfscope}%
\end{pgfscope}%
\begin{pgfscope}%
\definecolor{textcolor}{rgb}{0.000000,0.000000,0.000000}%
\pgfsetstrokecolor{textcolor}%
\pgfsetfillcolor{textcolor}%
\pgftext[x=0.282305in, y=1.322460in, left, base]{\color{textcolor}{\sffamily\fontsize{8.000000}{9.600000}\selectfont\catcode`\^=\active\def^{\ifmmode\sp\else\^{}\fi}\catcode`\%=\active\def%{\%}0.5}}%
\end{pgfscope}%
\begin{pgfscope}%
\pgfsetbuttcap%
\pgfsetroundjoin%
\definecolor{currentfill}{rgb}{0.000000,0.000000,0.000000}%
\pgfsetfillcolor{currentfill}%
\pgfsetlinewidth{0.501875pt}%
\definecolor{currentstroke}{rgb}{0.000000,0.000000,0.000000}%
\pgfsetstrokecolor{currentstroke}%
\pgfsetdash{}{0pt}%
\pgfsys@defobject{currentmarker}{\pgfqpoint{0.000000in}{0.000000in}}{\pgfqpoint{0.041667in}{0.000000in}}{%
\pgfpathmoveto{\pgfqpoint{0.000000in}{0.000000in}}%
\pgfpathlineto{\pgfqpoint{0.041667in}{0.000000in}}%
\pgfusepath{stroke,fill}%
}%
\begin{pgfscope}%
\pgfsys@transformshift{0.507620in}{1.617669in}%
\pgfsys@useobject{currentmarker}{}%
\end{pgfscope}%
\end{pgfscope}%
\begin{pgfscope}%
\definecolor{textcolor}{rgb}{0.000000,0.000000,0.000000}%
\pgfsetstrokecolor{textcolor}%
\pgfsetfillcolor{textcolor}%
\pgftext[x=0.282305in, y=1.575460in, left, base]{\color{textcolor}{\sffamily\fontsize{8.000000}{9.600000}\selectfont\catcode`\^=\active\def^{\ifmmode\sp\else\^{}\fi}\catcode`\%=\active\def%{\%}1.0}}%
\end{pgfscope}%
\begin{pgfscope}%
\pgfsetbuttcap%
\pgfsetroundjoin%
\definecolor{currentfill}{rgb}{0.000000,0.000000,0.000000}%
\pgfsetfillcolor{currentfill}%
\pgfsetlinewidth{0.501875pt}%
\definecolor{currentstroke}{rgb}{0.000000,0.000000,0.000000}%
\pgfsetstrokecolor{currentstroke}%
\pgfsetdash{}{0pt}%
\pgfsys@defobject{currentmarker}{\pgfqpoint{0.000000in}{0.000000in}}{\pgfqpoint{0.041667in}{0.000000in}}{%
\pgfpathmoveto{\pgfqpoint{0.000000in}{0.000000in}}%
\pgfpathlineto{\pgfqpoint{0.041667in}{0.000000in}}%
\pgfusepath{stroke,fill}%
}%
\begin{pgfscope}%
\pgfsys@transformshift{0.507620in}{1.870669in}%
\pgfsys@useobject{currentmarker}{}%
\end{pgfscope}%
\end{pgfscope}%
\begin{pgfscope}%
\definecolor{textcolor}{rgb}{0.000000,0.000000,0.000000}%
\pgfsetstrokecolor{textcolor}%
\pgfsetfillcolor{textcolor}%
\pgftext[x=0.282305in, y=1.828460in, left, base]{\color{textcolor}{\sffamily\fontsize{8.000000}{9.600000}\selectfont\catcode`\^=\active\def^{\ifmmode\sp\else\^{}\fi}\catcode`\%=\active\def%{\%}1.5}}%
\end{pgfscope}%
\begin{pgfscope}%
\pgfsetbuttcap%
\pgfsetroundjoin%
\definecolor{currentfill}{rgb}{0.000000,0.000000,0.000000}%
\pgfsetfillcolor{currentfill}%
\pgfsetlinewidth{0.501875pt}%
\definecolor{currentstroke}{rgb}{0.000000,0.000000,0.000000}%
\pgfsetstrokecolor{currentstroke}%
\pgfsetdash{}{0pt}%
\pgfsys@defobject{currentmarker}{\pgfqpoint{0.000000in}{0.000000in}}{\pgfqpoint{0.041667in}{0.000000in}}{%
\pgfpathmoveto{\pgfqpoint{0.000000in}{0.000000in}}%
\pgfpathlineto{\pgfqpoint{0.041667in}{0.000000in}}%
\pgfusepath{stroke,fill}%
}%
\begin{pgfscope}%
\pgfsys@transformshift{0.507620in}{2.123669in}%
\pgfsys@useobject{currentmarker}{}%
\end{pgfscope}%
\end{pgfscope}%
\begin{pgfscope}%
\definecolor{textcolor}{rgb}{0.000000,0.000000,0.000000}%
\pgfsetstrokecolor{textcolor}%
\pgfsetfillcolor{textcolor}%
\pgftext[x=0.282305in, y=2.081460in, left, base]{\color{textcolor}{\sffamily\fontsize{8.000000}{9.600000}\selectfont\catcode`\^=\active\def^{\ifmmode\sp\else\^{}\fi}\catcode`\%=\active\def%{\%}2.0}}%
\end{pgfscope}%
\begin{pgfscope}%
\pgfsetbuttcap%
\pgfsetroundjoin%
\definecolor{currentfill}{rgb}{0.000000,0.000000,0.000000}%
\pgfsetfillcolor{currentfill}%
\pgfsetlinewidth{0.301125pt}%
\definecolor{currentstroke}{rgb}{0.000000,0.000000,0.000000}%
\pgfsetstrokecolor{currentstroke}%
\pgfsetdash{}{0pt}%
\pgfsys@defobject{currentmarker}{\pgfqpoint{0.000000in}{0.000000in}}{\pgfqpoint{0.027778in}{0.000000in}}{%
\pgfpathmoveto{\pgfqpoint{0.000000in}{0.000000in}}%
\pgfpathlineto{\pgfqpoint{0.027778in}{0.000000in}}%
\pgfusepath{stroke,fill}%
}%
\begin{pgfscope}%
\pgfsys@transformshift{0.507620in}{0.403269in}%
\pgfsys@useobject{currentmarker}{}%
\end{pgfscope}%
\end{pgfscope}%
\begin{pgfscope}%
\pgfsetbuttcap%
\pgfsetroundjoin%
\definecolor{currentfill}{rgb}{0.000000,0.000000,0.000000}%
\pgfsetfillcolor{currentfill}%
\pgfsetlinewidth{0.301125pt}%
\definecolor{currentstroke}{rgb}{0.000000,0.000000,0.000000}%
\pgfsetstrokecolor{currentstroke}%
\pgfsetdash{}{0pt}%
\pgfsys@defobject{currentmarker}{\pgfqpoint{0.000000in}{0.000000in}}{\pgfqpoint{0.027778in}{0.000000in}}{%
\pgfpathmoveto{\pgfqpoint{0.000000in}{0.000000in}}%
\pgfpathlineto{\pgfqpoint{0.027778in}{0.000000in}}%
\pgfusepath{stroke,fill}%
}%
\begin{pgfscope}%
\pgfsys@transformshift{0.507620in}{0.453869in}%
\pgfsys@useobject{currentmarker}{}%
\end{pgfscope}%
\end{pgfscope}%
\begin{pgfscope}%
\pgfsetbuttcap%
\pgfsetroundjoin%
\definecolor{currentfill}{rgb}{0.000000,0.000000,0.000000}%
\pgfsetfillcolor{currentfill}%
\pgfsetlinewidth{0.301125pt}%
\definecolor{currentstroke}{rgb}{0.000000,0.000000,0.000000}%
\pgfsetstrokecolor{currentstroke}%
\pgfsetdash{}{0pt}%
\pgfsys@defobject{currentmarker}{\pgfqpoint{0.000000in}{0.000000in}}{\pgfqpoint{0.027778in}{0.000000in}}{%
\pgfpathmoveto{\pgfqpoint{0.000000in}{0.000000in}}%
\pgfpathlineto{\pgfqpoint{0.027778in}{0.000000in}}%
\pgfusepath{stroke,fill}%
}%
\begin{pgfscope}%
\pgfsys@transformshift{0.507620in}{0.504469in}%
\pgfsys@useobject{currentmarker}{}%
\end{pgfscope}%
\end{pgfscope}%
\begin{pgfscope}%
\pgfsetbuttcap%
\pgfsetroundjoin%
\definecolor{currentfill}{rgb}{0.000000,0.000000,0.000000}%
\pgfsetfillcolor{currentfill}%
\pgfsetlinewidth{0.301125pt}%
\definecolor{currentstroke}{rgb}{0.000000,0.000000,0.000000}%
\pgfsetstrokecolor{currentstroke}%
\pgfsetdash{}{0pt}%
\pgfsys@defobject{currentmarker}{\pgfqpoint{0.000000in}{0.000000in}}{\pgfqpoint{0.027778in}{0.000000in}}{%
\pgfpathmoveto{\pgfqpoint{0.000000in}{0.000000in}}%
\pgfpathlineto{\pgfqpoint{0.027778in}{0.000000in}}%
\pgfusepath{stroke,fill}%
}%
\begin{pgfscope}%
\pgfsys@transformshift{0.507620in}{0.555069in}%
\pgfsys@useobject{currentmarker}{}%
\end{pgfscope}%
\end{pgfscope}%
\begin{pgfscope}%
\pgfsetbuttcap%
\pgfsetroundjoin%
\definecolor{currentfill}{rgb}{0.000000,0.000000,0.000000}%
\pgfsetfillcolor{currentfill}%
\pgfsetlinewidth{0.301125pt}%
\definecolor{currentstroke}{rgb}{0.000000,0.000000,0.000000}%
\pgfsetstrokecolor{currentstroke}%
\pgfsetdash{}{0pt}%
\pgfsys@defobject{currentmarker}{\pgfqpoint{0.000000in}{0.000000in}}{\pgfqpoint{0.027778in}{0.000000in}}{%
\pgfpathmoveto{\pgfqpoint{0.000000in}{0.000000in}}%
\pgfpathlineto{\pgfqpoint{0.027778in}{0.000000in}}%
\pgfusepath{stroke,fill}%
}%
\begin{pgfscope}%
\pgfsys@transformshift{0.507620in}{0.656269in}%
\pgfsys@useobject{currentmarker}{}%
\end{pgfscope}%
\end{pgfscope}%
\begin{pgfscope}%
\pgfsetbuttcap%
\pgfsetroundjoin%
\definecolor{currentfill}{rgb}{0.000000,0.000000,0.000000}%
\pgfsetfillcolor{currentfill}%
\pgfsetlinewidth{0.301125pt}%
\definecolor{currentstroke}{rgb}{0.000000,0.000000,0.000000}%
\pgfsetstrokecolor{currentstroke}%
\pgfsetdash{}{0pt}%
\pgfsys@defobject{currentmarker}{\pgfqpoint{0.000000in}{0.000000in}}{\pgfqpoint{0.027778in}{0.000000in}}{%
\pgfpathmoveto{\pgfqpoint{0.000000in}{0.000000in}}%
\pgfpathlineto{\pgfqpoint{0.027778in}{0.000000in}}%
\pgfusepath{stroke,fill}%
}%
\begin{pgfscope}%
\pgfsys@transformshift{0.507620in}{0.706869in}%
\pgfsys@useobject{currentmarker}{}%
\end{pgfscope}%
\end{pgfscope}%
\begin{pgfscope}%
\pgfsetbuttcap%
\pgfsetroundjoin%
\definecolor{currentfill}{rgb}{0.000000,0.000000,0.000000}%
\pgfsetfillcolor{currentfill}%
\pgfsetlinewidth{0.301125pt}%
\definecolor{currentstroke}{rgb}{0.000000,0.000000,0.000000}%
\pgfsetstrokecolor{currentstroke}%
\pgfsetdash{}{0pt}%
\pgfsys@defobject{currentmarker}{\pgfqpoint{0.000000in}{0.000000in}}{\pgfqpoint{0.027778in}{0.000000in}}{%
\pgfpathmoveto{\pgfqpoint{0.000000in}{0.000000in}}%
\pgfpathlineto{\pgfqpoint{0.027778in}{0.000000in}}%
\pgfusepath{stroke,fill}%
}%
\begin{pgfscope}%
\pgfsys@transformshift{0.507620in}{0.757469in}%
\pgfsys@useobject{currentmarker}{}%
\end{pgfscope}%
\end{pgfscope}%
\begin{pgfscope}%
\pgfsetbuttcap%
\pgfsetroundjoin%
\definecolor{currentfill}{rgb}{0.000000,0.000000,0.000000}%
\pgfsetfillcolor{currentfill}%
\pgfsetlinewidth{0.301125pt}%
\definecolor{currentstroke}{rgb}{0.000000,0.000000,0.000000}%
\pgfsetstrokecolor{currentstroke}%
\pgfsetdash{}{0pt}%
\pgfsys@defobject{currentmarker}{\pgfqpoint{0.000000in}{0.000000in}}{\pgfqpoint{0.027778in}{0.000000in}}{%
\pgfpathmoveto{\pgfqpoint{0.000000in}{0.000000in}}%
\pgfpathlineto{\pgfqpoint{0.027778in}{0.000000in}}%
\pgfusepath{stroke,fill}%
}%
\begin{pgfscope}%
\pgfsys@transformshift{0.507620in}{0.808069in}%
\pgfsys@useobject{currentmarker}{}%
\end{pgfscope}%
\end{pgfscope}%
\begin{pgfscope}%
\pgfsetbuttcap%
\pgfsetroundjoin%
\definecolor{currentfill}{rgb}{0.000000,0.000000,0.000000}%
\pgfsetfillcolor{currentfill}%
\pgfsetlinewidth{0.301125pt}%
\definecolor{currentstroke}{rgb}{0.000000,0.000000,0.000000}%
\pgfsetstrokecolor{currentstroke}%
\pgfsetdash{}{0pt}%
\pgfsys@defobject{currentmarker}{\pgfqpoint{0.000000in}{0.000000in}}{\pgfqpoint{0.027778in}{0.000000in}}{%
\pgfpathmoveto{\pgfqpoint{0.000000in}{0.000000in}}%
\pgfpathlineto{\pgfqpoint{0.027778in}{0.000000in}}%
\pgfusepath{stroke,fill}%
}%
\begin{pgfscope}%
\pgfsys@transformshift{0.507620in}{0.909269in}%
\pgfsys@useobject{currentmarker}{}%
\end{pgfscope}%
\end{pgfscope}%
\begin{pgfscope}%
\pgfsetbuttcap%
\pgfsetroundjoin%
\definecolor{currentfill}{rgb}{0.000000,0.000000,0.000000}%
\pgfsetfillcolor{currentfill}%
\pgfsetlinewidth{0.301125pt}%
\definecolor{currentstroke}{rgb}{0.000000,0.000000,0.000000}%
\pgfsetstrokecolor{currentstroke}%
\pgfsetdash{}{0pt}%
\pgfsys@defobject{currentmarker}{\pgfqpoint{0.000000in}{0.000000in}}{\pgfqpoint{0.027778in}{0.000000in}}{%
\pgfpathmoveto{\pgfqpoint{0.000000in}{0.000000in}}%
\pgfpathlineto{\pgfqpoint{0.027778in}{0.000000in}}%
\pgfusepath{stroke,fill}%
}%
\begin{pgfscope}%
\pgfsys@transformshift{0.507620in}{0.959869in}%
\pgfsys@useobject{currentmarker}{}%
\end{pgfscope}%
\end{pgfscope}%
\begin{pgfscope}%
\pgfsetbuttcap%
\pgfsetroundjoin%
\definecolor{currentfill}{rgb}{0.000000,0.000000,0.000000}%
\pgfsetfillcolor{currentfill}%
\pgfsetlinewidth{0.301125pt}%
\definecolor{currentstroke}{rgb}{0.000000,0.000000,0.000000}%
\pgfsetstrokecolor{currentstroke}%
\pgfsetdash{}{0pt}%
\pgfsys@defobject{currentmarker}{\pgfqpoint{0.000000in}{0.000000in}}{\pgfqpoint{0.027778in}{0.000000in}}{%
\pgfpathmoveto{\pgfqpoint{0.000000in}{0.000000in}}%
\pgfpathlineto{\pgfqpoint{0.027778in}{0.000000in}}%
\pgfusepath{stroke,fill}%
}%
\begin{pgfscope}%
\pgfsys@transformshift{0.507620in}{1.010469in}%
\pgfsys@useobject{currentmarker}{}%
\end{pgfscope}%
\end{pgfscope}%
\begin{pgfscope}%
\pgfsetbuttcap%
\pgfsetroundjoin%
\definecolor{currentfill}{rgb}{0.000000,0.000000,0.000000}%
\pgfsetfillcolor{currentfill}%
\pgfsetlinewidth{0.301125pt}%
\definecolor{currentstroke}{rgb}{0.000000,0.000000,0.000000}%
\pgfsetstrokecolor{currentstroke}%
\pgfsetdash{}{0pt}%
\pgfsys@defobject{currentmarker}{\pgfqpoint{0.000000in}{0.000000in}}{\pgfqpoint{0.027778in}{0.000000in}}{%
\pgfpathmoveto{\pgfqpoint{0.000000in}{0.000000in}}%
\pgfpathlineto{\pgfqpoint{0.027778in}{0.000000in}}%
\pgfusepath{stroke,fill}%
}%
\begin{pgfscope}%
\pgfsys@transformshift{0.507620in}{1.061069in}%
\pgfsys@useobject{currentmarker}{}%
\end{pgfscope}%
\end{pgfscope}%
\begin{pgfscope}%
\pgfsetbuttcap%
\pgfsetroundjoin%
\definecolor{currentfill}{rgb}{0.000000,0.000000,0.000000}%
\pgfsetfillcolor{currentfill}%
\pgfsetlinewidth{0.301125pt}%
\definecolor{currentstroke}{rgb}{0.000000,0.000000,0.000000}%
\pgfsetstrokecolor{currentstroke}%
\pgfsetdash{}{0pt}%
\pgfsys@defobject{currentmarker}{\pgfqpoint{0.000000in}{0.000000in}}{\pgfqpoint{0.027778in}{0.000000in}}{%
\pgfpathmoveto{\pgfqpoint{0.000000in}{0.000000in}}%
\pgfpathlineto{\pgfqpoint{0.027778in}{0.000000in}}%
\pgfusepath{stroke,fill}%
}%
\begin{pgfscope}%
\pgfsys@transformshift{0.507620in}{1.162269in}%
\pgfsys@useobject{currentmarker}{}%
\end{pgfscope}%
\end{pgfscope}%
\begin{pgfscope}%
\pgfsetbuttcap%
\pgfsetroundjoin%
\definecolor{currentfill}{rgb}{0.000000,0.000000,0.000000}%
\pgfsetfillcolor{currentfill}%
\pgfsetlinewidth{0.301125pt}%
\definecolor{currentstroke}{rgb}{0.000000,0.000000,0.000000}%
\pgfsetstrokecolor{currentstroke}%
\pgfsetdash{}{0pt}%
\pgfsys@defobject{currentmarker}{\pgfqpoint{0.000000in}{0.000000in}}{\pgfqpoint{0.027778in}{0.000000in}}{%
\pgfpathmoveto{\pgfqpoint{0.000000in}{0.000000in}}%
\pgfpathlineto{\pgfqpoint{0.027778in}{0.000000in}}%
\pgfusepath{stroke,fill}%
}%
\begin{pgfscope}%
\pgfsys@transformshift{0.507620in}{1.212869in}%
\pgfsys@useobject{currentmarker}{}%
\end{pgfscope}%
\end{pgfscope}%
\begin{pgfscope}%
\pgfsetbuttcap%
\pgfsetroundjoin%
\definecolor{currentfill}{rgb}{0.000000,0.000000,0.000000}%
\pgfsetfillcolor{currentfill}%
\pgfsetlinewidth{0.301125pt}%
\definecolor{currentstroke}{rgb}{0.000000,0.000000,0.000000}%
\pgfsetstrokecolor{currentstroke}%
\pgfsetdash{}{0pt}%
\pgfsys@defobject{currentmarker}{\pgfqpoint{0.000000in}{0.000000in}}{\pgfqpoint{0.027778in}{0.000000in}}{%
\pgfpathmoveto{\pgfqpoint{0.000000in}{0.000000in}}%
\pgfpathlineto{\pgfqpoint{0.027778in}{0.000000in}}%
\pgfusepath{stroke,fill}%
}%
\begin{pgfscope}%
\pgfsys@transformshift{0.507620in}{1.263469in}%
\pgfsys@useobject{currentmarker}{}%
\end{pgfscope}%
\end{pgfscope}%
\begin{pgfscope}%
\pgfsetbuttcap%
\pgfsetroundjoin%
\definecolor{currentfill}{rgb}{0.000000,0.000000,0.000000}%
\pgfsetfillcolor{currentfill}%
\pgfsetlinewidth{0.301125pt}%
\definecolor{currentstroke}{rgb}{0.000000,0.000000,0.000000}%
\pgfsetstrokecolor{currentstroke}%
\pgfsetdash{}{0pt}%
\pgfsys@defobject{currentmarker}{\pgfqpoint{0.000000in}{0.000000in}}{\pgfqpoint{0.027778in}{0.000000in}}{%
\pgfpathmoveto{\pgfqpoint{0.000000in}{0.000000in}}%
\pgfpathlineto{\pgfqpoint{0.027778in}{0.000000in}}%
\pgfusepath{stroke,fill}%
}%
\begin{pgfscope}%
\pgfsys@transformshift{0.507620in}{1.314069in}%
\pgfsys@useobject{currentmarker}{}%
\end{pgfscope}%
\end{pgfscope}%
\begin{pgfscope}%
\pgfsetbuttcap%
\pgfsetroundjoin%
\definecolor{currentfill}{rgb}{0.000000,0.000000,0.000000}%
\pgfsetfillcolor{currentfill}%
\pgfsetlinewidth{0.301125pt}%
\definecolor{currentstroke}{rgb}{0.000000,0.000000,0.000000}%
\pgfsetstrokecolor{currentstroke}%
\pgfsetdash{}{0pt}%
\pgfsys@defobject{currentmarker}{\pgfqpoint{0.000000in}{0.000000in}}{\pgfqpoint{0.027778in}{0.000000in}}{%
\pgfpathmoveto{\pgfqpoint{0.000000in}{0.000000in}}%
\pgfpathlineto{\pgfqpoint{0.027778in}{0.000000in}}%
\pgfusepath{stroke,fill}%
}%
\begin{pgfscope}%
\pgfsys@transformshift{0.507620in}{1.415269in}%
\pgfsys@useobject{currentmarker}{}%
\end{pgfscope}%
\end{pgfscope}%
\begin{pgfscope}%
\pgfsetbuttcap%
\pgfsetroundjoin%
\definecolor{currentfill}{rgb}{0.000000,0.000000,0.000000}%
\pgfsetfillcolor{currentfill}%
\pgfsetlinewidth{0.301125pt}%
\definecolor{currentstroke}{rgb}{0.000000,0.000000,0.000000}%
\pgfsetstrokecolor{currentstroke}%
\pgfsetdash{}{0pt}%
\pgfsys@defobject{currentmarker}{\pgfqpoint{0.000000in}{0.000000in}}{\pgfqpoint{0.027778in}{0.000000in}}{%
\pgfpathmoveto{\pgfqpoint{0.000000in}{0.000000in}}%
\pgfpathlineto{\pgfqpoint{0.027778in}{0.000000in}}%
\pgfusepath{stroke,fill}%
}%
\begin{pgfscope}%
\pgfsys@transformshift{0.507620in}{1.465869in}%
\pgfsys@useobject{currentmarker}{}%
\end{pgfscope}%
\end{pgfscope}%
\begin{pgfscope}%
\pgfsetbuttcap%
\pgfsetroundjoin%
\definecolor{currentfill}{rgb}{0.000000,0.000000,0.000000}%
\pgfsetfillcolor{currentfill}%
\pgfsetlinewidth{0.301125pt}%
\definecolor{currentstroke}{rgb}{0.000000,0.000000,0.000000}%
\pgfsetstrokecolor{currentstroke}%
\pgfsetdash{}{0pt}%
\pgfsys@defobject{currentmarker}{\pgfqpoint{0.000000in}{0.000000in}}{\pgfqpoint{0.027778in}{0.000000in}}{%
\pgfpathmoveto{\pgfqpoint{0.000000in}{0.000000in}}%
\pgfpathlineto{\pgfqpoint{0.027778in}{0.000000in}}%
\pgfusepath{stroke,fill}%
}%
\begin{pgfscope}%
\pgfsys@transformshift{0.507620in}{1.516469in}%
\pgfsys@useobject{currentmarker}{}%
\end{pgfscope}%
\end{pgfscope}%
\begin{pgfscope}%
\pgfsetbuttcap%
\pgfsetroundjoin%
\definecolor{currentfill}{rgb}{0.000000,0.000000,0.000000}%
\pgfsetfillcolor{currentfill}%
\pgfsetlinewidth{0.301125pt}%
\definecolor{currentstroke}{rgb}{0.000000,0.000000,0.000000}%
\pgfsetstrokecolor{currentstroke}%
\pgfsetdash{}{0pt}%
\pgfsys@defobject{currentmarker}{\pgfqpoint{0.000000in}{0.000000in}}{\pgfqpoint{0.027778in}{0.000000in}}{%
\pgfpathmoveto{\pgfqpoint{0.000000in}{0.000000in}}%
\pgfpathlineto{\pgfqpoint{0.027778in}{0.000000in}}%
\pgfusepath{stroke,fill}%
}%
\begin{pgfscope}%
\pgfsys@transformshift{0.507620in}{1.567069in}%
\pgfsys@useobject{currentmarker}{}%
\end{pgfscope}%
\end{pgfscope}%
\begin{pgfscope}%
\pgfsetbuttcap%
\pgfsetroundjoin%
\definecolor{currentfill}{rgb}{0.000000,0.000000,0.000000}%
\pgfsetfillcolor{currentfill}%
\pgfsetlinewidth{0.301125pt}%
\definecolor{currentstroke}{rgb}{0.000000,0.000000,0.000000}%
\pgfsetstrokecolor{currentstroke}%
\pgfsetdash{}{0pt}%
\pgfsys@defobject{currentmarker}{\pgfqpoint{0.000000in}{0.000000in}}{\pgfqpoint{0.027778in}{0.000000in}}{%
\pgfpathmoveto{\pgfqpoint{0.000000in}{0.000000in}}%
\pgfpathlineto{\pgfqpoint{0.027778in}{0.000000in}}%
\pgfusepath{stroke,fill}%
}%
\begin{pgfscope}%
\pgfsys@transformshift{0.507620in}{1.668269in}%
\pgfsys@useobject{currentmarker}{}%
\end{pgfscope}%
\end{pgfscope}%
\begin{pgfscope}%
\pgfsetbuttcap%
\pgfsetroundjoin%
\definecolor{currentfill}{rgb}{0.000000,0.000000,0.000000}%
\pgfsetfillcolor{currentfill}%
\pgfsetlinewidth{0.301125pt}%
\definecolor{currentstroke}{rgb}{0.000000,0.000000,0.000000}%
\pgfsetstrokecolor{currentstroke}%
\pgfsetdash{}{0pt}%
\pgfsys@defobject{currentmarker}{\pgfqpoint{0.000000in}{0.000000in}}{\pgfqpoint{0.027778in}{0.000000in}}{%
\pgfpathmoveto{\pgfqpoint{0.000000in}{0.000000in}}%
\pgfpathlineto{\pgfqpoint{0.027778in}{0.000000in}}%
\pgfusepath{stroke,fill}%
}%
\begin{pgfscope}%
\pgfsys@transformshift{0.507620in}{1.718869in}%
\pgfsys@useobject{currentmarker}{}%
\end{pgfscope}%
\end{pgfscope}%
\begin{pgfscope}%
\pgfsetbuttcap%
\pgfsetroundjoin%
\definecolor{currentfill}{rgb}{0.000000,0.000000,0.000000}%
\pgfsetfillcolor{currentfill}%
\pgfsetlinewidth{0.301125pt}%
\definecolor{currentstroke}{rgb}{0.000000,0.000000,0.000000}%
\pgfsetstrokecolor{currentstroke}%
\pgfsetdash{}{0pt}%
\pgfsys@defobject{currentmarker}{\pgfqpoint{0.000000in}{0.000000in}}{\pgfqpoint{0.027778in}{0.000000in}}{%
\pgfpathmoveto{\pgfqpoint{0.000000in}{0.000000in}}%
\pgfpathlineto{\pgfqpoint{0.027778in}{0.000000in}}%
\pgfusepath{stroke,fill}%
}%
\begin{pgfscope}%
\pgfsys@transformshift{0.507620in}{1.769469in}%
\pgfsys@useobject{currentmarker}{}%
\end{pgfscope}%
\end{pgfscope}%
\begin{pgfscope}%
\pgfsetbuttcap%
\pgfsetroundjoin%
\definecolor{currentfill}{rgb}{0.000000,0.000000,0.000000}%
\pgfsetfillcolor{currentfill}%
\pgfsetlinewidth{0.301125pt}%
\definecolor{currentstroke}{rgb}{0.000000,0.000000,0.000000}%
\pgfsetstrokecolor{currentstroke}%
\pgfsetdash{}{0pt}%
\pgfsys@defobject{currentmarker}{\pgfqpoint{0.000000in}{0.000000in}}{\pgfqpoint{0.027778in}{0.000000in}}{%
\pgfpathmoveto{\pgfqpoint{0.000000in}{0.000000in}}%
\pgfpathlineto{\pgfqpoint{0.027778in}{0.000000in}}%
\pgfusepath{stroke,fill}%
}%
\begin{pgfscope}%
\pgfsys@transformshift{0.507620in}{1.820069in}%
\pgfsys@useobject{currentmarker}{}%
\end{pgfscope}%
\end{pgfscope}%
\begin{pgfscope}%
\pgfsetbuttcap%
\pgfsetroundjoin%
\definecolor{currentfill}{rgb}{0.000000,0.000000,0.000000}%
\pgfsetfillcolor{currentfill}%
\pgfsetlinewidth{0.301125pt}%
\definecolor{currentstroke}{rgb}{0.000000,0.000000,0.000000}%
\pgfsetstrokecolor{currentstroke}%
\pgfsetdash{}{0pt}%
\pgfsys@defobject{currentmarker}{\pgfqpoint{0.000000in}{0.000000in}}{\pgfqpoint{0.027778in}{0.000000in}}{%
\pgfpathmoveto{\pgfqpoint{0.000000in}{0.000000in}}%
\pgfpathlineto{\pgfqpoint{0.027778in}{0.000000in}}%
\pgfusepath{stroke,fill}%
}%
\begin{pgfscope}%
\pgfsys@transformshift{0.507620in}{1.921269in}%
\pgfsys@useobject{currentmarker}{}%
\end{pgfscope}%
\end{pgfscope}%
\begin{pgfscope}%
\pgfsetbuttcap%
\pgfsetroundjoin%
\definecolor{currentfill}{rgb}{0.000000,0.000000,0.000000}%
\pgfsetfillcolor{currentfill}%
\pgfsetlinewidth{0.301125pt}%
\definecolor{currentstroke}{rgb}{0.000000,0.000000,0.000000}%
\pgfsetstrokecolor{currentstroke}%
\pgfsetdash{}{0pt}%
\pgfsys@defobject{currentmarker}{\pgfqpoint{0.000000in}{0.000000in}}{\pgfqpoint{0.027778in}{0.000000in}}{%
\pgfpathmoveto{\pgfqpoint{0.000000in}{0.000000in}}%
\pgfpathlineto{\pgfqpoint{0.027778in}{0.000000in}}%
\pgfusepath{stroke,fill}%
}%
\begin{pgfscope}%
\pgfsys@transformshift{0.507620in}{1.971869in}%
\pgfsys@useobject{currentmarker}{}%
\end{pgfscope}%
\end{pgfscope}%
\begin{pgfscope}%
\pgfsetbuttcap%
\pgfsetroundjoin%
\definecolor{currentfill}{rgb}{0.000000,0.000000,0.000000}%
\pgfsetfillcolor{currentfill}%
\pgfsetlinewidth{0.301125pt}%
\definecolor{currentstroke}{rgb}{0.000000,0.000000,0.000000}%
\pgfsetstrokecolor{currentstroke}%
\pgfsetdash{}{0pt}%
\pgfsys@defobject{currentmarker}{\pgfqpoint{0.000000in}{0.000000in}}{\pgfqpoint{0.027778in}{0.000000in}}{%
\pgfpathmoveto{\pgfqpoint{0.000000in}{0.000000in}}%
\pgfpathlineto{\pgfqpoint{0.027778in}{0.000000in}}%
\pgfusepath{stroke,fill}%
}%
\begin{pgfscope}%
\pgfsys@transformshift{0.507620in}{2.022469in}%
\pgfsys@useobject{currentmarker}{}%
\end{pgfscope}%
\end{pgfscope}%
\begin{pgfscope}%
\pgfsetbuttcap%
\pgfsetroundjoin%
\definecolor{currentfill}{rgb}{0.000000,0.000000,0.000000}%
\pgfsetfillcolor{currentfill}%
\pgfsetlinewidth{0.301125pt}%
\definecolor{currentstroke}{rgb}{0.000000,0.000000,0.000000}%
\pgfsetstrokecolor{currentstroke}%
\pgfsetdash{}{0pt}%
\pgfsys@defobject{currentmarker}{\pgfqpoint{0.000000in}{0.000000in}}{\pgfqpoint{0.027778in}{0.000000in}}{%
\pgfpathmoveto{\pgfqpoint{0.000000in}{0.000000in}}%
\pgfpathlineto{\pgfqpoint{0.027778in}{0.000000in}}%
\pgfusepath{stroke,fill}%
}%
\begin{pgfscope}%
\pgfsys@transformshift{0.507620in}{2.073069in}%
\pgfsys@useobject{currentmarker}{}%
\end{pgfscope}%
\end{pgfscope}%
\begin{pgfscope}%
\definecolor{textcolor}{rgb}{0.000000,0.000000,0.000000}%
\pgfsetstrokecolor{textcolor}%
\pgfsetfillcolor{textcolor}%
\pgftext[x=0.140972in,y=1.238169in,,bottom,rotate=90.000000]{\color{textcolor}{\sffamily\fontsize{9.000000}{10.800000}\selectfont\catcode`\^=\active\def^{\ifmmode\sp\else\^{}\fi}\catcode`\%=\active\def%{\%}$p_N(x)$}}%
\end{pgfscope}%
\begin{pgfscope}%
\pgfsetrectcap%
\pgfsetmiterjoin%
\pgfsetlinewidth{0.602250pt}%
\definecolor{currentstroke}{rgb}{0.000000,0.000000,0.000000}%
\pgfsetstrokecolor{currentstroke}%
\pgfsetdash{}{0pt}%
\pgfpathmoveto{\pgfqpoint{0.507620in}{0.352669in}}%
\pgfpathlineto{\pgfqpoint{0.507620in}{2.123669in}}%
\pgfusepath{stroke}%
\end{pgfscope}%
\begin{pgfscope}%
\pgfsetrectcap%
\pgfsetmiterjoin%
\pgfsetlinewidth{0.602250pt}%
\definecolor{currentstroke}{rgb}{0.000000,0.000000,0.000000}%
\pgfsetstrokecolor{currentstroke}%
\pgfsetdash{}{0pt}%
\pgfpathmoveto{\pgfqpoint{0.507620in}{0.352669in}}%
\pgfpathlineto{\pgfqpoint{1.857769in}{0.352669in}}%
\pgfusepath{stroke}%
\end{pgfscope}%
\begin{pgfscope}%
\pgfpathrectangle{\pgfqpoint{0.507620in}{0.352669in}}{\pgfqpoint{1.350149in}{1.771000in}}%
\pgfusepath{clip}%
\pgfsetrectcap%
\pgfsetroundjoin%
\pgfsetlinewidth{0.903375pt}%
\definecolor{currentstroke}{rgb}{0.250980,0.400000,0.600000}%
\pgfsetstrokecolor{currentstroke}%
\pgfsetdash{}{0pt}%
\pgfpathmoveto{\pgfqpoint{0.507620in}{1.131130in}}%
\pgfpathlineto{\pgfqpoint{0.521135in}{1.118586in}}%
\pgfpathlineto{\pgfqpoint{0.534650in}{1.108148in}}%
\pgfpathlineto{\pgfqpoint{0.548165in}{1.099705in}}%
\pgfpathlineto{\pgfqpoint{0.561680in}{1.093144in}}%
\pgfpathlineto{\pgfqpoint{0.575195in}{1.088356in}}%
\pgfpathlineto{\pgfqpoint{0.588710in}{1.085235in}}%
\pgfpathlineto{\pgfqpoint{0.602225in}{1.083677in}}%
\pgfpathlineto{\pgfqpoint{0.615740in}{1.083580in}}%
\pgfpathlineto{\pgfqpoint{0.630606in}{1.085042in}}%
\pgfpathlineto{\pgfqpoint{0.645473in}{1.088022in}}%
\pgfpathlineto{\pgfqpoint{0.661691in}{1.092856in}}%
\pgfpathlineto{\pgfqpoint{0.679260in}{1.099780in}}%
\pgfpathlineto{\pgfqpoint{0.698181in}{1.108978in}}%
\pgfpathlineto{\pgfqpoint{0.718454in}{1.120579in}}%
\pgfpathlineto{\pgfqpoint{0.741429in}{1.135567in}}%
\pgfpathlineto{\pgfqpoint{0.767108in}{1.154176in}}%
\pgfpathlineto{\pgfqpoint{0.798192in}{1.178624in}}%
\pgfpathlineto{\pgfqpoint{0.842792in}{1.215793in}}%
\pgfpathlineto{\pgfqpoint{0.917124in}{1.277760in}}%
\pgfpathlineto{\pgfqpoint{0.952263in}{1.304943in}}%
\pgfpathlineto{\pgfqpoint{0.981996in}{1.326072in}}%
\pgfpathlineto{\pgfqpoint{1.009026in}{1.343431in}}%
\pgfpathlineto{\pgfqpoint{1.033353in}{1.357324in}}%
\pgfpathlineto{\pgfqpoint{1.056329in}{1.368787in}}%
\pgfpathlineto{\pgfqpoint{1.077953in}{1.377996in}}%
\pgfpathlineto{\pgfqpoint{1.099577in}{1.385590in}}%
\pgfpathlineto{\pgfqpoint{1.119849in}{1.391182in}}%
\pgfpathlineto{\pgfqpoint{1.140122in}{1.395253in}}%
\pgfpathlineto{\pgfqpoint{1.160394in}{1.397772in}}%
\pgfpathlineto{\pgfqpoint{1.180667in}{1.398719in}}%
\pgfpathlineto{\pgfqpoint{1.200939in}{1.398087in}}%
\pgfpathlineto{\pgfqpoint{1.221212in}{1.395882in}}%
\pgfpathlineto{\pgfqpoint{1.241484in}{1.392119in}}%
\pgfpathlineto{\pgfqpoint{1.261757in}{1.386828in}}%
\pgfpathlineto{\pgfqpoint{1.282029in}{1.380049in}}%
\pgfpathlineto{\pgfqpoint{1.303653in}{1.371237in}}%
\pgfpathlineto{\pgfqpoint{1.326629in}{1.360171in}}%
\pgfpathlineto{\pgfqpoint{1.349604in}{1.347462in}}%
\pgfpathlineto{\pgfqpoint{1.373931in}{1.332363in}}%
\pgfpathlineto{\pgfqpoint{1.400961in}{1.313822in}}%
\pgfpathlineto{\pgfqpoint{1.430694in}{1.291611in}}%
\pgfpathlineto{\pgfqpoint{1.467185in}{1.262375in}}%
\pgfpathlineto{\pgfqpoint{1.523948in}{1.214649in}}%
\pgfpathlineto{\pgfqpoint{1.576656in}{1.171007in}}%
\pgfpathlineto{\pgfqpoint{1.607741in}{1.147122in}}%
\pgfpathlineto{\pgfqpoint{1.633419in}{1.129182in}}%
\pgfpathlineto{\pgfqpoint{1.656395in}{1.114958in}}%
\pgfpathlineto{\pgfqpoint{1.676667in}{1.104167in}}%
\pgfpathlineto{\pgfqpoint{1.694237in}{1.096378in}}%
\pgfpathlineto{\pgfqpoint{1.711806in}{1.090242in}}%
\pgfpathlineto{\pgfqpoint{1.728024in}{1.086216in}}%
\pgfpathlineto{\pgfqpoint{1.742891in}{1.084048in}}%
\pgfpathlineto{\pgfqpoint{1.757757in}{1.083469in}}%
\pgfpathlineto{\pgfqpoint{1.771272in}{1.084430in}}%
\pgfpathlineto{\pgfqpoint{1.784787in}{1.086914in}}%
\pgfpathlineto{\pgfqpoint{1.798303in}{1.091022in}}%
\pgfpathlineto{\pgfqpoint{1.810466in}{1.096196in}}%
\pgfpathlineto{\pgfqpoint{1.822630in}{1.102850in}}%
\pgfpathlineto{\pgfqpoint{1.834793in}{1.111065in}}%
\pgfpathlineto{\pgfqpoint{1.846957in}{1.120922in}}%
\pgfpathlineto{\pgfqpoint{1.857769in}{1.131130in}}%
\pgfpathlineto{\pgfqpoint{1.857769in}{1.131130in}}%
\pgfusepath{stroke}%
\end{pgfscope}%
\begin{pgfscope}%
\pgfpathrectangle{\pgfqpoint{0.507620in}{0.352669in}}{\pgfqpoint{1.350149in}{1.771000in}}%
\pgfusepath{clip}%
\pgfsetrectcap%
\pgfsetroundjoin%
\pgfsetlinewidth{0.903375pt}%
\definecolor{currentstroke}{rgb}{0.768627,0.556863,0.211765}%
\pgfsetstrokecolor{currentstroke}%
\pgfsetdash{}{0pt}%
\pgfpathmoveto{\pgfqpoint{0.507620in}{1.131130in}}%
\pgfpathlineto{\pgfqpoint{0.514377in}{1.479360in}}%
\pgfpathlineto{\pgfqpoint{0.519783in}{1.690266in}}%
\pgfpathlineto{\pgfqpoint{0.525189in}{1.849550in}}%
\pgfpathlineto{\pgfqpoint{0.530595in}{1.964305in}}%
\pgfpathlineto{\pgfqpoint{0.534650in}{2.025027in}}%
\pgfpathlineto{\pgfqpoint{0.538704in}{2.066783in}}%
\pgfpathlineto{\pgfqpoint{0.542759in}{2.091846in}}%
\pgfpathlineto{\pgfqpoint{0.545462in}{2.100324in}}%
\pgfpathlineto{\pgfqpoint{0.546813in}{2.102313in}}%
\pgfpathlineto{\pgfqpoint{0.548165in}{2.102896in}}%
\pgfpathlineto{\pgfqpoint{0.549516in}{2.102140in}}%
\pgfpathlineto{\pgfqpoint{0.552219in}{2.096869in}}%
\pgfpathlineto{\pgfqpoint{0.554922in}{2.087001in}}%
\pgfpathlineto{\pgfqpoint{0.558977in}{2.064605in}}%
\pgfpathlineto{\pgfqpoint{0.564383in}{2.022825in}}%
\pgfpathlineto{\pgfqpoint{0.571140in}{1.955785in}}%
\pgfpathlineto{\pgfqpoint{0.580601in}{1.843512in}}%
\pgfpathlineto{\pgfqpoint{0.622497in}{1.319672in}}%
\pgfpathlineto{\pgfqpoint{0.631958in}{1.227454in}}%
\pgfpathlineto{\pgfqpoint{0.641418in}{1.150241in}}%
\pgfpathlineto{\pgfqpoint{0.649527in}{1.096344in}}%
\pgfpathlineto{\pgfqpoint{0.656285in}{1.059984in}}%
\pgfpathlineto{\pgfqpoint{0.663042in}{1.031122in}}%
\pgfpathlineto{\pgfqpoint{0.669800in}{1.009370in}}%
\pgfpathlineto{\pgfqpoint{0.675206in}{0.996767in}}%
\pgfpathlineto{\pgfqpoint{0.680612in}{0.988129in}}%
\pgfpathlineto{\pgfqpoint{0.684666in}{0.984073in}}%
\pgfpathlineto{\pgfqpoint{0.688721in}{0.981946in}}%
\pgfpathlineto{\pgfqpoint{0.692775in}{0.981613in}}%
\pgfpathlineto{\pgfqpoint{0.696830in}{0.982936in}}%
\pgfpathlineto{\pgfqpoint{0.700884in}{0.985777in}}%
\pgfpathlineto{\pgfqpoint{0.706290in}{0.991687in}}%
\pgfpathlineto{\pgfqpoint{0.713048in}{1.002027in}}%
\pgfpathlineto{\pgfqpoint{0.721157in}{1.017950in}}%
\pgfpathlineto{\pgfqpoint{0.730617in}{1.040092in}}%
\pgfpathlineto{\pgfqpoint{0.745484in}{1.079240in}}%
\pgfpathlineto{\pgfqpoint{0.771162in}{1.146955in}}%
\pgfpathlineto{\pgfqpoint{0.783326in}{1.174729in}}%
\pgfpathlineto{\pgfqpoint{0.794138in}{1.195683in}}%
\pgfpathlineto{\pgfqpoint{0.803598in}{1.210743in}}%
\pgfpathlineto{\pgfqpoint{0.811707in}{1.221104in}}%
\pgfpathlineto{\pgfqpoint{0.819816in}{1.229109in}}%
\pgfpathlineto{\pgfqpoint{0.827925in}{1.234830in}}%
\pgfpathlineto{\pgfqpoint{0.834683in}{1.237945in}}%
\pgfpathlineto{\pgfqpoint{0.841440in}{1.239669in}}%
\pgfpathlineto{\pgfqpoint{0.849549in}{1.240080in}}%
\pgfpathlineto{\pgfqpoint{0.857658in}{1.238923in}}%
\pgfpathlineto{\pgfqpoint{0.867119in}{1.235967in}}%
\pgfpathlineto{\pgfqpoint{0.879282in}{1.230371in}}%
\pgfpathlineto{\pgfqpoint{0.918476in}{1.210610in}}%
\pgfpathlineto{\pgfqpoint{0.927936in}{1.208029in}}%
\pgfpathlineto{\pgfqpoint{0.936045in}{1.207138in}}%
\pgfpathlineto{\pgfqpoint{0.944154in}{1.207675in}}%
\pgfpathlineto{\pgfqpoint{0.952263in}{1.209803in}}%
\pgfpathlineto{\pgfqpoint{0.960372in}{1.213651in}}%
\pgfpathlineto{\pgfqpoint{0.968481in}{1.219310in}}%
\pgfpathlineto{\pgfqpoint{0.976590in}{1.226833in}}%
\pgfpathlineto{\pgfqpoint{0.984699in}{1.236236in}}%
\pgfpathlineto{\pgfqpoint{0.994160in}{1.249552in}}%
\pgfpathlineto{\pgfqpoint{1.004972in}{1.267744in}}%
\pgfpathlineto{\pgfqpoint{1.015784in}{1.288872in}}%
\pgfpathlineto{\pgfqpoint{1.029299in}{1.318888in}}%
\pgfpathlineto{\pgfqpoint{1.045517in}{1.359056in}}%
\pgfpathlineto{\pgfqpoint{1.069844in}{1.424018in}}%
\pgfpathlineto{\pgfqpoint{1.096874in}{1.495339in}}%
\pgfpathlineto{\pgfqpoint{1.111740in}{1.530812in}}%
\pgfpathlineto{\pgfqpoint{1.123904in}{1.556492in}}%
\pgfpathlineto{\pgfqpoint{1.134716in}{1.576157in}}%
\pgfpathlineto{\pgfqpoint{1.144176in}{1.590558in}}%
\pgfpathlineto{\pgfqpoint{1.152285in}{1.600619in}}%
\pgfpathlineto{\pgfqpoint{1.160394in}{1.608436in}}%
\pgfpathlineto{\pgfqpoint{1.167152in}{1.613166in}}%
\pgfpathlineto{\pgfqpoint{1.173909in}{1.616226in}}%
\pgfpathlineto{\pgfqpoint{1.180667in}{1.617592in}}%
\pgfpathlineto{\pgfqpoint{1.187424in}{1.617250in}}%
\pgfpathlineto{\pgfqpoint{1.194182in}{1.615204in}}%
\pgfpathlineto{\pgfqpoint{1.200939in}{1.611472in}}%
\pgfpathlineto{\pgfqpoint{1.207697in}{1.606086in}}%
\pgfpathlineto{\pgfqpoint{1.215806in}{1.597509in}}%
\pgfpathlineto{\pgfqpoint{1.223915in}{1.586728in}}%
\pgfpathlineto{\pgfqpoint{1.233375in}{1.571547in}}%
\pgfpathlineto{\pgfqpoint{1.244187in}{1.551089in}}%
\pgfpathlineto{\pgfqpoint{1.256351in}{1.524662in}}%
\pgfpathlineto{\pgfqpoint{1.271217in}{1.488511in}}%
\pgfpathlineto{\pgfqpoint{1.292841in}{1.431337in}}%
\pgfpathlineto{\pgfqpoint{1.326629in}{1.341866in}}%
\pgfpathlineto{\pgfqpoint{1.342847in}{1.303422in}}%
\pgfpathlineto{\pgfqpoint{1.356362in}{1.275339in}}%
\pgfpathlineto{\pgfqpoint{1.367174in}{1.256012in}}%
\pgfpathlineto{\pgfqpoint{1.377986in}{1.239785in}}%
\pgfpathlineto{\pgfqpoint{1.387446in}{1.228269in}}%
\pgfpathlineto{\pgfqpoint{1.395555in}{1.220433in}}%
\pgfpathlineto{\pgfqpoint{1.403664in}{1.214466in}}%
\pgfpathlineto{\pgfqpoint{1.411773in}{1.210322in}}%
\pgfpathlineto{\pgfqpoint{1.419882in}{1.207915in}}%
\pgfpathlineto{\pgfqpoint{1.427991in}{1.207123in}}%
\pgfpathlineto{\pgfqpoint{1.436100in}{1.207790in}}%
\pgfpathlineto{\pgfqpoint{1.445561in}{1.210152in}}%
\pgfpathlineto{\pgfqpoint{1.456373in}{1.214498in}}%
\pgfpathlineto{\pgfqpoint{1.471239in}{1.222238in}}%
\pgfpathlineto{\pgfqpoint{1.494215in}{1.234283in}}%
\pgfpathlineto{\pgfqpoint{1.505027in}{1.238237in}}%
\pgfpathlineto{\pgfqpoint{1.513136in}{1.239855in}}%
\pgfpathlineto{\pgfqpoint{1.521245in}{1.239997in}}%
\pgfpathlineto{\pgfqpoint{1.528002in}{1.238795in}}%
\pgfpathlineto{\pgfqpoint{1.534760in}{1.236250in}}%
\pgfpathlineto{\pgfqpoint{1.541517in}{1.232248in}}%
\pgfpathlineto{\pgfqpoint{1.548275in}{1.226700in}}%
\pgfpathlineto{\pgfqpoint{1.556384in}{1.217914in}}%
\pgfpathlineto{\pgfqpoint{1.564493in}{1.206764in}}%
\pgfpathlineto{\pgfqpoint{1.573953in}{1.190807in}}%
\pgfpathlineto{\pgfqpoint{1.583414in}{1.171848in}}%
\pgfpathlineto{\pgfqpoint{1.594226in}{1.146955in}}%
\pgfpathlineto{\pgfqpoint{1.609092in}{1.108585in}}%
\pgfpathlineto{\pgfqpoint{1.638825in}{1.030223in}}%
\pgfpathlineto{\pgfqpoint{1.648286in}{1.009567in}}%
\pgfpathlineto{\pgfqpoint{1.656395in}{0.995459in}}%
\pgfpathlineto{\pgfqpoint{1.661801in}{0.988446in}}%
\pgfpathlineto{\pgfqpoint{1.667207in}{0.983722in}}%
\pgfpathlineto{\pgfqpoint{1.671261in}{0.981877in}}%
\pgfpathlineto{\pgfqpoint{1.675316in}{0.981642in}}%
\pgfpathlineto{\pgfqpoint{1.679370in}{0.983156in}}%
\pgfpathlineto{\pgfqpoint{1.683425in}{0.986554in}}%
\pgfpathlineto{\pgfqpoint{1.687479in}{0.991971in}}%
\pgfpathlineto{\pgfqpoint{1.691534in}{0.999535in}}%
\pgfpathlineto{\pgfqpoint{1.696940in}{1.013174in}}%
\pgfpathlineto{\pgfqpoint{1.702346in}{1.031122in}}%
\pgfpathlineto{\pgfqpoint{1.709103in}{1.059984in}}%
\pgfpathlineto{\pgfqpoint{1.715861in}{1.096344in}}%
\pgfpathlineto{\pgfqpoint{1.722618in}{1.140471in}}%
\pgfpathlineto{\pgfqpoint{1.730727in}{1.203828in}}%
\pgfpathlineto{\pgfqpoint{1.740188in}{1.291853in}}%
\pgfpathlineto{\pgfqpoint{1.749648in}{1.394012in}}%
\pgfpathlineto{\pgfqpoint{1.761812in}{1.542427in}}%
\pgfpathlineto{\pgfqpoint{1.795599in}{1.970275in}}%
\pgfpathlineto{\pgfqpoint{1.802357in}{2.034396in}}%
\pgfpathlineto{\pgfqpoint{1.807763in}{2.073009in}}%
\pgfpathlineto{\pgfqpoint{1.811818in}{2.092480in}}%
\pgfpathlineto{\pgfqpoint{1.814521in}{2.100109in}}%
\pgfpathlineto{\pgfqpoint{1.817224in}{2.102896in}}%
\pgfpathlineto{\pgfqpoint{1.818575in}{2.102313in}}%
\pgfpathlineto{\pgfqpoint{1.819927in}{2.100324in}}%
\pgfpathlineto{\pgfqpoint{1.822630in}{2.091846in}}%
\pgfpathlineto{\pgfqpoint{1.825333in}{2.076885in}}%
\pgfpathlineto{\pgfqpoint{1.828036in}{2.054827in}}%
\pgfpathlineto{\pgfqpoint{1.832090in}{2.007012in}}%
\pgfpathlineto{\pgfqpoint{1.836145in}{1.939430in}}%
\pgfpathlineto{\pgfqpoint{1.840199in}{1.849550in}}%
\pgfpathlineto{\pgfqpoint{1.845605in}{1.690266in}}%
\pgfpathlineto{\pgfqpoint{1.851011in}{1.479360in}}%
\pgfpathlineto{\pgfqpoint{1.856417in}{1.209002in}}%
\pgfpathlineto{\pgfqpoint{1.857769in}{1.131130in}}%
\pgfpathlineto{\pgfqpoint{1.857769in}{1.131130in}}%
\pgfusepath{stroke}%
\end{pgfscope}%
\begin{pgfscope}%
\pgfpathrectangle{\pgfqpoint{0.507620in}{0.352669in}}{\pgfqpoint{1.350149in}{1.771000in}}%
\pgfusepath{clip}%
\pgfsetrectcap%
\pgfsetroundjoin%
\pgfsetlinewidth{0.903375pt}%
\definecolor{currentstroke}{rgb}{0.619608,0.188235,0.172549}%
\pgfsetstrokecolor{currentstroke}%
\pgfsetdash{}{0pt}%
\pgfpathmoveto{\pgfqpoint{0.507620in}{1.131130in}}%
\pgfpathlineto{\pgfqpoint{0.513026in}{1.611858in}}%
\pgfpathlineto{\pgfqpoint{0.517080in}{1.859589in}}%
\pgfpathlineto{\pgfqpoint{0.521135in}{2.028887in}}%
\pgfpathlineto{\pgfqpoint{0.525021in}{2.130335in}}%
\pgfpathmoveto{\pgfqpoint{0.541541in}{2.130335in}}%
\pgfpathlineto{\pgfqpoint{0.546813in}{2.044592in}}%
\pgfpathlineto{\pgfqpoint{0.554922in}{1.879641in}}%
\pgfpathlineto{\pgfqpoint{0.575195in}{1.452385in}}%
\pgfpathlineto{\pgfqpoint{0.583304in}{1.313992in}}%
\pgfpathlineto{\pgfqpoint{0.590061in}{1.219667in}}%
\pgfpathlineto{\pgfqpoint{0.596819in}{1.144933in}}%
\pgfpathlineto{\pgfqpoint{0.602225in}{1.098736in}}%
\pgfpathlineto{\pgfqpoint{0.607631in}{1.063769in}}%
\pgfpathlineto{\pgfqpoint{0.613037in}{1.039004in}}%
\pgfpathlineto{\pgfqpoint{0.617091in}{1.026430in}}%
\pgfpathlineto{\pgfqpoint{0.621146in}{1.018423in}}%
\pgfpathlineto{\pgfqpoint{0.623849in}{1.015362in}}%
\pgfpathlineto{\pgfqpoint{0.626552in}{1.013942in}}%
\pgfpathlineto{\pgfqpoint{0.629255in}{1.014008in}}%
\pgfpathlineto{\pgfqpoint{0.631958in}{1.015411in}}%
\pgfpathlineto{\pgfqpoint{0.636012in}{1.019701in}}%
\pgfpathlineto{\pgfqpoint{0.640067in}{1.026196in}}%
\pgfpathlineto{\pgfqpoint{0.645473in}{1.037519in}}%
\pgfpathlineto{\pgfqpoint{0.653582in}{1.058295in}}%
\pgfpathlineto{\pgfqpoint{0.681963in}{1.134682in}}%
\pgfpathlineto{\pgfqpoint{0.690072in}{1.151326in}}%
\pgfpathlineto{\pgfqpoint{0.696830in}{1.162445in}}%
\pgfpathlineto{\pgfqpoint{0.703587in}{1.170989in}}%
\pgfpathlineto{\pgfqpoint{0.710345in}{1.177034in}}%
\pgfpathlineto{\pgfqpoint{0.717102in}{1.180765in}}%
\pgfpathlineto{\pgfqpoint{0.723860in}{1.182445in}}%
\pgfpathlineto{\pgfqpoint{0.730617in}{1.182393in}}%
\pgfpathlineto{\pgfqpoint{0.738726in}{1.180534in}}%
\pgfpathlineto{\pgfqpoint{0.748187in}{1.176668in}}%
\pgfpathlineto{\pgfqpoint{0.784677in}{1.159901in}}%
\pgfpathlineto{\pgfqpoint{0.794138in}{1.157984in}}%
\pgfpathlineto{\pgfqpoint{0.803598in}{1.157694in}}%
\pgfpathlineto{\pgfqpoint{0.813059in}{1.159072in}}%
\pgfpathlineto{\pgfqpoint{0.822519in}{1.162025in}}%
\pgfpathlineto{\pgfqpoint{0.833331in}{1.167075in}}%
\pgfpathlineto{\pgfqpoint{0.846846in}{1.175303in}}%
\pgfpathlineto{\pgfqpoint{0.867119in}{1.189787in}}%
\pgfpathlineto{\pgfqpoint{0.895500in}{1.209895in}}%
\pgfpathlineto{\pgfqpoint{0.915773in}{1.222249in}}%
\pgfpathlineto{\pgfqpoint{0.967130in}{1.251968in}}%
\pgfpathlineto{\pgfqpoint{0.979293in}{1.261624in}}%
\pgfpathlineto{\pgfqpoint{0.990105in}{1.272013in}}%
\pgfpathlineto{\pgfqpoint{1.000917in}{1.284434in}}%
\pgfpathlineto{\pgfqpoint{1.011729in}{1.299114in}}%
\pgfpathlineto{\pgfqpoint{1.022541in}{1.316165in}}%
\pgfpathlineto{\pgfqpoint{1.034705in}{1.338164in}}%
\pgfpathlineto{\pgfqpoint{1.048220in}{1.365834in}}%
\pgfpathlineto{\pgfqpoint{1.064438in}{1.402667in}}%
\pgfpathlineto{\pgfqpoint{1.094171in}{1.475018in}}%
\pgfpathlineto{\pgfqpoint{1.113092in}{1.519175in}}%
\pgfpathlineto{\pgfqpoint{1.126607in}{1.547258in}}%
\pgfpathlineto{\pgfqpoint{1.137419in}{1.566626in}}%
\pgfpathlineto{\pgfqpoint{1.146879in}{1.580793in}}%
\pgfpathlineto{\pgfqpoint{1.154988in}{1.590596in}}%
\pgfpathlineto{\pgfqpoint{1.163097in}{1.598051in}}%
\pgfpathlineto{\pgfqpoint{1.169855in}{1.602371in}}%
\pgfpathlineto{\pgfqpoint{1.176612in}{1.604912in}}%
\pgfpathlineto{\pgfqpoint{1.183370in}{1.605641in}}%
\pgfpathlineto{\pgfqpoint{1.190127in}{1.604548in}}%
\pgfpathlineto{\pgfqpoint{1.196885in}{1.601648in}}%
\pgfpathlineto{\pgfqpoint{1.203642in}{1.596978in}}%
\pgfpathlineto{\pgfqpoint{1.210400in}{1.590596in}}%
\pgfpathlineto{\pgfqpoint{1.218509in}{1.580793in}}%
\pgfpathlineto{\pgfqpoint{1.227969in}{1.566626in}}%
\pgfpathlineto{\pgfqpoint{1.238781in}{1.547258in}}%
\pgfpathlineto{\pgfqpoint{1.250945in}{1.522144in}}%
\pgfpathlineto{\pgfqpoint{1.267163in}{1.484794in}}%
\pgfpathlineto{\pgfqpoint{1.319871in}{1.360052in}}%
\pgfpathlineto{\pgfqpoint{1.333386in}{1.333025in}}%
\pgfpathlineto{\pgfqpoint{1.345550in}{1.311678in}}%
\pgfpathlineto{\pgfqpoint{1.357713in}{1.293335in}}%
\pgfpathlineto{\pgfqpoint{1.368525in}{1.279522in}}%
\pgfpathlineto{\pgfqpoint{1.379337in}{1.267894in}}%
\pgfpathlineto{\pgfqpoint{1.391501in}{1.257095in}}%
\pgfpathlineto{\pgfqpoint{1.403664in}{1.248232in}}%
\pgfpathlineto{\pgfqpoint{1.419882in}{1.238431in}}%
\pgfpathlineto{\pgfqpoint{1.469888in}{1.209895in}}%
\pgfpathlineto{\pgfqpoint{1.491512in}{1.194761in}}%
\pgfpathlineto{\pgfqpoint{1.523948in}{1.171806in}}%
\pgfpathlineto{\pgfqpoint{1.536111in}{1.164994in}}%
\pgfpathlineto{\pgfqpoint{1.546923in}{1.160577in}}%
\pgfpathlineto{\pgfqpoint{1.556384in}{1.158280in}}%
\pgfpathlineto{\pgfqpoint{1.565844in}{1.157613in}}%
\pgfpathlineto{\pgfqpoint{1.575305in}{1.158614in}}%
\pgfpathlineto{\pgfqpoint{1.584765in}{1.161178in}}%
\pgfpathlineto{\pgfqpoint{1.595577in}{1.165669in}}%
\pgfpathlineto{\pgfqpoint{1.629365in}{1.181337in}}%
\pgfpathlineto{\pgfqpoint{1.637474in}{1.182600in}}%
\pgfpathlineto{\pgfqpoint{1.644231in}{1.181999in}}%
\pgfpathlineto{\pgfqpoint{1.649637in}{1.180192in}}%
\pgfpathlineto{\pgfqpoint{1.655043in}{1.177034in}}%
\pgfpathlineto{\pgfqpoint{1.660449in}{1.172393in}}%
\pgfpathlineto{\pgfqpoint{1.667207in}{1.164359in}}%
\pgfpathlineto{\pgfqpoint{1.673964in}{1.153755in}}%
\pgfpathlineto{\pgfqpoint{1.682073in}{1.137693in}}%
\pgfpathlineto{\pgfqpoint{1.691534in}{1.114861in}}%
\pgfpathlineto{\pgfqpoint{1.705049in}{1.077328in}}%
\pgfpathlineto{\pgfqpoint{1.718564in}{1.040724in}}%
\pgfpathlineto{\pgfqpoint{1.725321in}{1.026196in}}%
\pgfpathlineto{\pgfqpoint{1.730727in}{1.018003in}}%
\pgfpathlineto{\pgfqpoint{1.734782in}{1.014552in}}%
\pgfpathlineto{\pgfqpoint{1.737485in}{1.013799in}}%
\pgfpathlineto{\pgfqpoint{1.740188in}{1.014457in}}%
\pgfpathlineto{\pgfqpoint{1.742891in}{1.016678in}}%
\pgfpathlineto{\pgfqpoint{1.745594in}{1.020617in}}%
\pgfpathlineto{\pgfqpoint{1.749648in}{1.030088in}}%
\pgfpathlineto{\pgfqpoint{1.753703in}{1.044301in}}%
\pgfpathlineto{\pgfqpoint{1.757757in}{1.063769in}}%
\pgfpathlineto{\pgfqpoint{1.763163in}{1.098736in}}%
\pgfpathlineto{\pgfqpoint{1.768569in}{1.144933in}}%
\pgfpathlineto{\pgfqpoint{1.773975in}{1.203164in}}%
\pgfpathlineto{\pgfqpoint{1.780733in}{1.293558in}}%
\pgfpathlineto{\pgfqpoint{1.787490in}{1.403351in}}%
\pgfpathlineto{\pgfqpoint{1.795599in}{1.558169in}}%
\pgfpathlineto{\pgfqpoint{1.810466in}{1.879641in}}%
\pgfpathlineto{\pgfqpoint{1.819927in}{2.068784in}}%
\pgfpathlineto{\pgfqpoint{1.823847in}{2.130335in}}%
\pgfpathmoveto{\pgfqpoint{1.840367in}{2.130335in}}%
\pgfpathlineto{\pgfqpoint{1.842902in}{2.070392in}}%
\pgfpathlineto{\pgfqpoint{1.846957in}{1.923997in}}%
\pgfpathlineto{\pgfqpoint{1.851011in}{1.704107in}}%
\pgfpathlineto{\pgfqpoint{1.855066in}{1.395200in}}%
\pgfpathlineto{\pgfqpoint{1.857769in}{1.131130in}}%
\pgfpathlineto{\pgfqpoint{1.857769in}{1.131130in}}%
\pgfusepath{stroke}%
\end{pgfscope}%
\begin{pgfscope}%
\definecolor{textcolor}{rgb}{0.000000,0.000000,0.000000}%
\pgfsetstrokecolor{textcolor}%
\pgfsetfillcolor{textcolor}%
\pgftext[x=1.182694in,y=2.207002in,,base]{\color{textcolor}{\sffamily\fontsize{8.000000}{9.600000}\selectfont\catcode`\^=\active\def^{\ifmmode\sp\else\^{}\fi}\catcode`\%=\active\def%{\%}(a) Equidistant nodes}}%
\end{pgfscope}%
\begin{pgfscope}%
\pgfpathrectangle{\pgfqpoint{0.507620in}{0.352669in}}{\pgfqpoint{1.350149in}{1.771000in}}%
\pgfusepath{clip}%
\pgfsetbuttcap%
\pgfsetroundjoin%
\definecolor{currentfill}{rgb}{0.250980,0.400000,0.600000}%
\pgfsetfillcolor{currentfill}%
\pgfsetlinewidth{1.003750pt}%
\definecolor{currentstroke}{rgb}{0.250980,0.400000,0.600000}%
\pgfsetstrokecolor{currentstroke}%
\pgfsetdash{}{0pt}%
\pgfsys@defobject{currentmarker}{\pgfqpoint{-0.013889in}{-0.013889in}}{\pgfqpoint{0.013889in}{0.013889in}}{%
\pgfpathmoveto{\pgfqpoint{0.000000in}{-0.013889in}}%
\pgfpathcurveto{\pgfqpoint{0.003683in}{-0.013889in}}{\pgfqpoint{0.007216in}{-0.012425in}}{\pgfqpoint{0.009821in}{-0.009821in}}%
\pgfpathcurveto{\pgfqpoint{0.012425in}{-0.007216in}}{\pgfqpoint{0.013889in}{-0.003683in}}{\pgfqpoint{0.013889in}{0.000000in}}%
\pgfpathcurveto{\pgfqpoint{0.013889in}{0.003683in}}{\pgfqpoint{0.012425in}{0.007216in}}{\pgfqpoint{0.009821in}{0.009821in}}%
\pgfpathcurveto{\pgfqpoint{0.007216in}{0.012425in}}{\pgfqpoint{0.003683in}{0.013889in}}{\pgfqpoint{0.000000in}{0.013889in}}%
\pgfpathcurveto{\pgfqpoint{-0.003683in}{0.013889in}}{\pgfqpoint{-0.007216in}{0.012425in}}{\pgfqpoint{-0.009821in}{0.009821in}}%
\pgfpathcurveto{\pgfqpoint{-0.012425in}{0.007216in}}{\pgfqpoint{-0.013889in}{0.003683in}}{\pgfqpoint{-0.013889in}{0.000000in}}%
\pgfpathcurveto{\pgfqpoint{-0.013889in}{-0.003683in}}{\pgfqpoint{-0.012425in}{-0.007216in}}{\pgfqpoint{-0.009821in}{-0.009821in}}%
\pgfpathcurveto{\pgfqpoint{-0.007216in}{-0.012425in}}{\pgfqpoint{-0.003683in}{-0.013889in}}{\pgfqpoint{0.000000in}{-0.013889in}}%
\pgfpathlineto{\pgfqpoint{0.000000in}{-0.013889in}}%
\pgfpathclose%
\pgfusepath{stroke,fill}%
}%
\begin{pgfscope}%
\pgfsys@transformshift{0.507620in}{1.131130in}%
\pgfsys@useobject{currentmarker}{}%
\end{pgfscope}%
\begin{pgfscope}%
\pgfsys@transformshift{0.777649in}{1.162269in}%
\pgfsys@useobject{currentmarker}{}%
\end{pgfscope}%
\begin{pgfscope}%
\pgfsys@transformshift{1.047679in}{1.364669in}%
\pgfsys@useobject{currentmarker}{}%
\end{pgfscope}%
\begin{pgfscope}%
\pgfsys@transformshift{1.317709in}{1.364669in}%
\pgfsys@useobject{currentmarker}{}%
\end{pgfscope}%
\begin{pgfscope}%
\pgfsys@transformshift{1.587739in}{1.162269in}%
\pgfsys@useobject{currentmarker}{}%
\end{pgfscope}%
\begin{pgfscope}%
\pgfsys@transformshift{1.857769in}{1.131130in}%
\pgfsys@useobject{currentmarker}{}%
\end{pgfscope}%
\end{pgfscope}%
\begin{pgfscope}%
\pgfpathrectangle{\pgfqpoint{0.507620in}{0.352669in}}{\pgfqpoint{1.350149in}{1.771000in}}%
\pgfusepath{clip}%
\pgfsetbuttcap%
\pgfsetroundjoin%
\definecolor{currentfill}{rgb}{0.768627,0.556863,0.211765}%
\pgfsetfillcolor{currentfill}%
\pgfsetlinewidth{1.003750pt}%
\definecolor{currentstroke}{rgb}{0.768627,0.556863,0.211765}%
\pgfsetstrokecolor{currentstroke}%
\pgfsetdash{}{0pt}%
\pgfsys@defobject{currentmarker}{\pgfqpoint{-0.013889in}{-0.013889in}}{\pgfqpoint{0.013889in}{0.013889in}}{%
\pgfpathmoveto{\pgfqpoint{0.000000in}{-0.013889in}}%
\pgfpathcurveto{\pgfqpoint{0.003683in}{-0.013889in}}{\pgfqpoint{0.007216in}{-0.012425in}}{\pgfqpoint{0.009821in}{-0.009821in}}%
\pgfpathcurveto{\pgfqpoint{0.012425in}{-0.007216in}}{\pgfqpoint{0.013889in}{-0.003683in}}{\pgfqpoint{0.013889in}{0.000000in}}%
\pgfpathcurveto{\pgfqpoint{0.013889in}{0.003683in}}{\pgfqpoint{0.012425in}{0.007216in}}{\pgfqpoint{0.009821in}{0.009821in}}%
\pgfpathcurveto{\pgfqpoint{0.007216in}{0.012425in}}{\pgfqpoint{0.003683in}{0.013889in}}{\pgfqpoint{0.000000in}{0.013889in}}%
\pgfpathcurveto{\pgfqpoint{-0.003683in}{0.013889in}}{\pgfqpoint{-0.007216in}{0.012425in}}{\pgfqpoint{-0.009821in}{0.009821in}}%
\pgfpathcurveto{\pgfqpoint{-0.012425in}{0.007216in}}{\pgfqpoint{-0.013889in}{0.003683in}}{\pgfqpoint{-0.013889in}{0.000000in}}%
\pgfpathcurveto{\pgfqpoint{-0.013889in}{-0.003683in}}{\pgfqpoint{-0.012425in}{-0.007216in}}{\pgfqpoint{-0.009821in}{-0.009821in}}%
\pgfpathcurveto{\pgfqpoint{-0.007216in}{-0.012425in}}{\pgfqpoint{-0.003683in}{-0.013889in}}{\pgfqpoint{0.000000in}{-0.013889in}}%
\pgfpathlineto{\pgfqpoint{0.000000in}{-0.013889in}}%
\pgfpathclose%
\pgfusepath{stroke,fill}%
}%
\begin{pgfscope}%
\pgfsys@transformshift{0.507620in}{1.131130in}%
\pgfsys@useobject{currentmarker}{}%
\end{pgfscope}%
\begin{pgfscope}%
\pgfsys@transformshift{0.642634in}{1.141433in}%
\pgfsys@useobject{currentmarker}{}%
\end{pgfscope}%
\begin{pgfscope}%
\pgfsys@transformshift{0.777649in}{1.162269in}%
\pgfsys@useobject{currentmarker}{}%
\end{pgfscope}%
\begin{pgfscope}%
\pgfsys@transformshift{0.912664in}{1.212869in}%
\pgfsys@useobject{currentmarker}{}%
\end{pgfscope}%
\begin{pgfscope}%
\pgfsys@transformshift{1.047679in}{1.364669in}%
\pgfsys@useobject{currentmarker}{}%
\end{pgfscope}%
\begin{pgfscope}%
\pgfsys@transformshift{1.182694in}{1.617669in}%
\pgfsys@useobject{currentmarker}{}%
\end{pgfscope}%
\begin{pgfscope}%
\pgfsys@transformshift{1.317709in}{1.364669in}%
\pgfsys@useobject{currentmarker}{}%
\end{pgfscope}%
\begin{pgfscope}%
\pgfsys@transformshift{1.452724in}{1.212869in}%
\pgfsys@useobject{currentmarker}{}%
\end{pgfscope}%
\begin{pgfscope}%
\pgfsys@transformshift{1.587739in}{1.162269in}%
\pgfsys@useobject{currentmarker}{}%
\end{pgfscope}%
\begin{pgfscope}%
\pgfsys@transformshift{1.722754in}{1.141433in}%
\pgfsys@useobject{currentmarker}{}%
\end{pgfscope}%
\begin{pgfscope}%
\pgfsys@transformshift{1.857769in}{1.131130in}%
\pgfsys@useobject{currentmarker}{}%
\end{pgfscope}%
\end{pgfscope}%
\begin{pgfscope}%
\pgfpathrectangle{\pgfqpoint{0.507620in}{0.352669in}}{\pgfqpoint{1.350149in}{1.771000in}}%
\pgfusepath{clip}%
\pgfsetbuttcap%
\pgfsetroundjoin%
\definecolor{currentfill}{rgb}{0.619608,0.188235,0.172549}%
\pgfsetfillcolor{currentfill}%
\pgfsetlinewidth{1.003750pt}%
\definecolor{currentstroke}{rgb}{0.619608,0.188235,0.172549}%
\pgfsetstrokecolor{currentstroke}%
\pgfsetdash{}{0pt}%
\pgfsys@defobject{currentmarker}{\pgfqpoint{-0.013889in}{-0.013889in}}{\pgfqpoint{0.013889in}{0.013889in}}{%
\pgfpathmoveto{\pgfqpoint{0.000000in}{-0.013889in}}%
\pgfpathcurveto{\pgfqpoint{0.003683in}{-0.013889in}}{\pgfqpoint{0.007216in}{-0.012425in}}{\pgfqpoint{0.009821in}{-0.009821in}}%
\pgfpathcurveto{\pgfqpoint{0.012425in}{-0.007216in}}{\pgfqpoint{0.013889in}{-0.003683in}}{\pgfqpoint{0.013889in}{0.000000in}}%
\pgfpathcurveto{\pgfqpoint{0.013889in}{0.003683in}}{\pgfqpoint{0.012425in}{0.007216in}}{\pgfqpoint{0.009821in}{0.009821in}}%
\pgfpathcurveto{\pgfqpoint{0.007216in}{0.012425in}}{\pgfqpoint{0.003683in}{0.013889in}}{\pgfqpoint{0.000000in}{0.013889in}}%
\pgfpathcurveto{\pgfqpoint{-0.003683in}{0.013889in}}{\pgfqpoint{-0.007216in}{0.012425in}}{\pgfqpoint{-0.009821in}{0.009821in}}%
\pgfpathcurveto{\pgfqpoint{-0.012425in}{0.007216in}}{\pgfqpoint{-0.013889in}{0.003683in}}{\pgfqpoint{-0.013889in}{0.000000in}}%
\pgfpathcurveto{\pgfqpoint{-0.013889in}{-0.003683in}}{\pgfqpoint{-0.012425in}{-0.007216in}}{\pgfqpoint{-0.009821in}{-0.009821in}}%
\pgfpathcurveto{\pgfqpoint{-0.007216in}{-0.012425in}}{\pgfqpoint{-0.003683in}{-0.013889in}}{\pgfqpoint{0.000000in}{-0.013889in}}%
\pgfpathlineto{\pgfqpoint{0.000000in}{-0.013889in}}%
\pgfpathclose%
\pgfusepath{stroke,fill}%
}%
\begin{pgfscope}%
\pgfsys@transformshift{0.507620in}{1.131130in}%
\pgfsys@useobject{currentmarker}{}%
\end{pgfscope}%
\begin{pgfscope}%
\pgfsys@transformshift{0.597630in}{1.137253in}%
\pgfsys@useobject{currentmarker}{}%
\end{pgfscope}%
\begin{pgfscope}%
\pgfsys@transformshift{0.687639in}{1.146700in}%
\pgfsys@useobject{currentmarker}{}%
\end{pgfscope}%
\begin{pgfscope}%
\pgfsys@transformshift{0.777649in}{1.162269in}%
\pgfsys@useobject{currentmarker}{}%
\end{pgfscope}%
\begin{pgfscope}%
\pgfsys@transformshift{0.867659in}{1.190186in}%
\pgfsys@useobject{currentmarker}{}%
\end{pgfscope}%
\begin{pgfscope}%
\pgfsys@transformshift{0.957669in}{1.245610in}%
\pgfsys@useobject{currentmarker}{}%
\end{pgfscope}%
\begin{pgfscope}%
\pgfsys@transformshift{1.047679in}{1.364669in}%
\pgfsys@useobject{currentmarker}{}%
\end{pgfscope}%
\begin{pgfscope}%
\pgfsys@transformshift{1.137689in}{1.567069in}%
\pgfsys@useobject{currentmarker}{}%
\end{pgfscope}%
\begin{pgfscope}%
\pgfsys@transformshift{1.227699in}{1.567069in}%
\pgfsys@useobject{currentmarker}{}%
\end{pgfscope}%
\begin{pgfscope}%
\pgfsys@transformshift{1.317709in}{1.364669in}%
\pgfsys@useobject{currentmarker}{}%
\end{pgfscope}%
\begin{pgfscope}%
\pgfsys@transformshift{1.407719in}{1.245610in}%
\pgfsys@useobject{currentmarker}{}%
\end{pgfscope}%
\begin{pgfscope}%
\pgfsys@transformshift{1.497729in}{1.190186in}%
\pgfsys@useobject{currentmarker}{}%
\end{pgfscope}%
\begin{pgfscope}%
\pgfsys@transformshift{1.587739in}{1.162269in}%
\pgfsys@useobject{currentmarker}{}%
\end{pgfscope}%
\begin{pgfscope}%
\pgfsys@transformshift{1.677749in}{1.146700in}%
\pgfsys@useobject{currentmarker}{}%
\end{pgfscope}%
\begin{pgfscope}%
\pgfsys@transformshift{1.767759in}{1.137253in}%
\pgfsys@useobject{currentmarker}{}%
\end{pgfscope}%
\begin{pgfscope}%
\pgfsys@transformshift{1.857769in}{1.131130in}%
\pgfsys@useobject{currentmarker}{}%
\end{pgfscope}%
\end{pgfscope}%
\begin{pgfscope}%
\pgfpathrectangle{\pgfqpoint{0.507620in}{0.352669in}}{\pgfqpoint{1.350149in}{1.771000in}}%
\pgfusepath{clip}%
\pgfsetrectcap%
\pgfsetroundjoin%
\pgfsetlinewidth{1.405250pt}%
\definecolor{currentstroke}{rgb}{0.000000,0.000000,0.000000}%
\pgfsetstrokecolor{currentstroke}%
\pgfsetdash{}{0pt}%
\pgfpathmoveto{\pgfqpoint{0.507620in}{1.131130in}}%
\pgfpathlineto{\pgfqpoint{0.581952in}{1.135998in}}%
\pgfpathlineto{\pgfqpoint{0.641418in}{1.141308in}}%
\pgfpathlineto{\pgfqpoint{0.691424in}{1.147203in}}%
\pgfpathlineto{\pgfqpoint{0.733320in}{1.153564in}}%
\pgfpathlineto{\pgfqpoint{0.769811in}{1.160550in}}%
\pgfpathlineto{\pgfqpoint{0.802247in}{1.168268in}}%
\pgfpathlineto{\pgfqpoint{0.830628in}{1.176544in}}%
\pgfpathlineto{\pgfqpoint{0.856307in}{1.185603in}}%
\pgfpathlineto{\pgfqpoint{0.879282in}{1.195303in}}%
\pgfpathlineto{\pgfqpoint{0.899555in}{1.205410in}}%
\pgfpathlineto{\pgfqpoint{0.918476in}{1.216438in}}%
\pgfpathlineto{\pgfqpoint{0.936045in}{1.228331in}}%
\pgfpathlineto{\pgfqpoint{0.952263in}{1.240986in}}%
\pgfpathlineto{\pgfqpoint{0.968481in}{1.255531in}}%
\pgfpathlineto{\pgfqpoint{0.983348in}{1.270789in}}%
\pgfpathlineto{\pgfqpoint{0.998214in}{1.288162in}}%
\pgfpathlineto{\pgfqpoint{1.013081in}{1.307931in}}%
\pgfpathlineto{\pgfqpoint{1.027947in}{1.330371in}}%
\pgfpathlineto{\pgfqpoint{1.042814in}{1.355716in}}%
\pgfpathlineto{\pgfqpoint{1.057680in}{1.384102in}}%
\pgfpathlineto{\pgfqpoint{1.073898in}{1.418461in}}%
\pgfpathlineto{\pgfqpoint{1.092819in}{1.462300in}}%
\pgfpathlineto{\pgfqpoint{1.134716in}{1.560936in}}%
\pgfpathlineto{\pgfqpoint{1.145528in}{1.582027in}}%
\pgfpathlineto{\pgfqpoint{1.154988in}{1.597222in}}%
\pgfpathlineto{\pgfqpoint{1.163097in}{1.607229in}}%
\pgfpathlineto{\pgfqpoint{1.169855in}{1.613134in}}%
\pgfpathlineto{\pgfqpoint{1.175261in}{1.616140in}}%
\pgfpathlineto{\pgfqpoint{1.180667in}{1.617555in}}%
\pgfpathlineto{\pgfqpoint{1.186073in}{1.617352in}}%
\pgfpathlineto{\pgfqpoint{1.191479in}{1.615536in}}%
\pgfpathlineto{\pgfqpoint{1.196885in}{1.612140in}}%
\pgfpathlineto{\pgfqpoint{1.203642in}{1.605774in}}%
\pgfpathlineto{\pgfqpoint{1.210400in}{1.597222in}}%
\pgfpathlineto{\pgfqpoint{1.218509in}{1.584404in}}%
\pgfpathlineto{\pgfqpoint{1.227969in}{1.566521in}}%
\pgfpathlineto{\pgfqpoint{1.240133in}{1.540124in}}%
\pgfpathlineto{\pgfqpoint{1.260405in}{1.491752in}}%
\pgfpathlineto{\pgfqpoint{1.288787in}{1.424505in}}%
\pgfpathlineto{\pgfqpoint{1.306356in}{1.386833in}}%
\pgfpathlineto{\pgfqpoint{1.322574in}{1.355716in}}%
\pgfpathlineto{\pgfqpoint{1.337441in}{1.330371in}}%
\pgfpathlineto{\pgfqpoint{1.352307in}{1.307931in}}%
\pgfpathlineto{\pgfqpoint{1.367174in}{1.288162in}}%
\pgfpathlineto{\pgfqpoint{1.382040in}{1.270789in}}%
\pgfpathlineto{\pgfqpoint{1.396907in}{1.255531in}}%
\pgfpathlineto{\pgfqpoint{1.413125in}{1.240986in}}%
\pgfpathlineto{\pgfqpoint{1.430694in}{1.227353in}}%
\pgfpathlineto{\pgfqpoint{1.448264in}{1.215592in}}%
\pgfpathlineto{\pgfqpoint{1.467185in}{1.204685in}}%
\pgfpathlineto{\pgfqpoint{1.487457in}{1.194685in}}%
\pgfpathlineto{\pgfqpoint{1.510433in}{1.185083in}}%
\pgfpathlineto{\pgfqpoint{1.534760in}{1.176544in}}%
\pgfpathlineto{\pgfqpoint{1.561790in}{1.168626in}}%
\pgfpathlineto{\pgfqpoint{1.592874in}{1.161133in}}%
\pgfpathlineto{\pgfqpoint{1.628013in}{1.154266in}}%
\pgfpathlineto{\pgfqpoint{1.668558in}{1.147941in}}%
\pgfpathlineto{\pgfqpoint{1.714509in}{1.142307in}}%
\pgfpathlineto{\pgfqpoint{1.769921in}{1.137074in}}%
\pgfpathlineto{\pgfqpoint{1.836145in}{1.132386in}}%
\pgfpathlineto{\pgfqpoint{1.857769in}{1.131130in}}%
\pgfpathlineto{\pgfqpoint{1.857769in}{1.131130in}}%
\pgfusepath{stroke}%
\end{pgfscope}%
\begin{pgfscope}%
\pgfsetrectcap%
\pgfsetroundjoin%
\pgfsetlinewidth{1.405250pt}%
\definecolor{currentstroke}{rgb}{0.000000,0.000000,0.000000}%
\pgfsetstrokecolor{currentstroke}%
\pgfsetdash{}{0pt}%
\pgfpathmoveto{\pgfqpoint{0.574286in}{2.016419in}}%
\pgfpathlineto{\pgfqpoint{0.657620in}{2.016419in}}%
\pgfpathlineto{\pgfqpoint{0.740953in}{2.016419in}}%
\pgfusepath{stroke}%
\end{pgfscope}%
\begin{pgfscope}%
\definecolor{textcolor}{rgb}{0.000000,0.000000,0.000000}%
\pgfsetstrokecolor{textcolor}%
\pgfsetfillcolor{textcolor}%
\pgftext[x=0.807620in,y=1.987252in,left,base]{\color{textcolor}{\sffamily\fontsize{6.000000}{7.200000}\selectfont\catcode`\^=\active\def^{\ifmmode\sp\else\^{}\fi}\catcode`\%=\active\def%{\%}$f(x)$}}%
\end{pgfscope}%
\begin{pgfscope}%
\pgfsetrectcap%
\pgfsetroundjoin%
\pgfsetlinewidth{0.903375pt}%
\definecolor{currentstroke}{rgb}{0.250980,0.400000,0.600000}%
\pgfsetstrokecolor{currentstroke}%
\pgfsetdash{}{0pt}%
\pgfpathmoveto{\pgfqpoint{0.574286in}{1.892022in}}%
\pgfpathlineto{\pgfqpoint{0.657620in}{1.892022in}}%
\pgfpathlineto{\pgfqpoint{0.740953in}{1.892022in}}%
\pgfusepath{stroke}%
\end{pgfscope}%
\begin{pgfscope}%
\definecolor{textcolor}{rgb}{0.000000,0.000000,0.000000}%
\pgfsetstrokecolor{textcolor}%
\pgfsetfillcolor{textcolor}%
\pgftext[x=0.807620in,y=1.862855in,left,base]{\color{textcolor}{\sffamily\fontsize{6.000000}{7.200000}\selectfont\catcode`\^=\active\def^{\ifmmode\sp\else\^{}\fi}\catcode`\%=\active\def%{\%}$N = 5$}}%
\end{pgfscope}%
\begin{pgfscope}%
\pgfsetrectcap%
\pgfsetroundjoin%
\pgfsetlinewidth{0.903375pt}%
\definecolor{currentstroke}{rgb}{0.768627,0.556863,0.211765}%
\pgfsetstrokecolor{currentstroke}%
\pgfsetdash{}{0pt}%
\pgfpathmoveto{\pgfqpoint{1.214832in}{2.016419in}}%
\pgfpathlineto{\pgfqpoint{1.298165in}{2.016419in}}%
\pgfpathlineto{\pgfqpoint{1.381498in}{2.016419in}}%
\pgfusepath{stroke}%
\end{pgfscope}%
\begin{pgfscope}%
\definecolor{textcolor}{rgb}{0.000000,0.000000,0.000000}%
\pgfsetstrokecolor{textcolor}%
\pgfsetfillcolor{textcolor}%
\pgftext[x=1.448165in,y=1.987252in,left,base]{\color{textcolor}{\sffamily\fontsize{6.000000}{7.200000}\selectfont\catcode`\^=\active\def^{\ifmmode\sp\else\^{}\fi}\catcode`\%=\active\def%{\%}$N = 10$}}%
\end{pgfscope}%
\begin{pgfscope}%
\pgfsetrectcap%
\pgfsetroundjoin%
\pgfsetlinewidth{0.903375pt}%
\definecolor{currentstroke}{rgb}{0.619608,0.188235,0.172549}%
\pgfsetstrokecolor{currentstroke}%
\pgfsetdash{}{0pt}%
\pgfpathmoveto{\pgfqpoint{1.214832in}{1.894104in}}%
\pgfpathlineto{\pgfqpoint{1.298165in}{1.894104in}}%
\pgfpathlineto{\pgfqpoint{1.381498in}{1.894104in}}%
\pgfusepath{stroke}%
\end{pgfscope}%
\begin{pgfscope}%
\definecolor{textcolor}{rgb}{0.000000,0.000000,0.000000}%
\pgfsetstrokecolor{textcolor}%
\pgfsetfillcolor{textcolor}%
\pgftext[x=1.448165in,y=1.864938in,left,base]{\color{textcolor}{\sffamily\fontsize{6.000000}{7.200000}\selectfont\catcode`\^=\active\def^{\ifmmode\sp\else\^{}\fi}\catcode`\%=\active\def%{\%}$N = 15$}}%
\end{pgfscope}%
\begin{pgfscope}%
\pgfsetbuttcap%
\pgfsetmiterjoin%
\definecolor{currentfill}{rgb}{1.000000,1.000000,1.000000}%
\pgfsetfillcolor{currentfill}%
\pgfsetlinewidth{0.000000pt}%
\definecolor{currentstroke}{rgb}{0.000000,0.000000,0.000000}%
\pgfsetstrokecolor{currentstroke}%
\pgfsetstrokeopacity{0.000000}%
\pgfsetdash{}{0pt}%
\pgfpathmoveto{\pgfqpoint{2.330321in}{0.352669in}}%
\pgfpathlineto{\pgfqpoint{3.680470in}{0.352669in}}%
\pgfpathlineto{\pgfqpoint{3.680470in}{2.123669in}}%
\pgfpathlineto{\pgfqpoint{2.330321in}{2.123669in}}%
\pgfpathlineto{\pgfqpoint{2.330321in}{0.352669in}}%
\pgfpathclose%
\pgfusepath{fill}%
\end{pgfscope}%
\begin{pgfscope}%
\pgfsetbuttcap%
\pgfsetroundjoin%
\definecolor{currentfill}{rgb}{0.000000,0.000000,0.000000}%
\pgfsetfillcolor{currentfill}%
\pgfsetlinewidth{0.501875pt}%
\definecolor{currentstroke}{rgb}{0.000000,0.000000,0.000000}%
\pgfsetstrokecolor{currentstroke}%
\pgfsetdash{}{0pt}%
\pgfsys@defobject{currentmarker}{\pgfqpoint{0.000000in}{0.000000in}}{\pgfqpoint{0.000000in}{0.041667in}}{%
\pgfpathmoveto{\pgfqpoint{0.000000in}{0.000000in}}%
\pgfpathlineto{\pgfqpoint{0.000000in}{0.041667in}}%
\pgfusepath{stroke,fill}%
}%
\begin{pgfscope}%
\pgfsys@transformshift{2.330321in}{0.352669in}%
\pgfsys@useobject{currentmarker}{}%
\end{pgfscope}%
\end{pgfscope}%
\begin{pgfscope}%
\definecolor{textcolor}{rgb}{0.000000,0.000000,0.000000}%
\pgfsetstrokecolor{textcolor}%
\pgfsetfillcolor{textcolor}%
\pgftext[x=2.330321in,y=0.304058in,,top]{\color{textcolor}{\sffamily\fontsize{8.000000}{9.600000}\selectfont\catcode`\^=\active\def^{\ifmmode\sp\else\^{}\fi}\catcode`\%=\active\def%{\%}\ensuremath{-}1.0}}%
\end{pgfscope}%
\begin{pgfscope}%
\pgfsetbuttcap%
\pgfsetroundjoin%
\definecolor{currentfill}{rgb}{0.000000,0.000000,0.000000}%
\pgfsetfillcolor{currentfill}%
\pgfsetlinewidth{0.501875pt}%
\definecolor{currentstroke}{rgb}{0.000000,0.000000,0.000000}%
\pgfsetstrokecolor{currentstroke}%
\pgfsetdash{}{0pt}%
\pgfsys@defobject{currentmarker}{\pgfqpoint{0.000000in}{0.000000in}}{\pgfqpoint{0.000000in}{0.041667in}}{%
\pgfpathmoveto{\pgfqpoint{0.000000in}{0.000000in}}%
\pgfpathlineto{\pgfqpoint{0.000000in}{0.041667in}}%
\pgfusepath{stroke,fill}%
}%
\begin{pgfscope}%
\pgfsys@transformshift{2.667858in}{0.352669in}%
\pgfsys@useobject{currentmarker}{}%
\end{pgfscope}%
\end{pgfscope}%
\begin{pgfscope}%
\definecolor{textcolor}{rgb}{0.000000,0.000000,0.000000}%
\pgfsetstrokecolor{textcolor}%
\pgfsetfillcolor{textcolor}%
\pgftext[x=2.667858in,y=0.304058in,,top]{\color{textcolor}{\sffamily\fontsize{8.000000}{9.600000}\selectfont\catcode`\^=\active\def^{\ifmmode\sp\else\^{}\fi}\catcode`\%=\active\def%{\%}\ensuremath{-}0.5}}%
\end{pgfscope}%
\begin{pgfscope}%
\pgfsetbuttcap%
\pgfsetroundjoin%
\definecolor{currentfill}{rgb}{0.000000,0.000000,0.000000}%
\pgfsetfillcolor{currentfill}%
\pgfsetlinewidth{0.501875pt}%
\definecolor{currentstroke}{rgb}{0.000000,0.000000,0.000000}%
\pgfsetstrokecolor{currentstroke}%
\pgfsetdash{}{0pt}%
\pgfsys@defobject{currentmarker}{\pgfqpoint{0.000000in}{0.000000in}}{\pgfqpoint{0.000000in}{0.041667in}}{%
\pgfpathmoveto{\pgfqpoint{0.000000in}{0.000000in}}%
\pgfpathlineto{\pgfqpoint{0.000000in}{0.041667in}}%
\pgfusepath{stroke,fill}%
}%
\begin{pgfscope}%
\pgfsys@transformshift{3.005395in}{0.352669in}%
\pgfsys@useobject{currentmarker}{}%
\end{pgfscope}%
\end{pgfscope}%
\begin{pgfscope}%
\definecolor{textcolor}{rgb}{0.000000,0.000000,0.000000}%
\pgfsetstrokecolor{textcolor}%
\pgfsetfillcolor{textcolor}%
\pgftext[x=3.005395in,y=0.304058in,,top]{\color{textcolor}{\sffamily\fontsize{8.000000}{9.600000}\selectfont\catcode`\^=\active\def^{\ifmmode\sp\else\^{}\fi}\catcode`\%=\active\def%{\%}0.0}}%
\end{pgfscope}%
\begin{pgfscope}%
\pgfsetbuttcap%
\pgfsetroundjoin%
\definecolor{currentfill}{rgb}{0.000000,0.000000,0.000000}%
\pgfsetfillcolor{currentfill}%
\pgfsetlinewidth{0.501875pt}%
\definecolor{currentstroke}{rgb}{0.000000,0.000000,0.000000}%
\pgfsetstrokecolor{currentstroke}%
\pgfsetdash{}{0pt}%
\pgfsys@defobject{currentmarker}{\pgfqpoint{0.000000in}{0.000000in}}{\pgfqpoint{0.000000in}{0.041667in}}{%
\pgfpathmoveto{\pgfqpoint{0.000000in}{0.000000in}}%
\pgfpathlineto{\pgfqpoint{0.000000in}{0.041667in}}%
\pgfusepath{stroke,fill}%
}%
\begin{pgfscope}%
\pgfsys@transformshift{3.342932in}{0.352669in}%
\pgfsys@useobject{currentmarker}{}%
\end{pgfscope}%
\end{pgfscope}%
\begin{pgfscope}%
\definecolor{textcolor}{rgb}{0.000000,0.000000,0.000000}%
\pgfsetstrokecolor{textcolor}%
\pgfsetfillcolor{textcolor}%
\pgftext[x=3.342932in,y=0.304058in,,top]{\color{textcolor}{\sffamily\fontsize{8.000000}{9.600000}\selectfont\catcode`\^=\active\def^{\ifmmode\sp\else\^{}\fi}\catcode`\%=\active\def%{\%}0.5}}%
\end{pgfscope}%
\begin{pgfscope}%
\pgfsetbuttcap%
\pgfsetroundjoin%
\definecolor{currentfill}{rgb}{0.000000,0.000000,0.000000}%
\pgfsetfillcolor{currentfill}%
\pgfsetlinewidth{0.501875pt}%
\definecolor{currentstroke}{rgb}{0.000000,0.000000,0.000000}%
\pgfsetstrokecolor{currentstroke}%
\pgfsetdash{}{0pt}%
\pgfsys@defobject{currentmarker}{\pgfqpoint{0.000000in}{0.000000in}}{\pgfqpoint{0.000000in}{0.041667in}}{%
\pgfpathmoveto{\pgfqpoint{0.000000in}{0.000000in}}%
\pgfpathlineto{\pgfqpoint{0.000000in}{0.041667in}}%
\pgfusepath{stroke,fill}%
}%
\begin{pgfscope}%
\pgfsys@transformshift{3.680470in}{0.352669in}%
\pgfsys@useobject{currentmarker}{}%
\end{pgfscope}%
\end{pgfscope}%
\begin{pgfscope}%
\definecolor{textcolor}{rgb}{0.000000,0.000000,0.000000}%
\pgfsetstrokecolor{textcolor}%
\pgfsetfillcolor{textcolor}%
\pgftext[x=3.680470in,y=0.304058in,,top]{\color{textcolor}{\sffamily\fontsize{8.000000}{9.600000}\selectfont\catcode`\^=\active\def^{\ifmmode\sp\else\^{}\fi}\catcode`\%=\active\def%{\%}1.0}}%
\end{pgfscope}%
\begin{pgfscope}%
\pgfsetbuttcap%
\pgfsetroundjoin%
\definecolor{currentfill}{rgb}{0.000000,0.000000,0.000000}%
\pgfsetfillcolor{currentfill}%
\pgfsetlinewidth{0.301125pt}%
\definecolor{currentstroke}{rgb}{0.000000,0.000000,0.000000}%
\pgfsetstrokecolor{currentstroke}%
\pgfsetdash{}{0pt}%
\pgfsys@defobject{currentmarker}{\pgfqpoint{0.000000in}{0.000000in}}{\pgfqpoint{0.000000in}{0.027778in}}{%
\pgfpathmoveto{\pgfqpoint{0.000000in}{0.000000in}}%
\pgfpathlineto{\pgfqpoint{0.000000in}{0.027778in}}%
\pgfusepath{stroke,fill}%
}%
\begin{pgfscope}%
\pgfsys@transformshift{2.397828in}{0.352669in}%
\pgfsys@useobject{currentmarker}{}%
\end{pgfscope}%
\end{pgfscope}%
\begin{pgfscope}%
\pgfsetbuttcap%
\pgfsetroundjoin%
\definecolor{currentfill}{rgb}{0.000000,0.000000,0.000000}%
\pgfsetfillcolor{currentfill}%
\pgfsetlinewidth{0.301125pt}%
\definecolor{currentstroke}{rgb}{0.000000,0.000000,0.000000}%
\pgfsetstrokecolor{currentstroke}%
\pgfsetdash{}{0pt}%
\pgfsys@defobject{currentmarker}{\pgfqpoint{0.000000in}{0.000000in}}{\pgfqpoint{0.000000in}{0.027778in}}{%
\pgfpathmoveto{\pgfqpoint{0.000000in}{0.000000in}}%
\pgfpathlineto{\pgfqpoint{0.000000in}{0.027778in}}%
\pgfusepath{stroke,fill}%
}%
\begin{pgfscope}%
\pgfsys@transformshift{2.465336in}{0.352669in}%
\pgfsys@useobject{currentmarker}{}%
\end{pgfscope}%
\end{pgfscope}%
\begin{pgfscope}%
\pgfsetbuttcap%
\pgfsetroundjoin%
\definecolor{currentfill}{rgb}{0.000000,0.000000,0.000000}%
\pgfsetfillcolor{currentfill}%
\pgfsetlinewidth{0.301125pt}%
\definecolor{currentstroke}{rgb}{0.000000,0.000000,0.000000}%
\pgfsetstrokecolor{currentstroke}%
\pgfsetdash{}{0pt}%
\pgfsys@defobject{currentmarker}{\pgfqpoint{0.000000in}{0.000000in}}{\pgfqpoint{0.000000in}{0.027778in}}{%
\pgfpathmoveto{\pgfqpoint{0.000000in}{0.000000in}}%
\pgfpathlineto{\pgfqpoint{0.000000in}{0.027778in}}%
\pgfusepath{stroke,fill}%
}%
\begin{pgfscope}%
\pgfsys@transformshift{2.532843in}{0.352669in}%
\pgfsys@useobject{currentmarker}{}%
\end{pgfscope}%
\end{pgfscope}%
\begin{pgfscope}%
\pgfsetbuttcap%
\pgfsetroundjoin%
\definecolor{currentfill}{rgb}{0.000000,0.000000,0.000000}%
\pgfsetfillcolor{currentfill}%
\pgfsetlinewidth{0.301125pt}%
\definecolor{currentstroke}{rgb}{0.000000,0.000000,0.000000}%
\pgfsetstrokecolor{currentstroke}%
\pgfsetdash{}{0pt}%
\pgfsys@defobject{currentmarker}{\pgfqpoint{0.000000in}{0.000000in}}{\pgfqpoint{0.000000in}{0.027778in}}{%
\pgfpathmoveto{\pgfqpoint{0.000000in}{0.000000in}}%
\pgfpathlineto{\pgfqpoint{0.000000in}{0.027778in}}%
\pgfusepath{stroke,fill}%
}%
\begin{pgfscope}%
\pgfsys@transformshift{2.600350in}{0.352669in}%
\pgfsys@useobject{currentmarker}{}%
\end{pgfscope}%
\end{pgfscope}%
\begin{pgfscope}%
\pgfsetbuttcap%
\pgfsetroundjoin%
\definecolor{currentfill}{rgb}{0.000000,0.000000,0.000000}%
\pgfsetfillcolor{currentfill}%
\pgfsetlinewidth{0.301125pt}%
\definecolor{currentstroke}{rgb}{0.000000,0.000000,0.000000}%
\pgfsetstrokecolor{currentstroke}%
\pgfsetdash{}{0pt}%
\pgfsys@defobject{currentmarker}{\pgfqpoint{0.000000in}{0.000000in}}{\pgfqpoint{0.000000in}{0.027778in}}{%
\pgfpathmoveto{\pgfqpoint{0.000000in}{0.000000in}}%
\pgfpathlineto{\pgfqpoint{0.000000in}{0.027778in}}%
\pgfusepath{stroke,fill}%
}%
\begin{pgfscope}%
\pgfsys@transformshift{2.735365in}{0.352669in}%
\pgfsys@useobject{currentmarker}{}%
\end{pgfscope}%
\end{pgfscope}%
\begin{pgfscope}%
\pgfsetbuttcap%
\pgfsetroundjoin%
\definecolor{currentfill}{rgb}{0.000000,0.000000,0.000000}%
\pgfsetfillcolor{currentfill}%
\pgfsetlinewidth{0.301125pt}%
\definecolor{currentstroke}{rgb}{0.000000,0.000000,0.000000}%
\pgfsetstrokecolor{currentstroke}%
\pgfsetdash{}{0pt}%
\pgfsys@defobject{currentmarker}{\pgfqpoint{0.000000in}{0.000000in}}{\pgfqpoint{0.000000in}{0.027778in}}{%
\pgfpathmoveto{\pgfqpoint{0.000000in}{0.000000in}}%
\pgfpathlineto{\pgfqpoint{0.000000in}{0.027778in}}%
\pgfusepath{stroke,fill}%
}%
\begin{pgfscope}%
\pgfsys@transformshift{2.802873in}{0.352669in}%
\pgfsys@useobject{currentmarker}{}%
\end{pgfscope}%
\end{pgfscope}%
\begin{pgfscope}%
\pgfsetbuttcap%
\pgfsetroundjoin%
\definecolor{currentfill}{rgb}{0.000000,0.000000,0.000000}%
\pgfsetfillcolor{currentfill}%
\pgfsetlinewidth{0.301125pt}%
\definecolor{currentstroke}{rgb}{0.000000,0.000000,0.000000}%
\pgfsetstrokecolor{currentstroke}%
\pgfsetdash{}{0pt}%
\pgfsys@defobject{currentmarker}{\pgfqpoint{0.000000in}{0.000000in}}{\pgfqpoint{0.000000in}{0.027778in}}{%
\pgfpathmoveto{\pgfqpoint{0.000000in}{0.000000in}}%
\pgfpathlineto{\pgfqpoint{0.000000in}{0.027778in}}%
\pgfusepath{stroke,fill}%
}%
\begin{pgfscope}%
\pgfsys@transformshift{2.870380in}{0.352669in}%
\pgfsys@useobject{currentmarker}{}%
\end{pgfscope}%
\end{pgfscope}%
\begin{pgfscope}%
\pgfsetbuttcap%
\pgfsetroundjoin%
\definecolor{currentfill}{rgb}{0.000000,0.000000,0.000000}%
\pgfsetfillcolor{currentfill}%
\pgfsetlinewidth{0.301125pt}%
\definecolor{currentstroke}{rgb}{0.000000,0.000000,0.000000}%
\pgfsetstrokecolor{currentstroke}%
\pgfsetdash{}{0pt}%
\pgfsys@defobject{currentmarker}{\pgfqpoint{0.000000in}{0.000000in}}{\pgfqpoint{0.000000in}{0.027778in}}{%
\pgfpathmoveto{\pgfqpoint{0.000000in}{0.000000in}}%
\pgfpathlineto{\pgfqpoint{0.000000in}{0.027778in}}%
\pgfusepath{stroke,fill}%
}%
\begin{pgfscope}%
\pgfsys@transformshift{2.937888in}{0.352669in}%
\pgfsys@useobject{currentmarker}{}%
\end{pgfscope}%
\end{pgfscope}%
\begin{pgfscope}%
\pgfsetbuttcap%
\pgfsetroundjoin%
\definecolor{currentfill}{rgb}{0.000000,0.000000,0.000000}%
\pgfsetfillcolor{currentfill}%
\pgfsetlinewidth{0.301125pt}%
\definecolor{currentstroke}{rgb}{0.000000,0.000000,0.000000}%
\pgfsetstrokecolor{currentstroke}%
\pgfsetdash{}{0pt}%
\pgfsys@defobject{currentmarker}{\pgfqpoint{0.000000in}{0.000000in}}{\pgfqpoint{0.000000in}{0.027778in}}{%
\pgfpathmoveto{\pgfqpoint{0.000000in}{0.000000in}}%
\pgfpathlineto{\pgfqpoint{0.000000in}{0.027778in}}%
\pgfusepath{stroke,fill}%
}%
\begin{pgfscope}%
\pgfsys@transformshift{3.072903in}{0.352669in}%
\pgfsys@useobject{currentmarker}{}%
\end{pgfscope}%
\end{pgfscope}%
\begin{pgfscope}%
\pgfsetbuttcap%
\pgfsetroundjoin%
\definecolor{currentfill}{rgb}{0.000000,0.000000,0.000000}%
\pgfsetfillcolor{currentfill}%
\pgfsetlinewidth{0.301125pt}%
\definecolor{currentstroke}{rgb}{0.000000,0.000000,0.000000}%
\pgfsetstrokecolor{currentstroke}%
\pgfsetdash{}{0pt}%
\pgfsys@defobject{currentmarker}{\pgfqpoint{0.000000in}{0.000000in}}{\pgfqpoint{0.000000in}{0.027778in}}{%
\pgfpathmoveto{\pgfqpoint{0.000000in}{0.000000in}}%
\pgfpathlineto{\pgfqpoint{0.000000in}{0.027778in}}%
\pgfusepath{stroke,fill}%
}%
\begin{pgfscope}%
\pgfsys@transformshift{3.140410in}{0.352669in}%
\pgfsys@useobject{currentmarker}{}%
\end{pgfscope}%
\end{pgfscope}%
\begin{pgfscope}%
\pgfsetbuttcap%
\pgfsetroundjoin%
\definecolor{currentfill}{rgb}{0.000000,0.000000,0.000000}%
\pgfsetfillcolor{currentfill}%
\pgfsetlinewidth{0.301125pt}%
\definecolor{currentstroke}{rgb}{0.000000,0.000000,0.000000}%
\pgfsetstrokecolor{currentstroke}%
\pgfsetdash{}{0pt}%
\pgfsys@defobject{currentmarker}{\pgfqpoint{0.000000in}{0.000000in}}{\pgfqpoint{0.000000in}{0.027778in}}{%
\pgfpathmoveto{\pgfqpoint{0.000000in}{0.000000in}}%
\pgfpathlineto{\pgfqpoint{0.000000in}{0.027778in}}%
\pgfusepath{stroke,fill}%
}%
\begin{pgfscope}%
\pgfsys@transformshift{3.207917in}{0.352669in}%
\pgfsys@useobject{currentmarker}{}%
\end{pgfscope}%
\end{pgfscope}%
\begin{pgfscope}%
\pgfsetbuttcap%
\pgfsetroundjoin%
\definecolor{currentfill}{rgb}{0.000000,0.000000,0.000000}%
\pgfsetfillcolor{currentfill}%
\pgfsetlinewidth{0.301125pt}%
\definecolor{currentstroke}{rgb}{0.000000,0.000000,0.000000}%
\pgfsetstrokecolor{currentstroke}%
\pgfsetdash{}{0pt}%
\pgfsys@defobject{currentmarker}{\pgfqpoint{0.000000in}{0.000000in}}{\pgfqpoint{0.000000in}{0.027778in}}{%
\pgfpathmoveto{\pgfqpoint{0.000000in}{0.000000in}}%
\pgfpathlineto{\pgfqpoint{0.000000in}{0.027778in}}%
\pgfusepath{stroke,fill}%
}%
\begin{pgfscope}%
\pgfsys@transformshift{3.275425in}{0.352669in}%
\pgfsys@useobject{currentmarker}{}%
\end{pgfscope}%
\end{pgfscope}%
\begin{pgfscope}%
\pgfsetbuttcap%
\pgfsetroundjoin%
\definecolor{currentfill}{rgb}{0.000000,0.000000,0.000000}%
\pgfsetfillcolor{currentfill}%
\pgfsetlinewidth{0.301125pt}%
\definecolor{currentstroke}{rgb}{0.000000,0.000000,0.000000}%
\pgfsetstrokecolor{currentstroke}%
\pgfsetdash{}{0pt}%
\pgfsys@defobject{currentmarker}{\pgfqpoint{0.000000in}{0.000000in}}{\pgfqpoint{0.000000in}{0.027778in}}{%
\pgfpathmoveto{\pgfqpoint{0.000000in}{0.000000in}}%
\pgfpathlineto{\pgfqpoint{0.000000in}{0.027778in}}%
\pgfusepath{stroke,fill}%
}%
\begin{pgfscope}%
\pgfsys@transformshift{3.410440in}{0.352669in}%
\pgfsys@useobject{currentmarker}{}%
\end{pgfscope}%
\end{pgfscope}%
\begin{pgfscope}%
\pgfsetbuttcap%
\pgfsetroundjoin%
\definecolor{currentfill}{rgb}{0.000000,0.000000,0.000000}%
\pgfsetfillcolor{currentfill}%
\pgfsetlinewidth{0.301125pt}%
\definecolor{currentstroke}{rgb}{0.000000,0.000000,0.000000}%
\pgfsetstrokecolor{currentstroke}%
\pgfsetdash{}{0pt}%
\pgfsys@defobject{currentmarker}{\pgfqpoint{0.000000in}{0.000000in}}{\pgfqpoint{0.000000in}{0.027778in}}{%
\pgfpathmoveto{\pgfqpoint{0.000000in}{0.000000in}}%
\pgfpathlineto{\pgfqpoint{0.000000in}{0.027778in}}%
\pgfusepath{stroke,fill}%
}%
\begin{pgfscope}%
\pgfsys@transformshift{3.477947in}{0.352669in}%
\pgfsys@useobject{currentmarker}{}%
\end{pgfscope}%
\end{pgfscope}%
\begin{pgfscope}%
\pgfsetbuttcap%
\pgfsetroundjoin%
\definecolor{currentfill}{rgb}{0.000000,0.000000,0.000000}%
\pgfsetfillcolor{currentfill}%
\pgfsetlinewidth{0.301125pt}%
\definecolor{currentstroke}{rgb}{0.000000,0.000000,0.000000}%
\pgfsetstrokecolor{currentstroke}%
\pgfsetdash{}{0pt}%
\pgfsys@defobject{currentmarker}{\pgfqpoint{0.000000in}{0.000000in}}{\pgfqpoint{0.000000in}{0.027778in}}{%
\pgfpathmoveto{\pgfqpoint{0.000000in}{0.000000in}}%
\pgfpathlineto{\pgfqpoint{0.000000in}{0.027778in}}%
\pgfusepath{stroke,fill}%
}%
\begin{pgfscope}%
\pgfsys@transformshift{3.545455in}{0.352669in}%
\pgfsys@useobject{currentmarker}{}%
\end{pgfscope}%
\end{pgfscope}%
\begin{pgfscope}%
\pgfsetbuttcap%
\pgfsetroundjoin%
\definecolor{currentfill}{rgb}{0.000000,0.000000,0.000000}%
\pgfsetfillcolor{currentfill}%
\pgfsetlinewidth{0.301125pt}%
\definecolor{currentstroke}{rgb}{0.000000,0.000000,0.000000}%
\pgfsetstrokecolor{currentstroke}%
\pgfsetdash{}{0pt}%
\pgfsys@defobject{currentmarker}{\pgfqpoint{0.000000in}{0.000000in}}{\pgfqpoint{0.000000in}{0.027778in}}{%
\pgfpathmoveto{\pgfqpoint{0.000000in}{0.000000in}}%
\pgfpathlineto{\pgfqpoint{0.000000in}{0.027778in}}%
\pgfusepath{stroke,fill}%
}%
\begin{pgfscope}%
\pgfsys@transformshift{3.612962in}{0.352669in}%
\pgfsys@useobject{currentmarker}{}%
\end{pgfscope}%
\end{pgfscope}%
\begin{pgfscope}%
\definecolor{textcolor}{rgb}{0.000000,0.000000,0.000000}%
\pgfsetstrokecolor{textcolor}%
\pgfsetfillcolor{textcolor}%
\pgftext[x=3.005395in,y=0.140972in,,top]{\color{textcolor}{\sffamily\fontsize{9.000000}{10.800000}\selectfont\catcode`\^=\active\def^{\ifmmode\sp\else\^{}\fi}\catcode`\%=\active\def%{\%}$x$}}%
\end{pgfscope}%
\begin{pgfscope}%
\pgfsetbuttcap%
\pgfsetroundjoin%
\definecolor{currentfill}{rgb}{0.000000,0.000000,0.000000}%
\pgfsetfillcolor{currentfill}%
\pgfsetlinewidth{0.501875pt}%
\definecolor{currentstroke}{rgb}{0.000000,0.000000,0.000000}%
\pgfsetstrokecolor{currentstroke}%
\pgfsetdash{}{0pt}%
\pgfsys@defobject{currentmarker}{\pgfqpoint{0.000000in}{0.000000in}}{\pgfqpoint{0.041667in}{0.000000in}}{%
\pgfpathmoveto{\pgfqpoint{0.000000in}{0.000000in}}%
\pgfpathlineto{\pgfqpoint{0.041667in}{0.000000in}}%
\pgfusepath{stroke,fill}%
}%
\begin{pgfscope}%
\pgfsys@transformshift{2.330321in}{0.352669in}%
\pgfsys@useobject{currentmarker}{}%
\end{pgfscope}%
\end{pgfscope}%
\begin{pgfscope}%
\definecolor{textcolor}{rgb}{0.000000,0.000000,0.000000}%
\pgfsetstrokecolor{textcolor}%
\pgfsetfillcolor{textcolor}%
\pgftext[x=2.019228in, y=0.310460in, left, base]{\color{textcolor}{\sffamily\fontsize{8.000000}{9.600000}\selectfont\catcode`\^=\active\def^{\ifmmode\sp\else\^{}\fi}\catcode`\%=\active\def%{\%}\ensuremath{-}0.2}}%
\end{pgfscope}%
\begin{pgfscope}%
\pgfsetbuttcap%
\pgfsetroundjoin%
\definecolor{currentfill}{rgb}{0.000000,0.000000,0.000000}%
\pgfsetfillcolor{currentfill}%
\pgfsetlinewidth{0.501875pt}%
\definecolor{currentstroke}{rgb}{0.000000,0.000000,0.000000}%
\pgfsetstrokecolor{currentstroke}%
\pgfsetdash{}{0pt}%
\pgfsys@defobject{currentmarker}{\pgfqpoint{0.000000in}{0.000000in}}{\pgfqpoint{0.041667in}{0.000000in}}{%
\pgfpathmoveto{\pgfqpoint{0.000000in}{0.000000in}}%
\pgfpathlineto{\pgfqpoint{0.041667in}{0.000000in}}%
\pgfusepath{stroke,fill}%
}%
\begin{pgfscope}%
\pgfsys@transformshift{2.330321in}{0.615039in}%
\pgfsys@useobject{currentmarker}{}%
\end{pgfscope}%
\end{pgfscope}%
\begin{pgfscope}%
\definecolor{textcolor}{rgb}{0.000000,0.000000,0.000000}%
\pgfsetstrokecolor{textcolor}%
\pgfsetfillcolor{textcolor}%
\pgftext[x=2.105006in, y=0.572830in, left, base]{\color{textcolor}{\sffamily\fontsize{8.000000}{9.600000}\selectfont\catcode`\^=\active\def^{\ifmmode\sp\else\^{}\fi}\catcode`\%=\active\def%{\%}0.0}}%
\end{pgfscope}%
\begin{pgfscope}%
\pgfsetbuttcap%
\pgfsetroundjoin%
\definecolor{currentfill}{rgb}{0.000000,0.000000,0.000000}%
\pgfsetfillcolor{currentfill}%
\pgfsetlinewidth{0.501875pt}%
\definecolor{currentstroke}{rgb}{0.000000,0.000000,0.000000}%
\pgfsetstrokecolor{currentstroke}%
\pgfsetdash{}{0pt}%
\pgfsys@defobject{currentmarker}{\pgfqpoint{0.000000in}{0.000000in}}{\pgfqpoint{0.041667in}{0.000000in}}{%
\pgfpathmoveto{\pgfqpoint{0.000000in}{0.000000in}}%
\pgfpathlineto{\pgfqpoint{0.041667in}{0.000000in}}%
\pgfusepath{stroke,fill}%
}%
\begin{pgfscope}%
\pgfsys@transformshift{2.330321in}{0.877409in}%
\pgfsys@useobject{currentmarker}{}%
\end{pgfscope}%
\end{pgfscope}%
\begin{pgfscope}%
\definecolor{textcolor}{rgb}{0.000000,0.000000,0.000000}%
\pgfsetstrokecolor{textcolor}%
\pgfsetfillcolor{textcolor}%
\pgftext[x=2.105006in, y=0.835200in, left, base]{\color{textcolor}{\sffamily\fontsize{8.000000}{9.600000}\selectfont\catcode`\^=\active\def^{\ifmmode\sp\else\^{}\fi}\catcode`\%=\active\def%{\%}0.2}}%
\end{pgfscope}%
\begin{pgfscope}%
\pgfsetbuttcap%
\pgfsetroundjoin%
\definecolor{currentfill}{rgb}{0.000000,0.000000,0.000000}%
\pgfsetfillcolor{currentfill}%
\pgfsetlinewidth{0.501875pt}%
\definecolor{currentstroke}{rgb}{0.000000,0.000000,0.000000}%
\pgfsetstrokecolor{currentstroke}%
\pgfsetdash{}{0pt}%
\pgfsys@defobject{currentmarker}{\pgfqpoint{0.000000in}{0.000000in}}{\pgfqpoint{0.041667in}{0.000000in}}{%
\pgfpathmoveto{\pgfqpoint{0.000000in}{0.000000in}}%
\pgfpathlineto{\pgfqpoint{0.041667in}{0.000000in}}%
\pgfusepath{stroke,fill}%
}%
\begin{pgfscope}%
\pgfsys@transformshift{2.330321in}{1.139780in}%
\pgfsys@useobject{currentmarker}{}%
\end{pgfscope}%
\end{pgfscope}%
\begin{pgfscope}%
\definecolor{textcolor}{rgb}{0.000000,0.000000,0.000000}%
\pgfsetstrokecolor{textcolor}%
\pgfsetfillcolor{textcolor}%
\pgftext[x=2.105006in, y=1.097571in, left, base]{\color{textcolor}{\sffamily\fontsize{8.000000}{9.600000}\selectfont\catcode`\^=\active\def^{\ifmmode\sp\else\^{}\fi}\catcode`\%=\active\def%{\%}0.4}}%
\end{pgfscope}%
\begin{pgfscope}%
\pgfsetbuttcap%
\pgfsetroundjoin%
\definecolor{currentfill}{rgb}{0.000000,0.000000,0.000000}%
\pgfsetfillcolor{currentfill}%
\pgfsetlinewidth{0.501875pt}%
\definecolor{currentstroke}{rgb}{0.000000,0.000000,0.000000}%
\pgfsetstrokecolor{currentstroke}%
\pgfsetdash{}{0pt}%
\pgfsys@defobject{currentmarker}{\pgfqpoint{0.000000in}{0.000000in}}{\pgfqpoint{0.041667in}{0.000000in}}{%
\pgfpathmoveto{\pgfqpoint{0.000000in}{0.000000in}}%
\pgfpathlineto{\pgfqpoint{0.041667in}{0.000000in}}%
\pgfusepath{stroke,fill}%
}%
\begin{pgfscope}%
\pgfsys@transformshift{2.330321in}{1.402150in}%
\pgfsys@useobject{currentmarker}{}%
\end{pgfscope}%
\end{pgfscope}%
\begin{pgfscope}%
\definecolor{textcolor}{rgb}{0.000000,0.000000,0.000000}%
\pgfsetstrokecolor{textcolor}%
\pgfsetfillcolor{textcolor}%
\pgftext[x=2.105006in, y=1.359941in, left, base]{\color{textcolor}{\sffamily\fontsize{8.000000}{9.600000}\selectfont\catcode`\^=\active\def^{\ifmmode\sp\else\^{}\fi}\catcode`\%=\active\def%{\%}0.6}}%
\end{pgfscope}%
\begin{pgfscope}%
\pgfsetbuttcap%
\pgfsetroundjoin%
\definecolor{currentfill}{rgb}{0.000000,0.000000,0.000000}%
\pgfsetfillcolor{currentfill}%
\pgfsetlinewidth{0.501875pt}%
\definecolor{currentstroke}{rgb}{0.000000,0.000000,0.000000}%
\pgfsetstrokecolor{currentstroke}%
\pgfsetdash{}{0pt}%
\pgfsys@defobject{currentmarker}{\pgfqpoint{0.000000in}{0.000000in}}{\pgfqpoint{0.041667in}{0.000000in}}{%
\pgfpathmoveto{\pgfqpoint{0.000000in}{0.000000in}}%
\pgfpathlineto{\pgfqpoint{0.041667in}{0.000000in}}%
\pgfusepath{stroke,fill}%
}%
\begin{pgfscope}%
\pgfsys@transformshift{2.330321in}{1.664521in}%
\pgfsys@useobject{currentmarker}{}%
\end{pgfscope}%
\end{pgfscope}%
\begin{pgfscope}%
\definecolor{textcolor}{rgb}{0.000000,0.000000,0.000000}%
\pgfsetstrokecolor{textcolor}%
\pgfsetfillcolor{textcolor}%
\pgftext[x=2.105006in, y=1.622311in, left, base]{\color{textcolor}{\sffamily\fontsize{8.000000}{9.600000}\selectfont\catcode`\^=\active\def^{\ifmmode\sp\else\^{}\fi}\catcode`\%=\active\def%{\%}0.8}}%
\end{pgfscope}%
\begin{pgfscope}%
\pgfsetbuttcap%
\pgfsetroundjoin%
\definecolor{currentfill}{rgb}{0.000000,0.000000,0.000000}%
\pgfsetfillcolor{currentfill}%
\pgfsetlinewidth{0.501875pt}%
\definecolor{currentstroke}{rgb}{0.000000,0.000000,0.000000}%
\pgfsetstrokecolor{currentstroke}%
\pgfsetdash{}{0pt}%
\pgfsys@defobject{currentmarker}{\pgfqpoint{0.000000in}{0.000000in}}{\pgfqpoint{0.041667in}{0.000000in}}{%
\pgfpathmoveto{\pgfqpoint{0.000000in}{0.000000in}}%
\pgfpathlineto{\pgfqpoint{0.041667in}{0.000000in}}%
\pgfusepath{stroke,fill}%
}%
\begin{pgfscope}%
\pgfsys@transformshift{2.330321in}{1.926891in}%
\pgfsys@useobject{currentmarker}{}%
\end{pgfscope}%
\end{pgfscope}%
\begin{pgfscope}%
\definecolor{textcolor}{rgb}{0.000000,0.000000,0.000000}%
\pgfsetstrokecolor{textcolor}%
\pgfsetfillcolor{textcolor}%
\pgftext[x=2.105006in, y=1.884682in, left, base]{\color{textcolor}{\sffamily\fontsize{8.000000}{9.600000}\selectfont\catcode`\^=\active\def^{\ifmmode\sp\else\^{}\fi}\catcode`\%=\active\def%{\%}1.0}}%
\end{pgfscope}%
\begin{pgfscope}%
\pgfsetbuttcap%
\pgfsetroundjoin%
\definecolor{currentfill}{rgb}{0.000000,0.000000,0.000000}%
\pgfsetfillcolor{currentfill}%
\pgfsetlinewidth{0.301125pt}%
\definecolor{currentstroke}{rgb}{0.000000,0.000000,0.000000}%
\pgfsetstrokecolor{currentstroke}%
\pgfsetdash{}{0pt}%
\pgfsys@defobject{currentmarker}{\pgfqpoint{0.000000in}{0.000000in}}{\pgfqpoint{0.027778in}{0.000000in}}{%
\pgfpathmoveto{\pgfqpoint{0.000000in}{0.000000in}}%
\pgfpathlineto{\pgfqpoint{0.027778in}{0.000000in}}%
\pgfusepath{stroke,fill}%
}%
\begin{pgfscope}%
\pgfsys@transformshift{2.330321in}{0.418261in}%
\pgfsys@useobject{currentmarker}{}%
\end{pgfscope}%
\end{pgfscope}%
\begin{pgfscope}%
\pgfsetbuttcap%
\pgfsetroundjoin%
\definecolor{currentfill}{rgb}{0.000000,0.000000,0.000000}%
\pgfsetfillcolor{currentfill}%
\pgfsetlinewidth{0.301125pt}%
\definecolor{currentstroke}{rgb}{0.000000,0.000000,0.000000}%
\pgfsetstrokecolor{currentstroke}%
\pgfsetdash{}{0pt}%
\pgfsys@defobject{currentmarker}{\pgfqpoint{0.000000in}{0.000000in}}{\pgfqpoint{0.027778in}{0.000000in}}{%
\pgfpathmoveto{\pgfqpoint{0.000000in}{0.000000in}}%
\pgfpathlineto{\pgfqpoint{0.027778in}{0.000000in}}%
\pgfusepath{stroke,fill}%
}%
\begin{pgfscope}%
\pgfsys@transformshift{2.330321in}{0.483854in}%
\pgfsys@useobject{currentmarker}{}%
\end{pgfscope}%
\end{pgfscope}%
\begin{pgfscope}%
\pgfsetbuttcap%
\pgfsetroundjoin%
\definecolor{currentfill}{rgb}{0.000000,0.000000,0.000000}%
\pgfsetfillcolor{currentfill}%
\pgfsetlinewidth{0.301125pt}%
\definecolor{currentstroke}{rgb}{0.000000,0.000000,0.000000}%
\pgfsetstrokecolor{currentstroke}%
\pgfsetdash{}{0pt}%
\pgfsys@defobject{currentmarker}{\pgfqpoint{0.000000in}{0.000000in}}{\pgfqpoint{0.027778in}{0.000000in}}{%
\pgfpathmoveto{\pgfqpoint{0.000000in}{0.000000in}}%
\pgfpathlineto{\pgfqpoint{0.027778in}{0.000000in}}%
\pgfusepath{stroke,fill}%
}%
\begin{pgfscope}%
\pgfsys@transformshift{2.330321in}{0.549447in}%
\pgfsys@useobject{currentmarker}{}%
\end{pgfscope}%
\end{pgfscope}%
\begin{pgfscope}%
\pgfsetbuttcap%
\pgfsetroundjoin%
\definecolor{currentfill}{rgb}{0.000000,0.000000,0.000000}%
\pgfsetfillcolor{currentfill}%
\pgfsetlinewidth{0.301125pt}%
\definecolor{currentstroke}{rgb}{0.000000,0.000000,0.000000}%
\pgfsetstrokecolor{currentstroke}%
\pgfsetdash{}{0pt}%
\pgfsys@defobject{currentmarker}{\pgfqpoint{0.000000in}{0.000000in}}{\pgfqpoint{0.027778in}{0.000000in}}{%
\pgfpathmoveto{\pgfqpoint{0.000000in}{0.000000in}}%
\pgfpathlineto{\pgfqpoint{0.027778in}{0.000000in}}%
\pgfusepath{stroke,fill}%
}%
\begin{pgfscope}%
\pgfsys@transformshift{2.330321in}{0.680632in}%
\pgfsys@useobject{currentmarker}{}%
\end{pgfscope}%
\end{pgfscope}%
\begin{pgfscope}%
\pgfsetbuttcap%
\pgfsetroundjoin%
\definecolor{currentfill}{rgb}{0.000000,0.000000,0.000000}%
\pgfsetfillcolor{currentfill}%
\pgfsetlinewidth{0.301125pt}%
\definecolor{currentstroke}{rgb}{0.000000,0.000000,0.000000}%
\pgfsetstrokecolor{currentstroke}%
\pgfsetdash{}{0pt}%
\pgfsys@defobject{currentmarker}{\pgfqpoint{0.000000in}{0.000000in}}{\pgfqpoint{0.027778in}{0.000000in}}{%
\pgfpathmoveto{\pgfqpoint{0.000000in}{0.000000in}}%
\pgfpathlineto{\pgfqpoint{0.027778in}{0.000000in}}%
\pgfusepath{stroke,fill}%
}%
\begin{pgfscope}%
\pgfsys@transformshift{2.330321in}{0.746224in}%
\pgfsys@useobject{currentmarker}{}%
\end{pgfscope}%
\end{pgfscope}%
\begin{pgfscope}%
\pgfsetbuttcap%
\pgfsetroundjoin%
\definecolor{currentfill}{rgb}{0.000000,0.000000,0.000000}%
\pgfsetfillcolor{currentfill}%
\pgfsetlinewidth{0.301125pt}%
\definecolor{currentstroke}{rgb}{0.000000,0.000000,0.000000}%
\pgfsetstrokecolor{currentstroke}%
\pgfsetdash{}{0pt}%
\pgfsys@defobject{currentmarker}{\pgfqpoint{0.000000in}{0.000000in}}{\pgfqpoint{0.027778in}{0.000000in}}{%
\pgfpathmoveto{\pgfqpoint{0.000000in}{0.000000in}}%
\pgfpathlineto{\pgfqpoint{0.027778in}{0.000000in}}%
\pgfusepath{stroke,fill}%
}%
\begin{pgfscope}%
\pgfsys@transformshift{2.330321in}{0.811817in}%
\pgfsys@useobject{currentmarker}{}%
\end{pgfscope}%
\end{pgfscope}%
\begin{pgfscope}%
\pgfsetbuttcap%
\pgfsetroundjoin%
\definecolor{currentfill}{rgb}{0.000000,0.000000,0.000000}%
\pgfsetfillcolor{currentfill}%
\pgfsetlinewidth{0.301125pt}%
\definecolor{currentstroke}{rgb}{0.000000,0.000000,0.000000}%
\pgfsetstrokecolor{currentstroke}%
\pgfsetdash{}{0pt}%
\pgfsys@defobject{currentmarker}{\pgfqpoint{0.000000in}{0.000000in}}{\pgfqpoint{0.027778in}{0.000000in}}{%
\pgfpathmoveto{\pgfqpoint{0.000000in}{0.000000in}}%
\pgfpathlineto{\pgfqpoint{0.027778in}{0.000000in}}%
\pgfusepath{stroke,fill}%
}%
\begin{pgfscope}%
\pgfsys@transformshift{2.330321in}{0.943002in}%
\pgfsys@useobject{currentmarker}{}%
\end{pgfscope}%
\end{pgfscope}%
\begin{pgfscope}%
\pgfsetbuttcap%
\pgfsetroundjoin%
\definecolor{currentfill}{rgb}{0.000000,0.000000,0.000000}%
\pgfsetfillcolor{currentfill}%
\pgfsetlinewidth{0.301125pt}%
\definecolor{currentstroke}{rgb}{0.000000,0.000000,0.000000}%
\pgfsetstrokecolor{currentstroke}%
\pgfsetdash{}{0pt}%
\pgfsys@defobject{currentmarker}{\pgfqpoint{0.000000in}{0.000000in}}{\pgfqpoint{0.027778in}{0.000000in}}{%
\pgfpathmoveto{\pgfqpoint{0.000000in}{0.000000in}}%
\pgfpathlineto{\pgfqpoint{0.027778in}{0.000000in}}%
\pgfusepath{stroke,fill}%
}%
\begin{pgfscope}%
\pgfsys@transformshift{2.330321in}{1.008595in}%
\pgfsys@useobject{currentmarker}{}%
\end{pgfscope}%
\end{pgfscope}%
\begin{pgfscope}%
\pgfsetbuttcap%
\pgfsetroundjoin%
\definecolor{currentfill}{rgb}{0.000000,0.000000,0.000000}%
\pgfsetfillcolor{currentfill}%
\pgfsetlinewidth{0.301125pt}%
\definecolor{currentstroke}{rgb}{0.000000,0.000000,0.000000}%
\pgfsetstrokecolor{currentstroke}%
\pgfsetdash{}{0pt}%
\pgfsys@defobject{currentmarker}{\pgfqpoint{0.000000in}{0.000000in}}{\pgfqpoint{0.027778in}{0.000000in}}{%
\pgfpathmoveto{\pgfqpoint{0.000000in}{0.000000in}}%
\pgfpathlineto{\pgfqpoint{0.027778in}{0.000000in}}%
\pgfusepath{stroke,fill}%
}%
\begin{pgfscope}%
\pgfsys@transformshift{2.330321in}{1.074187in}%
\pgfsys@useobject{currentmarker}{}%
\end{pgfscope}%
\end{pgfscope}%
\begin{pgfscope}%
\pgfsetbuttcap%
\pgfsetroundjoin%
\definecolor{currentfill}{rgb}{0.000000,0.000000,0.000000}%
\pgfsetfillcolor{currentfill}%
\pgfsetlinewidth{0.301125pt}%
\definecolor{currentstroke}{rgb}{0.000000,0.000000,0.000000}%
\pgfsetstrokecolor{currentstroke}%
\pgfsetdash{}{0pt}%
\pgfsys@defobject{currentmarker}{\pgfqpoint{0.000000in}{0.000000in}}{\pgfqpoint{0.027778in}{0.000000in}}{%
\pgfpathmoveto{\pgfqpoint{0.000000in}{0.000000in}}%
\pgfpathlineto{\pgfqpoint{0.027778in}{0.000000in}}%
\pgfusepath{stroke,fill}%
}%
\begin{pgfscope}%
\pgfsys@transformshift{2.330321in}{1.205372in}%
\pgfsys@useobject{currentmarker}{}%
\end{pgfscope}%
\end{pgfscope}%
\begin{pgfscope}%
\pgfsetbuttcap%
\pgfsetroundjoin%
\definecolor{currentfill}{rgb}{0.000000,0.000000,0.000000}%
\pgfsetfillcolor{currentfill}%
\pgfsetlinewidth{0.301125pt}%
\definecolor{currentstroke}{rgb}{0.000000,0.000000,0.000000}%
\pgfsetstrokecolor{currentstroke}%
\pgfsetdash{}{0pt}%
\pgfsys@defobject{currentmarker}{\pgfqpoint{0.000000in}{0.000000in}}{\pgfqpoint{0.027778in}{0.000000in}}{%
\pgfpathmoveto{\pgfqpoint{0.000000in}{0.000000in}}%
\pgfpathlineto{\pgfqpoint{0.027778in}{0.000000in}}%
\pgfusepath{stroke,fill}%
}%
\begin{pgfscope}%
\pgfsys@transformshift{2.330321in}{1.270965in}%
\pgfsys@useobject{currentmarker}{}%
\end{pgfscope}%
\end{pgfscope}%
\begin{pgfscope}%
\pgfsetbuttcap%
\pgfsetroundjoin%
\definecolor{currentfill}{rgb}{0.000000,0.000000,0.000000}%
\pgfsetfillcolor{currentfill}%
\pgfsetlinewidth{0.301125pt}%
\definecolor{currentstroke}{rgb}{0.000000,0.000000,0.000000}%
\pgfsetstrokecolor{currentstroke}%
\pgfsetdash{}{0pt}%
\pgfsys@defobject{currentmarker}{\pgfqpoint{0.000000in}{0.000000in}}{\pgfqpoint{0.027778in}{0.000000in}}{%
\pgfpathmoveto{\pgfqpoint{0.000000in}{0.000000in}}%
\pgfpathlineto{\pgfqpoint{0.027778in}{0.000000in}}%
\pgfusepath{stroke,fill}%
}%
\begin{pgfscope}%
\pgfsys@transformshift{2.330321in}{1.336558in}%
\pgfsys@useobject{currentmarker}{}%
\end{pgfscope}%
\end{pgfscope}%
\begin{pgfscope}%
\pgfsetbuttcap%
\pgfsetroundjoin%
\definecolor{currentfill}{rgb}{0.000000,0.000000,0.000000}%
\pgfsetfillcolor{currentfill}%
\pgfsetlinewidth{0.301125pt}%
\definecolor{currentstroke}{rgb}{0.000000,0.000000,0.000000}%
\pgfsetstrokecolor{currentstroke}%
\pgfsetdash{}{0pt}%
\pgfsys@defobject{currentmarker}{\pgfqpoint{0.000000in}{0.000000in}}{\pgfqpoint{0.027778in}{0.000000in}}{%
\pgfpathmoveto{\pgfqpoint{0.000000in}{0.000000in}}%
\pgfpathlineto{\pgfqpoint{0.027778in}{0.000000in}}%
\pgfusepath{stroke,fill}%
}%
\begin{pgfscope}%
\pgfsys@transformshift{2.330321in}{1.467743in}%
\pgfsys@useobject{currentmarker}{}%
\end{pgfscope}%
\end{pgfscope}%
\begin{pgfscope}%
\pgfsetbuttcap%
\pgfsetroundjoin%
\definecolor{currentfill}{rgb}{0.000000,0.000000,0.000000}%
\pgfsetfillcolor{currentfill}%
\pgfsetlinewidth{0.301125pt}%
\definecolor{currentstroke}{rgb}{0.000000,0.000000,0.000000}%
\pgfsetstrokecolor{currentstroke}%
\pgfsetdash{}{0pt}%
\pgfsys@defobject{currentmarker}{\pgfqpoint{0.000000in}{0.000000in}}{\pgfqpoint{0.027778in}{0.000000in}}{%
\pgfpathmoveto{\pgfqpoint{0.000000in}{0.000000in}}%
\pgfpathlineto{\pgfqpoint{0.027778in}{0.000000in}}%
\pgfusepath{stroke,fill}%
}%
\begin{pgfscope}%
\pgfsys@transformshift{2.330321in}{1.533335in}%
\pgfsys@useobject{currentmarker}{}%
\end{pgfscope}%
\end{pgfscope}%
\begin{pgfscope}%
\pgfsetbuttcap%
\pgfsetroundjoin%
\definecolor{currentfill}{rgb}{0.000000,0.000000,0.000000}%
\pgfsetfillcolor{currentfill}%
\pgfsetlinewidth{0.301125pt}%
\definecolor{currentstroke}{rgb}{0.000000,0.000000,0.000000}%
\pgfsetstrokecolor{currentstroke}%
\pgfsetdash{}{0pt}%
\pgfsys@defobject{currentmarker}{\pgfqpoint{0.000000in}{0.000000in}}{\pgfqpoint{0.027778in}{0.000000in}}{%
\pgfpathmoveto{\pgfqpoint{0.000000in}{0.000000in}}%
\pgfpathlineto{\pgfqpoint{0.027778in}{0.000000in}}%
\pgfusepath{stroke,fill}%
}%
\begin{pgfscope}%
\pgfsys@transformshift{2.330321in}{1.598928in}%
\pgfsys@useobject{currentmarker}{}%
\end{pgfscope}%
\end{pgfscope}%
\begin{pgfscope}%
\pgfsetbuttcap%
\pgfsetroundjoin%
\definecolor{currentfill}{rgb}{0.000000,0.000000,0.000000}%
\pgfsetfillcolor{currentfill}%
\pgfsetlinewidth{0.301125pt}%
\definecolor{currentstroke}{rgb}{0.000000,0.000000,0.000000}%
\pgfsetstrokecolor{currentstroke}%
\pgfsetdash{}{0pt}%
\pgfsys@defobject{currentmarker}{\pgfqpoint{0.000000in}{0.000000in}}{\pgfqpoint{0.027778in}{0.000000in}}{%
\pgfpathmoveto{\pgfqpoint{0.000000in}{0.000000in}}%
\pgfpathlineto{\pgfqpoint{0.027778in}{0.000000in}}%
\pgfusepath{stroke,fill}%
}%
\begin{pgfscope}%
\pgfsys@transformshift{2.330321in}{1.730113in}%
\pgfsys@useobject{currentmarker}{}%
\end{pgfscope}%
\end{pgfscope}%
\begin{pgfscope}%
\pgfsetbuttcap%
\pgfsetroundjoin%
\definecolor{currentfill}{rgb}{0.000000,0.000000,0.000000}%
\pgfsetfillcolor{currentfill}%
\pgfsetlinewidth{0.301125pt}%
\definecolor{currentstroke}{rgb}{0.000000,0.000000,0.000000}%
\pgfsetstrokecolor{currentstroke}%
\pgfsetdash{}{0pt}%
\pgfsys@defobject{currentmarker}{\pgfqpoint{0.000000in}{0.000000in}}{\pgfqpoint{0.027778in}{0.000000in}}{%
\pgfpathmoveto{\pgfqpoint{0.000000in}{0.000000in}}%
\pgfpathlineto{\pgfqpoint{0.027778in}{0.000000in}}%
\pgfusepath{stroke,fill}%
}%
\begin{pgfscope}%
\pgfsys@transformshift{2.330321in}{1.795706in}%
\pgfsys@useobject{currentmarker}{}%
\end{pgfscope}%
\end{pgfscope}%
\begin{pgfscope}%
\pgfsetbuttcap%
\pgfsetroundjoin%
\definecolor{currentfill}{rgb}{0.000000,0.000000,0.000000}%
\pgfsetfillcolor{currentfill}%
\pgfsetlinewidth{0.301125pt}%
\definecolor{currentstroke}{rgb}{0.000000,0.000000,0.000000}%
\pgfsetstrokecolor{currentstroke}%
\pgfsetdash{}{0pt}%
\pgfsys@defobject{currentmarker}{\pgfqpoint{0.000000in}{0.000000in}}{\pgfqpoint{0.027778in}{0.000000in}}{%
\pgfpathmoveto{\pgfqpoint{0.000000in}{0.000000in}}%
\pgfpathlineto{\pgfqpoint{0.027778in}{0.000000in}}%
\pgfusepath{stroke,fill}%
}%
\begin{pgfscope}%
\pgfsys@transformshift{2.330321in}{1.861298in}%
\pgfsys@useobject{currentmarker}{}%
\end{pgfscope}%
\end{pgfscope}%
\begin{pgfscope}%
\pgfsetbuttcap%
\pgfsetroundjoin%
\definecolor{currentfill}{rgb}{0.000000,0.000000,0.000000}%
\pgfsetfillcolor{currentfill}%
\pgfsetlinewidth{0.301125pt}%
\definecolor{currentstroke}{rgb}{0.000000,0.000000,0.000000}%
\pgfsetstrokecolor{currentstroke}%
\pgfsetdash{}{0pt}%
\pgfsys@defobject{currentmarker}{\pgfqpoint{0.000000in}{0.000000in}}{\pgfqpoint{0.027778in}{0.000000in}}{%
\pgfpathmoveto{\pgfqpoint{0.000000in}{0.000000in}}%
\pgfpathlineto{\pgfqpoint{0.027778in}{0.000000in}}%
\pgfusepath{stroke,fill}%
}%
\begin{pgfscope}%
\pgfsys@transformshift{2.330321in}{1.992484in}%
\pgfsys@useobject{currentmarker}{}%
\end{pgfscope}%
\end{pgfscope}%
\begin{pgfscope}%
\pgfsetbuttcap%
\pgfsetroundjoin%
\definecolor{currentfill}{rgb}{0.000000,0.000000,0.000000}%
\pgfsetfillcolor{currentfill}%
\pgfsetlinewidth{0.301125pt}%
\definecolor{currentstroke}{rgb}{0.000000,0.000000,0.000000}%
\pgfsetstrokecolor{currentstroke}%
\pgfsetdash{}{0pt}%
\pgfsys@defobject{currentmarker}{\pgfqpoint{0.000000in}{0.000000in}}{\pgfqpoint{0.027778in}{0.000000in}}{%
\pgfpathmoveto{\pgfqpoint{0.000000in}{0.000000in}}%
\pgfpathlineto{\pgfqpoint{0.027778in}{0.000000in}}%
\pgfusepath{stroke,fill}%
}%
\begin{pgfscope}%
\pgfsys@transformshift{2.330321in}{2.058076in}%
\pgfsys@useobject{currentmarker}{}%
\end{pgfscope}%
\end{pgfscope}%
\begin{pgfscope}%
\pgfsetbuttcap%
\pgfsetroundjoin%
\definecolor{currentfill}{rgb}{0.000000,0.000000,0.000000}%
\pgfsetfillcolor{currentfill}%
\pgfsetlinewidth{0.301125pt}%
\definecolor{currentstroke}{rgb}{0.000000,0.000000,0.000000}%
\pgfsetstrokecolor{currentstroke}%
\pgfsetdash{}{0pt}%
\pgfsys@defobject{currentmarker}{\pgfqpoint{0.000000in}{0.000000in}}{\pgfqpoint{0.027778in}{0.000000in}}{%
\pgfpathmoveto{\pgfqpoint{0.000000in}{0.000000in}}%
\pgfpathlineto{\pgfqpoint{0.027778in}{0.000000in}}%
\pgfusepath{stroke,fill}%
}%
\begin{pgfscope}%
\pgfsys@transformshift{2.330321in}{2.123669in}%
\pgfsys@useobject{currentmarker}{}%
\end{pgfscope}%
\end{pgfscope}%
\begin{pgfscope}%
\pgfsetrectcap%
\pgfsetmiterjoin%
\pgfsetlinewidth{0.602250pt}%
\definecolor{currentstroke}{rgb}{0.000000,0.000000,0.000000}%
\pgfsetstrokecolor{currentstroke}%
\pgfsetdash{}{0pt}%
\pgfpathmoveto{\pgfqpoint{2.330321in}{0.352669in}}%
\pgfpathlineto{\pgfqpoint{2.330321in}{2.123669in}}%
\pgfusepath{stroke}%
\end{pgfscope}%
\begin{pgfscope}%
\pgfsetrectcap%
\pgfsetmiterjoin%
\pgfsetlinewidth{0.602250pt}%
\definecolor{currentstroke}{rgb}{0.000000,0.000000,0.000000}%
\pgfsetstrokecolor{currentstroke}%
\pgfsetdash{}{0pt}%
\pgfpathmoveto{\pgfqpoint{2.330321in}{0.352669in}}%
\pgfpathlineto{\pgfqpoint{3.680470in}{0.352669in}}%
\pgfusepath{stroke}%
\end{pgfscope}%
\begin{pgfscope}%
\pgfpathrectangle{\pgfqpoint{2.330321in}{0.352669in}}{\pgfqpoint{1.350149in}{1.771000in}}%
\pgfusepath{clip}%
\pgfsetrectcap%
\pgfsetroundjoin%
\pgfsetlinewidth{0.903375pt}%
\definecolor{currentstroke}{rgb}{0.250980,0.400000,0.600000}%
\pgfsetstrokecolor{currentstroke}%
\pgfsetdash{}{0pt}%
\pgfpathmoveto{\pgfqpoint{2.330321in}{0.665495in}}%
\pgfpathlineto{\pgfqpoint{2.342484in}{0.662158in}}%
\pgfpathlineto{\pgfqpoint{2.355999in}{0.660078in}}%
\pgfpathlineto{\pgfqpoint{2.369514in}{0.659616in}}%
\pgfpathlineto{\pgfqpoint{2.383029in}{0.660672in}}%
\pgfpathlineto{\pgfqpoint{2.396544in}{0.663151in}}%
\pgfpathlineto{\pgfqpoint{2.411411in}{0.667409in}}%
\pgfpathlineto{\pgfqpoint{2.426277in}{0.673153in}}%
\pgfpathlineto{\pgfqpoint{2.442495in}{0.680971in}}%
\pgfpathlineto{\pgfqpoint{2.460065in}{0.691101in}}%
\pgfpathlineto{\pgfqpoint{2.478986in}{0.703732in}}%
\pgfpathlineto{\pgfqpoint{2.500610in}{0.720072in}}%
\pgfpathlineto{\pgfqpoint{2.523585in}{0.739314in}}%
\pgfpathlineto{\pgfqpoint{2.550615in}{0.763940in}}%
\pgfpathlineto{\pgfqpoint{2.583051in}{0.795565in}}%
\pgfpathlineto{\pgfqpoint{2.631705in}{0.845334in}}%
\pgfpathlineto{\pgfqpoint{2.703335in}{0.918423in}}%
\pgfpathlineto{\pgfqpoint{2.739825in}{0.953395in}}%
\pgfpathlineto{\pgfqpoint{2.769558in}{0.979938in}}%
\pgfpathlineto{\pgfqpoint{2.796588in}{1.002171in}}%
\pgfpathlineto{\pgfqpoint{2.822267in}{1.021363in}}%
\pgfpathlineto{\pgfqpoint{2.845242in}{1.036768in}}%
\pgfpathlineto{\pgfqpoint{2.868218in}{1.050371in}}%
\pgfpathlineto{\pgfqpoint{2.889842in}{1.061427in}}%
\pgfpathlineto{\pgfqpoint{2.910114in}{1.070181in}}%
\pgfpathlineto{\pgfqpoint{2.930387in}{1.077321in}}%
\pgfpathlineto{\pgfqpoint{2.949308in}{1.082488in}}%
\pgfpathlineto{\pgfqpoint{2.968229in}{1.086184in}}%
\pgfpathlineto{\pgfqpoint{2.987150in}{1.088387in}}%
\pgfpathlineto{\pgfqpoint{3.006071in}{1.089087in}}%
\pgfpathlineto{\pgfqpoint{3.024992in}{1.088280in}}%
\pgfpathlineto{\pgfqpoint{3.043913in}{1.085969in}}%
\pgfpathlineto{\pgfqpoint{3.062834in}{1.082168in}}%
\pgfpathlineto{\pgfqpoint{3.081755in}{1.076896in}}%
\pgfpathlineto{\pgfqpoint{3.102027in}{1.069647in}}%
\pgfpathlineto{\pgfqpoint{3.122300in}{1.060787in}}%
\pgfpathlineto{\pgfqpoint{3.143924in}{1.049623in}}%
\pgfpathlineto{\pgfqpoint{3.165548in}{1.036768in}}%
\pgfpathlineto{\pgfqpoint{3.188523in}{1.021363in}}%
\pgfpathlineto{\pgfqpoint{3.212850in}{1.003230in}}%
\pgfpathlineto{\pgfqpoint{3.238529in}{0.982247in}}%
\pgfpathlineto{\pgfqpoint{3.266910in}{0.957129in}}%
\pgfpathlineto{\pgfqpoint{3.299346in}{0.926385in}}%
\pgfpathlineto{\pgfqpoint{3.339891in}{0.885766in}}%
\pgfpathlineto{\pgfqpoint{3.457472in}{0.766501in}}%
\pgfpathlineto{\pgfqpoint{3.487205in}{0.739314in}}%
\pgfpathlineto{\pgfqpoint{3.511532in}{0.718997in}}%
\pgfpathlineto{\pgfqpoint{3.533156in}{0.702775in}}%
\pgfpathlineto{\pgfqpoint{3.553429in}{0.689438in}}%
\pgfpathlineto{\pgfqpoint{3.570998in}{0.679562in}}%
\pgfpathlineto{\pgfqpoint{3.587216in}{0.672003in}}%
\pgfpathlineto{\pgfqpoint{3.603434in}{0.666096in}}%
\pgfpathlineto{\pgfqpoint{3.618301in}{0.662264in}}%
\pgfpathlineto{\pgfqpoint{3.631816in}{0.660201in}}%
\pgfpathlineto{\pgfqpoint{3.645331in}{0.659590in}}%
\pgfpathlineto{\pgfqpoint{3.658846in}{0.660528in}}%
\pgfpathlineto{\pgfqpoint{3.671009in}{0.662776in}}%
\pgfpathlineto{\pgfqpoint{3.680470in}{0.665495in}}%
\pgfpathlineto{\pgfqpoint{3.680470in}{0.665495in}}%
\pgfusepath{stroke}%
\end{pgfscope}%
\begin{pgfscope}%
\pgfpathrectangle{\pgfqpoint{2.330321in}{0.352669in}}{\pgfqpoint{1.350149in}{1.771000in}}%
\pgfusepath{clip}%
\pgfsetrectcap%
\pgfsetroundjoin%
\pgfsetlinewidth{0.903375pt}%
\definecolor{currentstroke}{rgb}{0.768627,0.556863,0.211765}%
\pgfsetstrokecolor{currentstroke}%
\pgfsetdash{}{0pt}%
\pgfpathmoveto{\pgfqpoint{2.330321in}{0.665495in}}%
\pgfpathlineto{\pgfqpoint{2.334375in}{0.672852in}}%
\pgfpathlineto{\pgfqpoint{2.338430in}{0.677483in}}%
\pgfpathlineto{\pgfqpoint{2.342484in}{0.679857in}}%
\pgfpathlineto{\pgfqpoint{2.346539in}{0.680396in}}%
\pgfpathlineto{\pgfqpoint{2.350593in}{0.679477in}}%
\pgfpathlineto{\pgfqpoint{2.355999in}{0.676558in}}%
\pgfpathlineto{\pgfqpoint{2.364108in}{0.669914in}}%
\pgfpathlineto{\pgfqpoint{2.385732in}{0.650561in}}%
\pgfpathlineto{\pgfqpoint{2.393841in}{0.645787in}}%
\pgfpathlineto{\pgfqpoint{2.400599in}{0.643549in}}%
\pgfpathlineto{\pgfqpoint{2.407356in}{0.643054in}}%
\pgfpathlineto{\pgfqpoint{2.414114in}{0.644346in}}%
\pgfpathlineto{\pgfqpoint{2.420871in}{0.647391in}}%
\pgfpathlineto{\pgfqpoint{2.427629in}{0.652086in}}%
\pgfpathlineto{\pgfqpoint{2.435738in}{0.659682in}}%
\pgfpathlineto{\pgfqpoint{2.445198in}{0.670820in}}%
\pgfpathlineto{\pgfqpoint{2.457362in}{0.687780in}}%
\pgfpathlineto{\pgfqpoint{2.504664in}{0.757066in}}%
\pgfpathlineto{\pgfqpoint{2.515476in}{0.768751in}}%
\pgfpathlineto{\pgfqpoint{2.524937in}{0.776690in}}%
\pgfpathlineto{\pgfqpoint{2.533046in}{0.781633in}}%
\pgfpathlineto{\pgfqpoint{2.541155in}{0.784784in}}%
\pgfpathlineto{\pgfqpoint{2.549264in}{0.786130in}}%
\pgfpathlineto{\pgfqpoint{2.557373in}{0.785709in}}%
\pgfpathlineto{\pgfqpoint{2.565482in}{0.783607in}}%
\pgfpathlineto{\pgfqpoint{2.573591in}{0.779954in}}%
\pgfpathlineto{\pgfqpoint{2.583051in}{0.773962in}}%
\pgfpathlineto{\pgfqpoint{2.593863in}{0.765239in}}%
\pgfpathlineto{\pgfqpoint{2.608730in}{0.751027in}}%
\pgfpathlineto{\pgfqpoint{2.641166in}{0.719103in}}%
\pgfpathlineto{\pgfqpoint{2.651978in}{0.710776in}}%
\pgfpathlineto{\pgfqpoint{2.661438in}{0.705396in}}%
\pgfpathlineto{\pgfqpoint{2.669547in}{0.702560in}}%
\pgfpathlineto{\pgfqpoint{2.676305in}{0.701652in}}%
\pgfpathlineto{\pgfqpoint{2.683062in}{0.702220in}}%
\pgfpathlineto{\pgfqpoint{2.689820in}{0.704395in}}%
\pgfpathlineto{\pgfqpoint{2.696577in}{0.708295in}}%
\pgfpathlineto{\pgfqpoint{2.703335in}{0.714025in}}%
\pgfpathlineto{\pgfqpoint{2.710092in}{0.721676in}}%
\pgfpathlineto{\pgfqpoint{2.718201in}{0.733499in}}%
\pgfpathlineto{\pgfqpoint{2.726310in}{0.748296in}}%
\pgfpathlineto{\pgfqpoint{2.734419in}{0.766136in}}%
\pgfpathlineto{\pgfqpoint{2.743880in}{0.790835in}}%
\pgfpathlineto{\pgfqpoint{2.753340in}{0.819710in}}%
\pgfpathlineto{\pgfqpoint{2.764152in}{0.857727in}}%
\pgfpathlineto{\pgfqpoint{2.776316in}{0.906637in}}%
\pgfpathlineto{\pgfqpoint{2.789831in}{0.968066in}}%
\pgfpathlineto{\pgfqpoint{2.804697in}{1.043265in}}%
\pgfpathlineto{\pgfqpoint{2.822267in}{1.140529in}}%
\pgfpathlineto{\pgfqpoint{2.845242in}{1.276914in}}%
\pgfpathlineto{\pgfqpoint{2.893896in}{1.568400in}}%
\pgfpathlineto{\pgfqpoint{2.911466in}{1.663406in}}%
\pgfpathlineto{\pgfqpoint{2.926332in}{1.735505in}}%
\pgfpathlineto{\pgfqpoint{2.938496in}{1.787493in}}%
\pgfpathlineto{\pgfqpoint{2.949308in}{1.827645in}}%
\pgfpathlineto{\pgfqpoint{2.958768in}{1.857657in}}%
\pgfpathlineto{\pgfqpoint{2.968229in}{1.882565in}}%
\pgfpathlineto{\pgfqpoint{2.976338in}{1.899657in}}%
\pgfpathlineto{\pgfqpoint{2.983095in}{1.910797in}}%
\pgfpathlineto{\pgfqpoint{2.989853in}{1.919054in}}%
\pgfpathlineto{\pgfqpoint{2.995259in}{1.923553in}}%
\pgfpathlineto{\pgfqpoint{3.000665in}{1.926164in}}%
\pgfpathlineto{\pgfqpoint{3.006071in}{1.926876in}}%
\pgfpathlineto{\pgfqpoint{3.011477in}{1.925689in}}%
\pgfpathlineto{\pgfqpoint{3.016883in}{1.922605in}}%
\pgfpathlineto{\pgfqpoint{3.022289in}{1.917636in}}%
\pgfpathlineto{\pgfqpoint{3.027695in}{1.910797in}}%
\pgfpathlineto{\pgfqpoint{3.034452in}{1.899657in}}%
\pgfpathlineto{\pgfqpoint{3.041210in}{1.885691in}}%
\pgfpathlineto{\pgfqpoint{3.049319in}{1.865306in}}%
\pgfpathlineto{\pgfqpoint{3.058779in}{1.836724in}}%
\pgfpathlineto{\pgfqpoint{3.068240in}{1.803251in}}%
\pgfpathlineto{\pgfqpoint{3.079052in}{1.759452in}}%
\pgfpathlineto{\pgfqpoint{3.091215in}{1.703809in}}%
\pgfpathlineto{\pgfqpoint{3.106082in}{1.627967in}}%
\pgfpathlineto{\pgfqpoint{3.123651in}{1.529676in}}%
\pgfpathlineto{\pgfqpoint{3.147978in}{1.384124in}}%
\pgfpathlineto{\pgfqpoint{3.189875in}{1.132781in}}%
\pgfpathlineto{\pgfqpoint{3.208796in}{1.029055in}}%
\pgfpathlineto{\pgfqpoint{3.223662in}{0.955219in}}%
\pgfpathlineto{\pgfqpoint{3.237177in}{0.895226in}}%
\pgfpathlineto{\pgfqpoint{3.249341in}{0.847730in}}%
\pgfpathlineto{\pgfqpoint{3.260153in}{0.811037in}}%
\pgfpathlineto{\pgfqpoint{3.270965in}{0.779737in}}%
\pgfpathlineto{\pgfqpoint{3.280425in}{0.756833in}}%
\pgfpathlineto{\pgfqpoint{3.288534in}{0.740521in}}%
\pgfpathlineto{\pgfqpoint{3.296643in}{0.727221in}}%
\pgfpathlineto{\pgfqpoint{3.304752in}{0.716850in}}%
\pgfpathlineto{\pgfqpoint{3.311510in}{0.710362in}}%
\pgfpathlineto{\pgfqpoint{3.318267in}{0.705741in}}%
\pgfpathlineto{\pgfqpoint{3.325025in}{0.702890in}}%
\pgfpathlineto{\pgfqpoint{3.331782in}{0.701695in}}%
\pgfpathlineto{\pgfqpoint{3.338540in}{0.702029in}}%
\pgfpathlineto{\pgfqpoint{3.345297in}{0.703754in}}%
\pgfpathlineto{\pgfqpoint{3.353406in}{0.707452in}}%
\pgfpathlineto{\pgfqpoint{3.362867in}{0.713660in}}%
\pgfpathlineto{\pgfqpoint{3.373679in}{0.722688in}}%
\pgfpathlineto{\pgfqpoint{3.388545in}{0.737192in}}%
\pgfpathlineto{\pgfqpoint{3.416927in}{0.765239in}}%
\pgfpathlineto{\pgfqpoint{3.429091in}{0.774921in}}%
\pgfpathlineto{\pgfqpoint{3.438551in}{0.780664in}}%
\pgfpathlineto{\pgfqpoint{3.446660in}{0.784070in}}%
\pgfpathlineto{\pgfqpoint{3.454769in}{0.785899in}}%
\pgfpathlineto{\pgfqpoint{3.462878in}{0.786030in}}%
\pgfpathlineto{\pgfqpoint{3.469636in}{0.784784in}}%
\pgfpathlineto{\pgfqpoint{3.477745in}{0.781633in}}%
\pgfpathlineto{\pgfqpoint{3.485854in}{0.776690in}}%
\pgfpathlineto{\pgfqpoint{3.493963in}{0.770022in}}%
\pgfpathlineto{\pgfqpoint{3.503423in}{0.760227in}}%
\pgfpathlineto{\pgfqpoint{3.514235in}{0.746734in}}%
\pgfpathlineto{\pgfqpoint{3.527750in}{0.727296in}}%
\pgfpathlineto{\pgfqpoint{3.566944in}{0.669100in}}%
\pgfpathlineto{\pgfqpoint{3.576404in}{0.658280in}}%
\pgfpathlineto{\pgfqpoint{3.584513in}{0.651021in}}%
\pgfpathlineto{\pgfqpoint{3.591271in}{0.646646in}}%
\pgfpathlineto{\pgfqpoint{3.598028in}{0.643946in}}%
\pgfpathlineto{\pgfqpoint{3.604786in}{0.643010in}}%
\pgfpathlineto{\pgfqpoint{3.611543in}{0.643860in}}%
\pgfpathlineto{\pgfqpoint{3.618301in}{0.646433in}}%
\pgfpathlineto{\pgfqpoint{3.625058in}{0.650561in}}%
\pgfpathlineto{\pgfqpoint{3.634519in}{0.658379in}}%
\pgfpathlineto{\pgfqpoint{3.656143in}{0.677435in}}%
\pgfpathlineto{\pgfqpoint{3.661549in}{0.679924in}}%
\pgfpathlineto{\pgfqpoint{3.665603in}{0.680396in}}%
\pgfpathlineto{\pgfqpoint{3.669658in}{0.679290in}}%
\pgfpathlineto{\pgfqpoint{3.673712in}{0.676212in}}%
\pgfpathlineto{\pgfqpoint{3.677767in}{0.670727in}}%
\pgfpathlineto{\pgfqpoint{3.680470in}{0.665495in}}%
\pgfpathlineto{\pgfqpoint{3.680470in}{0.665495in}}%
\pgfusepath{stroke}%
\end{pgfscope}%
\begin{pgfscope}%
\pgfpathrectangle{\pgfqpoint{2.330321in}{0.352669in}}{\pgfqpoint{1.350149in}{1.771000in}}%
\pgfusepath{clip}%
\pgfsetrectcap%
\pgfsetroundjoin%
\pgfsetlinewidth{0.903375pt}%
\definecolor{currentstroke}{rgb}{0.619608,0.188235,0.172549}%
\pgfsetstrokecolor{currentstroke}%
\pgfsetdash{}{0pt}%
\pgfpathmoveto{\pgfqpoint{2.330321in}{0.665495in}}%
\pgfpathlineto{\pgfqpoint{2.337078in}{0.667094in}}%
\pgfpathlineto{\pgfqpoint{2.349242in}{0.667822in}}%
\pgfpathlineto{\pgfqpoint{2.364108in}{0.668995in}}%
\pgfpathlineto{\pgfqpoint{2.376272in}{0.671432in}}%
\pgfpathlineto{\pgfqpoint{2.393841in}{0.676707in}}%
\pgfpathlineto{\pgfqpoint{2.420871in}{0.684787in}}%
\pgfpathlineto{\pgfqpoint{2.437089in}{0.688006in}}%
\pgfpathlineto{\pgfqpoint{2.457362in}{0.690416in}}%
\pgfpathlineto{\pgfqpoint{2.487095in}{0.693918in}}%
\pgfpathlineto{\pgfqpoint{2.501961in}{0.697248in}}%
\pgfpathlineto{\pgfqpoint{2.515476in}{0.701774in}}%
\pgfpathlineto{\pgfqpoint{2.528991in}{0.707864in}}%
\pgfpathlineto{\pgfqpoint{2.543858in}{0.716263in}}%
\pgfpathlineto{\pgfqpoint{2.561427in}{0.727986in}}%
\pgfpathlineto{\pgfqpoint{2.629002in}{0.775085in}}%
\pgfpathlineto{\pgfqpoint{2.650626in}{0.787053in}}%
\pgfpathlineto{\pgfqpoint{2.687117in}{0.806863in}}%
\pgfpathlineto{\pgfqpoint{2.700632in}{0.816215in}}%
\pgfpathlineto{\pgfqpoint{2.711444in}{0.825405in}}%
\pgfpathlineto{\pgfqpoint{2.722256in}{0.836637in}}%
\pgfpathlineto{\pgfqpoint{2.731716in}{0.848511in}}%
\pgfpathlineto{\pgfqpoint{2.741177in}{0.862614in}}%
\pgfpathlineto{\pgfqpoint{2.750637in}{0.879220in}}%
\pgfpathlineto{\pgfqpoint{2.760098in}{0.898567in}}%
\pgfpathlineto{\pgfqpoint{2.770910in}{0.924279in}}%
\pgfpathlineto{\pgfqpoint{2.781722in}{0.954019in}}%
\pgfpathlineto{\pgfqpoint{2.792534in}{0.987881in}}%
\pgfpathlineto{\pgfqpoint{2.804697in}{1.030871in}}%
\pgfpathlineto{\pgfqpoint{2.818212in}{1.084450in}}%
\pgfpathlineto{\pgfqpoint{2.833079in}{1.149761in}}%
\pgfpathlineto{\pgfqpoint{2.850648in}{1.233944in}}%
\pgfpathlineto{\pgfqpoint{2.877678in}{1.372147in}}%
\pgfpathlineto{\pgfqpoint{2.908763in}{1.529522in}}%
\pgfpathlineto{\pgfqpoint{2.924981in}{1.604248in}}%
\pgfpathlineto{\pgfqpoint{2.938496in}{1.659787in}}%
\pgfpathlineto{\pgfqpoint{2.949308in}{1.698662in}}%
\pgfpathlineto{\pgfqpoint{2.958768in}{1.728003in}}%
\pgfpathlineto{\pgfqpoint{2.968229in}{1.752536in}}%
\pgfpathlineto{\pgfqpoint{2.976338in}{1.769464in}}%
\pgfpathlineto{\pgfqpoint{2.983095in}{1.780539in}}%
\pgfpathlineto{\pgfqpoint{2.989853in}{1.788768in}}%
\pgfpathlineto{\pgfqpoint{2.995259in}{1.793260in}}%
\pgfpathlineto{\pgfqpoint{3.000665in}{1.795868in}}%
\pgfpathlineto{\pgfqpoint{3.006071in}{1.796581in}}%
\pgfpathlineto{\pgfqpoint{3.011477in}{1.795394in}}%
\pgfpathlineto{\pgfqpoint{3.016883in}{1.792313in}}%
\pgfpathlineto{\pgfqpoint{3.022289in}{1.787354in}}%
\pgfpathlineto{\pgfqpoint{3.027695in}{1.780539in}}%
\pgfpathlineto{\pgfqpoint{3.034452in}{1.769464in}}%
\pgfpathlineto{\pgfqpoint{3.041210in}{1.755626in}}%
\pgfpathlineto{\pgfqpoint{3.049319in}{1.735520in}}%
\pgfpathlineto{\pgfqpoint{3.058779in}{1.707512in}}%
\pgfpathlineto{\pgfqpoint{3.069591in}{1.670004in}}%
\pgfpathlineto{\pgfqpoint{3.081755in}{1.621639in}}%
\pgfpathlineto{\pgfqpoint{3.095270in}{1.561562in}}%
\pgfpathlineto{\pgfqpoint{3.112839in}{1.476181in}}%
\pgfpathlineto{\pgfqpoint{3.147978in}{1.295357in}}%
\pgfpathlineto{\pgfqpoint{3.170954in}{1.181359in}}%
\pgfpathlineto{\pgfqpoint{3.188523in}{1.101640in}}%
\pgfpathlineto{\pgfqpoint{3.203390in}{1.041111in}}%
\pgfpathlineto{\pgfqpoint{3.216905in}{0.992404in}}%
\pgfpathlineto{\pgfqpoint{3.229068in}{0.954019in}}%
\pgfpathlineto{\pgfqpoint{3.239880in}{0.924279in}}%
\pgfpathlineto{\pgfqpoint{3.250692in}{0.898567in}}%
\pgfpathlineto{\pgfqpoint{3.261504in}{0.876685in}}%
\pgfpathlineto{\pgfqpoint{3.272316in}{0.858341in}}%
\pgfpathlineto{\pgfqpoint{3.283128in}{0.843166in}}%
\pgfpathlineto{\pgfqpoint{3.293940in}{0.830736in}}%
\pgfpathlineto{\pgfqpoint{3.304752in}{0.820585in}}%
\pgfpathlineto{\pgfqpoint{3.316916in}{0.811292in}}%
\pgfpathlineto{\pgfqpoint{3.331782in}{0.802043in}}%
\pgfpathlineto{\pgfqpoint{3.356109in}{0.789162in}}%
\pgfpathlineto{\pgfqpoint{3.381788in}{0.775085in}}%
\pgfpathlineto{\pgfqpoint{3.402060in}{0.762190in}}%
\pgfpathlineto{\pgfqpoint{3.426388in}{0.744794in}}%
\pgfpathlineto{\pgfqpoint{3.461527in}{0.719696in}}%
\pgfpathlineto{\pgfqpoint{3.477745in}{0.709988in}}%
\pgfpathlineto{\pgfqpoint{3.492611in}{0.702867in}}%
\pgfpathlineto{\pgfqpoint{3.506126in}{0.698031in}}%
\pgfpathlineto{\pgfqpoint{3.520993in}{0.694411in}}%
\pgfpathlineto{\pgfqpoint{3.538562in}{0.691869in}}%
\pgfpathlineto{\pgfqpoint{3.581810in}{0.686574in}}%
\pgfpathlineto{\pgfqpoint{3.598028in}{0.682646in}}%
\pgfpathlineto{\pgfqpoint{3.623707in}{0.674534in}}%
\pgfpathlineto{\pgfqpoint{3.639925in}{0.670182in}}%
\pgfpathlineto{\pgfqpoint{3.652088in}{0.668364in}}%
\pgfpathlineto{\pgfqpoint{3.680470in}{0.665495in}}%
\pgfpathlineto{\pgfqpoint{3.680470in}{0.665495in}}%
\pgfusepath{stroke}%
\end{pgfscope}%
\begin{pgfscope}%
\definecolor{textcolor}{rgb}{0.000000,0.000000,0.000000}%
\pgfsetstrokecolor{textcolor}%
\pgfsetfillcolor{textcolor}%
\pgftext[x=3.005395in,y=2.207002in,,base]{\color{textcolor}{\sffamily\fontsize{8.000000}{9.600000}\selectfont\catcode`\^=\active\def^{\ifmmode\sp\else\^{}\fi}\catcode`\%=\active\def%{\%}(b) Chebyshev nodes}}%
\end{pgfscope}%
\begin{pgfscope}%
\pgfpathrectangle{\pgfqpoint{2.330321in}{0.352669in}}{\pgfqpoint{1.350149in}{1.771000in}}%
\pgfusepath{clip}%
\pgfsetbuttcap%
\pgfsetroundjoin%
\definecolor{currentfill}{rgb}{0.250980,0.400000,0.600000}%
\pgfsetfillcolor{currentfill}%
\pgfsetlinewidth{1.003750pt}%
\definecolor{currentstroke}{rgb}{0.250980,0.400000,0.600000}%
\pgfsetstrokecolor{currentstroke}%
\pgfsetdash{}{0pt}%
\pgfsys@defobject{currentmarker}{\pgfqpoint{-0.013889in}{-0.013889in}}{\pgfqpoint{0.013889in}{0.013889in}}{%
\pgfpathmoveto{\pgfqpoint{0.000000in}{-0.013889in}}%
\pgfpathcurveto{\pgfqpoint{0.003683in}{-0.013889in}}{\pgfqpoint{0.007216in}{-0.012425in}}{\pgfqpoint{0.009821in}{-0.009821in}}%
\pgfpathcurveto{\pgfqpoint{0.012425in}{-0.007216in}}{\pgfqpoint{0.013889in}{-0.003683in}}{\pgfqpoint{0.013889in}{0.000000in}}%
\pgfpathcurveto{\pgfqpoint{0.013889in}{0.003683in}}{\pgfqpoint{0.012425in}{0.007216in}}{\pgfqpoint{0.009821in}{0.009821in}}%
\pgfpathcurveto{\pgfqpoint{0.007216in}{0.012425in}}{\pgfqpoint{0.003683in}{0.013889in}}{\pgfqpoint{0.000000in}{0.013889in}}%
\pgfpathcurveto{\pgfqpoint{-0.003683in}{0.013889in}}{\pgfqpoint{-0.007216in}{0.012425in}}{\pgfqpoint{-0.009821in}{0.009821in}}%
\pgfpathcurveto{\pgfqpoint{-0.012425in}{0.007216in}}{\pgfqpoint{-0.013889in}{0.003683in}}{\pgfqpoint{-0.013889in}{0.000000in}}%
\pgfpathcurveto{\pgfqpoint{-0.013889in}{-0.003683in}}{\pgfqpoint{-0.012425in}{-0.007216in}}{\pgfqpoint{-0.009821in}{-0.009821in}}%
\pgfpathcurveto{\pgfqpoint{-0.007216in}{-0.012425in}}{\pgfqpoint{-0.003683in}{-0.013889in}}{\pgfqpoint{0.000000in}{-0.013889in}}%
\pgfpathlineto{\pgfqpoint{0.000000in}{-0.013889in}}%
\pgfpathclose%
\pgfusepath{stroke,fill}%
}%
\begin{pgfscope}%
\pgfsys@transformshift{3.680470in}{0.665495in}%
\pgfsys@useobject{currentmarker}{}%
\end{pgfscope}%
\begin{pgfscope}%
\pgfsys@transformshift{3.551542in}{0.690595in}%
\pgfsys@useobject{currentmarker}{}%
\end{pgfscope}%
\begin{pgfscope}%
\pgfsys@transformshift{3.214005in}{1.002326in}%
\pgfsys@useobject{currentmarker}{}%
\end{pgfscope}%
\begin{pgfscope}%
\pgfsys@transformshift{2.796786in}{1.002326in}%
\pgfsys@useobject{currentmarker}{}%
\end{pgfscope}%
\begin{pgfscope}%
\pgfsys@transformshift{2.459248in}{0.690595in}%
\pgfsys@useobject{currentmarker}{}%
\end{pgfscope}%
\begin{pgfscope}%
\pgfsys@transformshift{2.330321in}{0.665495in}%
\pgfsys@useobject{currentmarker}{}%
\end{pgfscope}%
\end{pgfscope}%
\begin{pgfscope}%
\pgfpathrectangle{\pgfqpoint{2.330321in}{0.352669in}}{\pgfqpoint{1.350149in}{1.771000in}}%
\pgfusepath{clip}%
\pgfsetbuttcap%
\pgfsetroundjoin%
\definecolor{currentfill}{rgb}{0.768627,0.556863,0.211765}%
\pgfsetfillcolor{currentfill}%
\pgfsetlinewidth{1.003750pt}%
\definecolor{currentstroke}{rgb}{0.768627,0.556863,0.211765}%
\pgfsetstrokecolor{currentstroke}%
\pgfsetdash{}{0pt}%
\pgfsys@defobject{currentmarker}{\pgfqpoint{-0.013889in}{-0.013889in}}{\pgfqpoint{0.013889in}{0.013889in}}{%
\pgfpathmoveto{\pgfqpoint{0.000000in}{-0.013889in}}%
\pgfpathcurveto{\pgfqpoint{0.003683in}{-0.013889in}}{\pgfqpoint{0.007216in}{-0.012425in}}{\pgfqpoint{0.009821in}{-0.009821in}}%
\pgfpathcurveto{\pgfqpoint{0.012425in}{-0.007216in}}{\pgfqpoint{0.013889in}{-0.003683in}}{\pgfqpoint{0.013889in}{0.000000in}}%
\pgfpathcurveto{\pgfqpoint{0.013889in}{0.003683in}}{\pgfqpoint{0.012425in}{0.007216in}}{\pgfqpoint{0.009821in}{0.009821in}}%
\pgfpathcurveto{\pgfqpoint{0.007216in}{0.012425in}}{\pgfqpoint{0.003683in}{0.013889in}}{\pgfqpoint{0.000000in}{0.013889in}}%
\pgfpathcurveto{\pgfqpoint{-0.003683in}{0.013889in}}{\pgfqpoint{-0.007216in}{0.012425in}}{\pgfqpoint{-0.009821in}{0.009821in}}%
\pgfpathcurveto{\pgfqpoint{-0.012425in}{0.007216in}}{\pgfqpoint{-0.013889in}{0.003683in}}{\pgfqpoint{-0.013889in}{0.000000in}}%
\pgfpathcurveto{\pgfqpoint{-0.013889in}{-0.003683in}}{\pgfqpoint{-0.012425in}{-0.007216in}}{\pgfqpoint{-0.009821in}{-0.009821in}}%
\pgfpathcurveto{\pgfqpoint{-0.007216in}{-0.012425in}}{\pgfqpoint{-0.003683in}{-0.013889in}}{\pgfqpoint{0.000000in}{-0.013889in}}%
\pgfpathlineto{\pgfqpoint{0.000000in}{-0.013889in}}%
\pgfpathclose%
\pgfusepath{stroke,fill}%
}%
\begin{pgfscope}%
\pgfsys@transformshift{3.680470in}{0.665495in}%
\pgfsys@useobject{currentmarker}{}%
\end{pgfscope}%
\begin{pgfscope}%
\pgfsys@transformshift{3.647429in}{0.670596in}%
\pgfsys@useobject{currentmarker}{}%
\end{pgfscope}%
\begin{pgfscope}%
\pgfsys@transformshift{3.551542in}{0.690595in}%
\pgfsys@useobject{currentmarker}{}%
\end{pgfscope}%
\begin{pgfscope}%
\pgfsys@transformshift{3.402194in}{0.751162in}%
\pgfsys@useobject{currentmarker}{}%
\end{pgfscope}%
\begin{pgfscope}%
\pgfsys@transformshift{3.214005in}{1.002326in}%
\pgfsys@useobject{currentmarker}{}%
\end{pgfscope}%
\begin{pgfscope}%
\pgfsys@transformshift{3.005395in}{1.926891in}%
\pgfsys@useobject{currentmarker}{}%
\end{pgfscope}%
\begin{pgfscope}%
\pgfsys@transformshift{2.796786in}{1.002326in}%
\pgfsys@useobject{currentmarker}{}%
\end{pgfscope}%
\begin{pgfscope}%
\pgfsys@transformshift{2.608596in}{0.751162in}%
\pgfsys@useobject{currentmarker}{}%
\end{pgfscope}%
\begin{pgfscope}%
\pgfsys@transformshift{2.459248in}{0.690595in}%
\pgfsys@useobject{currentmarker}{}%
\end{pgfscope}%
\begin{pgfscope}%
\pgfsys@transformshift{2.363361in}{0.670596in}%
\pgfsys@useobject{currentmarker}{}%
\end{pgfscope}%
\begin{pgfscope}%
\pgfsys@transformshift{2.330321in}{0.665495in}%
\pgfsys@useobject{currentmarker}{}%
\end{pgfscope}%
\end{pgfscope}%
\begin{pgfscope}%
\pgfpathrectangle{\pgfqpoint{2.330321in}{0.352669in}}{\pgfqpoint{1.350149in}{1.771000in}}%
\pgfusepath{clip}%
\pgfsetbuttcap%
\pgfsetroundjoin%
\definecolor{currentfill}{rgb}{0.619608,0.188235,0.172549}%
\pgfsetfillcolor{currentfill}%
\pgfsetlinewidth{1.003750pt}%
\definecolor{currentstroke}{rgb}{0.619608,0.188235,0.172549}%
\pgfsetstrokecolor{currentstroke}%
\pgfsetdash{}{0pt}%
\pgfsys@defobject{currentmarker}{\pgfqpoint{-0.013889in}{-0.013889in}}{\pgfqpoint{0.013889in}{0.013889in}}{%
\pgfpathmoveto{\pgfqpoint{0.000000in}{-0.013889in}}%
\pgfpathcurveto{\pgfqpoint{0.003683in}{-0.013889in}}{\pgfqpoint{0.007216in}{-0.012425in}}{\pgfqpoint{0.009821in}{-0.009821in}}%
\pgfpathcurveto{\pgfqpoint{0.012425in}{-0.007216in}}{\pgfqpoint{0.013889in}{-0.003683in}}{\pgfqpoint{0.013889in}{0.000000in}}%
\pgfpathcurveto{\pgfqpoint{0.013889in}{0.003683in}}{\pgfqpoint{0.012425in}{0.007216in}}{\pgfqpoint{0.009821in}{0.009821in}}%
\pgfpathcurveto{\pgfqpoint{0.007216in}{0.012425in}}{\pgfqpoint{0.003683in}{0.013889in}}{\pgfqpoint{0.000000in}{0.013889in}}%
\pgfpathcurveto{\pgfqpoint{-0.003683in}{0.013889in}}{\pgfqpoint{-0.007216in}{0.012425in}}{\pgfqpoint{-0.009821in}{0.009821in}}%
\pgfpathcurveto{\pgfqpoint{-0.012425in}{0.007216in}}{\pgfqpoint{-0.013889in}{0.003683in}}{\pgfqpoint{-0.013889in}{0.000000in}}%
\pgfpathcurveto{\pgfqpoint{-0.013889in}{-0.003683in}}{\pgfqpoint{-0.012425in}{-0.007216in}}{\pgfqpoint{-0.009821in}{-0.009821in}}%
\pgfpathcurveto{\pgfqpoint{-0.007216in}{-0.012425in}}{\pgfqpoint{-0.003683in}{-0.013889in}}{\pgfqpoint{0.000000in}{-0.013889in}}%
\pgfpathlineto{\pgfqpoint{0.000000in}{-0.013889in}}%
\pgfpathclose%
\pgfusepath{stroke,fill}%
}%
\begin{pgfscope}%
\pgfsys@transformshift{3.680470in}{0.665495in}%
\pgfsys@useobject{currentmarker}{}%
\end{pgfscope}%
\begin{pgfscope}%
\pgfsys@transformshift{3.665718in}{0.667683in}%
\pgfsys@useobject{currentmarker}{}%
\end{pgfscope}%
\begin{pgfscope}%
\pgfsys@transformshift{3.622106in}{0.675039in}%
\pgfsys@useobject{currentmarker}{}%
\end{pgfscope}%
\begin{pgfscope}%
\pgfsys@transformshift{3.551542in}{0.690595in}%
\pgfsys@useobject{currentmarker}{}%
\end{pgfscope}%
\begin{pgfscope}%
\pgfsys@transformshift{3.457108in}{0.722626in}%
\pgfsys@useobject{currentmarker}{}%
\end{pgfscope}%
\begin{pgfscope}%
\pgfsys@transformshift{3.342932in}{0.795984in}%
\pgfsys@useobject{currentmarker}{}%
\end{pgfscope}%
\begin{pgfscope}%
\pgfsys@transformshift{3.214005in}{1.002326in}%
\pgfsys@useobject{currentmarker}{}%
\end{pgfscope}%
\begin{pgfscope}%
\pgfsys@transformshift{3.075960in}{1.645434in}%
\pgfsys@useobject{currentmarker}{}%
\end{pgfscope}%
\begin{pgfscope}%
\pgfsys@transformshift{2.934831in}{1.645434in}%
\pgfsys@useobject{currentmarker}{}%
\end{pgfscope}%
\begin{pgfscope}%
\pgfsys@transformshift{2.796786in}{1.002326in}%
\pgfsys@useobject{currentmarker}{}%
\end{pgfscope}%
\begin{pgfscope}%
\pgfsys@transformshift{2.667858in}{0.795984in}%
\pgfsys@useobject{currentmarker}{}%
\end{pgfscope}%
\begin{pgfscope}%
\pgfsys@transformshift{2.553682in}{0.722626in}%
\pgfsys@useobject{currentmarker}{}%
\end{pgfscope}%
\begin{pgfscope}%
\pgfsys@transformshift{2.459248in}{0.690595in}%
\pgfsys@useobject{currentmarker}{}%
\end{pgfscope}%
\begin{pgfscope}%
\pgfsys@transformshift{2.388684in}{0.675039in}%
\pgfsys@useobject{currentmarker}{}%
\end{pgfscope}%
\begin{pgfscope}%
\pgfsys@transformshift{2.345073in}{0.667683in}%
\pgfsys@useobject{currentmarker}{}%
\end{pgfscope}%
\begin{pgfscope}%
\pgfsys@transformshift{2.330321in}{0.665495in}%
\pgfsys@useobject{currentmarker}{}%
\end{pgfscope}%
\end{pgfscope}%
\begin{pgfscope}%
\pgfpathrectangle{\pgfqpoint{2.330321in}{0.352669in}}{\pgfqpoint{1.350149in}{1.771000in}}%
\pgfusepath{clip}%
\pgfsetrectcap%
\pgfsetroundjoin%
\pgfsetlinewidth{1.405250pt}%
\definecolor{currentstroke}{rgb}{0.000000,0.000000,0.000000}%
\pgfsetstrokecolor{currentstroke}%
\pgfsetdash{}{0pt}%
\pgfpathmoveto{\pgfqpoint{2.330321in}{0.665495in}}%
\pgfpathlineto{\pgfqpoint{2.377623in}{0.673036in}}%
\pgfpathlineto{\pgfqpoint{2.419520in}{0.681194in}}%
\pgfpathlineto{\pgfqpoint{2.456010in}{0.689757in}}%
\pgfpathlineto{\pgfqpoint{2.488446in}{0.698810in}}%
\pgfpathlineto{\pgfqpoint{2.518179in}{0.708595in}}%
\pgfpathlineto{\pgfqpoint{2.545209in}{0.719012in}}%
\pgfpathlineto{\pgfqpoint{2.569536in}{0.729898in}}%
\pgfpathlineto{\pgfqpoint{2.592512in}{0.741767in}}%
\pgfpathlineto{\pgfqpoint{2.614136in}{0.754630in}}%
\pgfpathlineto{\pgfqpoint{2.634408in}{0.768470in}}%
\pgfpathlineto{\pgfqpoint{2.653329in}{0.783234in}}%
\pgfpathlineto{\pgfqpoint{2.670899in}{0.798826in}}%
\pgfpathlineto{\pgfqpoint{2.687117in}{0.815104in}}%
\pgfpathlineto{\pgfqpoint{2.703335in}{0.833491in}}%
\pgfpathlineto{\pgfqpoint{2.718201in}{0.852493in}}%
\pgfpathlineto{\pgfqpoint{2.733068in}{0.873871in}}%
\pgfpathlineto{\pgfqpoint{2.747934in}{0.897992in}}%
\pgfpathlineto{\pgfqpoint{2.761449in}{0.922657in}}%
\pgfpathlineto{\pgfqpoint{2.774964in}{0.950307in}}%
\pgfpathlineto{\pgfqpoint{2.788479in}{0.981357in}}%
\pgfpathlineto{\pgfqpoint{2.801994in}{1.016271in}}%
\pgfpathlineto{\pgfqpoint{2.815509in}{1.055557in}}%
\pgfpathlineto{\pgfqpoint{2.829024in}{1.099755in}}%
\pgfpathlineto{\pgfqpoint{2.842539in}{1.149413in}}%
\pgfpathlineto{\pgfqpoint{2.856054in}{1.205041in}}%
\pgfpathlineto{\pgfqpoint{2.869569in}{1.267037in}}%
\pgfpathlineto{\pgfqpoint{2.884436in}{1.342783in}}%
\pgfpathlineto{\pgfqpoint{2.900654in}{1.434011in}}%
\pgfpathlineto{\pgfqpoint{2.920926in}{1.557863in}}%
\pgfpathlineto{\pgfqpoint{2.953362in}{1.757248in}}%
\pgfpathlineto{\pgfqpoint{2.965526in}{1.821673in}}%
\pgfpathlineto{\pgfqpoint{2.974986in}{1.863558in}}%
\pgfpathlineto{\pgfqpoint{2.981744in}{1.887834in}}%
\pgfpathlineto{\pgfqpoint{2.988501in}{1.906669in}}%
\pgfpathlineto{\pgfqpoint{2.993907in}{1.917462in}}%
\pgfpathlineto{\pgfqpoint{2.997962in}{1.922927in}}%
\pgfpathlineto{\pgfqpoint{3.002016in}{1.926070in}}%
\pgfpathlineto{\pgfqpoint{3.006071in}{1.926858in}}%
\pgfpathlineto{\pgfqpoint{3.008774in}{1.926070in}}%
\pgfpathlineto{\pgfqpoint{3.012828in}{1.922927in}}%
\pgfpathlineto{\pgfqpoint{3.016883in}{1.917462in}}%
\pgfpathlineto{\pgfqpoint{3.020937in}{1.909734in}}%
\pgfpathlineto{\pgfqpoint{3.026343in}{1.896053in}}%
\pgfpathlineto{\pgfqpoint{3.033101in}{1.873882in}}%
\pgfpathlineto{\pgfqpoint{3.041210in}{1.840650in}}%
\pgfpathlineto{\pgfqpoint{3.050670in}{1.794285in}}%
\pgfpathlineto{\pgfqpoint{3.062834in}{1.725849in}}%
\pgfpathlineto{\pgfqpoint{3.081755in}{1.608968in}}%
\pgfpathlineto{\pgfqpoint{3.110136in}{1.434011in}}%
\pgfpathlineto{\pgfqpoint{3.126354in}{1.342783in}}%
\pgfpathlineto{\pgfqpoint{3.141221in}{1.267037in}}%
\pgfpathlineto{\pgfqpoint{3.156087in}{1.199197in}}%
\pgfpathlineto{\pgfqpoint{3.169602in}{1.144186in}}%
\pgfpathlineto{\pgfqpoint{3.183117in}{1.095098in}}%
\pgfpathlineto{\pgfqpoint{3.196632in}{1.051416in}}%
\pgfpathlineto{\pgfqpoint{3.210147in}{1.012591in}}%
\pgfpathlineto{\pgfqpoint{3.223662in}{0.978086in}}%
\pgfpathlineto{\pgfqpoint{3.237177in}{0.947396in}}%
\pgfpathlineto{\pgfqpoint{3.250692in}{0.920063in}}%
\pgfpathlineto{\pgfqpoint{3.265559in}{0.893384in}}%
\pgfpathlineto{\pgfqpoint{3.280425in}{0.869791in}}%
\pgfpathlineto{\pgfqpoint{3.295292in}{0.848871in}}%
\pgfpathlineto{\pgfqpoint{3.311510in}{0.828678in}}%
\pgfpathlineto{\pgfqpoint{3.327728in}{0.810849in}}%
\pgfpathlineto{\pgfqpoint{3.345297in}{0.793817in}}%
\pgfpathlineto{\pgfqpoint{3.362867in}{0.778815in}}%
\pgfpathlineto{\pgfqpoint{3.381788in}{0.764593in}}%
\pgfpathlineto{\pgfqpoint{3.402060in}{0.751244in}}%
\pgfpathlineto{\pgfqpoint{3.423685in}{0.738820in}}%
\pgfpathlineto{\pgfqpoint{3.446660in}{0.727340in}}%
\pgfpathlineto{\pgfqpoint{3.470987in}{0.716798in}}%
\pgfpathlineto{\pgfqpoint{3.498017in}{0.706696in}}%
\pgfpathlineto{\pgfqpoint{3.527750in}{0.697193in}}%
\pgfpathlineto{\pgfqpoint{3.560186in}{0.688389in}}%
\pgfpathlineto{\pgfqpoint{3.596677in}{0.680050in}}%
\pgfpathlineto{\pgfqpoint{3.637222in}{0.672327in}}%
\pgfpathlineto{\pgfqpoint{3.680470in}{0.665495in}}%
\pgfpathlineto{\pgfqpoint{3.680470in}{0.665495in}}%
\pgfusepath{stroke}%
\end{pgfscope}%
\begin{pgfscope}%
\pgfsetrectcap%
\pgfsetroundjoin%
\pgfsetlinewidth{1.405250pt}%
\definecolor{currentstroke}{rgb}{0.000000,0.000000,0.000000}%
\pgfsetstrokecolor{currentstroke}%
\pgfsetdash{}{0pt}%
\pgfpathmoveto{\pgfqpoint{2.396987in}{2.016419in}}%
\pgfpathlineto{\pgfqpoint{2.480321in}{2.016419in}}%
\pgfpathlineto{\pgfqpoint{2.563654in}{2.016419in}}%
\pgfusepath{stroke}%
\end{pgfscope}%
\begin{pgfscope}%
\definecolor{textcolor}{rgb}{0.000000,0.000000,0.000000}%
\pgfsetstrokecolor{textcolor}%
\pgfsetfillcolor{textcolor}%
\pgftext[x=2.630321in,y=1.987252in,left,base]{\color{textcolor}{\sffamily\fontsize{6.000000}{7.200000}\selectfont\catcode`\^=\active\def^{\ifmmode\sp\else\^{}\fi}\catcode`\%=\active\def%{\%}$f(x)$}}%
\end{pgfscope}%
\begin{pgfscope}%
\pgfsetrectcap%
\pgfsetroundjoin%
\pgfsetlinewidth{0.903375pt}%
\definecolor{currentstroke}{rgb}{0.250980,0.400000,0.600000}%
\pgfsetstrokecolor{currentstroke}%
\pgfsetdash{}{0pt}%
\pgfpathmoveto{\pgfqpoint{2.396987in}{1.892022in}}%
\pgfpathlineto{\pgfqpoint{2.480321in}{1.892022in}}%
\pgfpathlineto{\pgfqpoint{2.563654in}{1.892022in}}%
\pgfusepath{stroke}%
\end{pgfscope}%
\begin{pgfscope}%
\definecolor{textcolor}{rgb}{0.000000,0.000000,0.000000}%
\pgfsetstrokecolor{textcolor}%
\pgfsetfillcolor{textcolor}%
\pgftext[x=2.630321in,y=1.862855in,left,base]{\color{textcolor}{\sffamily\fontsize{6.000000}{7.200000}\selectfont\catcode`\^=\active\def^{\ifmmode\sp\else\^{}\fi}\catcode`\%=\active\def%{\%}$N = 5$}}%
\end{pgfscope}%
\begin{pgfscope}%
\pgfsetrectcap%
\pgfsetroundjoin%
\pgfsetlinewidth{0.903375pt}%
\definecolor{currentstroke}{rgb}{0.768627,0.556863,0.211765}%
\pgfsetstrokecolor{currentstroke}%
\pgfsetdash{}{0pt}%
\pgfpathmoveto{\pgfqpoint{3.037533in}{2.016419in}}%
\pgfpathlineto{\pgfqpoint{3.120866in}{2.016419in}}%
\pgfpathlineto{\pgfqpoint{3.204199in}{2.016419in}}%
\pgfusepath{stroke}%
\end{pgfscope}%
\begin{pgfscope}%
\definecolor{textcolor}{rgb}{0.000000,0.000000,0.000000}%
\pgfsetstrokecolor{textcolor}%
\pgfsetfillcolor{textcolor}%
\pgftext[x=3.270866in,y=1.987252in,left,base]{\color{textcolor}{\sffamily\fontsize{6.000000}{7.200000}\selectfont\catcode`\^=\active\def^{\ifmmode\sp\else\^{}\fi}\catcode`\%=\active\def%{\%}$N = 10$}}%
\end{pgfscope}%
\begin{pgfscope}%
\pgfsetrectcap%
\pgfsetroundjoin%
\pgfsetlinewidth{0.903375pt}%
\definecolor{currentstroke}{rgb}{0.619608,0.188235,0.172549}%
\pgfsetstrokecolor{currentstroke}%
\pgfsetdash{}{0pt}%
\pgfpathmoveto{\pgfqpoint{3.037533in}{1.894104in}}%
\pgfpathlineto{\pgfqpoint{3.120866in}{1.894104in}}%
\pgfpathlineto{\pgfqpoint{3.204199in}{1.894104in}}%
\pgfusepath{stroke}%
\end{pgfscope}%
\begin{pgfscope}%
\definecolor{textcolor}{rgb}{0.000000,0.000000,0.000000}%
\pgfsetstrokecolor{textcolor}%
\pgfsetfillcolor{textcolor}%
\pgftext[x=3.270866in,y=1.864938in,left,base]{\color{textcolor}{\sffamily\fontsize{6.000000}{7.200000}\selectfont\catcode`\^=\active\def^{\ifmmode\sp\else\^{}\fi}\catcode`\%=\active\def%{\%}$N = 15$}}%
\end{pgfscope}%
\end{pgfpicture}%
\makeatother%
\endgroup%

\caption[Runge phenomenon on equidistant vs.\ Chebyshev nodes]{%
  Polynomial interpolation of $f(x) = 1/(1 + 25x^2)$ on $[-1,1]$.
  (a) Equidistant nodes: the interpolant converges near the centre
  but develops violent oscillations near the endpoints that grow
  with increasing polynomial degree --- the Runge phenomenon.
  (b) Chebyshev Gauss--Lobatto nodes (\cref{sec:spec-chebyshev}):
  the interpolant converges uniformly across the entire interval,
  with the error decreasing monotonically as $N$ increases.  The
  clustered node placement near the endpoints suppresses the node
  polynomial $\omega_{N+1}$, eliminating the instability.}
\label{fig:runge-phenomenon}
\end{figure}

\paragraph{The cure: clustered nodes.}

The Runge phenomenon is not a failure of polynomial interpolation
itself but of the node placement.  The fix is to cluster nodes toward
the endpoints of the interval, where the equidistant node polynomial
is largest.  The idea is to make $\abs{\omega_{N+1}(x)}$ as uniform as
possible across $[a, b]$, so that no region of the interval suffers
disproportionately large error.

The optimal node distribution on $[-1, 1]$ --- the one that minimises
$\max_{x} \abs{\omega_{N+1}(x)}$ --- turns out to be the zeros of the
$(N+1)$-th Chebyshev polynomial of the first kind:
%
\begin{equation}
  x_j = \cos\!\left(\frac{2j + 1}{2(N+1)}\,\pi\right) ,
  \qquad j = 0, 1, \ldots, N \,.
  \label{eq:chebyshev-zeros}
\end{equation}
%
These nodes cluster quadratically near $\pm 1$, with spacing
$\order{1/N^2}$ at the endpoints compared to $\order{1/N}$ in the
interior.  For this distribution, the node polynomial satisfies
$\max_{x \in [-1,1]} \abs{\omega_{N+1}(x)} = 2^{-N}$, the smallest
supremum norm on $[-1,1]$ achievable by any monic polynomial of
degree $N+1$ --- a classical result of Chebyshev \cite{Trefethen2000}.  The Lebesgue constant
grows only logarithmically, $\Lambda_N \sim (2/\pi)\ln N$, compared
to the exponential growth for equidistant nodes.

A closely related distribution --- the \emph{Chebyshev--Gauss--Lobatto}
(CGL) nodes --- includes the endpoints $\pm 1$:
%
\begin{equation}
  x_j = \cos\!\left(\frac{j}{N}\,\pi\right) ,
  \qquad j = 0, 1, \ldots, N \,.
  \label{eq:cgl-nodes}
\end{equation}
%
The CGL nodes are the extrema of the $N$-th Chebyshev polynomial
(which include the endpoints $\pm 1$).  They share the logarithmic Lebesgue
growth and the quadratic endpoint clustering of the Chebyshev zeros,
with the practical advantage that including $\pm 1$ simplifies
boundary conditions and inter-element coupling in a spectral-element
or DG context.
\xrefnote{The next section (\cref{sec:spec-chebyshev}) develops the
Chebyshev polynomial theory that explains why these particular nodes
are optimal.  \Cref{sec:spec-lgl} derives the Legendre--Gauss--Lobatto
(LGL) nodes used by the codebase's DG solver, which share the
essential clustering property.}

\paragraph{From interpolation to spectral methods.}

The Runge phenomenon teaches that polynomial approximation can
achieve arbitrarily high accuracy for smooth functions, but only if
the nodes are chosen correctly.  Spectral methods build on this
insight: they represent the numerical solution as a polynomial on
clustered nodes and use the polynomial structure to compute
derivatives, integrals, and inter-element fluxes.  The rest of this
chapter develops the machinery: Chebyshev and Legendre polynomial
bases that make clustered nodes natural, differentiation and
quadrature operators built directly from them, and the Vandermonde
matrices that convert between nodal and modal representations in the
DG solver.

% sec:spec-chebyshev — Chebyshev Polynomials and Gauss-Lobatto Nodes
\section{Chebyshev Polynomials and Gauss-Lobatto Nodes}
\label{sec:spec-chebyshev}

The previous section showed that clustered nodes cure the Runge
phenomenon, and that the Chebyshev zeros minimise the node polynomial
over all monic polynomials of a given degree.  But where do these
special nodes come from?  The answer lies in a single change of
variable: the map $x = \cos\theta$ transforms polynomial
approximation on $[-1, 1]$ into trigonometric approximation on
$[0, \pi]$, and the rich theory of Fourier series carries over
wholesale.  The Chebyshev polynomials that emerge from this
correspondence carry algebraic, orthogonality, and minimax properties
that explain why the Chebyshev--Gauss--Lobatto (CGL) nodes of
\cref{eq:cgl-nodes} --- the extrema of these polynomials --- are the
natural choice for spectral collocation.
\coderef{spectral}{rs}

\paragraph{The cosine definition.}

The $n$-th Chebyshev polynomial of the first kind is defined by the
trigonometric identity
%
\begin{equation}
  T_n(x) = \cos(n\,\arccos x) \,,
  \qquad x \in [-1, 1] \,.
  \label{eq:chebyshev-def}
\end{equation}
%
Setting $\theta = \arccos x$ reveals the essential idea: $T_n$ is
simply $\cos(n\theta)$ written in the variable
$x = \cos\theta$.  The first few are $T_0(x) = 1$,
$T_1(x) = x$, and $T_2(x) = 2x^2 - 1$ (the double-angle
formula).  Despite the transcendental appearance of the definition,
$T_n$ is a polynomial of degree exactly $n$ with integer
coefficients and leading coefficient $2^{n-1}$ for $n \geq 1$.
\histnote{Chebyshev introduced these polynomials in the 1850s in the
context of mechanism design --- specifically, approximating a
rectilinear linkage --- long before their role in numerical analysis
was recognised.}

\paragraph{Three-term recurrence.}

The trigonometric identity
$\cos((n+1)\theta) = 2\cos\theta\,\cos(n\theta) - \cos((n-1)\theta)$
translates directly into the polynomial recurrence
%
\begin{equation}
  T_{n+1}(x) = 2x\,T_n(x) - T_{n-1}(x) \,,
  \label{eq:chebyshev-recurrence}
\end{equation}
%
with initial conditions $T_0(x) = 1$, $T_1(x) = x$.  This
recurrence generates every Chebyshev polynomial from the two
simplest, and is the standard method for evaluating $T_n(x)$ at a
given point --- each step costs one multiply and one subtract, with no
trigonometric function calls.  The recurrence also reveals that the
leading coefficient of $T_n$ doubles at each step, confirming the
factor $2^{n-1}$ for $n \geq 1$.
\implnote{The codebase's \texttt{ChebyshevBasis} in
\codemodule{spectral} uses the cosine formula
(\cref{eq:chebyshev-def}) to place the CGL nodes directly, rather
than evaluating the recurrence.  The recurrence appears in the
Legendre polynomial evaluation (\cref{sec:spec-legendre}), where the
analogous three-term relation is the primary computational tool.}

\paragraph{Zeros and extrema.}

The zeros and extrema of $T_n$ are direct consequences of the
cosine definition.  Setting $T_n(x) = \cos(n\theta) = 0$ gives
$n\theta = (2k+1)\pi/2$, so the $n$ zeros are
%
\begin{equation}
  x_k = \cos\!\left(\frac{2k + 1}{2n}\,\pi\right) ,
  \qquad k = 0, 1, \ldots, n-1 \,.
  \label{eq:chebyshev-zeros-Tn}
\end{equation}
%
These are exactly the Chebyshev zeros of $T_n$ --- the same nodes
that appeared in \cref{eq:chebyshev-zeros} of the previous section
with $n = N + 1$.  The extrema of $T_n$ occur where
$\cos(n\theta) = \pm 1$, i.e.\ $n\theta = j\pi$, giving
the $n + 1$ points
%
\begin{equation}
  x_j = \cos\!\left(\frac{j}{n}\,\pi\right) ,
  \qquad j = 0, 1, \ldots, n \,,
  \label{eq:chebyshev-extrema}
\end{equation}
%
which are the CGL nodes of \cref{eq:cgl-nodes} with $n = N$.
At these points, $T_n(x_j) = \cos(j\pi) = (-1)^j$: the polynomial
alternates between $+1$ and $-1$ at its extrema, touching both
bounds at every one of $n + 1$ equally spaced points in
$\theta$-space.

\paragraph{The minimax property.}

The alternation $T_n(x_j) = (-1)^j$ at the extrema is the key to
the minimax optimality claimed in the previous section.  Among all monic
polynomials of degree $n$ on $[-1, 1]$, the scaled Chebyshev
polynomial $\tilde{T}_n(x) = T_n(x) / 2^{n-1}$ has the smallest
possible supremum norm:
%
\begin{equation}
  \max_{x \in [-1,1]} \abs{\tilde{T}_n(x)}
    = \frac{1}{2^{n-1}} \,.
  \label{eq:chebyshev-minimax}
\end{equation}
%
Any other monic polynomial of degree $n$ must have a larger
supremum norm.  The proof is a short contradiction argument using
the equioscillation theorem: suppose a monic polynomial $q_n$
satisfies $\norm{q_n}_\infty < 2^{-(n-1)}$.  Then
$\tilde{T}_n - q_n$ is a polynomial of degree at most $n - 1$
(the leading terms cancel), but
$\tilde{T}_n(x_j) = (-1)^j / 2^{n-1}$ has magnitude exceeding
$\norm{q_n}_\infty$ at each of the $n + 1$ extremal points.
Since $\abs{\tilde{T}_n(x_j)} > \norm{q_n}_\infty \geq \abs{q_n(x_j)}$,
the sign of $\tilde{T}_n(x_j) - q_n(x_j)$ equals the sign of
$\tilde{T}_n(x_j) = (-1)^j$.  The difference $\tilde{T}_n - q_n$
therefore changes sign between each consecutive pair of extremal
points, giving at least $n$ zeros --- impossible for a polynomial of
degree $n - 1$.

The minimax property explains why the Chebyshev zeros produce the
smallest node polynomial: the monic polynomial
$\omega_{N+1}(x) = \prod_{j=0}^{N}(x - x_j)$
of \cref{eq:node-polynomial} is minimised in supremum norm precisely
when it equals $\tilde{T}_{N+1}$, i.e.\ when the nodes are the
zeros of $T_{N+1}$.
\physnote{The equioscillation property is also the reason Chebyshev
nodes suppress the Runge phenomenon: the node polynomial oscillates
\emph{uniformly} across $[-1,1]$, so no region of the interval
suffers disproportionately large interpolation error.  Equidistant
nodes lack this balance --- the node polynomial is orders of magnitude
larger near the endpoints.}

\paragraph{Orthogonality.}

The Chebyshev polynomials are orthogonal with respect to the
weighted inner product
%
\begin{equation}
  \int_{-1}^{1} T_m(x)\,T_n(x)\,
    \frac{\d x}{\sqrt{1 - x^2}}
  = \begin{cases}
      0 & m \neq n \\
      \pi & m = n = 0 \\
      \pi/2 & m = n \geq 1
    \end{cases}
  \label{eq:chebyshev-ortho}
\end{equation}
%
This follows directly from the substitution $x = \cos\theta$: the
weight function $(1 - x^2)^{-1/2}$ cancels the Jacobian of the
change of variable, and the integral reduces to
$\int_0^\pi \cos(m\theta)\,\cos(n\theta)\,\d\theta$, which is a
standard Fourier orthogonality integral.  The weight
$(1 - x^2)^{-1/2}$ diverges at $x = \pm 1$, reflecting the
clustering of the $\theta$-grid near the endpoints of $[-1,1]$
--- the same clustering that suppresses the Runge phenomenon.
\warnnote{The Chebyshev weight $(1-x^2)^{-1/2}$ differs from the
Legendre weight, which is simply $1$.  The two families are
orthogonal with respect to different inner products, and the choice
of weight function determines which nodes (CGL vs LGL) are natural
for a given application.  \Cref{sec:spec-legendre} develops the
Legendre case.}

\paragraph{Chebyshev expansion.}

Because the Chebyshev polynomials are orthogonal
(\cref{eq:chebyshev-ortho}), any function
$f \in L^2_w[-1,1]$ (square-integrable with respect to
the Chebyshev weight) in a Chebyshev series, converging in the
weighted $L^2$ norm:
%
\begin{equation}
  f(x) = \sum_{n=0}^{\infty} a_n\,T_n(x) \,,
  \qquad
  a_n = \frac{2}{\pi\,c_n}
    \int_{-1}^{1} f(x)\,T_n(x)\,
    \frac{\d x}{\sqrt{1 - x^2}} \,,
  \label{eq:chebyshev-expansion}
\end{equation}
%
where $c_0 = 2$ and $c_n = 1$ for $n \geq 1$.  The coefficients
$a_n$ are Fourier cosine coefficients in disguise: under
$x = \cos\theta$, the expansion becomes
$f(\cos\theta) = \sum a_n \cos(n\theta)$, and the $a_n$ are
computed by the standard cosine transform.  For smooth functions,
the coefficients $a_n$ decay faster than any power of $1/n$ ---
this is the spectral convergence that gives spectral methods their
name, developed quantitatively in \cref{sec:spec-convergence}.

\paragraph{Discrete orthogonality and aliasing.}

Continuous Chebyshev expansions are elegant but impractical: the
integral in \cref{eq:chebyshev-expansion} is rarely available in
closed form.  The CGL nodes provide an exact quadrature rule for
polynomials: for any polynomials $T_m$ and $T_n$ with
$m + n \leq 2N - 1$, the Gauss--Lobatto quadrature on the $N + 1$
CGL nodes reproduces the continuous orthogonality integral exactly.
This discrete orthogonality means that the Chebyshev expansion of a
polynomial of degree $N$ can be computed exactly from its values at
the CGL nodes, using the discrete cosine transform (DCT).

When the function is not a polynomial --- or is a polynomial of degree
higher than $N$ --- the discrete expansion aliases high modes onto low
ones, just as sampling a continuous signal at a finite rate folds
frequencies above the Nyquist limit.  Concretely, the discrete
coefficient $\hat{a}_k$ computed from $N + 1$ samples equals the true
$a_k$ plus contributions from $a_{2N-k}$, $a_{2N+k}$, and so on.
For smooth functions these aliased contributions are negligible
because the high-mode coefficients decay rapidly; for non-smooth
functions, aliasing introduces persistent errors that limit
convergence to algebraic rates.
\xrefnote{The spectral convergence theorem
(\cref{sec:spec-convergence}) quantifies the decay rate of Chebyshev
coefficients as a function of the smoothness of $f$.}

\paragraph{The codebase implementation.}

The \texttt{ChebyshevBasis} struct in \codemodule{spectral}
encapsulates the CGL nodes and the associated differentiation
matrices for $n = N + 1$ collocation points (the struct field
\texttt{n} counts points, not polynomial degree):
%
\begin{lstlisting}[language=Rust]
pub struct ChebyshevBasis {
    pub n: usize,
    pub nodes: Vec<f64>,   // CGL: cos(pi*k/(n-1))
    pub d1: Vec<f64>,      // first derivative matrix
    pub d2: Vec<f64>,      // second derivative matrix
}
\end{lstlisting}
%
The nodes are computed directly from the cosine formula
$x_k = \cos(\pi k / (n-1))$ for $k = 0, \ldots, n-1$, where
$n = N + 1$ is the number of collocation points.  The
differentiation matrices \texttt{d1} and \texttt{d2} --- constructed
using the barycentric formula of Trefethen \cite{Trefethen2000} ---
are the subject of \cref{sec:spec-diffmat}.  The XCTS initial-data
solver uses \texttt{ChebyshevBasis} for elliptic problems on
multi-domain radial grids: each radial subdomain carries its own
set of CGL nodes, and the exponential convergence established above
means that modest polynomial degrees ($N \sim 20$--$30$) suffice to
resolve the metric fields to machine precision.  The Chebyshev
framework thus converts the abstract minimax and orthogonality
results of this section into a concrete computational tool ---
one that underpins every spectral solve in the codebase.

% sec:spec-diffmat — Differentiation Matrices
\section{Differentiation Matrices}
\label{sec:spec-diffmat}

Once a function is represented by its values at collocation nodes,
differentiation reduces to a single matrix-vector multiply --- no
symbolic calculus, no finite-difference stencils, just linear
algebra.  A spectral solver performs this multiply thousands of
times per time step.  Given the function values
$f(x_0), \ldots, f(x_N)$ at the
collocation nodes, we need the derivative values
$f'(x_0), \ldots, f'(x_N)$ at the same nodes.  Because the
interpolating polynomial $p_N$ of \cref{eq:lagrange-interp} is a
linear combination of the function values, its derivative at any node
is also a linear combination --- the map from values to derivatives is
a matrix-vector product.  The matrix encoding this map is the
\emph{spectral differentiation matrix}, and constructing it
efficiently and accurately is the subject of this section.
\coderef{spectral}{rs}

\paragraph{Differentiating the Lagrange basis.}

The interpolant $p_N(x) = \sum_{j=0}^{N} f(x_j)\,\ell_j(x)$ from
\cref{eq:lagrange-interp} is built from the Lagrange basis polynomials
$\ell_j$ of \cref{eq:lagrange-basis}.  Differentiating and evaluating
at a node $x_i$ gives
%
\begin{equation}
  p_N'(x_i) = \sum_{j=0}^{N} f(x_j)\,\ell_j'(x_i) \,.
  \label{eq:diffmat-deriv}
\end{equation}
%
This is a matrix-vector product $\mathbf{f}' = D\,\mathbf{f}$, where
the vectors $\mathbf{f}$ and $\mathbf{f}'$ collect the function values
and derivative values at the nodes, and the differentiation matrix $D$
has entries
%
\begin{equation}
  D_{ij} = \ell_j'(x_i) \,.
  \label{eq:diffmat-def}
\end{equation}
%
Each row of $D$ encodes the differentiation stencil at a single node,
and each column records how the function value at $x_j$ influences
every derivative.  For $N + 1$ nodes, $D$ is an
$(N+1) \times (N+1)$ dense matrix --- every node contributes to
every derivative, reflecting the global nature of polynomial
interpolation.
\physnote{The density of $D$ is the spectral analogue of a wide
finite-difference stencil.  A second-order centred stencil uses 3
points; a spectral stencil uses all $N + 1$.  The extra information is
what buys exponential convergence.}

\paragraph{Explicit formula for CGL nodes.}

Differentiating the Lagrange basis product
(\cref{eq:lagrange-basis}) directly is possible but numerically
fragile: each $\ell_j'(x_i)$ involves a sum of $N$ terms, each a
product of $N - 1$ factors, with catastrophic cancellation at high
$N$.  A stable alternative uses the \emph{barycentric weights}
$w_j$ defined by
%
\begin{equation}
  w_j = \frac{1}{\prod_{\substack{k=0 \\ k \neq j}}^{N}
    (x_j - x_k)} \,,
  \label{eq:bary-weights}
\end{equation}
%
which encode the node geometry in a single factor per node.
Differentiating the barycentric interpolation formula and
evaluating at $x_i$ produces the off-diagonal entries of $D$:
%
\begin{equation}
  D_{ij} = \frac{w_j / w_i}{x_i - x_j} \,, \qquad i \neq j \,.
  \label{eq:diffmat-offdiag}
\end{equation}
%
Each entry is a ratio of two barycentric weights divided by a
node spacing --- a single division and a single multiplication,
with no sums over $N$ terms and no catastrophic cancellation.
For CGL nodes $x_k = \cos(k\pi/N)$, the barycentric weights take the
closed form $w_k = (-1)^k / c_k$, where $c_0 = c_N = 2$ and
$c_k = 1$ for $1 \leq k \leq N - 1$.  Substituting into
\cref{eq:diffmat-offdiag} gives the explicit entry
%
\begin{equation}
  D_{ij} = \frac{c_i}{c_j}\,
    \frac{(-1)^{i+j}}{x_i - x_j} \,, \qquad i \neq j \,,
  \label{eq:diffmat-cgl}
\end{equation}
%
which requires only a sign, a ratio of endpoint flags, and a node
difference --- no products over $N$ terms, and no cancellation.
\histnote{The barycentric form of the differentiation matrix appears
in Trefethen's \emph{Spectral Methods in MATLAB}
\cite{Trefethen2000}, building on earlier work by
Fornberg \cite{Fornberg1988}.}

\paragraph{The row-sum property and diagonal entries.}

Differentiating the constant function $f(x) = 1$ must yield
$f'(x) = 0$, so each row of $D$ sums to zero:
%
\begin{equation}
  \sum_{j=0}^{N} D_{ij} = 0 \qquad
  \text{for every } i \,.
  \label{eq:diffmat-rowsum}
\end{equation}
%
Since the off-diagonal entries are already determined by
\cref{eq:diffmat-cgl}, this identity pins down the remaining
unknowns --- the diagonal entries --- as the negative row sums of
the off-diagonal part:
%
\begin{equation}
  D_{ii} = -\sum_{\substack{j=0 \\ j \neq i}}^{N} D_{ij} \,.
  \label{eq:diffmat-diag}
\end{equation}
%
Computing the diagonal via \cref{eq:diffmat-diag} rather than from
the direct formula avoids a separate derivation and,
more importantly, guarantees that the row-sum property holds to
machine precision regardless of roundoff in the off-diagonal
entries.  Both the Chebyshev and LGL implementations in the
codebase use this approach.
\warnnote{An equivalent closed-form expression exists for the
CGL diagonal entries:
$D_{ii} = -x_i / [2(1 - x_i^2)]$ for interior nodes,
$D_{00} = (2N^2 + 1)/6$, $D_{NN} = -(2N^2 + 1)/6$.
Using this formula directly can violate the row-sum condition by
$\order{N\,\epsilon_{\text{mach}}}$, which degrades derivative
accuracy for large $N$.}

\paragraph{Higher-order derivatives.}

The second-derivative matrix $D^{(2)}$ maps function values to second
derivatives: $\mathbf{f}'' = D^{(2)}\,\mathbf{f}$.  The simplest
construction is the matrix square
%
\begin{equation}
  D^{(2)} = D \cdot D \,,
  \label{eq:diffmat-d2}
\end{equation}
%
since applying $D$ twice differentiates the interpolant twice ---
crucially, both applications use the same collocation grid, so the
composition is exact for polynomials of degree $N$.  The matrix
$D^{(2)}$ inherits the conditioning of $D$ squared:
$\norm{D^{(2)}}_2 = \order{N^4}$, which is severe but manageable
at the moderate polynomial degrees ($N \sim 20$--$30$) used in
practice.  The codebase's \texttt{ChebyshevBasis} stores both
\texttt{d1} and \texttt{d2} as flat row-major arrays, computing
\texttt{d2} from the matrix product at construction time.  For
elliptic problems --- where the second derivative appears in the
Laplacian --- having $D^{(2)}$ precomputed avoids a redundant
matrix multiply at every solver iteration.
\implnote{Higher derivatives ($D^{(3)}$, $D^{(4)}$) can be computed
by further matrix products, but roundoff accumulates: each
multiplication amplifies errors by a factor of $\order{N^2}$.
In practice, third and fourth derivatives are rarely needed directly;
the XCTS solver requires only first and second.}

\paragraph{Conditioning.}

The spectral differentiation matrix $D$ is dense and has large
entries near the endpoints, reflecting the tight node clustering
there.  Its operator norm scales as
%
\begin{equation}
  \norm{D}_2 = \order{N^2} \,,
  \label{eq:diffmat-norm}
\end{equation}
%
meaning that an $\order{\epsilon}$ perturbation in the function values
produces an $\order{N^2\,\epsilon}$ perturbation in the
derivatives \cite{Trefethen2000}.  For $N = 20$, the amplification
factor is $\order{400}$; for $N = 50$, it exceeds $\order{2500}$.
This growth is not a deficiency of the spectral approach --- it is
intrinsic to numerical differentiation.  The finite-difference
analogue is no better: a centred difference
$f'(x) \approx (f(x+h) - f(x-h))/(2h)$ amplifies data errors by
$1/h$, and achieving $\order{h^p}$ accuracy with $N$ points on a
fixed interval forces $h = \order{1/N}$, giving $\order{N}$
amplification --- less than spectral's $\order{N^2}$ --- but with
far worse truncation error that more than offsets this advantage.

The crucial difference is in what happens to the \emph{truncation}
error.  A $p$-th order finite-difference stencil has truncation
error $\order{h^p}$, which improves only polynomially as $N$
grows.  The spectral method has truncation error that decreases
\emph{exponentially} with $N$ for smooth functions
(\cref{sec:spec-convergence}): the rapid decay of the Chebyshev
coefficients ensures that the interpolation error shrinks faster
than $D$ can amplify it.  At $N = 20$, the spectral derivative of an
analytic function is typically accurate to $10^{-12}$ or
better --- twelve orders of magnitude beyond what a second-order
finite difference achieves on the same grid.
\xrefnote{The spectral convergence theorem
(\cref{sec:spec-convergence}) makes this trade-off precise: for
analytic functions, the Chebyshev interpolation error decreases
geometrically at rate $\rho^{-N}$, where $\rho > 1$ depends on
the function's analyticity strip.}

\begin{figure}[htbp]
  \centering
  %% Creator: Matplotlib, PGF backend
%%
%% To include the figure in your LaTeX document, write
%%   \input{<filename>.pgf}
%%
%% Make sure the required packages are loaded in your preamble
%%   \usepackage{pgf}
%%
%% Also ensure that all the required font packages are loaded; for instance,
%% the lmodern package is sometimes necessary when using math font.
%%   \usepackage{lmodern}
%%
%% Figures using additional raster images can only be included by \input if
%% they are in the same directory as the main LaTeX file. For loading figures
%% from other directories you can use the `import` package
%%   \usepackage{import}
%%
%% and then include the figures with
%%   \import{<path to file>}{<filename>.pgf}
%%
%% Matplotlib used the following preamble
%%   \def\mathdefault#1{#1}
%%   \everymath=\expandafter{\the\everymath\displaystyle}
%%   \IfFileExists{scrextend.sty}{
%%     \usepackage[fontsize=10.000000pt]{scrextend}
%%   }{
%%     \renewcommand{\normalsize}{\fontsize{10.000000}{12.000000}\selectfont}
%%     \normalsize
%%   }
%%   \usepackage{fontspec}
%%   \usepackage{unicode-math}
%%   \setmainfont{STIX Two Text}[Numbers=OldStyle, Ligatures=TeX]
%%   \setmathfont{STIX Two Math}
%%   \setmonofont{Source Code Pro}[Scale=MatchLowercase]
%%   \ifdefined\pdftexversion\else  % non-pdftex case.
%%     \usepackage{fontspec}
%%     \setmainfont{DejaVuSerif.ttf}[Path=\detokenize{/Users/gvn/.pyenv/versions/3.12.11/lib/python3.12/site-packages/matplotlib/mpl-data/fonts/ttf/}]
%%     \setsansfont{DejaVuSans.ttf}[Path=\detokenize{/Users/gvn/.pyenv/versions/3.12.11/lib/python3.12/site-packages/matplotlib/mpl-data/fonts/ttf/}]
%%     \setmonofont{DejaVuSansMono.ttf}[Path=\detokenize{/Users/gvn/.pyenv/versions/3.12.11/lib/python3.12/site-packages/matplotlib/mpl-data/fonts/ttf/}]
%%   \fi
%%   \makeatletter\@ifpackageloaded{underscore}{}{\usepackage[strings]{underscore}}\makeatother
%%
\begingroup%
\makeatletter%
\begin{pgfpicture}%
\pgfpathrectangle{\pgfpointorigin}{\pgfqpoint{3.736961in}{2.361438in}}%
\pgfusepath{use as bounding box, clip}%
\begin{pgfscope}%
\pgfsetbuttcap%
\pgfsetmiterjoin%
\definecolor{currentfill}{rgb}{1.000000,1.000000,1.000000}%
\pgfsetfillcolor{currentfill}%
\pgfsetlinewidth{0.000000pt}%
\definecolor{currentstroke}{rgb}{1.000000,1.000000,1.000000}%
\pgfsetstrokecolor{currentstroke}%
\pgfsetdash{}{0pt}%
\pgfpathmoveto{\pgfqpoint{-0.000000in}{0.000000in}}%
\pgfpathlineto{\pgfqpoint{3.736961in}{0.000000in}}%
\pgfpathlineto{\pgfqpoint{3.736961in}{2.361438in}}%
\pgfpathlineto{\pgfqpoint{-0.000000in}{2.361438in}}%
\pgfpathlineto{\pgfqpoint{-0.000000in}{0.000000in}}%
\pgfpathclose%
\pgfusepath{fill}%
\end{pgfscope}%
\begin{pgfscope}%
\pgfsetbuttcap%
\pgfsetmiterjoin%
\definecolor{currentfill}{rgb}{1.000000,1.000000,1.000000}%
\pgfsetfillcolor{currentfill}%
\pgfsetlinewidth{0.000000pt}%
\definecolor{currentstroke}{rgb}{0.000000,0.000000,0.000000}%
\pgfsetstrokecolor{currentstroke}%
\pgfsetstrokeopacity{0.000000}%
\pgfsetdash{}{0pt}%
\pgfpathmoveto{\pgfqpoint{0.544111in}{0.352669in}}%
\pgfpathlineto{\pgfqpoint{3.716961in}{0.352669in}}%
\pgfpathlineto{\pgfqpoint{3.716961in}{2.299229in}}%
\pgfpathlineto{\pgfqpoint{0.544111in}{2.299229in}}%
\pgfpathlineto{\pgfqpoint{0.544111in}{0.352669in}}%
\pgfpathclose%
\pgfusepath{fill}%
\end{pgfscope}%
\begin{pgfscope}%
\pgfpathrectangle{\pgfqpoint{0.544111in}{0.352669in}}{\pgfqpoint{3.172850in}{1.946560in}}%
\pgfusepath{clip}%
\pgfsetbuttcap%
\pgfsetroundjoin%
\pgfsetlinewidth{0.501875pt}%
\definecolor{currentstroke}{rgb}{0.431373,0.423529,0.462745}%
\pgfsetstrokecolor{currentstroke}%
\pgfsetstrokeopacity{0.500000}%
\pgfsetdash{{0.500000pt}{0.825000pt}}{0.000000pt}%
\pgfpathmoveto{\pgfqpoint{0.544111in}{0.390134in}}%
\pgfpathlineto{\pgfqpoint{3.716961in}{0.390134in}}%
\pgfusepath{stroke}%
\end{pgfscope}%
\begin{pgfscope}%
\pgfsetbuttcap%
\pgfsetroundjoin%
\definecolor{currentfill}{rgb}{0.000000,0.000000,0.000000}%
\pgfsetfillcolor{currentfill}%
\pgfsetlinewidth{0.501875pt}%
\definecolor{currentstroke}{rgb}{0.000000,0.000000,0.000000}%
\pgfsetstrokecolor{currentstroke}%
\pgfsetdash{}{0pt}%
\pgfsys@defobject{currentmarker}{\pgfqpoint{0.000000in}{0.000000in}}{\pgfqpoint{0.000000in}{0.041667in}}{%
\pgfpathmoveto{\pgfqpoint{0.000000in}{0.000000in}}%
\pgfpathlineto{\pgfqpoint{0.000000in}{0.041667in}}%
\pgfusepath{stroke,fill}%
}%
\begin{pgfscope}%
\pgfsys@transformshift{1.507909in}{0.352669in}%
\pgfsys@useobject{currentmarker}{}%
\end{pgfscope}%
\end{pgfscope}%
\begin{pgfscope}%
\definecolor{textcolor}{rgb}{0.000000,0.000000,0.000000}%
\pgfsetstrokecolor{textcolor}%
\pgfsetfillcolor{textcolor}%
\pgftext[x=1.507909in,y=0.304058in,,top]{\color{textcolor}{\sffamily\fontsize{8.000000}{9.600000}\selectfont\catcode`\^=\active\def^{\ifmmode\sp\else\^{}\fi}\catcode`\%=\active\def%{\%}$\mathdefault{10^{1}}$}}%
\end{pgfscope}%
\begin{pgfscope}%
\pgfsetbuttcap%
\pgfsetroundjoin%
\definecolor{currentfill}{rgb}{0.000000,0.000000,0.000000}%
\pgfsetfillcolor{currentfill}%
\pgfsetlinewidth{0.501875pt}%
\definecolor{currentstroke}{rgb}{0.000000,0.000000,0.000000}%
\pgfsetstrokecolor{currentstroke}%
\pgfsetdash{}{0pt}%
\pgfsys@defobject{currentmarker}{\pgfqpoint{0.000000in}{0.000000in}}{\pgfqpoint{0.000000in}{0.041667in}}{%
\pgfpathmoveto{\pgfqpoint{0.000000in}{0.000000in}}%
\pgfpathlineto{\pgfqpoint{0.000000in}{0.041667in}}%
\pgfusepath{stroke,fill}%
}%
\begin{pgfscope}%
\pgfsys@transformshift{2.886791in}{0.352669in}%
\pgfsys@useobject{currentmarker}{}%
\end{pgfscope}%
\end{pgfscope}%
\begin{pgfscope}%
\definecolor{textcolor}{rgb}{0.000000,0.000000,0.000000}%
\pgfsetstrokecolor{textcolor}%
\pgfsetfillcolor{textcolor}%
\pgftext[x=2.886791in,y=0.304058in,,top]{\color{textcolor}{\sffamily\fontsize{8.000000}{9.600000}\selectfont\catcode`\^=\active\def^{\ifmmode\sp\else\^{}\fi}\catcode`\%=\active\def%{\%}$\mathdefault{10^{2}}$}}%
\end{pgfscope}%
\begin{pgfscope}%
\pgfsetbuttcap%
\pgfsetroundjoin%
\definecolor{currentfill}{rgb}{0.000000,0.000000,0.000000}%
\pgfsetfillcolor{currentfill}%
\pgfsetlinewidth{0.301125pt}%
\definecolor{currentstroke}{rgb}{0.000000,0.000000,0.000000}%
\pgfsetstrokecolor{currentstroke}%
\pgfsetdash{}{0pt}%
\pgfsys@defobject{currentmarker}{\pgfqpoint{0.000000in}{0.000000in}}{\pgfqpoint{0.000000in}{0.027778in}}{%
\pgfpathmoveto{\pgfqpoint{0.000000in}{0.000000in}}%
\pgfpathlineto{\pgfqpoint{0.000000in}{0.027778in}}%
\pgfusepath{stroke,fill}%
}%
\begin{pgfscope}%
\pgfsys@transformshift{0.544111in}{0.352669in}%
\pgfsys@useobject{currentmarker}{}%
\end{pgfscope}%
\end{pgfscope}%
\begin{pgfscope}%
\pgfsetbuttcap%
\pgfsetroundjoin%
\definecolor{currentfill}{rgb}{0.000000,0.000000,0.000000}%
\pgfsetfillcolor{currentfill}%
\pgfsetlinewidth{0.301125pt}%
\definecolor{currentstroke}{rgb}{0.000000,0.000000,0.000000}%
\pgfsetstrokecolor{currentstroke}%
\pgfsetdash{}{0pt}%
\pgfsys@defobject{currentmarker}{\pgfqpoint{0.000000in}{0.000000in}}{\pgfqpoint{0.000000in}{0.027778in}}{%
\pgfpathmoveto{\pgfqpoint{0.000000in}{0.000000in}}%
\pgfpathlineto{\pgfqpoint{0.000000in}{0.027778in}}%
\pgfusepath{stroke,fill}%
}%
\begin{pgfscope}%
\pgfsys@transformshift{0.786920in}{0.352669in}%
\pgfsys@useobject{currentmarker}{}%
\end{pgfscope}%
\end{pgfscope}%
\begin{pgfscope}%
\pgfsetbuttcap%
\pgfsetroundjoin%
\definecolor{currentfill}{rgb}{0.000000,0.000000,0.000000}%
\pgfsetfillcolor{currentfill}%
\pgfsetlinewidth{0.301125pt}%
\definecolor{currentstroke}{rgb}{0.000000,0.000000,0.000000}%
\pgfsetstrokecolor{currentstroke}%
\pgfsetdash{}{0pt}%
\pgfsys@defobject{currentmarker}{\pgfqpoint{0.000000in}{0.000000in}}{\pgfqpoint{0.000000in}{0.027778in}}{%
\pgfpathmoveto{\pgfqpoint{0.000000in}{0.000000in}}%
\pgfpathlineto{\pgfqpoint{0.000000in}{0.027778in}}%
\pgfusepath{stroke,fill}%
}%
\begin{pgfscope}%
\pgfsys@transformshift{0.959196in}{0.352669in}%
\pgfsys@useobject{currentmarker}{}%
\end{pgfscope}%
\end{pgfscope}%
\begin{pgfscope}%
\pgfsetbuttcap%
\pgfsetroundjoin%
\definecolor{currentfill}{rgb}{0.000000,0.000000,0.000000}%
\pgfsetfillcolor{currentfill}%
\pgfsetlinewidth{0.301125pt}%
\definecolor{currentstroke}{rgb}{0.000000,0.000000,0.000000}%
\pgfsetstrokecolor{currentstroke}%
\pgfsetdash{}{0pt}%
\pgfsys@defobject{currentmarker}{\pgfqpoint{0.000000in}{0.000000in}}{\pgfqpoint{0.000000in}{0.027778in}}{%
\pgfpathmoveto{\pgfqpoint{0.000000in}{0.000000in}}%
\pgfpathlineto{\pgfqpoint{0.000000in}{0.027778in}}%
\pgfusepath{stroke,fill}%
}%
\begin{pgfscope}%
\pgfsys@transformshift{1.092824in}{0.352669in}%
\pgfsys@useobject{currentmarker}{}%
\end{pgfscope}%
\end{pgfscope}%
\begin{pgfscope}%
\pgfsetbuttcap%
\pgfsetroundjoin%
\definecolor{currentfill}{rgb}{0.000000,0.000000,0.000000}%
\pgfsetfillcolor{currentfill}%
\pgfsetlinewidth{0.301125pt}%
\definecolor{currentstroke}{rgb}{0.000000,0.000000,0.000000}%
\pgfsetstrokecolor{currentstroke}%
\pgfsetdash{}{0pt}%
\pgfsys@defobject{currentmarker}{\pgfqpoint{0.000000in}{0.000000in}}{\pgfqpoint{0.000000in}{0.027778in}}{%
\pgfpathmoveto{\pgfqpoint{0.000000in}{0.000000in}}%
\pgfpathlineto{\pgfqpoint{0.000000in}{0.027778in}}%
\pgfusepath{stroke,fill}%
}%
\begin{pgfscope}%
\pgfsys@transformshift{1.202005in}{0.352669in}%
\pgfsys@useobject{currentmarker}{}%
\end{pgfscope}%
\end{pgfscope}%
\begin{pgfscope}%
\pgfsetbuttcap%
\pgfsetroundjoin%
\definecolor{currentfill}{rgb}{0.000000,0.000000,0.000000}%
\pgfsetfillcolor{currentfill}%
\pgfsetlinewidth{0.301125pt}%
\definecolor{currentstroke}{rgb}{0.000000,0.000000,0.000000}%
\pgfsetstrokecolor{currentstroke}%
\pgfsetdash{}{0pt}%
\pgfsys@defobject{currentmarker}{\pgfqpoint{0.000000in}{0.000000in}}{\pgfqpoint{0.000000in}{0.027778in}}{%
\pgfpathmoveto{\pgfqpoint{0.000000in}{0.000000in}}%
\pgfpathlineto{\pgfqpoint{0.000000in}{0.027778in}}%
\pgfusepath{stroke,fill}%
}%
\begin{pgfscope}%
\pgfsys@transformshift{1.294317in}{0.352669in}%
\pgfsys@useobject{currentmarker}{}%
\end{pgfscope}%
\end{pgfscope}%
\begin{pgfscope}%
\pgfsetbuttcap%
\pgfsetroundjoin%
\definecolor{currentfill}{rgb}{0.000000,0.000000,0.000000}%
\pgfsetfillcolor{currentfill}%
\pgfsetlinewidth{0.301125pt}%
\definecolor{currentstroke}{rgb}{0.000000,0.000000,0.000000}%
\pgfsetstrokecolor{currentstroke}%
\pgfsetdash{}{0pt}%
\pgfsys@defobject{currentmarker}{\pgfqpoint{0.000000in}{0.000000in}}{\pgfqpoint{0.000000in}{0.027778in}}{%
\pgfpathmoveto{\pgfqpoint{0.000000in}{0.000000in}}%
\pgfpathlineto{\pgfqpoint{0.000000in}{0.027778in}}%
\pgfusepath{stroke,fill}%
}%
\begin{pgfscope}%
\pgfsys@transformshift{1.374281in}{0.352669in}%
\pgfsys@useobject{currentmarker}{}%
\end{pgfscope}%
\end{pgfscope}%
\begin{pgfscope}%
\pgfsetbuttcap%
\pgfsetroundjoin%
\definecolor{currentfill}{rgb}{0.000000,0.000000,0.000000}%
\pgfsetfillcolor{currentfill}%
\pgfsetlinewidth{0.301125pt}%
\definecolor{currentstroke}{rgb}{0.000000,0.000000,0.000000}%
\pgfsetstrokecolor{currentstroke}%
\pgfsetdash{}{0pt}%
\pgfsys@defobject{currentmarker}{\pgfqpoint{0.000000in}{0.000000in}}{\pgfqpoint{0.000000in}{0.027778in}}{%
\pgfpathmoveto{\pgfqpoint{0.000000in}{0.000000in}}%
\pgfpathlineto{\pgfqpoint{0.000000in}{0.027778in}}%
\pgfusepath{stroke,fill}%
}%
\begin{pgfscope}%
\pgfsys@transformshift{1.444815in}{0.352669in}%
\pgfsys@useobject{currentmarker}{}%
\end{pgfscope}%
\end{pgfscope}%
\begin{pgfscope}%
\pgfsetbuttcap%
\pgfsetroundjoin%
\definecolor{currentfill}{rgb}{0.000000,0.000000,0.000000}%
\pgfsetfillcolor{currentfill}%
\pgfsetlinewidth{0.301125pt}%
\definecolor{currentstroke}{rgb}{0.000000,0.000000,0.000000}%
\pgfsetstrokecolor{currentstroke}%
\pgfsetdash{}{0pt}%
\pgfsys@defobject{currentmarker}{\pgfqpoint{0.000000in}{0.000000in}}{\pgfqpoint{0.000000in}{0.027778in}}{%
\pgfpathmoveto{\pgfqpoint{0.000000in}{0.000000in}}%
\pgfpathlineto{\pgfqpoint{0.000000in}{0.027778in}}%
\pgfusepath{stroke,fill}%
}%
\begin{pgfscope}%
\pgfsys@transformshift{1.922994in}{0.352669in}%
\pgfsys@useobject{currentmarker}{}%
\end{pgfscope}%
\end{pgfscope}%
\begin{pgfscope}%
\pgfsetbuttcap%
\pgfsetroundjoin%
\definecolor{currentfill}{rgb}{0.000000,0.000000,0.000000}%
\pgfsetfillcolor{currentfill}%
\pgfsetlinewidth{0.301125pt}%
\definecolor{currentstroke}{rgb}{0.000000,0.000000,0.000000}%
\pgfsetstrokecolor{currentstroke}%
\pgfsetdash{}{0pt}%
\pgfsys@defobject{currentmarker}{\pgfqpoint{0.000000in}{0.000000in}}{\pgfqpoint{0.000000in}{0.027778in}}{%
\pgfpathmoveto{\pgfqpoint{0.000000in}{0.000000in}}%
\pgfpathlineto{\pgfqpoint{0.000000in}{0.027778in}}%
\pgfusepath{stroke,fill}%
}%
\begin{pgfscope}%
\pgfsys@transformshift{2.165803in}{0.352669in}%
\pgfsys@useobject{currentmarker}{}%
\end{pgfscope}%
\end{pgfscope}%
\begin{pgfscope}%
\pgfsetbuttcap%
\pgfsetroundjoin%
\definecolor{currentfill}{rgb}{0.000000,0.000000,0.000000}%
\pgfsetfillcolor{currentfill}%
\pgfsetlinewidth{0.301125pt}%
\definecolor{currentstroke}{rgb}{0.000000,0.000000,0.000000}%
\pgfsetstrokecolor{currentstroke}%
\pgfsetdash{}{0pt}%
\pgfsys@defobject{currentmarker}{\pgfqpoint{0.000000in}{0.000000in}}{\pgfqpoint{0.000000in}{0.027778in}}{%
\pgfpathmoveto{\pgfqpoint{0.000000in}{0.000000in}}%
\pgfpathlineto{\pgfqpoint{0.000000in}{0.027778in}}%
\pgfusepath{stroke,fill}%
}%
\begin{pgfscope}%
\pgfsys@transformshift{2.338079in}{0.352669in}%
\pgfsys@useobject{currentmarker}{}%
\end{pgfscope}%
\end{pgfscope}%
\begin{pgfscope}%
\pgfsetbuttcap%
\pgfsetroundjoin%
\definecolor{currentfill}{rgb}{0.000000,0.000000,0.000000}%
\pgfsetfillcolor{currentfill}%
\pgfsetlinewidth{0.301125pt}%
\definecolor{currentstroke}{rgb}{0.000000,0.000000,0.000000}%
\pgfsetstrokecolor{currentstroke}%
\pgfsetdash{}{0pt}%
\pgfsys@defobject{currentmarker}{\pgfqpoint{0.000000in}{0.000000in}}{\pgfqpoint{0.000000in}{0.027778in}}{%
\pgfpathmoveto{\pgfqpoint{0.000000in}{0.000000in}}%
\pgfpathlineto{\pgfqpoint{0.000000in}{0.027778in}}%
\pgfusepath{stroke,fill}%
}%
\begin{pgfscope}%
\pgfsys@transformshift{2.471706in}{0.352669in}%
\pgfsys@useobject{currentmarker}{}%
\end{pgfscope}%
\end{pgfscope}%
\begin{pgfscope}%
\pgfsetbuttcap%
\pgfsetroundjoin%
\definecolor{currentfill}{rgb}{0.000000,0.000000,0.000000}%
\pgfsetfillcolor{currentfill}%
\pgfsetlinewidth{0.301125pt}%
\definecolor{currentstroke}{rgb}{0.000000,0.000000,0.000000}%
\pgfsetstrokecolor{currentstroke}%
\pgfsetdash{}{0pt}%
\pgfsys@defobject{currentmarker}{\pgfqpoint{0.000000in}{0.000000in}}{\pgfqpoint{0.000000in}{0.027778in}}{%
\pgfpathmoveto{\pgfqpoint{0.000000in}{0.000000in}}%
\pgfpathlineto{\pgfqpoint{0.000000in}{0.027778in}}%
\pgfusepath{stroke,fill}%
}%
\begin{pgfscope}%
\pgfsys@transformshift{2.580888in}{0.352669in}%
\pgfsys@useobject{currentmarker}{}%
\end{pgfscope}%
\end{pgfscope}%
\begin{pgfscope}%
\pgfsetbuttcap%
\pgfsetroundjoin%
\definecolor{currentfill}{rgb}{0.000000,0.000000,0.000000}%
\pgfsetfillcolor{currentfill}%
\pgfsetlinewidth{0.301125pt}%
\definecolor{currentstroke}{rgb}{0.000000,0.000000,0.000000}%
\pgfsetstrokecolor{currentstroke}%
\pgfsetdash{}{0pt}%
\pgfsys@defobject{currentmarker}{\pgfqpoint{0.000000in}{0.000000in}}{\pgfqpoint{0.000000in}{0.027778in}}{%
\pgfpathmoveto{\pgfqpoint{0.000000in}{0.000000in}}%
\pgfpathlineto{\pgfqpoint{0.000000in}{0.027778in}}%
\pgfusepath{stroke,fill}%
}%
\begin{pgfscope}%
\pgfsys@transformshift{2.673200in}{0.352669in}%
\pgfsys@useobject{currentmarker}{}%
\end{pgfscope}%
\end{pgfscope}%
\begin{pgfscope}%
\pgfsetbuttcap%
\pgfsetroundjoin%
\definecolor{currentfill}{rgb}{0.000000,0.000000,0.000000}%
\pgfsetfillcolor{currentfill}%
\pgfsetlinewidth{0.301125pt}%
\definecolor{currentstroke}{rgb}{0.000000,0.000000,0.000000}%
\pgfsetstrokecolor{currentstroke}%
\pgfsetdash{}{0pt}%
\pgfsys@defobject{currentmarker}{\pgfqpoint{0.000000in}{0.000000in}}{\pgfqpoint{0.000000in}{0.027778in}}{%
\pgfpathmoveto{\pgfqpoint{0.000000in}{0.000000in}}%
\pgfpathlineto{\pgfqpoint{0.000000in}{0.027778in}}%
\pgfusepath{stroke,fill}%
}%
\begin{pgfscope}%
\pgfsys@transformshift{2.753164in}{0.352669in}%
\pgfsys@useobject{currentmarker}{}%
\end{pgfscope}%
\end{pgfscope}%
\begin{pgfscope}%
\pgfsetbuttcap%
\pgfsetroundjoin%
\definecolor{currentfill}{rgb}{0.000000,0.000000,0.000000}%
\pgfsetfillcolor{currentfill}%
\pgfsetlinewidth{0.301125pt}%
\definecolor{currentstroke}{rgb}{0.000000,0.000000,0.000000}%
\pgfsetstrokecolor{currentstroke}%
\pgfsetdash{}{0pt}%
\pgfsys@defobject{currentmarker}{\pgfqpoint{0.000000in}{0.000000in}}{\pgfqpoint{0.000000in}{0.027778in}}{%
\pgfpathmoveto{\pgfqpoint{0.000000in}{0.000000in}}%
\pgfpathlineto{\pgfqpoint{0.000000in}{0.027778in}}%
\pgfusepath{stroke,fill}%
}%
\begin{pgfscope}%
\pgfsys@transformshift{2.823697in}{0.352669in}%
\pgfsys@useobject{currentmarker}{}%
\end{pgfscope}%
\end{pgfscope}%
\begin{pgfscope}%
\pgfsetbuttcap%
\pgfsetroundjoin%
\definecolor{currentfill}{rgb}{0.000000,0.000000,0.000000}%
\pgfsetfillcolor{currentfill}%
\pgfsetlinewidth{0.301125pt}%
\definecolor{currentstroke}{rgb}{0.000000,0.000000,0.000000}%
\pgfsetstrokecolor{currentstroke}%
\pgfsetdash{}{0pt}%
\pgfsys@defobject{currentmarker}{\pgfqpoint{0.000000in}{0.000000in}}{\pgfqpoint{0.000000in}{0.027778in}}{%
\pgfpathmoveto{\pgfqpoint{0.000000in}{0.000000in}}%
\pgfpathlineto{\pgfqpoint{0.000000in}{0.027778in}}%
\pgfusepath{stroke,fill}%
}%
\begin{pgfscope}%
\pgfsys@transformshift{3.301876in}{0.352669in}%
\pgfsys@useobject{currentmarker}{}%
\end{pgfscope}%
\end{pgfscope}%
\begin{pgfscope}%
\pgfsetbuttcap%
\pgfsetroundjoin%
\definecolor{currentfill}{rgb}{0.000000,0.000000,0.000000}%
\pgfsetfillcolor{currentfill}%
\pgfsetlinewidth{0.301125pt}%
\definecolor{currentstroke}{rgb}{0.000000,0.000000,0.000000}%
\pgfsetstrokecolor{currentstroke}%
\pgfsetdash{}{0pt}%
\pgfsys@defobject{currentmarker}{\pgfqpoint{0.000000in}{0.000000in}}{\pgfqpoint{0.000000in}{0.027778in}}{%
\pgfpathmoveto{\pgfqpoint{0.000000in}{0.000000in}}%
\pgfpathlineto{\pgfqpoint{0.000000in}{0.027778in}}%
\pgfusepath{stroke,fill}%
}%
\begin{pgfscope}%
\pgfsys@transformshift{3.544685in}{0.352669in}%
\pgfsys@useobject{currentmarker}{}%
\end{pgfscope}%
\end{pgfscope}%
\begin{pgfscope}%
\pgfsetbuttcap%
\pgfsetroundjoin%
\definecolor{currentfill}{rgb}{0.000000,0.000000,0.000000}%
\pgfsetfillcolor{currentfill}%
\pgfsetlinewidth{0.301125pt}%
\definecolor{currentstroke}{rgb}{0.000000,0.000000,0.000000}%
\pgfsetstrokecolor{currentstroke}%
\pgfsetdash{}{0pt}%
\pgfsys@defobject{currentmarker}{\pgfqpoint{0.000000in}{0.000000in}}{\pgfqpoint{0.000000in}{0.027778in}}{%
\pgfpathmoveto{\pgfqpoint{0.000000in}{0.000000in}}%
\pgfpathlineto{\pgfqpoint{0.000000in}{0.027778in}}%
\pgfusepath{stroke,fill}%
}%
\begin{pgfscope}%
\pgfsys@transformshift{3.716961in}{0.352669in}%
\pgfsys@useobject{currentmarker}{}%
\end{pgfscope}%
\end{pgfscope}%
\begin{pgfscope}%
\definecolor{textcolor}{rgb}{0.000000,0.000000,0.000000}%
\pgfsetstrokecolor{textcolor}%
\pgfsetfillcolor{textcolor}%
\pgftext[x=2.130536in,y=0.140972in,,top]{\color{textcolor}{\sffamily\fontsize{9.000000}{10.800000}\selectfont\catcode`\^=\active\def^{\ifmmode\sp\else\^{}\fi}\catcode`\%=\active\def%{\%}Number of points $N$}}%
\end{pgfscope}%
\begin{pgfscope}%
\pgfsetbuttcap%
\pgfsetroundjoin%
\definecolor{currentfill}{rgb}{0.000000,0.000000,0.000000}%
\pgfsetfillcolor{currentfill}%
\pgfsetlinewidth{0.501875pt}%
\definecolor{currentstroke}{rgb}{0.000000,0.000000,0.000000}%
\pgfsetstrokecolor{currentstroke}%
\pgfsetdash{}{0pt}%
\pgfsys@defobject{currentmarker}{\pgfqpoint{0.000000in}{0.000000in}}{\pgfqpoint{0.041667in}{0.000000in}}{%
\pgfpathmoveto{\pgfqpoint{0.000000in}{0.000000in}}%
\pgfpathlineto{\pgfqpoint{0.041667in}{0.000000in}}%
\pgfusepath{stroke,fill}%
}%
\begin{pgfscope}%
\pgfsys@transformshift{0.544111in}{0.352669in}%
\pgfsys@useobject{currentmarker}{}%
\end{pgfscope}%
\end{pgfscope}%
\begin{pgfscope}%
\definecolor{textcolor}{rgb}{0.000000,0.000000,0.000000}%
\pgfsetstrokecolor{textcolor}%
\pgfsetfillcolor{textcolor}%
\pgftext[x=0.207722in, y=0.310460in, left, base]{\color{textcolor}{\sffamily\fontsize{8.000000}{9.600000}\selectfont\catcode`\^=\active\def^{\ifmmode\sp\else\^{}\fi}\catcode`\%=\active\def%{\%}$\mathdefault{10^{-16}}$}}%
\end{pgfscope}%
\begin{pgfscope}%
\pgfsetbuttcap%
\pgfsetroundjoin%
\definecolor{currentfill}{rgb}{0.000000,0.000000,0.000000}%
\pgfsetfillcolor{currentfill}%
\pgfsetlinewidth{0.501875pt}%
\definecolor{currentstroke}{rgb}{0.000000,0.000000,0.000000}%
\pgfsetstrokecolor{currentstroke}%
\pgfsetdash{}{0pt}%
\pgfsys@defobject{currentmarker}{\pgfqpoint{0.000000in}{0.000000in}}{\pgfqpoint{0.041667in}{0.000000in}}{%
\pgfpathmoveto{\pgfqpoint{0.000000in}{0.000000in}}%
\pgfpathlineto{\pgfqpoint{0.041667in}{0.000000in}}%
\pgfusepath{stroke,fill}%
}%
\begin{pgfscope}%
\pgfsys@transformshift{0.544111in}{0.677095in}%
\pgfsys@useobject{currentmarker}{}%
\end{pgfscope}%
\end{pgfscope}%
\begin{pgfscope}%
\definecolor{textcolor}{rgb}{0.000000,0.000000,0.000000}%
\pgfsetstrokecolor{textcolor}%
\pgfsetfillcolor{textcolor}%
\pgftext[x=0.209889in, y=0.634886in, left, base]{\color{textcolor}{\sffamily\fontsize{8.000000}{9.600000}\selectfont\catcode`\^=\active\def^{\ifmmode\sp\else\^{}\fi}\catcode`\%=\active\def%{\%}$\mathdefault{10^{-13}}$}}%
\end{pgfscope}%
\begin{pgfscope}%
\pgfsetbuttcap%
\pgfsetroundjoin%
\definecolor{currentfill}{rgb}{0.000000,0.000000,0.000000}%
\pgfsetfillcolor{currentfill}%
\pgfsetlinewidth{0.501875pt}%
\definecolor{currentstroke}{rgb}{0.000000,0.000000,0.000000}%
\pgfsetstrokecolor{currentstroke}%
\pgfsetdash{}{0pt}%
\pgfsys@defobject{currentmarker}{\pgfqpoint{0.000000in}{0.000000in}}{\pgfqpoint{0.041667in}{0.000000in}}{%
\pgfpathmoveto{\pgfqpoint{0.000000in}{0.000000in}}%
\pgfpathlineto{\pgfqpoint{0.041667in}{0.000000in}}%
\pgfusepath{stroke,fill}%
}%
\begin{pgfscope}%
\pgfsys@transformshift{0.544111in}{1.001522in}%
\pgfsys@useobject{currentmarker}{}%
\end{pgfscope}%
\end{pgfscope}%
\begin{pgfscope}%
\definecolor{textcolor}{rgb}{0.000000,0.000000,0.000000}%
\pgfsetstrokecolor{textcolor}%
\pgfsetfillcolor{textcolor}%
\pgftext[x=0.206139in, y=0.959313in, left, base]{\color{textcolor}{\sffamily\fontsize{8.000000}{9.600000}\selectfont\catcode`\^=\active\def^{\ifmmode\sp\else\^{}\fi}\catcode`\%=\active\def%{\%}$\mathdefault{10^{-10}}$}}%
\end{pgfscope}%
\begin{pgfscope}%
\pgfsetbuttcap%
\pgfsetroundjoin%
\definecolor{currentfill}{rgb}{0.000000,0.000000,0.000000}%
\pgfsetfillcolor{currentfill}%
\pgfsetlinewidth{0.501875pt}%
\definecolor{currentstroke}{rgb}{0.000000,0.000000,0.000000}%
\pgfsetstrokecolor{currentstroke}%
\pgfsetdash{}{0pt}%
\pgfsys@defobject{currentmarker}{\pgfqpoint{0.000000in}{0.000000in}}{\pgfqpoint{0.041667in}{0.000000in}}{%
\pgfpathmoveto{\pgfqpoint{0.000000in}{0.000000in}}%
\pgfpathlineto{\pgfqpoint{0.041667in}{0.000000in}}%
\pgfusepath{stroke,fill}%
}%
\begin{pgfscope}%
\pgfsys@transformshift{0.544111in}{1.325949in}%
\pgfsys@useobject{currentmarker}{}%
\end{pgfscope}%
\end{pgfscope}%
\begin{pgfscope}%
\definecolor{textcolor}{rgb}{0.000000,0.000000,0.000000}%
\pgfsetstrokecolor{textcolor}%
\pgfsetfillcolor{textcolor}%
\pgftext[x=0.258556in, y=1.283740in, left, base]{\color{textcolor}{\sffamily\fontsize{8.000000}{9.600000}\selectfont\catcode`\^=\active\def^{\ifmmode\sp\else\^{}\fi}\catcode`\%=\active\def%{\%}$\mathdefault{10^{-7}}$}}%
\end{pgfscope}%
\begin{pgfscope}%
\pgfsetbuttcap%
\pgfsetroundjoin%
\definecolor{currentfill}{rgb}{0.000000,0.000000,0.000000}%
\pgfsetfillcolor{currentfill}%
\pgfsetlinewidth{0.501875pt}%
\definecolor{currentstroke}{rgb}{0.000000,0.000000,0.000000}%
\pgfsetstrokecolor{currentstroke}%
\pgfsetdash{}{0pt}%
\pgfsys@defobject{currentmarker}{\pgfqpoint{0.000000in}{0.000000in}}{\pgfqpoint{0.041667in}{0.000000in}}{%
\pgfpathmoveto{\pgfqpoint{0.000000in}{0.000000in}}%
\pgfpathlineto{\pgfqpoint{0.041667in}{0.000000in}}%
\pgfusepath{stroke,fill}%
}%
\begin{pgfscope}%
\pgfsys@transformshift{0.544111in}{1.650375in}%
\pgfsys@useobject{currentmarker}{}%
\end{pgfscope}%
\end{pgfscope}%
\begin{pgfscope}%
\definecolor{textcolor}{rgb}{0.000000,0.000000,0.000000}%
\pgfsetstrokecolor{textcolor}%
\pgfsetfillcolor{textcolor}%
\pgftext[x=0.253472in, y=1.608166in, left, base]{\color{textcolor}{\sffamily\fontsize{8.000000}{9.600000}\selectfont\catcode`\^=\active\def^{\ifmmode\sp\else\^{}\fi}\catcode`\%=\active\def%{\%}$\mathdefault{10^{-4}}$}}%
\end{pgfscope}%
\begin{pgfscope}%
\pgfsetbuttcap%
\pgfsetroundjoin%
\definecolor{currentfill}{rgb}{0.000000,0.000000,0.000000}%
\pgfsetfillcolor{currentfill}%
\pgfsetlinewidth{0.501875pt}%
\definecolor{currentstroke}{rgb}{0.000000,0.000000,0.000000}%
\pgfsetstrokecolor{currentstroke}%
\pgfsetdash{}{0pt}%
\pgfsys@defobject{currentmarker}{\pgfqpoint{0.000000in}{0.000000in}}{\pgfqpoint{0.041667in}{0.000000in}}{%
\pgfpathmoveto{\pgfqpoint{0.000000in}{0.000000in}}%
\pgfpathlineto{\pgfqpoint{0.041667in}{0.000000in}}%
\pgfusepath{stroke,fill}%
}%
\begin{pgfscope}%
\pgfsys@transformshift{0.544111in}{1.974802in}%
\pgfsys@useobject{currentmarker}{}%
\end{pgfscope}%
\end{pgfscope}%
\begin{pgfscope}%
\definecolor{textcolor}{rgb}{0.000000,0.000000,0.000000}%
\pgfsetstrokecolor{textcolor}%
\pgfsetfillcolor{textcolor}%
\pgftext[x=0.256639in, y=1.932593in, left, base]{\color{textcolor}{\sffamily\fontsize{8.000000}{9.600000}\selectfont\catcode`\^=\active\def^{\ifmmode\sp\else\^{}\fi}\catcode`\%=\active\def%{\%}$\mathdefault{10^{-1}}$}}%
\end{pgfscope}%
\begin{pgfscope}%
\pgfsetbuttcap%
\pgfsetroundjoin%
\definecolor{currentfill}{rgb}{0.000000,0.000000,0.000000}%
\pgfsetfillcolor{currentfill}%
\pgfsetlinewidth{0.501875pt}%
\definecolor{currentstroke}{rgb}{0.000000,0.000000,0.000000}%
\pgfsetstrokecolor{currentstroke}%
\pgfsetdash{}{0pt}%
\pgfsys@defobject{currentmarker}{\pgfqpoint{0.000000in}{0.000000in}}{\pgfqpoint{0.041667in}{0.000000in}}{%
\pgfpathmoveto{\pgfqpoint{0.000000in}{0.000000in}}%
\pgfpathlineto{\pgfqpoint{0.041667in}{0.000000in}}%
\pgfusepath{stroke,fill}%
}%
\begin{pgfscope}%
\pgfsys@transformshift{0.544111in}{2.299229in}%
\pgfsys@useobject{currentmarker}{}%
\end{pgfscope}%
\end{pgfscope}%
\begin{pgfscope}%
\definecolor{textcolor}{rgb}{0.000000,0.000000,0.000000}%
\pgfsetstrokecolor{textcolor}%
\pgfsetfillcolor{textcolor}%
\pgftext[x=0.323472in, y=2.257020in, left, base]{\color{textcolor}{\sffamily\fontsize{8.000000}{9.600000}\selectfont\catcode`\^=\active\def^{\ifmmode\sp\else\^{}\fi}\catcode`\%=\active\def%{\%}$\mathdefault{10^{2}}$}}%
\end{pgfscope}%
\begin{pgfscope}%
\definecolor{textcolor}{rgb}{0.000000,0.000000,0.000000}%
\pgfsetstrokecolor{textcolor}%
\pgfsetfillcolor{textcolor}%
\pgftext[x=0.150583in,y=1.325949in,,bottom,rotate=90.000000]{\color{textcolor}{\sffamily\fontsize{9.000000}{10.800000}\selectfont\catcode`\^=\active\def^{\ifmmode\sp\else\^{}\fi}\catcode`\%=\active\def%{\%}$\|f_{\mathrm{numerical}}' - f_{\mathrm{exact}}'\|_\infty$}}%
\end{pgfscope}%
\begin{pgfscope}%
\pgfpathrectangle{\pgfqpoint{0.544111in}{0.352669in}}{\pgfqpoint{3.172850in}{1.946560in}}%
\pgfusepath{clip}%
\pgfsetrectcap%
\pgfsetroundjoin%
\pgfsetlinewidth{1.204500pt}%
\definecolor{currentstroke}{rgb}{0.250980,0.400000,0.600000}%
\pgfsetstrokecolor{currentstroke}%
\pgfsetdash{}{0pt}%
\pgfpathmoveto{\pgfqpoint{0.959196in}{2.112526in}}%
\pgfpathlineto{\pgfqpoint{1.092824in}{2.116285in}}%
\pgfpathlineto{\pgfqpoint{1.294317in}{2.098464in}}%
\pgfpathlineto{\pgfqpoint{1.374281in}{2.088009in}}%
\pgfpathlineto{\pgfqpoint{1.617090in}{2.071322in}}%
\pgfpathlineto{\pgfqpoint{1.952211in}{2.046585in}}%
\pgfpathlineto{\pgfqpoint{2.506600in}{2.003848in}}%
\pgfpathlineto{\pgfqpoint{3.716961in}{1.909075in}}%
\pgfpathlineto{\pgfqpoint{3.716961in}{1.909075in}}%
\pgfusepath{stroke}%
\end{pgfscope}%
\begin{pgfscope}%
\pgfpathrectangle{\pgfqpoint{0.544111in}{0.352669in}}{\pgfqpoint{3.172850in}{1.946560in}}%
\pgfusepath{clip}%
\pgfsetbuttcap%
\pgfsetroundjoin%
\pgfsetlinewidth{1.204500pt}%
\definecolor{currentstroke}{rgb}{0.431373,0.423529,0.462745}%
\pgfsetstrokecolor{currentstroke}%
\pgfsetdash{{4.440000pt}{1.920000pt}}{0.000000pt}%
\pgfpathmoveto{\pgfqpoint{0.959196in}{2.112526in}}%
\pgfpathlineto{\pgfqpoint{1.092824in}{2.105236in}}%
\pgfpathlineto{\pgfqpoint{1.202005in}{2.106982in}}%
\pgfpathlineto{\pgfqpoint{1.294317in}{2.098464in}}%
\pgfpathlineto{\pgfqpoint{1.374281in}{2.088009in}}%
\pgfpathlineto{\pgfqpoint{1.617090in}{2.071322in}}%
\pgfpathlineto{\pgfqpoint{1.952211in}{2.046585in}}%
\pgfpathlineto{\pgfqpoint{2.506600in}{2.003848in}}%
\pgfpathlineto{\pgfqpoint{3.716961in}{1.909075in}}%
\pgfpathlineto{\pgfqpoint{3.716961in}{1.909075in}}%
\pgfusepath{stroke}%
\end{pgfscope}%
\begin{pgfscope}%
\pgfsetrectcap%
\pgfsetmiterjoin%
\pgfsetlinewidth{0.602250pt}%
\definecolor{currentstroke}{rgb}{0.000000,0.000000,0.000000}%
\pgfsetstrokecolor{currentstroke}%
\pgfsetdash{}{0pt}%
\pgfpathmoveto{\pgfqpoint{0.544111in}{0.352669in}}%
\pgfpathlineto{\pgfqpoint{0.544111in}{2.299229in}}%
\pgfusepath{stroke}%
\end{pgfscope}%
\begin{pgfscope}%
\pgfsetrectcap%
\pgfsetmiterjoin%
\pgfsetlinewidth{0.602250pt}%
\definecolor{currentstroke}{rgb}{0.000000,0.000000,0.000000}%
\pgfsetstrokecolor{currentstroke}%
\pgfsetdash{}{0pt}%
\pgfpathmoveto{\pgfqpoint{3.716961in}{0.352669in}}%
\pgfpathlineto{\pgfqpoint{3.716961in}{2.299229in}}%
\pgfusepath{stroke}%
\end{pgfscope}%
\begin{pgfscope}%
\pgfsetrectcap%
\pgfsetmiterjoin%
\pgfsetlinewidth{0.602250pt}%
\definecolor{currentstroke}{rgb}{0.000000,0.000000,0.000000}%
\pgfsetstrokecolor{currentstroke}%
\pgfsetdash{}{0pt}%
\pgfpathmoveto{\pgfqpoint{0.544111in}{0.352669in}}%
\pgfpathlineto{\pgfqpoint{3.716961in}{0.352669in}}%
\pgfusepath{stroke}%
\end{pgfscope}%
\begin{pgfscope}%
\pgfsetrectcap%
\pgfsetmiterjoin%
\pgfsetlinewidth{0.602250pt}%
\definecolor{currentstroke}{rgb}{0.000000,0.000000,0.000000}%
\pgfsetstrokecolor{currentstroke}%
\pgfsetdash{}{0pt}%
\pgfpathmoveto{\pgfqpoint{0.544111in}{2.299229in}}%
\pgfpathlineto{\pgfqpoint{3.716961in}{2.299229in}}%
\pgfusepath{stroke}%
\end{pgfscope}%
\begin{pgfscope}%
\pgfpathrectangle{\pgfqpoint{0.544111in}{0.352669in}}{\pgfqpoint{3.172850in}{1.946560in}}%
\pgfusepath{clip}%
\pgfsetrectcap%
\pgfsetroundjoin%
\pgfsetlinewidth{1.405250pt}%
\definecolor{currentstroke}{rgb}{0.768627,0.556863,0.211765}%
\pgfsetstrokecolor{currentstroke}%
\pgfsetdash{}{0pt}%
\pgfpathmoveto{\pgfqpoint{0.544111in}{2.136707in}}%
\pgfpathlineto{\pgfqpoint{0.786920in}{2.154354in}}%
\pgfpathlineto{\pgfqpoint{0.959196in}{2.153779in}}%
\pgfpathlineto{\pgfqpoint{1.092824in}{2.120553in}}%
\pgfpathlineto{\pgfqpoint{1.202005in}{2.115440in}}%
\pgfpathlineto{\pgfqpoint{1.294317in}{2.097680in}}%
\pgfpathlineto{\pgfqpoint{1.374281in}{2.049832in}}%
\pgfpathlineto{\pgfqpoint{1.444815in}{2.056335in}}%
\pgfpathlineto{\pgfqpoint{1.507909in}{2.012255in}}%
\pgfpathlineto{\pgfqpoint{1.564984in}{1.976881in}}%
\pgfpathlineto{\pgfqpoint{1.617090in}{1.979265in}}%
\pgfpathlineto{\pgfqpoint{1.665023in}{1.923845in}}%
\pgfpathlineto{\pgfqpoint{1.709402in}{1.911399in}}%
\pgfpathlineto{\pgfqpoint{1.750718in}{1.883984in}}%
\pgfpathlineto{\pgfqpoint{1.789366in}{1.834586in}}%
\pgfpathlineto{\pgfqpoint{1.825671in}{1.828714in}}%
\pgfpathlineto{\pgfqpoint{1.859900in}{1.771555in}}%
\pgfpathlineto{\pgfqpoint{1.892277in}{1.743389in}}%
\pgfpathlineto{\pgfqpoint{1.922994in}{1.728351in}}%
\pgfpathlineto{\pgfqpoint{1.952211in}{1.672564in}}%
\pgfpathlineto{\pgfqpoint{1.980069in}{1.659977in}}%
\pgfpathlineto{\pgfqpoint{2.006689in}{1.607319in}}%
\pgfpathlineto{\pgfqpoint{2.032175in}{1.564019in}}%
\pgfpathlineto{\pgfqpoint{2.056621in}{1.556373in}}%
\pgfpathlineto{\pgfqpoint{2.080108in}{1.497232in}}%
\pgfpathlineto{\pgfqpoint{2.102709in}{1.479182in}}%
\pgfpathlineto{\pgfqpoint{2.124487in}{1.432769in}}%
\pgfpathlineto{\pgfqpoint{2.145501in}{1.384244in}}%
\pgfpathlineto{\pgfqpoint{2.165803in}{1.373003in}}%
\pgfpathlineto{\pgfqpoint{2.185439in}{1.312078in}}%
\pgfpathlineto{\pgfqpoint{2.204451in}{1.289684in}}%
\pgfpathlineto{\pgfqpoint{2.222879in}{1.246460in}}%
\pgfpathlineto{\pgfqpoint{2.240756in}{1.194898in}}%
\pgfpathlineto{\pgfqpoint{2.258115in}{1.180980in}}%
\pgfpathlineto{\pgfqpoint{2.274985in}{1.119185in}}%
\pgfpathlineto{\pgfqpoint{2.291392in}{1.093503in}}%
\pgfpathlineto{\pgfqpoint{2.307362in}{1.051152in}}%
\pgfpathlineto{\pgfqpoint{2.322917in}{0.997761in}}%
\pgfpathlineto{\pgfqpoint{2.338079in}{0.982177in}}%
\pgfusepath{stroke}%
\end{pgfscope}%
\begin{pgfscope}%
\definecolor{textcolor}{rgb}{0.431373,0.423529,0.462745}%
\pgfsetstrokecolor{textcolor}%
\pgfsetfillcolor{textcolor}%
\pgftext[x=2.258115in,y=0.441731in,right,bottom]{\color{textcolor}{\sffamily\fontsize{7.000000}{8.400000}\selectfont\catcode`\^=\active\def^{\ifmmode\sp\else\^{}\fi}\catcode`\%=\active\def%{\%}$\epsilon_{\mathrm{m}}$}}%
\end{pgfscope}%
\begin{pgfscope}%
\pgfsetrectcap%
\pgfsetroundjoin%
\pgfsetlinewidth{1.405250pt}%
\definecolor{currentstroke}{rgb}{0.768627,0.556863,0.211765}%
\pgfsetstrokecolor{currentstroke}%
\pgfsetdash{}{0pt}%
\pgfpathmoveto{\pgfqpoint{2.661436in}{2.181613in}}%
\pgfpathlineto{\pgfqpoint{2.758658in}{2.181613in}}%
\pgfpathlineto{\pgfqpoint{2.855880in}{2.181613in}}%
\pgfusepath{stroke}%
\end{pgfscope}%
\begin{pgfscope}%
\definecolor{textcolor}{rgb}{0.000000,0.000000,0.000000}%
\pgfsetstrokecolor{textcolor}%
\pgfsetfillcolor{textcolor}%
\pgftext[x=2.933658in,y=2.147585in,left,base]{\color{textcolor}{\sffamily\fontsize{7.000000}{8.400000}\selectfont\catcode`\^=\active\def^{\ifmmode\sp\else\^{}\fi}\catcode`\%=\active\def%{\%}Spectral (CGL)}}%
\end{pgfscope}%
\begin{pgfscope}%
\pgfsetrectcap%
\pgfsetroundjoin%
\pgfsetlinewidth{1.204500pt}%
\definecolor{currentstroke}{rgb}{0.250980,0.400000,0.600000}%
\pgfsetstrokecolor{currentstroke}%
\pgfsetdash{}{0pt}%
\pgfpathmoveto{\pgfqpoint{2.661436in}{2.029946in}}%
\pgfpathlineto{\pgfqpoint{2.758658in}{2.029946in}}%
\pgfpathlineto{\pgfqpoint{2.855880in}{2.029946in}}%
\pgfusepath{stroke}%
\end{pgfscope}%
\begin{pgfscope}%
\definecolor{textcolor}{rgb}{0.000000,0.000000,0.000000}%
\pgfsetstrokecolor{textcolor}%
\pgfsetfillcolor{textcolor}%
\pgftext[x=2.933658in,y=1.995918in,left,base]{\color{textcolor}{\sffamily\fontsize{7.000000}{8.400000}\selectfont\catcode`\^=\active\def^{\ifmmode\sp\else\^{}\fi}\catcode`\%=\active\def%{\%}FD, $\mathcal{O}(h^2)$}}%
\end{pgfscope}%
\begin{pgfscope}%
\pgfsetbuttcap%
\pgfsetroundjoin%
\pgfsetlinewidth{1.204500pt}%
\definecolor{currentstroke}{rgb}{0.431373,0.423529,0.462745}%
\pgfsetstrokecolor{currentstroke}%
\pgfsetdash{{4.440000pt}{1.920000pt}}{0.000000pt}%
\pgfpathmoveto{\pgfqpoint{2.661436in}{1.878557in}}%
\pgfpathlineto{\pgfqpoint{2.758658in}{1.878557in}}%
\pgfpathlineto{\pgfqpoint{2.855880in}{1.878557in}}%
\pgfusepath{stroke}%
\end{pgfscope}%
\begin{pgfscope}%
\definecolor{textcolor}{rgb}{0.000000,0.000000,0.000000}%
\pgfsetstrokecolor{textcolor}%
\pgfsetfillcolor{textcolor}%
\pgftext[x=2.933658in,y=1.844529in,left,base]{\color{textcolor}{\sffamily\fontsize{7.000000}{8.400000}\selectfont\catcode`\^=\active\def^{\ifmmode\sp\else\^{}\fi}\catcode`\%=\active\def%{\%}FD, $\mathcal{O}(h^4)$}}%
\end{pgfscope}%
\end{pgfpicture}%
\makeatother%
\endgroup%

  \caption[Spectral vs.\ finite-difference differentiation convergence]{%
    Supremum-norm error in the derivative of
    $f(x) = e^{\sin(\pi x)}$ on $[-1,1]$.  The spectral
    differentiation matrix on CGL nodes (gold) achieves exponential
    convergence, reaching machine precision near $N = 20$.
    Second-order (blue) and fourth-order (grey, dashed) centred
    finite differences on equidistant grids converge only
    algebraically at rates $\order{h^2}$ and $\order{h^4}$,
    requiring hundreds of grid points to match the accuracy the
    spectral method achieves with twenty.}
  \label{fig:diffmat-convergence}
\end{figure}

\paragraph{The codebase implementations.}

The \texttt{chebyshev\_basis} constructor in \codemodule{spectral}
builds the CGL differentiation matrix using
\cref{eq:diffmat-cgl} for the off-diagonal entries and
\cref{eq:diffmat-diag} for the diagonal.  The second-derivative
matrix is computed as the matrix square $D^{(2)} = D \cdot D$:
%
\begin{lstlisting}[language=Rust]
// Off-diagonal: D[i,j] = (c_i/c_j) * (-1)^{i+j} / (x_i - x_j)
let sign_ij = if (i + j) % 2 == 0 { 1.0 } else { -1.0 };
d1[i * n + j] = (c[i] / c[j]) * sign_ij / (nodes[i] - nodes[j]);

// Diagonal: negative row sum
d1[i * n + i] = -row_sum;

// D2 = D1 * D1
\end{lstlisting}
%
The DG solver in \codemodule{dg} uses a different but analogous
formula for the LGL differentiation matrix, with the Chebyshev
$c_i$ weights replaced by evaluations of the Legendre polynomial
$P_N$ at each node \cite{HesthavenWarburton2008}.  Despite the
different node distributions, both constructions share the same
row-sum diagonal, the same $\order{N^2}$ conditioning, and the
same exponential convergence for smooth data --- the spectral
differentiation matrix is a general tool, not tied to a particular
polynomial family.
\xrefnote{\Cref{sec:spec-legendre} develops the Legendre polynomial
theory behind the LGL nodes, and \cref{sec:spec-lgl} details the
quadrature weights that accompany them.}

% sec:spec-legendre — Legendre Polynomials and Gaussian Quadrature
\section{Legendre Polynomials and Gauss Quadrature}
\label{sec:spec-legendre}

The Chebyshev polynomials of \cref{sec:spec-chebyshev} carry a weight
function $(1 - x^2)^{-1/2}$ that diverges at the endpoints --- natural
for the cosine change of variable, but awkward for a DG solver that
needs unweighted $L^2$ inner products to build mass and stiffness
matrices.  The Legendre polynomials are the orthogonal family for the
simplest possible weight: $w(x) = 1$.  Their zeros and extrema define
the Gauss and Gauss--Lobatto quadrature rules that underpin every
integral evaluation in the DG codebase, and the exactness of these
rules for polynomial integrands is what makes spectral-element methods
work.
\coderef{dg}{rs}

\paragraph{The Bonnet recurrence.}

To use Legendre polynomials in a numerical code, we need a stable way
to evaluate $P_n(x)$ at arbitrary points without computing
coefficients explicitly.  The Legendre polynomials satisfy the
three-term recurrence
%
\begin{equation}
  (n+1)\,P_{n+1}(x) = (2n+1)\,x\,P_n(x) - n\,P_{n-1}(x) \,,
  \label{eq:legendre-recurrence}
\end{equation}
%
with $P_0(x) = 1$ and $P_1(x) = x$.  The first few are
$P_2(x) = (3x^2 - 1)/2$ and $P_3(x) = (5x^3 - 3x)/2$.  Each $P_n$
is a polynomial of degree exactly $n$ with $P_n(1) = 1$ and
$P_n(-1) = (-1)^n$; the normalisation at $x = 1$ fixes the leading
coefficient and distinguishes the Legendre convention from the monic
convention used in the Chebyshev minimax theorem
(\cref{eq:chebyshev-minimax}).  The recurrence is the standard method
for evaluating $P_n(x)$: starting from $P_0 = 1$ and $P_1 = x$, each
step costs two multiplies and a subtract, accumulating the polynomial
value from degree $0$ up to $n$.
\implnote{The codebase's \texttt{legendre\_and\_deriv} function in
\codemodule{dg} evaluates $P_n(x)$ and $P_n'(x)$ simultaneously
by differentiating the Bonnet recurrence
(\cref{eq:legendre-recurrence}): each step updates $P_k'$
alongside $P_k$, accumulating both values from degree $0$ up to $n$.
Both are needed at each Newton step when computing LGL nodes.}

\begin{figure}[t]
\centering
%% Creator: Matplotlib, PGF backend
%%
%% To include the figure in your LaTeX document, write
%%   \input{<filename>.pgf}
%%
%% Make sure the required packages are loaded in your preamble
%%   \usepackage{pgf}
%%
%% Also ensure that all the required font packages are loaded; for instance,
%% the lmodern package is sometimes necessary when using math font.
%%   \usepackage{lmodern}
%%
%% Figures using additional raster images can only be included by \input if
%% they are in the same directory as the main LaTeX file. For loading figures
%% from other directories you can use the `import` package
%%   \usepackage{import}
%%
%% and then include the figures with
%%   \import{<path to file>}{<filename>.pgf}
%%
%% Matplotlib used the following preamble
%%   \def\mathdefault#1{#1}
%%   \everymath=\expandafter{\the\everymath\displaystyle}
%%   \IfFileExists{scrextend.sty}{
%%     \usepackage[fontsize=10.000000pt]{scrextend}
%%   }{
%%     \renewcommand{\normalsize}{\fontsize{10.000000}{12.000000}\selectfont}
%%     \normalsize
%%   }
%%   \usepackage{fontspec}
%%   \usepackage{unicode-math}
%%   \setmainfont{STIX Two Text}[Numbers=OldStyle, Ligatures=TeX]
%%   \setmathfont{STIX Two Math}
%%   \setmonofont{Source Code Pro}[Scale=MatchLowercase]
%%   \ifdefined\pdftexversion\else  % non-pdftex case.
%%     \usepackage{fontspec}
%%     \setmainfont{DejaVuSerif.ttf}[Path=\detokenize{/Users/gvn/.pyenv/versions/3.12.11/lib/python3.12/site-packages/matplotlib/mpl-data/fonts/ttf/}]
%%     \setsansfont{DejaVuSans.ttf}[Path=\detokenize{/Users/gvn/.pyenv/versions/3.12.11/lib/python3.12/site-packages/matplotlib/mpl-data/fonts/ttf/}]
%%     \setmonofont{DejaVuSansMono.ttf}[Path=\detokenize{/Users/gvn/.pyenv/versions/3.12.11/lib/python3.12/site-packages/matplotlib/mpl-data/fonts/ttf/}]
%%   \fi
%%   \makeatletter\@ifpackageloaded{underscore}{}{\usepackage[strings]{underscore}}\makeatother
%%
\begingroup%
\makeatletter%
\begin{pgfpicture}%
\pgfpathrectangle{\pgfpointorigin}{\pgfqpoint{3.824168in}{2.297669in}}%
\pgfusepath{use as bounding box, clip}%
\begin{pgfscope}%
\pgfsetbuttcap%
\pgfsetmiterjoin%
\definecolor{currentfill}{rgb}{1.000000,1.000000,1.000000}%
\pgfsetfillcolor{currentfill}%
\pgfsetlinewidth{0.000000pt}%
\definecolor{currentstroke}{rgb}{1.000000,1.000000,1.000000}%
\pgfsetstrokecolor{currentstroke}%
\pgfsetdash{}{0pt}%
\pgfpathmoveto{\pgfqpoint{0.000000in}{0.000000in}}%
\pgfpathlineto{\pgfqpoint{3.824168in}{0.000000in}}%
\pgfpathlineto{\pgfqpoint{3.824168in}{2.297669in}}%
\pgfpathlineto{\pgfqpoint{0.000000in}{2.297669in}}%
\pgfpathlineto{\pgfqpoint{0.000000in}{0.000000in}}%
\pgfpathclose%
\pgfusepath{fill}%
\end{pgfscope}%
\begin{pgfscope}%
\pgfsetbuttcap%
\pgfsetmiterjoin%
\definecolor{currentfill}{rgb}{1.000000,1.000000,1.000000}%
\pgfsetfillcolor{currentfill}%
\pgfsetlinewidth{0.000000pt}%
\definecolor{currentstroke}{rgb}{0.000000,0.000000,0.000000}%
\pgfsetstrokecolor{currentstroke}%
\pgfsetstrokeopacity{0.000000}%
\pgfsetdash{}{0pt}%
\pgfpathmoveto{\pgfqpoint{0.507620in}{0.352669in}}%
\pgfpathlineto{\pgfqpoint{3.680470in}{0.352669in}}%
\pgfpathlineto{\pgfqpoint{3.680470in}{2.277669in}}%
\pgfpathlineto{\pgfqpoint{0.507620in}{2.277669in}}%
\pgfpathlineto{\pgfqpoint{0.507620in}{0.352669in}}%
\pgfpathclose%
\pgfusepath{fill}%
\end{pgfscope}%
\begin{pgfscope}%
\pgfpathrectangle{\pgfqpoint{0.507620in}{0.352669in}}{\pgfqpoint{3.172850in}{1.925000in}}%
\pgfusepath{clip}%
\pgfsetrectcap%
\pgfsetroundjoin%
\pgfsetlinewidth{0.301125pt}%
\definecolor{currentstroke}{rgb}{0.000000,0.000000,0.000000}%
\pgfsetstrokecolor{currentstroke}%
\pgfsetdash{}{0pt}%
\pgfpathmoveto{\pgfqpoint{0.507620in}{1.315169in}}%
\pgfpathlineto{\pgfqpoint{3.680470in}{1.315169in}}%
\pgfusepath{stroke}%
\end{pgfscope}%
\begin{pgfscope}%
\pgfsetbuttcap%
\pgfsetroundjoin%
\definecolor{currentfill}{rgb}{0.000000,0.000000,0.000000}%
\pgfsetfillcolor{currentfill}%
\pgfsetlinewidth{0.501875pt}%
\definecolor{currentstroke}{rgb}{0.000000,0.000000,0.000000}%
\pgfsetstrokecolor{currentstroke}%
\pgfsetdash{}{0pt}%
\pgfsys@defobject{currentmarker}{\pgfqpoint{0.000000in}{0.000000in}}{\pgfqpoint{0.000000in}{0.041667in}}{%
\pgfpathmoveto{\pgfqpoint{0.000000in}{0.000000in}}%
\pgfpathlineto{\pgfqpoint{0.000000in}{0.041667in}}%
\pgfusepath{stroke,fill}%
}%
\begin{pgfscope}%
\pgfsys@transformshift{0.507620in}{0.352669in}%
\pgfsys@useobject{currentmarker}{}%
\end{pgfscope}%
\end{pgfscope}%
\begin{pgfscope}%
\definecolor{textcolor}{rgb}{0.000000,0.000000,0.000000}%
\pgfsetstrokecolor{textcolor}%
\pgfsetfillcolor{textcolor}%
\pgftext[x=0.507620in,y=0.304058in,,top]{\color{textcolor}{\sffamily\fontsize{8.000000}{9.600000}\selectfont\catcode`\^=\active\def^{\ifmmode\sp\else\^{}\fi}\catcode`\%=\active\def%{\%}\ensuremath{-}1.00}}%
\end{pgfscope}%
\begin{pgfscope}%
\pgfsetbuttcap%
\pgfsetroundjoin%
\definecolor{currentfill}{rgb}{0.000000,0.000000,0.000000}%
\pgfsetfillcolor{currentfill}%
\pgfsetlinewidth{0.501875pt}%
\definecolor{currentstroke}{rgb}{0.000000,0.000000,0.000000}%
\pgfsetstrokecolor{currentstroke}%
\pgfsetdash{}{0pt}%
\pgfsys@defobject{currentmarker}{\pgfqpoint{0.000000in}{0.000000in}}{\pgfqpoint{0.000000in}{0.041667in}}{%
\pgfpathmoveto{\pgfqpoint{0.000000in}{0.000000in}}%
\pgfpathlineto{\pgfqpoint{0.000000in}{0.041667in}}%
\pgfusepath{stroke,fill}%
}%
\begin{pgfscope}%
\pgfsys@transformshift{0.904226in}{0.352669in}%
\pgfsys@useobject{currentmarker}{}%
\end{pgfscope}%
\end{pgfscope}%
\begin{pgfscope}%
\definecolor{textcolor}{rgb}{0.000000,0.000000,0.000000}%
\pgfsetstrokecolor{textcolor}%
\pgfsetfillcolor{textcolor}%
\pgftext[x=0.904226in,y=0.304058in,,top]{\color{textcolor}{\sffamily\fontsize{8.000000}{9.600000}\selectfont\catcode`\^=\active\def^{\ifmmode\sp\else\^{}\fi}\catcode`\%=\active\def%{\%}\ensuremath{-}0.75}}%
\end{pgfscope}%
\begin{pgfscope}%
\pgfsetbuttcap%
\pgfsetroundjoin%
\definecolor{currentfill}{rgb}{0.000000,0.000000,0.000000}%
\pgfsetfillcolor{currentfill}%
\pgfsetlinewidth{0.501875pt}%
\definecolor{currentstroke}{rgb}{0.000000,0.000000,0.000000}%
\pgfsetstrokecolor{currentstroke}%
\pgfsetdash{}{0pt}%
\pgfsys@defobject{currentmarker}{\pgfqpoint{0.000000in}{0.000000in}}{\pgfqpoint{0.000000in}{0.041667in}}{%
\pgfpathmoveto{\pgfqpoint{0.000000in}{0.000000in}}%
\pgfpathlineto{\pgfqpoint{0.000000in}{0.041667in}}%
\pgfusepath{stroke,fill}%
}%
\begin{pgfscope}%
\pgfsys@transformshift{1.300832in}{0.352669in}%
\pgfsys@useobject{currentmarker}{}%
\end{pgfscope}%
\end{pgfscope}%
\begin{pgfscope}%
\definecolor{textcolor}{rgb}{0.000000,0.000000,0.000000}%
\pgfsetstrokecolor{textcolor}%
\pgfsetfillcolor{textcolor}%
\pgftext[x=1.300832in,y=0.304058in,,top]{\color{textcolor}{\sffamily\fontsize{8.000000}{9.600000}\selectfont\catcode`\^=\active\def^{\ifmmode\sp\else\^{}\fi}\catcode`\%=\active\def%{\%}\ensuremath{-}0.50}}%
\end{pgfscope}%
\begin{pgfscope}%
\pgfsetbuttcap%
\pgfsetroundjoin%
\definecolor{currentfill}{rgb}{0.000000,0.000000,0.000000}%
\pgfsetfillcolor{currentfill}%
\pgfsetlinewidth{0.501875pt}%
\definecolor{currentstroke}{rgb}{0.000000,0.000000,0.000000}%
\pgfsetstrokecolor{currentstroke}%
\pgfsetdash{}{0pt}%
\pgfsys@defobject{currentmarker}{\pgfqpoint{0.000000in}{0.000000in}}{\pgfqpoint{0.000000in}{0.041667in}}{%
\pgfpathmoveto{\pgfqpoint{0.000000in}{0.000000in}}%
\pgfpathlineto{\pgfqpoint{0.000000in}{0.041667in}}%
\pgfusepath{stroke,fill}%
}%
\begin{pgfscope}%
\pgfsys@transformshift{1.697438in}{0.352669in}%
\pgfsys@useobject{currentmarker}{}%
\end{pgfscope}%
\end{pgfscope}%
\begin{pgfscope}%
\definecolor{textcolor}{rgb}{0.000000,0.000000,0.000000}%
\pgfsetstrokecolor{textcolor}%
\pgfsetfillcolor{textcolor}%
\pgftext[x=1.697438in,y=0.304058in,,top]{\color{textcolor}{\sffamily\fontsize{8.000000}{9.600000}\selectfont\catcode`\^=\active\def^{\ifmmode\sp\else\^{}\fi}\catcode`\%=\active\def%{\%}\ensuremath{-}0.25}}%
\end{pgfscope}%
\begin{pgfscope}%
\pgfsetbuttcap%
\pgfsetroundjoin%
\definecolor{currentfill}{rgb}{0.000000,0.000000,0.000000}%
\pgfsetfillcolor{currentfill}%
\pgfsetlinewidth{0.501875pt}%
\definecolor{currentstroke}{rgb}{0.000000,0.000000,0.000000}%
\pgfsetstrokecolor{currentstroke}%
\pgfsetdash{}{0pt}%
\pgfsys@defobject{currentmarker}{\pgfqpoint{0.000000in}{0.000000in}}{\pgfqpoint{0.000000in}{0.041667in}}{%
\pgfpathmoveto{\pgfqpoint{0.000000in}{0.000000in}}%
\pgfpathlineto{\pgfqpoint{0.000000in}{0.041667in}}%
\pgfusepath{stroke,fill}%
}%
\begin{pgfscope}%
\pgfsys@transformshift{2.094045in}{0.352669in}%
\pgfsys@useobject{currentmarker}{}%
\end{pgfscope}%
\end{pgfscope}%
\begin{pgfscope}%
\definecolor{textcolor}{rgb}{0.000000,0.000000,0.000000}%
\pgfsetstrokecolor{textcolor}%
\pgfsetfillcolor{textcolor}%
\pgftext[x=2.094045in,y=0.304058in,,top]{\color{textcolor}{\sffamily\fontsize{8.000000}{9.600000}\selectfont\catcode`\^=\active\def^{\ifmmode\sp\else\^{}\fi}\catcode`\%=\active\def%{\%}0.00}}%
\end{pgfscope}%
\begin{pgfscope}%
\pgfsetbuttcap%
\pgfsetroundjoin%
\definecolor{currentfill}{rgb}{0.000000,0.000000,0.000000}%
\pgfsetfillcolor{currentfill}%
\pgfsetlinewidth{0.501875pt}%
\definecolor{currentstroke}{rgb}{0.000000,0.000000,0.000000}%
\pgfsetstrokecolor{currentstroke}%
\pgfsetdash{}{0pt}%
\pgfsys@defobject{currentmarker}{\pgfqpoint{0.000000in}{0.000000in}}{\pgfqpoint{0.000000in}{0.041667in}}{%
\pgfpathmoveto{\pgfqpoint{0.000000in}{0.000000in}}%
\pgfpathlineto{\pgfqpoint{0.000000in}{0.041667in}}%
\pgfusepath{stroke,fill}%
}%
\begin{pgfscope}%
\pgfsys@transformshift{2.490651in}{0.352669in}%
\pgfsys@useobject{currentmarker}{}%
\end{pgfscope}%
\end{pgfscope}%
\begin{pgfscope}%
\definecolor{textcolor}{rgb}{0.000000,0.000000,0.000000}%
\pgfsetstrokecolor{textcolor}%
\pgfsetfillcolor{textcolor}%
\pgftext[x=2.490651in,y=0.304058in,,top]{\color{textcolor}{\sffamily\fontsize{8.000000}{9.600000}\selectfont\catcode`\^=\active\def^{\ifmmode\sp\else\^{}\fi}\catcode`\%=\active\def%{\%}0.25}}%
\end{pgfscope}%
\begin{pgfscope}%
\pgfsetbuttcap%
\pgfsetroundjoin%
\definecolor{currentfill}{rgb}{0.000000,0.000000,0.000000}%
\pgfsetfillcolor{currentfill}%
\pgfsetlinewidth{0.501875pt}%
\definecolor{currentstroke}{rgb}{0.000000,0.000000,0.000000}%
\pgfsetstrokecolor{currentstroke}%
\pgfsetdash{}{0pt}%
\pgfsys@defobject{currentmarker}{\pgfqpoint{0.000000in}{0.000000in}}{\pgfqpoint{0.000000in}{0.041667in}}{%
\pgfpathmoveto{\pgfqpoint{0.000000in}{0.000000in}}%
\pgfpathlineto{\pgfqpoint{0.000000in}{0.041667in}}%
\pgfusepath{stroke,fill}%
}%
\begin{pgfscope}%
\pgfsys@transformshift{2.887257in}{0.352669in}%
\pgfsys@useobject{currentmarker}{}%
\end{pgfscope}%
\end{pgfscope}%
\begin{pgfscope}%
\definecolor{textcolor}{rgb}{0.000000,0.000000,0.000000}%
\pgfsetstrokecolor{textcolor}%
\pgfsetfillcolor{textcolor}%
\pgftext[x=2.887257in,y=0.304058in,,top]{\color{textcolor}{\sffamily\fontsize{8.000000}{9.600000}\selectfont\catcode`\^=\active\def^{\ifmmode\sp\else\^{}\fi}\catcode`\%=\active\def%{\%}0.50}}%
\end{pgfscope}%
\begin{pgfscope}%
\pgfsetbuttcap%
\pgfsetroundjoin%
\definecolor{currentfill}{rgb}{0.000000,0.000000,0.000000}%
\pgfsetfillcolor{currentfill}%
\pgfsetlinewidth{0.501875pt}%
\definecolor{currentstroke}{rgb}{0.000000,0.000000,0.000000}%
\pgfsetstrokecolor{currentstroke}%
\pgfsetdash{}{0pt}%
\pgfsys@defobject{currentmarker}{\pgfqpoint{0.000000in}{0.000000in}}{\pgfqpoint{0.000000in}{0.041667in}}{%
\pgfpathmoveto{\pgfqpoint{0.000000in}{0.000000in}}%
\pgfpathlineto{\pgfqpoint{0.000000in}{0.041667in}}%
\pgfusepath{stroke,fill}%
}%
\begin{pgfscope}%
\pgfsys@transformshift{3.283863in}{0.352669in}%
\pgfsys@useobject{currentmarker}{}%
\end{pgfscope}%
\end{pgfscope}%
\begin{pgfscope}%
\definecolor{textcolor}{rgb}{0.000000,0.000000,0.000000}%
\pgfsetstrokecolor{textcolor}%
\pgfsetfillcolor{textcolor}%
\pgftext[x=3.283863in,y=0.304058in,,top]{\color{textcolor}{\sffamily\fontsize{8.000000}{9.600000}\selectfont\catcode`\^=\active\def^{\ifmmode\sp\else\^{}\fi}\catcode`\%=\active\def%{\%}0.75}}%
\end{pgfscope}%
\begin{pgfscope}%
\pgfsetbuttcap%
\pgfsetroundjoin%
\definecolor{currentfill}{rgb}{0.000000,0.000000,0.000000}%
\pgfsetfillcolor{currentfill}%
\pgfsetlinewidth{0.501875pt}%
\definecolor{currentstroke}{rgb}{0.000000,0.000000,0.000000}%
\pgfsetstrokecolor{currentstroke}%
\pgfsetdash{}{0pt}%
\pgfsys@defobject{currentmarker}{\pgfqpoint{0.000000in}{0.000000in}}{\pgfqpoint{0.000000in}{0.041667in}}{%
\pgfpathmoveto{\pgfqpoint{0.000000in}{0.000000in}}%
\pgfpathlineto{\pgfqpoint{0.000000in}{0.041667in}}%
\pgfusepath{stroke,fill}%
}%
\begin{pgfscope}%
\pgfsys@transformshift{3.680470in}{0.352669in}%
\pgfsys@useobject{currentmarker}{}%
\end{pgfscope}%
\end{pgfscope}%
\begin{pgfscope}%
\definecolor{textcolor}{rgb}{0.000000,0.000000,0.000000}%
\pgfsetstrokecolor{textcolor}%
\pgfsetfillcolor{textcolor}%
\pgftext[x=3.680470in,y=0.304058in,,top]{\color{textcolor}{\sffamily\fontsize{8.000000}{9.600000}\selectfont\catcode`\^=\active\def^{\ifmmode\sp\else\^{}\fi}\catcode`\%=\active\def%{\%}1.00}}%
\end{pgfscope}%
\begin{pgfscope}%
\pgfsetbuttcap%
\pgfsetroundjoin%
\definecolor{currentfill}{rgb}{0.000000,0.000000,0.000000}%
\pgfsetfillcolor{currentfill}%
\pgfsetlinewidth{0.301125pt}%
\definecolor{currentstroke}{rgb}{0.000000,0.000000,0.000000}%
\pgfsetstrokecolor{currentstroke}%
\pgfsetdash{}{0pt}%
\pgfsys@defobject{currentmarker}{\pgfqpoint{0.000000in}{0.000000in}}{\pgfqpoint{0.000000in}{0.027778in}}{%
\pgfpathmoveto{\pgfqpoint{0.000000in}{0.000000in}}%
\pgfpathlineto{\pgfqpoint{0.000000in}{0.027778in}}%
\pgfusepath{stroke,fill}%
}%
\begin{pgfscope}%
\pgfsys@transformshift{0.586941in}{0.352669in}%
\pgfsys@useobject{currentmarker}{}%
\end{pgfscope}%
\end{pgfscope}%
\begin{pgfscope}%
\pgfsetbuttcap%
\pgfsetroundjoin%
\definecolor{currentfill}{rgb}{0.000000,0.000000,0.000000}%
\pgfsetfillcolor{currentfill}%
\pgfsetlinewidth{0.301125pt}%
\definecolor{currentstroke}{rgb}{0.000000,0.000000,0.000000}%
\pgfsetstrokecolor{currentstroke}%
\pgfsetdash{}{0pt}%
\pgfsys@defobject{currentmarker}{\pgfqpoint{0.000000in}{0.000000in}}{\pgfqpoint{0.000000in}{0.027778in}}{%
\pgfpathmoveto{\pgfqpoint{0.000000in}{0.000000in}}%
\pgfpathlineto{\pgfqpoint{0.000000in}{0.027778in}}%
\pgfusepath{stroke,fill}%
}%
\begin{pgfscope}%
\pgfsys@transformshift{0.666262in}{0.352669in}%
\pgfsys@useobject{currentmarker}{}%
\end{pgfscope}%
\end{pgfscope}%
\begin{pgfscope}%
\pgfsetbuttcap%
\pgfsetroundjoin%
\definecolor{currentfill}{rgb}{0.000000,0.000000,0.000000}%
\pgfsetfillcolor{currentfill}%
\pgfsetlinewidth{0.301125pt}%
\definecolor{currentstroke}{rgb}{0.000000,0.000000,0.000000}%
\pgfsetstrokecolor{currentstroke}%
\pgfsetdash{}{0pt}%
\pgfsys@defobject{currentmarker}{\pgfqpoint{0.000000in}{0.000000in}}{\pgfqpoint{0.000000in}{0.027778in}}{%
\pgfpathmoveto{\pgfqpoint{0.000000in}{0.000000in}}%
\pgfpathlineto{\pgfqpoint{0.000000in}{0.027778in}}%
\pgfusepath{stroke,fill}%
}%
\begin{pgfscope}%
\pgfsys@transformshift{0.745583in}{0.352669in}%
\pgfsys@useobject{currentmarker}{}%
\end{pgfscope}%
\end{pgfscope}%
\begin{pgfscope}%
\pgfsetbuttcap%
\pgfsetroundjoin%
\definecolor{currentfill}{rgb}{0.000000,0.000000,0.000000}%
\pgfsetfillcolor{currentfill}%
\pgfsetlinewidth{0.301125pt}%
\definecolor{currentstroke}{rgb}{0.000000,0.000000,0.000000}%
\pgfsetstrokecolor{currentstroke}%
\pgfsetdash{}{0pt}%
\pgfsys@defobject{currentmarker}{\pgfqpoint{0.000000in}{0.000000in}}{\pgfqpoint{0.000000in}{0.027778in}}{%
\pgfpathmoveto{\pgfqpoint{0.000000in}{0.000000in}}%
\pgfpathlineto{\pgfqpoint{0.000000in}{0.027778in}}%
\pgfusepath{stroke,fill}%
}%
\begin{pgfscope}%
\pgfsys@transformshift{0.824905in}{0.352669in}%
\pgfsys@useobject{currentmarker}{}%
\end{pgfscope}%
\end{pgfscope}%
\begin{pgfscope}%
\pgfsetbuttcap%
\pgfsetroundjoin%
\definecolor{currentfill}{rgb}{0.000000,0.000000,0.000000}%
\pgfsetfillcolor{currentfill}%
\pgfsetlinewidth{0.301125pt}%
\definecolor{currentstroke}{rgb}{0.000000,0.000000,0.000000}%
\pgfsetstrokecolor{currentstroke}%
\pgfsetdash{}{0pt}%
\pgfsys@defobject{currentmarker}{\pgfqpoint{0.000000in}{0.000000in}}{\pgfqpoint{0.000000in}{0.027778in}}{%
\pgfpathmoveto{\pgfqpoint{0.000000in}{0.000000in}}%
\pgfpathlineto{\pgfqpoint{0.000000in}{0.027778in}}%
\pgfusepath{stroke,fill}%
}%
\begin{pgfscope}%
\pgfsys@transformshift{0.983547in}{0.352669in}%
\pgfsys@useobject{currentmarker}{}%
\end{pgfscope}%
\end{pgfscope}%
\begin{pgfscope}%
\pgfsetbuttcap%
\pgfsetroundjoin%
\definecolor{currentfill}{rgb}{0.000000,0.000000,0.000000}%
\pgfsetfillcolor{currentfill}%
\pgfsetlinewidth{0.301125pt}%
\definecolor{currentstroke}{rgb}{0.000000,0.000000,0.000000}%
\pgfsetstrokecolor{currentstroke}%
\pgfsetdash{}{0pt}%
\pgfsys@defobject{currentmarker}{\pgfqpoint{0.000000in}{0.000000in}}{\pgfqpoint{0.000000in}{0.027778in}}{%
\pgfpathmoveto{\pgfqpoint{0.000000in}{0.000000in}}%
\pgfpathlineto{\pgfqpoint{0.000000in}{0.027778in}}%
\pgfusepath{stroke,fill}%
}%
\begin{pgfscope}%
\pgfsys@transformshift{1.062868in}{0.352669in}%
\pgfsys@useobject{currentmarker}{}%
\end{pgfscope}%
\end{pgfscope}%
\begin{pgfscope}%
\pgfsetbuttcap%
\pgfsetroundjoin%
\definecolor{currentfill}{rgb}{0.000000,0.000000,0.000000}%
\pgfsetfillcolor{currentfill}%
\pgfsetlinewidth{0.301125pt}%
\definecolor{currentstroke}{rgb}{0.000000,0.000000,0.000000}%
\pgfsetstrokecolor{currentstroke}%
\pgfsetdash{}{0pt}%
\pgfsys@defobject{currentmarker}{\pgfqpoint{0.000000in}{0.000000in}}{\pgfqpoint{0.000000in}{0.027778in}}{%
\pgfpathmoveto{\pgfqpoint{0.000000in}{0.000000in}}%
\pgfpathlineto{\pgfqpoint{0.000000in}{0.027778in}}%
\pgfusepath{stroke,fill}%
}%
\begin{pgfscope}%
\pgfsys@transformshift{1.142190in}{0.352669in}%
\pgfsys@useobject{currentmarker}{}%
\end{pgfscope}%
\end{pgfscope}%
\begin{pgfscope}%
\pgfsetbuttcap%
\pgfsetroundjoin%
\definecolor{currentfill}{rgb}{0.000000,0.000000,0.000000}%
\pgfsetfillcolor{currentfill}%
\pgfsetlinewidth{0.301125pt}%
\definecolor{currentstroke}{rgb}{0.000000,0.000000,0.000000}%
\pgfsetstrokecolor{currentstroke}%
\pgfsetdash{}{0pt}%
\pgfsys@defobject{currentmarker}{\pgfqpoint{0.000000in}{0.000000in}}{\pgfqpoint{0.000000in}{0.027778in}}{%
\pgfpathmoveto{\pgfqpoint{0.000000in}{0.000000in}}%
\pgfpathlineto{\pgfqpoint{0.000000in}{0.027778in}}%
\pgfusepath{stroke,fill}%
}%
\begin{pgfscope}%
\pgfsys@transformshift{1.221511in}{0.352669in}%
\pgfsys@useobject{currentmarker}{}%
\end{pgfscope}%
\end{pgfscope}%
\begin{pgfscope}%
\pgfsetbuttcap%
\pgfsetroundjoin%
\definecolor{currentfill}{rgb}{0.000000,0.000000,0.000000}%
\pgfsetfillcolor{currentfill}%
\pgfsetlinewidth{0.301125pt}%
\definecolor{currentstroke}{rgb}{0.000000,0.000000,0.000000}%
\pgfsetstrokecolor{currentstroke}%
\pgfsetdash{}{0pt}%
\pgfsys@defobject{currentmarker}{\pgfqpoint{0.000000in}{0.000000in}}{\pgfqpoint{0.000000in}{0.027778in}}{%
\pgfpathmoveto{\pgfqpoint{0.000000in}{0.000000in}}%
\pgfpathlineto{\pgfqpoint{0.000000in}{0.027778in}}%
\pgfusepath{stroke,fill}%
}%
\begin{pgfscope}%
\pgfsys@transformshift{1.380153in}{0.352669in}%
\pgfsys@useobject{currentmarker}{}%
\end{pgfscope}%
\end{pgfscope}%
\begin{pgfscope}%
\pgfsetbuttcap%
\pgfsetroundjoin%
\definecolor{currentfill}{rgb}{0.000000,0.000000,0.000000}%
\pgfsetfillcolor{currentfill}%
\pgfsetlinewidth{0.301125pt}%
\definecolor{currentstroke}{rgb}{0.000000,0.000000,0.000000}%
\pgfsetstrokecolor{currentstroke}%
\pgfsetdash{}{0pt}%
\pgfsys@defobject{currentmarker}{\pgfqpoint{0.000000in}{0.000000in}}{\pgfqpoint{0.000000in}{0.027778in}}{%
\pgfpathmoveto{\pgfqpoint{0.000000in}{0.000000in}}%
\pgfpathlineto{\pgfqpoint{0.000000in}{0.027778in}}%
\pgfusepath{stroke,fill}%
}%
\begin{pgfscope}%
\pgfsys@transformshift{1.459475in}{0.352669in}%
\pgfsys@useobject{currentmarker}{}%
\end{pgfscope}%
\end{pgfscope}%
\begin{pgfscope}%
\pgfsetbuttcap%
\pgfsetroundjoin%
\definecolor{currentfill}{rgb}{0.000000,0.000000,0.000000}%
\pgfsetfillcolor{currentfill}%
\pgfsetlinewidth{0.301125pt}%
\definecolor{currentstroke}{rgb}{0.000000,0.000000,0.000000}%
\pgfsetstrokecolor{currentstroke}%
\pgfsetdash{}{0pt}%
\pgfsys@defobject{currentmarker}{\pgfqpoint{0.000000in}{0.000000in}}{\pgfqpoint{0.000000in}{0.027778in}}{%
\pgfpathmoveto{\pgfqpoint{0.000000in}{0.000000in}}%
\pgfpathlineto{\pgfqpoint{0.000000in}{0.027778in}}%
\pgfusepath{stroke,fill}%
}%
\begin{pgfscope}%
\pgfsys@transformshift{1.538796in}{0.352669in}%
\pgfsys@useobject{currentmarker}{}%
\end{pgfscope}%
\end{pgfscope}%
\begin{pgfscope}%
\pgfsetbuttcap%
\pgfsetroundjoin%
\definecolor{currentfill}{rgb}{0.000000,0.000000,0.000000}%
\pgfsetfillcolor{currentfill}%
\pgfsetlinewidth{0.301125pt}%
\definecolor{currentstroke}{rgb}{0.000000,0.000000,0.000000}%
\pgfsetstrokecolor{currentstroke}%
\pgfsetdash{}{0pt}%
\pgfsys@defobject{currentmarker}{\pgfqpoint{0.000000in}{0.000000in}}{\pgfqpoint{0.000000in}{0.027778in}}{%
\pgfpathmoveto{\pgfqpoint{0.000000in}{0.000000in}}%
\pgfpathlineto{\pgfqpoint{0.000000in}{0.027778in}}%
\pgfusepath{stroke,fill}%
}%
\begin{pgfscope}%
\pgfsys@transformshift{1.618117in}{0.352669in}%
\pgfsys@useobject{currentmarker}{}%
\end{pgfscope}%
\end{pgfscope}%
\begin{pgfscope}%
\pgfsetbuttcap%
\pgfsetroundjoin%
\definecolor{currentfill}{rgb}{0.000000,0.000000,0.000000}%
\pgfsetfillcolor{currentfill}%
\pgfsetlinewidth{0.301125pt}%
\definecolor{currentstroke}{rgb}{0.000000,0.000000,0.000000}%
\pgfsetstrokecolor{currentstroke}%
\pgfsetdash{}{0pt}%
\pgfsys@defobject{currentmarker}{\pgfqpoint{0.000000in}{0.000000in}}{\pgfqpoint{0.000000in}{0.027778in}}{%
\pgfpathmoveto{\pgfqpoint{0.000000in}{0.000000in}}%
\pgfpathlineto{\pgfqpoint{0.000000in}{0.027778in}}%
\pgfusepath{stroke,fill}%
}%
\begin{pgfscope}%
\pgfsys@transformshift{1.776760in}{0.352669in}%
\pgfsys@useobject{currentmarker}{}%
\end{pgfscope}%
\end{pgfscope}%
\begin{pgfscope}%
\pgfsetbuttcap%
\pgfsetroundjoin%
\definecolor{currentfill}{rgb}{0.000000,0.000000,0.000000}%
\pgfsetfillcolor{currentfill}%
\pgfsetlinewidth{0.301125pt}%
\definecolor{currentstroke}{rgb}{0.000000,0.000000,0.000000}%
\pgfsetstrokecolor{currentstroke}%
\pgfsetdash{}{0pt}%
\pgfsys@defobject{currentmarker}{\pgfqpoint{0.000000in}{0.000000in}}{\pgfqpoint{0.000000in}{0.027778in}}{%
\pgfpathmoveto{\pgfqpoint{0.000000in}{0.000000in}}%
\pgfpathlineto{\pgfqpoint{0.000000in}{0.027778in}}%
\pgfusepath{stroke,fill}%
}%
\begin{pgfscope}%
\pgfsys@transformshift{1.856081in}{0.352669in}%
\pgfsys@useobject{currentmarker}{}%
\end{pgfscope}%
\end{pgfscope}%
\begin{pgfscope}%
\pgfsetbuttcap%
\pgfsetroundjoin%
\definecolor{currentfill}{rgb}{0.000000,0.000000,0.000000}%
\pgfsetfillcolor{currentfill}%
\pgfsetlinewidth{0.301125pt}%
\definecolor{currentstroke}{rgb}{0.000000,0.000000,0.000000}%
\pgfsetstrokecolor{currentstroke}%
\pgfsetdash{}{0pt}%
\pgfsys@defobject{currentmarker}{\pgfqpoint{0.000000in}{0.000000in}}{\pgfqpoint{0.000000in}{0.027778in}}{%
\pgfpathmoveto{\pgfqpoint{0.000000in}{0.000000in}}%
\pgfpathlineto{\pgfqpoint{0.000000in}{0.027778in}}%
\pgfusepath{stroke,fill}%
}%
\begin{pgfscope}%
\pgfsys@transformshift{1.935402in}{0.352669in}%
\pgfsys@useobject{currentmarker}{}%
\end{pgfscope}%
\end{pgfscope}%
\begin{pgfscope}%
\pgfsetbuttcap%
\pgfsetroundjoin%
\definecolor{currentfill}{rgb}{0.000000,0.000000,0.000000}%
\pgfsetfillcolor{currentfill}%
\pgfsetlinewidth{0.301125pt}%
\definecolor{currentstroke}{rgb}{0.000000,0.000000,0.000000}%
\pgfsetstrokecolor{currentstroke}%
\pgfsetdash{}{0pt}%
\pgfsys@defobject{currentmarker}{\pgfqpoint{0.000000in}{0.000000in}}{\pgfqpoint{0.000000in}{0.027778in}}{%
\pgfpathmoveto{\pgfqpoint{0.000000in}{0.000000in}}%
\pgfpathlineto{\pgfqpoint{0.000000in}{0.027778in}}%
\pgfusepath{stroke,fill}%
}%
\begin{pgfscope}%
\pgfsys@transformshift{2.014723in}{0.352669in}%
\pgfsys@useobject{currentmarker}{}%
\end{pgfscope}%
\end{pgfscope}%
\begin{pgfscope}%
\pgfsetbuttcap%
\pgfsetroundjoin%
\definecolor{currentfill}{rgb}{0.000000,0.000000,0.000000}%
\pgfsetfillcolor{currentfill}%
\pgfsetlinewidth{0.301125pt}%
\definecolor{currentstroke}{rgb}{0.000000,0.000000,0.000000}%
\pgfsetstrokecolor{currentstroke}%
\pgfsetdash{}{0pt}%
\pgfsys@defobject{currentmarker}{\pgfqpoint{0.000000in}{0.000000in}}{\pgfqpoint{0.000000in}{0.027778in}}{%
\pgfpathmoveto{\pgfqpoint{0.000000in}{0.000000in}}%
\pgfpathlineto{\pgfqpoint{0.000000in}{0.027778in}}%
\pgfusepath{stroke,fill}%
}%
\begin{pgfscope}%
\pgfsys@transformshift{2.173366in}{0.352669in}%
\pgfsys@useobject{currentmarker}{}%
\end{pgfscope}%
\end{pgfscope}%
\begin{pgfscope}%
\pgfsetbuttcap%
\pgfsetroundjoin%
\definecolor{currentfill}{rgb}{0.000000,0.000000,0.000000}%
\pgfsetfillcolor{currentfill}%
\pgfsetlinewidth{0.301125pt}%
\definecolor{currentstroke}{rgb}{0.000000,0.000000,0.000000}%
\pgfsetstrokecolor{currentstroke}%
\pgfsetdash{}{0pt}%
\pgfsys@defobject{currentmarker}{\pgfqpoint{0.000000in}{0.000000in}}{\pgfqpoint{0.000000in}{0.027778in}}{%
\pgfpathmoveto{\pgfqpoint{0.000000in}{0.000000in}}%
\pgfpathlineto{\pgfqpoint{0.000000in}{0.027778in}}%
\pgfusepath{stroke,fill}%
}%
\begin{pgfscope}%
\pgfsys@transformshift{2.252687in}{0.352669in}%
\pgfsys@useobject{currentmarker}{}%
\end{pgfscope}%
\end{pgfscope}%
\begin{pgfscope}%
\pgfsetbuttcap%
\pgfsetroundjoin%
\definecolor{currentfill}{rgb}{0.000000,0.000000,0.000000}%
\pgfsetfillcolor{currentfill}%
\pgfsetlinewidth{0.301125pt}%
\definecolor{currentstroke}{rgb}{0.000000,0.000000,0.000000}%
\pgfsetstrokecolor{currentstroke}%
\pgfsetdash{}{0pt}%
\pgfsys@defobject{currentmarker}{\pgfqpoint{0.000000in}{0.000000in}}{\pgfqpoint{0.000000in}{0.027778in}}{%
\pgfpathmoveto{\pgfqpoint{0.000000in}{0.000000in}}%
\pgfpathlineto{\pgfqpoint{0.000000in}{0.027778in}}%
\pgfusepath{stroke,fill}%
}%
\begin{pgfscope}%
\pgfsys@transformshift{2.332008in}{0.352669in}%
\pgfsys@useobject{currentmarker}{}%
\end{pgfscope}%
\end{pgfscope}%
\begin{pgfscope}%
\pgfsetbuttcap%
\pgfsetroundjoin%
\definecolor{currentfill}{rgb}{0.000000,0.000000,0.000000}%
\pgfsetfillcolor{currentfill}%
\pgfsetlinewidth{0.301125pt}%
\definecolor{currentstroke}{rgb}{0.000000,0.000000,0.000000}%
\pgfsetstrokecolor{currentstroke}%
\pgfsetdash{}{0pt}%
\pgfsys@defobject{currentmarker}{\pgfqpoint{0.000000in}{0.000000in}}{\pgfqpoint{0.000000in}{0.027778in}}{%
\pgfpathmoveto{\pgfqpoint{0.000000in}{0.000000in}}%
\pgfpathlineto{\pgfqpoint{0.000000in}{0.027778in}}%
\pgfusepath{stroke,fill}%
}%
\begin{pgfscope}%
\pgfsys@transformshift{2.411330in}{0.352669in}%
\pgfsys@useobject{currentmarker}{}%
\end{pgfscope}%
\end{pgfscope}%
\begin{pgfscope}%
\pgfsetbuttcap%
\pgfsetroundjoin%
\definecolor{currentfill}{rgb}{0.000000,0.000000,0.000000}%
\pgfsetfillcolor{currentfill}%
\pgfsetlinewidth{0.301125pt}%
\definecolor{currentstroke}{rgb}{0.000000,0.000000,0.000000}%
\pgfsetstrokecolor{currentstroke}%
\pgfsetdash{}{0pt}%
\pgfsys@defobject{currentmarker}{\pgfqpoint{0.000000in}{0.000000in}}{\pgfqpoint{0.000000in}{0.027778in}}{%
\pgfpathmoveto{\pgfqpoint{0.000000in}{0.000000in}}%
\pgfpathlineto{\pgfqpoint{0.000000in}{0.027778in}}%
\pgfusepath{stroke,fill}%
}%
\begin{pgfscope}%
\pgfsys@transformshift{2.569972in}{0.352669in}%
\pgfsys@useobject{currentmarker}{}%
\end{pgfscope}%
\end{pgfscope}%
\begin{pgfscope}%
\pgfsetbuttcap%
\pgfsetroundjoin%
\definecolor{currentfill}{rgb}{0.000000,0.000000,0.000000}%
\pgfsetfillcolor{currentfill}%
\pgfsetlinewidth{0.301125pt}%
\definecolor{currentstroke}{rgb}{0.000000,0.000000,0.000000}%
\pgfsetstrokecolor{currentstroke}%
\pgfsetdash{}{0pt}%
\pgfsys@defobject{currentmarker}{\pgfqpoint{0.000000in}{0.000000in}}{\pgfqpoint{0.000000in}{0.027778in}}{%
\pgfpathmoveto{\pgfqpoint{0.000000in}{0.000000in}}%
\pgfpathlineto{\pgfqpoint{0.000000in}{0.027778in}}%
\pgfusepath{stroke,fill}%
}%
\begin{pgfscope}%
\pgfsys@transformshift{2.649293in}{0.352669in}%
\pgfsys@useobject{currentmarker}{}%
\end{pgfscope}%
\end{pgfscope}%
\begin{pgfscope}%
\pgfsetbuttcap%
\pgfsetroundjoin%
\definecolor{currentfill}{rgb}{0.000000,0.000000,0.000000}%
\pgfsetfillcolor{currentfill}%
\pgfsetlinewidth{0.301125pt}%
\definecolor{currentstroke}{rgb}{0.000000,0.000000,0.000000}%
\pgfsetstrokecolor{currentstroke}%
\pgfsetdash{}{0pt}%
\pgfsys@defobject{currentmarker}{\pgfqpoint{0.000000in}{0.000000in}}{\pgfqpoint{0.000000in}{0.027778in}}{%
\pgfpathmoveto{\pgfqpoint{0.000000in}{0.000000in}}%
\pgfpathlineto{\pgfqpoint{0.000000in}{0.027778in}}%
\pgfusepath{stroke,fill}%
}%
\begin{pgfscope}%
\pgfsys@transformshift{2.728615in}{0.352669in}%
\pgfsys@useobject{currentmarker}{}%
\end{pgfscope}%
\end{pgfscope}%
\begin{pgfscope}%
\pgfsetbuttcap%
\pgfsetroundjoin%
\definecolor{currentfill}{rgb}{0.000000,0.000000,0.000000}%
\pgfsetfillcolor{currentfill}%
\pgfsetlinewidth{0.301125pt}%
\definecolor{currentstroke}{rgb}{0.000000,0.000000,0.000000}%
\pgfsetstrokecolor{currentstroke}%
\pgfsetdash{}{0pt}%
\pgfsys@defobject{currentmarker}{\pgfqpoint{0.000000in}{0.000000in}}{\pgfqpoint{0.000000in}{0.027778in}}{%
\pgfpathmoveto{\pgfqpoint{0.000000in}{0.000000in}}%
\pgfpathlineto{\pgfqpoint{0.000000in}{0.027778in}}%
\pgfusepath{stroke,fill}%
}%
\begin{pgfscope}%
\pgfsys@transformshift{2.807936in}{0.352669in}%
\pgfsys@useobject{currentmarker}{}%
\end{pgfscope}%
\end{pgfscope}%
\begin{pgfscope}%
\pgfsetbuttcap%
\pgfsetroundjoin%
\definecolor{currentfill}{rgb}{0.000000,0.000000,0.000000}%
\pgfsetfillcolor{currentfill}%
\pgfsetlinewidth{0.301125pt}%
\definecolor{currentstroke}{rgb}{0.000000,0.000000,0.000000}%
\pgfsetstrokecolor{currentstroke}%
\pgfsetdash{}{0pt}%
\pgfsys@defobject{currentmarker}{\pgfqpoint{0.000000in}{0.000000in}}{\pgfqpoint{0.000000in}{0.027778in}}{%
\pgfpathmoveto{\pgfqpoint{0.000000in}{0.000000in}}%
\pgfpathlineto{\pgfqpoint{0.000000in}{0.027778in}}%
\pgfusepath{stroke,fill}%
}%
\begin{pgfscope}%
\pgfsys@transformshift{2.966578in}{0.352669in}%
\pgfsys@useobject{currentmarker}{}%
\end{pgfscope}%
\end{pgfscope}%
\begin{pgfscope}%
\pgfsetbuttcap%
\pgfsetroundjoin%
\definecolor{currentfill}{rgb}{0.000000,0.000000,0.000000}%
\pgfsetfillcolor{currentfill}%
\pgfsetlinewidth{0.301125pt}%
\definecolor{currentstroke}{rgb}{0.000000,0.000000,0.000000}%
\pgfsetstrokecolor{currentstroke}%
\pgfsetdash{}{0pt}%
\pgfsys@defobject{currentmarker}{\pgfqpoint{0.000000in}{0.000000in}}{\pgfqpoint{0.000000in}{0.027778in}}{%
\pgfpathmoveto{\pgfqpoint{0.000000in}{0.000000in}}%
\pgfpathlineto{\pgfqpoint{0.000000in}{0.027778in}}%
\pgfusepath{stroke,fill}%
}%
\begin{pgfscope}%
\pgfsys@transformshift{3.045900in}{0.352669in}%
\pgfsys@useobject{currentmarker}{}%
\end{pgfscope}%
\end{pgfscope}%
\begin{pgfscope}%
\pgfsetbuttcap%
\pgfsetroundjoin%
\definecolor{currentfill}{rgb}{0.000000,0.000000,0.000000}%
\pgfsetfillcolor{currentfill}%
\pgfsetlinewidth{0.301125pt}%
\definecolor{currentstroke}{rgb}{0.000000,0.000000,0.000000}%
\pgfsetstrokecolor{currentstroke}%
\pgfsetdash{}{0pt}%
\pgfsys@defobject{currentmarker}{\pgfqpoint{0.000000in}{0.000000in}}{\pgfqpoint{0.000000in}{0.027778in}}{%
\pgfpathmoveto{\pgfqpoint{0.000000in}{0.000000in}}%
\pgfpathlineto{\pgfqpoint{0.000000in}{0.027778in}}%
\pgfusepath{stroke,fill}%
}%
\begin{pgfscope}%
\pgfsys@transformshift{3.125221in}{0.352669in}%
\pgfsys@useobject{currentmarker}{}%
\end{pgfscope}%
\end{pgfscope}%
\begin{pgfscope}%
\pgfsetbuttcap%
\pgfsetroundjoin%
\definecolor{currentfill}{rgb}{0.000000,0.000000,0.000000}%
\pgfsetfillcolor{currentfill}%
\pgfsetlinewidth{0.301125pt}%
\definecolor{currentstroke}{rgb}{0.000000,0.000000,0.000000}%
\pgfsetstrokecolor{currentstroke}%
\pgfsetdash{}{0pt}%
\pgfsys@defobject{currentmarker}{\pgfqpoint{0.000000in}{0.000000in}}{\pgfqpoint{0.000000in}{0.027778in}}{%
\pgfpathmoveto{\pgfqpoint{0.000000in}{0.000000in}}%
\pgfpathlineto{\pgfqpoint{0.000000in}{0.027778in}}%
\pgfusepath{stroke,fill}%
}%
\begin{pgfscope}%
\pgfsys@transformshift{3.204542in}{0.352669in}%
\pgfsys@useobject{currentmarker}{}%
\end{pgfscope}%
\end{pgfscope}%
\begin{pgfscope}%
\pgfsetbuttcap%
\pgfsetroundjoin%
\definecolor{currentfill}{rgb}{0.000000,0.000000,0.000000}%
\pgfsetfillcolor{currentfill}%
\pgfsetlinewidth{0.301125pt}%
\definecolor{currentstroke}{rgb}{0.000000,0.000000,0.000000}%
\pgfsetstrokecolor{currentstroke}%
\pgfsetdash{}{0pt}%
\pgfsys@defobject{currentmarker}{\pgfqpoint{0.000000in}{0.000000in}}{\pgfqpoint{0.000000in}{0.027778in}}{%
\pgfpathmoveto{\pgfqpoint{0.000000in}{0.000000in}}%
\pgfpathlineto{\pgfqpoint{0.000000in}{0.027778in}}%
\pgfusepath{stroke,fill}%
}%
\begin{pgfscope}%
\pgfsys@transformshift{3.363185in}{0.352669in}%
\pgfsys@useobject{currentmarker}{}%
\end{pgfscope}%
\end{pgfscope}%
\begin{pgfscope}%
\pgfsetbuttcap%
\pgfsetroundjoin%
\definecolor{currentfill}{rgb}{0.000000,0.000000,0.000000}%
\pgfsetfillcolor{currentfill}%
\pgfsetlinewidth{0.301125pt}%
\definecolor{currentstroke}{rgb}{0.000000,0.000000,0.000000}%
\pgfsetstrokecolor{currentstroke}%
\pgfsetdash{}{0pt}%
\pgfsys@defobject{currentmarker}{\pgfqpoint{0.000000in}{0.000000in}}{\pgfqpoint{0.000000in}{0.027778in}}{%
\pgfpathmoveto{\pgfqpoint{0.000000in}{0.000000in}}%
\pgfpathlineto{\pgfqpoint{0.000000in}{0.027778in}}%
\pgfusepath{stroke,fill}%
}%
\begin{pgfscope}%
\pgfsys@transformshift{3.442506in}{0.352669in}%
\pgfsys@useobject{currentmarker}{}%
\end{pgfscope}%
\end{pgfscope}%
\begin{pgfscope}%
\pgfsetbuttcap%
\pgfsetroundjoin%
\definecolor{currentfill}{rgb}{0.000000,0.000000,0.000000}%
\pgfsetfillcolor{currentfill}%
\pgfsetlinewidth{0.301125pt}%
\definecolor{currentstroke}{rgb}{0.000000,0.000000,0.000000}%
\pgfsetstrokecolor{currentstroke}%
\pgfsetdash{}{0pt}%
\pgfsys@defobject{currentmarker}{\pgfqpoint{0.000000in}{0.000000in}}{\pgfqpoint{0.000000in}{0.027778in}}{%
\pgfpathmoveto{\pgfqpoint{0.000000in}{0.000000in}}%
\pgfpathlineto{\pgfqpoint{0.000000in}{0.027778in}}%
\pgfusepath{stroke,fill}%
}%
\begin{pgfscope}%
\pgfsys@transformshift{3.521827in}{0.352669in}%
\pgfsys@useobject{currentmarker}{}%
\end{pgfscope}%
\end{pgfscope}%
\begin{pgfscope}%
\pgfsetbuttcap%
\pgfsetroundjoin%
\definecolor{currentfill}{rgb}{0.000000,0.000000,0.000000}%
\pgfsetfillcolor{currentfill}%
\pgfsetlinewidth{0.301125pt}%
\definecolor{currentstroke}{rgb}{0.000000,0.000000,0.000000}%
\pgfsetstrokecolor{currentstroke}%
\pgfsetdash{}{0pt}%
\pgfsys@defobject{currentmarker}{\pgfqpoint{0.000000in}{0.000000in}}{\pgfqpoint{0.000000in}{0.027778in}}{%
\pgfpathmoveto{\pgfqpoint{0.000000in}{0.000000in}}%
\pgfpathlineto{\pgfqpoint{0.000000in}{0.027778in}}%
\pgfusepath{stroke,fill}%
}%
\begin{pgfscope}%
\pgfsys@transformshift{3.601148in}{0.352669in}%
\pgfsys@useobject{currentmarker}{}%
\end{pgfscope}%
\end{pgfscope}%
\begin{pgfscope}%
\definecolor{textcolor}{rgb}{0.000000,0.000000,0.000000}%
\pgfsetstrokecolor{textcolor}%
\pgfsetfillcolor{textcolor}%
\pgftext[x=2.094045in,y=0.140972in,,top]{\color{textcolor}{\sffamily\fontsize{9.000000}{10.800000}\selectfont\catcode`\^=\active\def^{\ifmmode\sp\else\^{}\fi}\catcode`\%=\active\def%{\%}$x$}}%
\end{pgfscope}%
\begin{pgfscope}%
\pgfsetbuttcap%
\pgfsetroundjoin%
\definecolor{currentfill}{rgb}{0.000000,0.000000,0.000000}%
\pgfsetfillcolor{currentfill}%
\pgfsetlinewidth{0.501875pt}%
\definecolor{currentstroke}{rgb}{0.000000,0.000000,0.000000}%
\pgfsetstrokecolor{currentstroke}%
\pgfsetdash{}{0pt}%
\pgfsys@defobject{currentmarker}{\pgfqpoint{0.000000in}{0.000000in}}{\pgfqpoint{0.041667in}{0.000000in}}{%
\pgfpathmoveto{\pgfqpoint{0.000000in}{0.000000in}}%
\pgfpathlineto{\pgfqpoint{0.041667in}{0.000000in}}%
\pgfusepath{stroke,fill}%
}%
\begin{pgfscope}%
\pgfsys@transformshift{0.507620in}{0.440169in}%
\pgfsys@useobject{currentmarker}{}%
\end{pgfscope}%
\end{pgfscope}%
\begin{pgfscope}%
\definecolor{textcolor}{rgb}{0.000000,0.000000,0.000000}%
\pgfsetstrokecolor{textcolor}%
\pgfsetfillcolor{textcolor}%
\pgftext[x=0.196527in, y=0.397960in, left, base]{\color{textcolor}{\sffamily\fontsize{8.000000}{9.600000}\selectfont\catcode`\^=\active\def^{\ifmmode\sp\else\^{}\fi}\catcode`\%=\active\def%{\%}\ensuremath{-}1.0}}%
\end{pgfscope}%
\begin{pgfscope}%
\pgfsetbuttcap%
\pgfsetroundjoin%
\definecolor{currentfill}{rgb}{0.000000,0.000000,0.000000}%
\pgfsetfillcolor{currentfill}%
\pgfsetlinewidth{0.501875pt}%
\definecolor{currentstroke}{rgb}{0.000000,0.000000,0.000000}%
\pgfsetstrokecolor{currentstroke}%
\pgfsetdash{}{0pt}%
\pgfsys@defobject{currentmarker}{\pgfqpoint{0.000000in}{0.000000in}}{\pgfqpoint{0.041667in}{0.000000in}}{%
\pgfpathmoveto{\pgfqpoint{0.000000in}{0.000000in}}%
\pgfpathlineto{\pgfqpoint{0.041667in}{0.000000in}}%
\pgfusepath{stroke,fill}%
}%
\begin{pgfscope}%
\pgfsys@transformshift{0.507620in}{0.877669in}%
\pgfsys@useobject{currentmarker}{}%
\end{pgfscope}%
\end{pgfscope}%
\begin{pgfscope}%
\definecolor{textcolor}{rgb}{0.000000,0.000000,0.000000}%
\pgfsetstrokecolor{textcolor}%
\pgfsetfillcolor{textcolor}%
\pgftext[x=0.196527in, y=0.835460in, left, base]{\color{textcolor}{\sffamily\fontsize{8.000000}{9.600000}\selectfont\catcode`\^=\active\def^{\ifmmode\sp\else\^{}\fi}\catcode`\%=\active\def%{\%}\ensuremath{-}0.5}}%
\end{pgfscope}%
\begin{pgfscope}%
\pgfsetbuttcap%
\pgfsetroundjoin%
\definecolor{currentfill}{rgb}{0.000000,0.000000,0.000000}%
\pgfsetfillcolor{currentfill}%
\pgfsetlinewidth{0.501875pt}%
\definecolor{currentstroke}{rgb}{0.000000,0.000000,0.000000}%
\pgfsetstrokecolor{currentstroke}%
\pgfsetdash{}{0pt}%
\pgfsys@defobject{currentmarker}{\pgfqpoint{0.000000in}{0.000000in}}{\pgfqpoint{0.041667in}{0.000000in}}{%
\pgfpathmoveto{\pgfqpoint{0.000000in}{0.000000in}}%
\pgfpathlineto{\pgfqpoint{0.041667in}{0.000000in}}%
\pgfusepath{stroke,fill}%
}%
\begin{pgfscope}%
\pgfsys@transformshift{0.507620in}{1.315169in}%
\pgfsys@useobject{currentmarker}{}%
\end{pgfscope}%
\end{pgfscope}%
\begin{pgfscope}%
\definecolor{textcolor}{rgb}{0.000000,0.000000,0.000000}%
\pgfsetstrokecolor{textcolor}%
\pgfsetfillcolor{textcolor}%
\pgftext[x=0.282305in, y=1.272960in, left, base]{\color{textcolor}{\sffamily\fontsize{8.000000}{9.600000}\selectfont\catcode`\^=\active\def^{\ifmmode\sp\else\^{}\fi}\catcode`\%=\active\def%{\%}0.0}}%
\end{pgfscope}%
\begin{pgfscope}%
\pgfsetbuttcap%
\pgfsetroundjoin%
\definecolor{currentfill}{rgb}{0.000000,0.000000,0.000000}%
\pgfsetfillcolor{currentfill}%
\pgfsetlinewidth{0.501875pt}%
\definecolor{currentstroke}{rgb}{0.000000,0.000000,0.000000}%
\pgfsetstrokecolor{currentstroke}%
\pgfsetdash{}{0pt}%
\pgfsys@defobject{currentmarker}{\pgfqpoint{0.000000in}{0.000000in}}{\pgfqpoint{0.041667in}{0.000000in}}{%
\pgfpathmoveto{\pgfqpoint{0.000000in}{0.000000in}}%
\pgfpathlineto{\pgfqpoint{0.041667in}{0.000000in}}%
\pgfusepath{stroke,fill}%
}%
\begin{pgfscope}%
\pgfsys@transformshift{0.507620in}{1.752669in}%
\pgfsys@useobject{currentmarker}{}%
\end{pgfscope}%
\end{pgfscope}%
\begin{pgfscope}%
\definecolor{textcolor}{rgb}{0.000000,0.000000,0.000000}%
\pgfsetstrokecolor{textcolor}%
\pgfsetfillcolor{textcolor}%
\pgftext[x=0.282305in, y=1.710460in, left, base]{\color{textcolor}{\sffamily\fontsize{8.000000}{9.600000}\selectfont\catcode`\^=\active\def^{\ifmmode\sp\else\^{}\fi}\catcode`\%=\active\def%{\%}0.5}}%
\end{pgfscope}%
\begin{pgfscope}%
\pgfsetbuttcap%
\pgfsetroundjoin%
\definecolor{currentfill}{rgb}{0.000000,0.000000,0.000000}%
\pgfsetfillcolor{currentfill}%
\pgfsetlinewidth{0.501875pt}%
\definecolor{currentstroke}{rgb}{0.000000,0.000000,0.000000}%
\pgfsetstrokecolor{currentstroke}%
\pgfsetdash{}{0pt}%
\pgfsys@defobject{currentmarker}{\pgfqpoint{0.000000in}{0.000000in}}{\pgfqpoint{0.041667in}{0.000000in}}{%
\pgfpathmoveto{\pgfqpoint{0.000000in}{0.000000in}}%
\pgfpathlineto{\pgfqpoint{0.041667in}{0.000000in}}%
\pgfusepath{stroke,fill}%
}%
\begin{pgfscope}%
\pgfsys@transformshift{0.507620in}{2.190169in}%
\pgfsys@useobject{currentmarker}{}%
\end{pgfscope}%
\end{pgfscope}%
\begin{pgfscope}%
\definecolor{textcolor}{rgb}{0.000000,0.000000,0.000000}%
\pgfsetstrokecolor{textcolor}%
\pgfsetfillcolor{textcolor}%
\pgftext[x=0.282305in, y=2.147960in, left, base]{\color{textcolor}{\sffamily\fontsize{8.000000}{9.600000}\selectfont\catcode`\^=\active\def^{\ifmmode\sp\else\^{}\fi}\catcode`\%=\active\def%{\%}1.0}}%
\end{pgfscope}%
\begin{pgfscope}%
\pgfsetbuttcap%
\pgfsetroundjoin%
\definecolor{currentfill}{rgb}{0.000000,0.000000,0.000000}%
\pgfsetfillcolor{currentfill}%
\pgfsetlinewidth{0.301125pt}%
\definecolor{currentstroke}{rgb}{0.000000,0.000000,0.000000}%
\pgfsetstrokecolor{currentstroke}%
\pgfsetdash{}{0pt}%
\pgfsys@defobject{currentmarker}{\pgfqpoint{0.000000in}{0.000000in}}{\pgfqpoint{0.027778in}{0.000000in}}{%
\pgfpathmoveto{\pgfqpoint{0.000000in}{0.000000in}}%
\pgfpathlineto{\pgfqpoint{0.027778in}{0.000000in}}%
\pgfusepath{stroke,fill}%
}%
\begin{pgfscope}%
\pgfsys@transformshift{0.507620in}{0.352669in}%
\pgfsys@useobject{currentmarker}{}%
\end{pgfscope}%
\end{pgfscope}%
\begin{pgfscope}%
\pgfsetbuttcap%
\pgfsetroundjoin%
\definecolor{currentfill}{rgb}{0.000000,0.000000,0.000000}%
\pgfsetfillcolor{currentfill}%
\pgfsetlinewidth{0.301125pt}%
\definecolor{currentstroke}{rgb}{0.000000,0.000000,0.000000}%
\pgfsetstrokecolor{currentstroke}%
\pgfsetdash{}{0pt}%
\pgfsys@defobject{currentmarker}{\pgfqpoint{0.000000in}{0.000000in}}{\pgfqpoint{0.027778in}{0.000000in}}{%
\pgfpathmoveto{\pgfqpoint{0.000000in}{0.000000in}}%
\pgfpathlineto{\pgfqpoint{0.027778in}{0.000000in}}%
\pgfusepath{stroke,fill}%
}%
\begin{pgfscope}%
\pgfsys@transformshift{0.507620in}{0.527669in}%
\pgfsys@useobject{currentmarker}{}%
\end{pgfscope}%
\end{pgfscope}%
\begin{pgfscope}%
\pgfsetbuttcap%
\pgfsetroundjoin%
\definecolor{currentfill}{rgb}{0.000000,0.000000,0.000000}%
\pgfsetfillcolor{currentfill}%
\pgfsetlinewidth{0.301125pt}%
\definecolor{currentstroke}{rgb}{0.000000,0.000000,0.000000}%
\pgfsetstrokecolor{currentstroke}%
\pgfsetdash{}{0pt}%
\pgfsys@defobject{currentmarker}{\pgfqpoint{0.000000in}{0.000000in}}{\pgfqpoint{0.027778in}{0.000000in}}{%
\pgfpathmoveto{\pgfqpoint{0.000000in}{0.000000in}}%
\pgfpathlineto{\pgfqpoint{0.027778in}{0.000000in}}%
\pgfusepath{stroke,fill}%
}%
\begin{pgfscope}%
\pgfsys@transformshift{0.507620in}{0.615169in}%
\pgfsys@useobject{currentmarker}{}%
\end{pgfscope}%
\end{pgfscope}%
\begin{pgfscope}%
\pgfsetbuttcap%
\pgfsetroundjoin%
\definecolor{currentfill}{rgb}{0.000000,0.000000,0.000000}%
\pgfsetfillcolor{currentfill}%
\pgfsetlinewidth{0.301125pt}%
\definecolor{currentstroke}{rgb}{0.000000,0.000000,0.000000}%
\pgfsetstrokecolor{currentstroke}%
\pgfsetdash{}{0pt}%
\pgfsys@defobject{currentmarker}{\pgfqpoint{0.000000in}{0.000000in}}{\pgfqpoint{0.027778in}{0.000000in}}{%
\pgfpathmoveto{\pgfqpoint{0.000000in}{0.000000in}}%
\pgfpathlineto{\pgfqpoint{0.027778in}{0.000000in}}%
\pgfusepath{stroke,fill}%
}%
\begin{pgfscope}%
\pgfsys@transformshift{0.507620in}{0.702669in}%
\pgfsys@useobject{currentmarker}{}%
\end{pgfscope}%
\end{pgfscope}%
\begin{pgfscope}%
\pgfsetbuttcap%
\pgfsetroundjoin%
\definecolor{currentfill}{rgb}{0.000000,0.000000,0.000000}%
\pgfsetfillcolor{currentfill}%
\pgfsetlinewidth{0.301125pt}%
\definecolor{currentstroke}{rgb}{0.000000,0.000000,0.000000}%
\pgfsetstrokecolor{currentstroke}%
\pgfsetdash{}{0pt}%
\pgfsys@defobject{currentmarker}{\pgfqpoint{0.000000in}{0.000000in}}{\pgfqpoint{0.027778in}{0.000000in}}{%
\pgfpathmoveto{\pgfqpoint{0.000000in}{0.000000in}}%
\pgfpathlineto{\pgfqpoint{0.027778in}{0.000000in}}%
\pgfusepath{stroke,fill}%
}%
\begin{pgfscope}%
\pgfsys@transformshift{0.507620in}{0.790169in}%
\pgfsys@useobject{currentmarker}{}%
\end{pgfscope}%
\end{pgfscope}%
\begin{pgfscope}%
\pgfsetbuttcap%
\pgfsetroundjoin%
\definecolor{currentfill}{rgb}{0.000000,0.000000,0.000000}%
\pgfsetfillcolor{currentfill}%
\pgfsetlinewidth{0.301125pt}%
\definecolor{currentstroke}{rgb}{0.000000,0.000000,0.000000}%
\pgfsetstrokecolor{currentstroke}%
\pgfsetdash{}{0pt}%
\pgfsys@defobject{currentmarker}{\pgfqpoint{0.000000in}{0.000000in}}{\pgfqpoint{0.027778in}{0.000000in}}{%
\pgfpathmoveto{\pgfqpoint{0.000000in}{0.000000in}}%
\pgfpathlineto{\pgfqpoint{0.027778in}{0.000000in}}%
\pgfusepath{stroke,fill}%
}%
\begin{pgfscope}%
\pgfsys@transformshift{0.507620in}{0.965169in}%
\pgfsys@useobject{currentmarker}{}%
\end{pgfscope}%
\end{pgfscope}%
\begin{pgfscope}%
\pgfsetbuttcap%
\pgfsetroundjoin%
\definecolor{currentfill}{rgb}{0.000000,0.000000,0.000000}%
\pgfsetfillcolor{currentfill}%
\pgfsetlinewidth{0.301125pt}%
\definecolor{currentstroke}{rgb}{0.000000,0.000000,0.000000}%
\pgfsetstrokecolor{currentstroke}%
\pgfsetdash{}{0pt}%
\pgfsys@defobject{currentmarker}{\pgfqpoint{0.000000in}{0.000000in}}{\pgfqpoint{0.027778in}{0.000000in}}{%
\pgfpathmoveto{\pgfqpoint{0.000000in}{0.000000in}}%
\pgfpathlineto{\pgfqpoint{0.027778in}{0.000000in}}%
\pgfusepath{stroke,fill}%
}%
\begin{pgfscope}%
\pgfsys@transformshift{0.507620in}{1.052669in}%
\pgfsys@useobject{currentmarker}{}%
\end{pgfscope}%
\end{pgfscope}%
\begin{pgfscope}%
\pgfsetbuttcap%
\pgfsetroundjoin%
\definecolor{currentfill}{rgb}{0.000000,0.000000,0.000000}%
\pgfsetfillcolor{currentfill}%
\pgfsetlinewidth{0.301125pt}%
\definecolor{currentstroke}{rgb}{0.000000,0.000000,0.000000}%
\pgfsetstrokecolor{currentstroke}%
\pgfsetdash{}{0pt}%
\pgfsys@defobject{currentmarker}{\pgfqpoint{0.000000in}{0.000000in}}{\pgfqpoint{0.027778in}{0.000000in}}{%
\pgfpathmoveto{\pgfqpoint{0.000000in}{0.000000in}}%
\pgfpathlineto{\pgfqpoint{0.027778in}{0.000000in}}%
\pgfusepath{stroke,fill}%
}%
\begin{pgfscope}%
\pgfsys@transformshift{0.507620in}{1.140169in}%
\pgfsys@useobject{currentmarker}{}%
\end{pgfscope}%
\end{pgfscope}%
\begin{pgfscope}%
\pgfsetbuttcap%
\pgfsetroundjoin%
\definecolor{currentfill}{rgb}{0.000000,0.000000,0.000000}%
\pgfsetfillcolor{currentfill}%
\pgfsetlinewidth{0.301125pt}%
\definecolor{currentstroke}{rgb}{0.000000,0.000000,0.000000}%
\pgfsetstrokecolor{currentstroke}%
\pgfsetdash{}{0pt}%
\pgfsys@defobject{currentmarker}{\pgfqpoint{0.000000in}{0.000000in}}{\pgfqpoint{0.027778in}{0.000000in}}{%
\pgfpathmoveto{\pgfqpoint{0.000000in}{0.000000in}}%
\pgfpathlineto{\pgfqpoint{0.027778in}{0.000000in}}%
\pgfusepath{stroke,fill}%
}%
\begin{pgfscope}%
\pgfsys@transformshift{0.507620in}{1.227669in}%
\pgfsys@useobject{currentmarker}{}%
\end{pgfscope}%
\end{pgfscope}%
\begin{pgfscope}%
\pgfsetbuttcap%
\pgfsetroundjoin%
\definecolor{currentfill}{rgb}{0.000000,0.000000,0.000000}%
\pgfsetfillcolor{currentfill}%
\pgfsetlinewidth{0.301125pt}%
\definecolor{currentstroke}{rgb}{0.000000,0.000000,0.000000}%
\pgfsetstrokecolor{currentstroke}%
\pgfsetdash{}{0pt}%
\pgfsys@defobject{currentmarker}{\pgfqpoint{0.000000in}{0.000000in}}{\pgfqpoint{0.027778in}{0.000000in}}{%
\pgfpathmoveto{\pgfqpoint{0.000000in}{0.000000in}}%
\pgfpathlineto{\pgfqpoint{0.027778in}{0.000000in}}%
\pgfusepath{stroke,fill}%
}%
\begin{pgfscope}%
\pgfsys@transformshift{0.507620in}{1.402669in}%
\pgfsys@useobject{currentmarker}{}%
\end{pgfscope}%
\end{pgfscope}%
\begin{pgfscope}%
\pgfsetbuttcap%
\pgfsetroundjoin%
\definecolor{currentfill}{rgb}{0.000000,0.000000,0.000000}%
\pgfsetfillcolor{currentfill}%
\pgfsetlinewidth{0.301125pt}%
\definecolor{currentstroke}{rgb}{0.000000,0.000000,0.000000}%
\pgfsetstrokecolor{currentstroke}%
\pgfsetdash{}{0pt}%
\pgfsys@defobject{currentmarker}{\pgfqpoint{0.000000in}{0.000000in}}{\pgfqpoint{0.027778in}{0.000000in}}{%
\pgfpathmoveto{\pgfqpoint{0.000000in}{0.000000in}}%
\pgfpathlineto{\pgfqpoint{0.027778in}{0.000000in}}%
\pgfusepath{stroke,fill}%
}%
\begin{pgfscope}%
\pgfsys@transformshift{0.507620in}{1.490169in}%
\pgfsys@useobject{currentmarker}{}%
\end{pgfscope}%
\end{pgfscope}%
\begin{pgfscope}%
\pgfsetbuttcap%
\pgfsetroundjoin%
\definecolor{currentfill}{rgb}{0.000000,0.000000,0.000000}%
\pgfsetfillcolor{currentfill}%
\pgfsetlinewidth{0.301125pt}%
\definecolor{currentstroke}{rgb}{0.000000,0.000000,0.000000}%
\pgfsetstrokecolor{currentstroke}%
\pgfsetdash{}{0pt}%
\pgfsys@defobject{currentmarker}{\pgfqpoint{0.000000in}{0.000000in}}{\pgfqpoint{0.027778in}{0.000000in}}{%
\pgfpathmoveto{\pgfqpoint{0.000000in}{0.000000in}}%
\pgfpathlineto{\pgfqpoint{0.027778in}{0.000000in}}%
\pgfusepath{stroke,fill}%
}%
\begin{pgfscope}%
\pgfsys@transformshift{0.507620in}{1.577669in}%
\pgfsys@useobject{currentmarker}{}%
\end{pgfscope}%
\end{pgfscope}%
\begin{pgfscope}%
\pgfsetbuttcap%
\pgfsetroundjoin%
\definecolor{currentfill}{rgb}{0.000000,0.000000,0.000000}%
\pgfsetfillcolor{currentfill}%
\pgfsetlinewidth{0.301125pt}%
\definecolor{currentstroke}{rgb}{0.000000,0.000000,0.000000}%
\pgfsetstrokecolor{currentstroke}%
\pgfsetdash{}{0pt}%
\pgfsys@defobject{currentmarker}{\pgfqpoint{0.000000in}{0.000000in}}{\pgfqpoint{0.027778in}{0.000000in}}{%
\pgfpathmoveto{\pgfqpoint{0.000000in}{0.000000in}}%
\pgfpathlineto{\pgfqpoint{0.027778in}{0.000000in}}%
\pgfusepath{stroke,fill}%
}%
\begin{pgfscope}%
\pgfsys@transformshift{0.507620in}{1.665169in}%
\pgfsys@useobject{currentmarker}{}%
\end{pgfscope}%
\end{pgfscope}%
\begin{pgfscope}%
\pgfsetbuttcap%
\pgfsetroundjoin%
\definecolor{currentfill}{rgb}{0.000000,0.000000,0.000000}%
\pgfsetfillcolor{currentfill}%
\pgfsetlinewidth{0.301125pt}%
\definecolor{currentstroke}{rgb}{0.000000,0.000000,0.000000}%
\pgfsetstrokecolor{currentstroke}%
\pgfsetdash{}{0pt}%
\pgfsys@defobject{currentmarker}{\pgfqpoint{0.000000in}{0.000000in}}{\pgfqpoint{0.027778in}{0.000000in}}{%
\pgfpathmoveto{\pgfqpoint{0.000000in}{0.000000in}}%
\pgfpathlineto{\pgfqpoint{0.027778in}{0.000000in}}%
\pgfusepath{stroke,fill}%
}%
\begin{pgfscope}%
\pgfsys@transformshift{0.507620in}{1.840169in}%
\pgfsys@useobject{currentmarker}{}%
\end{pgfscope}%
\end{pgfscope}%
\begin{pgfscope}%
\pgfsetbuttcap%
\pgfsetroundjoin%
\definecolor{currentfill}{rgb}{0.000000,0.000000,0.000000}%
\pgfsetfillcolor{currentfill}%
\pgfsetlinewidth{0.301125pt}%
\definecolor{currentstroke}{rgb}{0.000000,0.000000,0.000000}%
\pgfsetstrokecolor{currentstroke}%
\pgfsetdash{}{0pt}%
\pgfsys@defobject{currentmarker}{\pgfqpoint{0.000000in}{0.000000in}}{\pgfqpoint{0.027778in}{0.000000in}}{%
\pgfpathmoveto{\pgfqpoint{0.000000in}{0.000000in}}%
\pgfpathlineto{\pgfqpoint{0.027778in}{0.000000in}}%
\pgfusepath{stroke,fill}%
}%
\begin{pgfscope}%
\pgfsys@transformshift{0.507620in}{1.927669in}%
\pgfsys@useobject{currentmarker}{}%
\end{pgfscope}%
\end{pgfscope}%
\begin{pgfscope}%
\pgfsetbuttcap%
\pgfsetroundjoin%
\definecolor{currentfill}{rgb}{0.000000,0.000000,0.000000}%
\pgfsetfillcolor{currentfill}%
\pgfsetlinewidth{0.301125pt}%
\definecolor{currentstroke}{rgb}{0.000000,0.000000,0.000000}%
\pgfsetstrokecolor{currentstroke}%
\pgfsetdash{}{0pt}%
\pgfsys@defobject{currentmarker}{\pgfqpoint{0.000000in}{0.000000in}}{\pgfqpoint{0.027778in}{0.000000in}}{%
\pgfpathmoveto{\pgfqpoint{0.000000in}{0.000000in}}%
\pgfpathlineto{\pgfqpoint{0.027778in}{0.000000in}}%
\pgfusepath{stroke,fill}%
}%
\begin{pgfscope}%
\pgfsys@transformshift{0.507620in}{2.015169in}%
\pgfsys@useobject{currentmarker}{}%
\end{pgfscope}%
\end{pgfscope}%
\begin{pgfscope}%
\pgfsetbuttcap%
\pgfsetroundjoin%
\definecolor{currentfill}{rgb}{0.000000,0.000000,0.000000}%
\pgfsetfillcolor{currentfill}%
\pgfsetlinewidth{0.301125pt}%
\definecolor{currentstroke}{rgb}{0.000000,0.000000,0.000000}%
\pgfsetstrokecolor{currentstroke}%
\pgfsetdash{}{0pt}%
\pgfsys@defobject{currentmarker}{\pgfqpoint{0.000000in}{0.000000in}}{\pgfqpoint{0.027778in}{0.000000in}}{%
\pgfpathmoveto{\pgfqpoint{0.000000in}{0.000000in}}%
\pgfpathlineto{\pgfqpoint{0.027778in}{0.000000in}}%
\pgfusepath{stroke,fill}%
}%
\begin{pgfscope}%
\pgfsys@transformshift{0.507620in}{2.102669in}%
\pgfsys@useobject{currentmarker}{}%
\end{pgfscope}%
\end{pgfscope}%
\begin{pgfscope}%
\pgfsetbuttcap%
\pgfsetroundjoin%
\definecolor{currentfill}{rgb}{0.000000,0.000000,0.000000}%
\pgfsetfillcolor{currentfill}%
\pgfsetlinewidth{0.301125pt}%
\definecolor{currentstroke}{rgb}{0.000000,0.000000,0.000000}%
\pgfsetstrokecolor{currentstroke}%
\pgfsetdash{}{0pt}%
\pgfsys@defobject{currentmarker}{\pgfqpoint{0.000000in}{0.000000in}}{\pgfqpoint{0.027778in}{0.000000in}}{%
\pgfpathmoveto{\pgfqpoint{0.000000in}{0.000000in}}%
\pgfpathlineto{\pgfqpoint{0.027778in}{0.000000in}}%
\pgfusepath{stroke,fill}%
}%
\begin{pgfscope}%
\pgfsys@transformshift{0.507620in}{2.277669in}%
\pgfsys@useobject{currentmarker}{}%
\end{pgfscope}%
\end{pgfscope}%
\begin{pgfscope}%
\definecolor{textcolor}{rgb}{0.000000,0.000000,0.000000}%
\pgfsetstrokecolor{textcolor}%
\pgfsetfillcolor{textcolor}%
\pgftext[x=0.140972in,y=1.315169in,,bottom,rotate=90.000000]{\color{textcolor}{\sffamily\fontsize{9.000000}{10.800000}\selectfont\catcode`\^=\active\def^{\ifmmode\sp\else\^{}\fi}\catcode`\%=\active\def%{\%}$P_n(x)$}}%
\end{pgfscope}%
\begin{pgfscope}%
\pgfsetrectcap%
\pgfsetmiterjoin%
\pgfsetlinewidth{0.602250pt}%
\definecolor{currentstroke}{rgb}{0.000000,0.000000,0.000000}%
\pgfsetstrokecolor{currentstroke}%
\pgfsetdash{}{0pt}%
\pgfpathmoveto{\pgfqpoint{0.507620in}{0.352669in}}%
\pgfpathlineto{\pgfqpoint{0.507620in}{2.277669in}}%
\pgfusepath{stroke}%
\end{pgfscope}%
\begin{pgfscope}%
\pgfsetrectcap%
\pgfsetmiterjoin%
\pgfsetlinewidth{0.602250pt}%
\definecolor{currentstroke}{rgb}{0.000000,0.000000,0.000000}%
\pgfsetstrokecolor{currentstroke}%
\pgfsetdash{}{0pt}%
\pgfpathmoveto{\pgfqpoint{0.507620in}{0.352669in}}%
\pgfpathlineto{\pgfqpoint{3.680470in}{0.352669in}}%
\pgfusepath{stroke}%
\end{pgfscope}%
\begin{pgfscope}%
\pgfpathrectangle{\pgfqpoint{0.507620in}{0.352669in}}{\pgfqpoint{3.172850in}{1.925000in}}%
\pgfusepath{clip}%
\pgfsetbuttcap%
\pgfsetroundjoin%
\pgfsetlinewidth{0.803000pt}%
\definecolor{currentstroke}{rgb}{0.431373,0.423529,0.462745}%
\pgfsetstrokecolor{currentstroke}%
\pgfsetstrokeopacity{0.600000}%
\pgfsetdash{{3.200000pt}{2.400000pt}}{0.000000pt}%
\pgfpathmoveto{\pgfqpoint{0.507620in}{2.190169in}}%
\pgfpathlineto{\pgfqpoint{3.680470in}{2.190169in}}%
\pgfpathlineto{\pgfqpoint{3.680470in}{2.190169in}}%
\pgfusepath{stroke}%
\end{pgfscope}%
\begin{pgfscope}%
\pgfpathrectangle{\pgfqpoint{0.507620in}{0.352669in}}{\pgfqpoint{3.172850in}{1.925000in}}%
\pgfusepath{clip}%
\pgfsetbuttcap%
\pgfsetroundjoin%
\pgfsetlinewidth{0.803000pt}%
\definecolor{currentstroke}{rgb}{0.431373,0.423529,0.462745}%
\pgfsetstrokecolor{currentstroke}%
\pgfsetstrokeopacity{0.600000}%
\pgfsetdash{{1.600000pt}{1.600000pt}}{0.000000pt}%
\pgfpathmoveto{\pgfqpoint{0.507620in}{0.440169in}}%
\pgfpathlineto{\pgfqpoint{3.680470in}{2.190169in}}%
\pgfpathlineto{\pgfqpoint{3.680470in}{2.190169in}}%
\pgfusepath{stroke}%
\end{pgfscope}%
\begin{pgfscope}%
\pgfpathrectangle{\pgfqpoint{0.507620in}{0.352669in}}{\pgfqpoint{3.172850in}{1.925000in}}%
\pgfusepath{clip}%
\pgfsetrectcap%
\pgfsetroundjoin%
\pgfsetlinewidth{1.204500pt}%
\definecolor{currentstroke}{rgb}{0.768627,0.556863,0.211765}%
\pgfsetstrokecolor{currentstroke}%
\pgfsetdash{}{0pt}%
\pgfpathmoveto{\pgfqpoint{0.507620in}{2.190169in}}%
\pgfpathlineto{\pgfqpoint{0.558487in}{2.107350in}}%
\pgfpathlineto{\pgfqpoint{0.609354in}{2.027230in}}%
\pgfpathlineto{\pgfqpoint{0.660222in}{1.949808in}}%
\pgfpathlineto{\pgfqpoint{0.711089in}{1.875086in}}%
\pgfpathlineto{\pgfqpoint{0.761956in}{1.803062in}}%
\pgfpathlineto{\pgfqpoint{0.812824in}{1.733737in}}%
\pgfpathlineto{\pgfqpoint{0.863691in}{1.667111in}}%
\pgfpathlineto{\pgfqpoint{0.914558in}{1.603183in}}%
\pgfpathlineto{\pgfqpoint{0.959067in}{1.549460in}}%
\pgfpathlineto{\pgfqpoint{1.003576in}{1.497804in}}%
\pgfpathlineto{\pgfqpoint{1.048085in}{1.448214in}}%
\pgfpathlineto{\pgfqpoint{1.092594in}{1.400690in}}%
\pgfpathlineto{\pgfqpoint{1.137103in}{1.355232in}}%
\pgfpathlineto{\pgfqpoint{1.181612in}{1.311841in}}%
\pgfpathlineto{\pgfqpoint{1.226121in}{1.270516in}}%
\pgfpathlineto{\pgfqpoint{1.270630in}{1.231257in}}%
\pgfpathlineto{\pgfqpoint{1.315139in}{1.194064in}}%
\pgfpathlineto{\pgfqpoint{1.359647in}{1.158938in}}%
\pgfpathlineto{\pgfqpoint{1.404156in}{1.125878in}}%
\pgfpathlineto{\pgfqpoint{1.448665in}{1.094884in}}%
\pgfpathlineto{\pgfqpoint{1.493174in}{1.065956in}}%
\pgfpathlineto{\pgfqpoint{1.531325in}{1.042806in}}%
\pgfpathlineto{\pgfqpoint{1.569475in}{1.021173in}}%
\pgfpathlineto{\pgfqpoint{1.607626in}{1.001059in}}%
\pgfpathlineto{\pgfqpoint{1.645776in}{0.982463in}}%
\pgfpathlineto{\pgfqpoint{1.683927in}{0.965385in}}%
\pgfpathlineto{\pgfqpoint{1.722077in}{0.949824in}}%
\pgfpathlineto{\pgfqpoint{1.760228in}{0.935782in}}%
\pgfpathlineto{\pgfqpoint{1.798378in}{0.923258in}}%
\pgfpathlineto{\pgfqpoint{1.836529in}{0.912252in}}%
\pgfpathlineto{\pgfqpoint{1.874679in}{0.902764in}}%
\pgfpathlineto{\pgfqpoint{1.912830in}{0.894794in}}%
\pgfpathlineto{\pgfqpoint{1.950980in}{0.888343in}}%
\pgfpathlineto{\pgfqpoint{1.989131in}{0.883409in}}%
\pgfpathlineto{\pgfqpoint{2.027281in}{0.879993in}}%
\pgfpathlineto{\pgfqpoint{2.065432in}{0.878096in}}%
\pgfpathlineto{\pgfqpoint{2.103582in}{0.877716in}}%
\pgfpathlineto{\pgfqpoint{2.141733in}{0.878855in}}%
\pgfpathlineto{\pgfqpoint{2.179883in}{0.881511in}}%
\pgfpathlineto{\pgfqpoint{2.218034in}{0.885686in}}%
\pgfpathlineto{\pgfqpoint{2.256184in}{0.891379in}}%
\pgfpathlineto{\pgfqpoint{2.294335in}{0.898590in}}%
\pgfpathlineto{\pgfqpoint{2.332485in}{0.907318in}}%
\pgfpathlineto{\pgfqpoint{2.370636in}{0.917565in}}%
\pgfpathlineto{\pgfqpoint{2.408786in}{0.929330in}}%
\pgfpathlineto{\pgfqpoint{2.446937in}{0.942614in}}%
\pgfpathlineto{\pgfqpoint{2.485087in}{0.957415in}}%
\pgfpathlineto{\pgfqpoint{2.523238in}{0.973734in}}%
\pgfpathlineto{\pgfqpoint{2.561388in}{0.991571in}}%
\pgfpathlineto{\pgfqpoint{2.599539in}{1.010926in}}%
\pgfpathlineto{\pgfqpoint{2.637689in}{1.031800in}}%
\pgfpathlineto{\pgfqpoint{2.675840in}{1.054191in}}%
\pgfpathlineto{\pgfqpoint{2.713990in}{1.078101in}}%
\pgfpathlineto{\pgfqpoint{2.752141in}{1.103529in}}%
\pgfpathlineto{\pgfqpoint{2.796650in}{1.135113in}}%
\pgfpathlineto{\pgfqpoint{2.841159in}{1.168763in}}%
\pgfpathlineto{\pgfqpoint{2.885667in}{1.204480in}}%
\pgfpathlineto{\pgfqpoint{2.930176in}{1.242263in}}%
\pgfpathlineto{\pgfqpoint{2.974685in}{1.282112in}}%
\pgfpathlineto{\pgfqpoint{3.019194in}{1.324028in}}%
\pgfpathlineto{\pgfqpoint{3.063703in}{1.368009in}}%
\pgfpathlineto{\pgfqpoint{3.108212in}{1.414057in}}%
\pgfpathlineto{\pgfqpoint{3.152721in}{1.462172in}}%
\pgfpathlineto{\pgfqpoint{3.197230in}{1.512352in}}%
\pgfpathlineto{\pgfqpoint{3.241739in}{1.564599in}}%
\pgfpathlineto{\pgfqpoint{3.286248in}{1.618912in}}%
\pgfpathlineto{\pgfqpoint{3.337115in}{1.683514in}}%
\pgfpathlineto{\pgfqpoint{3.387982in}{1.750815in}}%
\pgfpathlineto{\pgfqpoint{3.438850in}{1.820815in}}%
\pgfpathlineto{\pgfqpoint{3.489717in}{1.893513in}}%
\pgfpathlineto{\pgfqpoint{3.540584in}{1.968911in}}%
\pgfpathlineto{\pgfqpoint{3.591452in}{2.047007in}}%
\pgfpathlineto{\pgfqpoint{3.642319in}{2.127802in}}%
\pgfpathlineto{\pgfqpoint{3.680470in}{2.190169in}}%
\pgfpathlineto{\pgfqpoint{3.680470in}{2.190169in}}%
\pgfusepath{stroke}%
\end{pgfscope}%
\begin{pgfscope}%
\pgfpathrectangle{\pgfqpoint{0.507620in}{0.352669in}}{\pgfqpoint{3.172850in}{1.925000in}}%
\pgfusepath{clip}%
\pgfsetrectcap%
\pgfsetroundjoin%
\pgfsetlinewidth{1.204500pt}%
\definecolor{currentstroke}{rgb}{0.250980,0.400000,0.600000}%
\pgfsetstrokecolor{currentstroke}%
\pgfsetdash{}{0pt}%
\pgfpathmoveto{\pgfqpoint{0.507620in}{0.440169in}}%
\pgfpathlineto{\pgfqpoint{0.539412in}{0.542761in}}%
\pgfpathlineto{\pgfqpoint{0.571204in}{0.640188in}}%
\pgfpathlineto{\pgfqpoint{0.602996in}{0.732556in}}%
\pgfpathlineto{\pgfqpoint{0.634788in}{0.819969in}}%
\pgfpathlineto{\pgfqpoint{0.666580in}{0.902533in}}%
\pgfpathlineto{\pgfqpoint{0.698372in}{0.980355in}}%
\pgfpathlineto{\pgfqpoint{0.730164in}{1.053539in}}%
\pgfpathlineto{\pgfqpoint{0.761956in}{1.122192in}}%
\pgfpathlineto{\pgfqpoint{0.787390in}{1.173922in}}%
\pgfpathlineto{\pgfqpoint{0.812824in}{1.222874in}}%
\pgfpathlineto{\pgfqpoint{0.838257in}{1.269102in}}%
\pgfpathlineto{\pgfqpoint{0.863691in}{1.312659in}}%
\pgfpathlineto{\pgfqpoint{0.889125in}{1.353599in}}%
\pgfpathlineto{\pgfqpoint{0.914558in}{1.391978in}}%
\pgfpathlineto{\pgfqpoint{0.939992in}{1.427848in}}%
\pgfpathlineto{\pgfqpoint{0.965426in}{1.461264in}}%
\pgfpathlineto{\pgfqpoint{0.990859in}{1.492281in}}%
\pgfpathlineto{\pgfqpoint{1.016293in}{1.520951in}}%
\pgfpathlineto{\pgfqpoint{1.041727in}{1.547330in}}%
\pgfpathlineto{\pgfqpoint{1.067160in}{1.571471in}}%
\pgfpathlineto{\pgfqpoint{1.092594in}{1.593428in}}%
\pgfpathlineto{\pgfqpoint{1.118028in}{1.613256in}}%
\pgfpathlineto{\pgfqpoint{1.143461in}{1.631008in}}%
\pgfpathlineto{\pgfqpoint{1.168895in}{1.646739in}}%
\pgfpathlineto{\pgfqpoint{1.194329in}{1.660502in}}%
\pgfpathlineto{\pgfqpoint{1.219762in}{1.672353in}}%
\pgfpathlineto{\pgfqpoint{1.245196in}{1.682344in}}%
\pgfpathlineto{\pgfqpoint{1.270630in}{1.690530in}}%
\pgfpathlineto{\pgfqpoint{1.296063in}{1.696965in}}%
\pgfpathlineto{\pgfqpoint{1.321497in}{1.701704in}}%
\pgfpathlineto{\pgfqpoint{1.346931in}{1.704799in}}%
\pgfpathlineto{\pgfqpoint{1.372364in}{1.706306in}}%
\pgfpathlineto{\pgfqpoint{1.397798in}{1.706278in}}%
\pgfpathlineto{\pgfqpoint{1.423232in}{1.704770in}}%
\pgfpathlineto{\pgfqpoint{1.448665in}{1.701835in}}%
\pgfpathlineto{\pgfqpoint{1.474099in}{1.697528in}}%
\pgfpathlineto{\pgfqpoint{1.499533in}{1.691902in}}%
\pgfpathlineto{\pgfqpoint{1.524966in}{1.685012in}}%
\pgfpathlineto{\pgfqpoint{1.556758in}{1.674705in}}%
\pgfpathlineto{\pgfqpoint{1.588550in}{1.662612in}}%
\pgfpathlineto{\pgfqpoint{1.620343in}{1.648840in}}%
\pgfpathlineto{\pgfqpoint{1.652135in}{1.633494in}}%
\pgfpathlineto{\pgfqpoint{1.683927in}{1.616679in}}%
\pgfpathlineto{\pgfqpoint{1.722077in}{1.594712in}}%
\pgfpathlineto{\pgfqpoint{1.760228in}{1.570966in}}%
\pgfpathlineto{\pgfqpoint{1.798378in}{1.545622in}}%
\pgfpathlineto{\pgfqpoint{1.842887in}{1.514279in}}%
\pgfpathlineto{\pgfqpoint{1.893754in}{1.476473in}}%
\pgfpathlineto{\pgfqpoint{1.950980in}{1.431926in}}%
\pgfpathlineto{\pgfqpoint{2.020923in}{1.375451in}}%
\pgfpathlineto{\pgfqpoint{2.262543in}{1.178386in}}%
\pgfpathlineto{\pgfqpoint{2.319768in}{1.134721in}}%
\pgfpathlineto{\pgfqpoint{2.370636in}{1.097929in}}%
\pgfpathlineto{\pgfqpoint{2.415145in}{1.067651in}}%
\pgfpathlineto{\pgfqpoint{2.453295in}{1.043352in}}%
\pgfpathlineto{\pgfqpoint{2.491446in}{1.020772in}}%
\pgfpathlineto{\pgfqpoint{2.529596in}{1.000093in}}%
\pgfpathlineto{\pgfqpoint{2.561388in}{0.984444in}}%
\pgfpathlineto{\pgfqpoint{2.593180in}{0.970349in}}%
\pgfpathlineto{\pgfqpoint{2.624972in}{0.957912in}}%
\pgfpathlineto{\pgfqpoint{2.656764in}{0.947239in}}%
\pgfpathlineto{\pgfqpoint{2.688557in}{0.938435in}}%
\pgfpathlineto{\pgfqpoint{2.713990in}{0.932810in}}%
\pgfpathlineto{\pgfqpoint{2.739424in}{0.928503in}}%
\pgfpathlineto{\pgfqpoint{2.764858in}{0.925568in}}%
\pgfpathlineto{\pgfqpoint{2.790291in}{0.924060in}}%
\pgfpathlineto{\pgfqpoint{2.815725in}{0.924032in}}%
\pgfpathlineto{\pgfqpoint{2.841159in}{0.925538in}}%
\pgfpathlineto{\pgfqpoint{2.866592in}{0.928634in}}%
\pgfpathlineto{\pgfqpoint{2.892026in}{0.933372in}}%
\pgfpathlineto{\pgfqpoint{2.917460in}{0.939807in}}%
\pgfpathlineto{\pgfqpoint{2.942893in}{0.947994in}}%
\pgfpathlineto{\pgfqpoint{2.968327in}{0.957985in}}%
\pgfpathlineto{\pgfqpoint{2.993761in}{0.969835in}}%
\pgfpathlineto{\pgfqpoint{3.019194in}{0.983599in}}%
\pgfpathlineto{\pgfqpoint{3.044628in}{0.999330in}}%
\pgfpathlineto{\pgfqpoint{3.070062in}{1.017082in}}%
\pgfpathlineto{\pgfqpoint{3.095495in}{1.036910in}}%
\pgfpathlineto{\pgfqpoint{3.120929in}{1.058867in}}%
\pgfpathlineto{\pgfqpoint{3.146363in}{1.083008in}}%
\pgfpathlineto{\pgfqpoint{3.171796in}{1.109386in}}%
\pgfpathlineto{\pgfqpoint{3.197230in}{1.138057in}}%
\pgfpathlineto{\pgfqpoint{3.222664in}{1.169073in}}%
\pgfpathlineto{\pgfqpoint{3.248097in}{1.202489in}}%
\pgfpathlineto{\pgfqpoint{3.273531in}{1.238360in}}%
\pgfpathlineto{\pgfqpoint{3.298965in}{1.276738in}}%
\pgfpathlineto{\pgfqpoint{3.324398in}{1.317679in}}%
\pgfpathlineto{\pgfqpoint{3.349832in}{1.361236in}}%
\pgfpathlineto{\pgfqpoint{3.375266in}{1.407463in}}%
\pgfpathlineto{\pgfqpoint{3.400699in}{1.456415in}}%
\pgfpathlineto{\pgfqpoint{3.426133in}{1.508145in}}%
\pgfpathlineto{\pgfqpoint{3.451567in}{1.562708in}}%
\pgfpathlineto{\pgfqpoint{3.477000in}{1.620158in}}%
\pgfpathlineto{\pgfqpoint{3.508792in}{1.696112in}}%
\pgfpathlineto{\pgfqpoint{3.540584in}{1.776767in}}%
\pgfpathlineto{\pgfqpoint{3.572376in}{1.862228in}}%
\pgfpathlineto{\pgfqpoint{3.604169in}{1.952601in}}%
\pgfpathlineto{\pgfqpoint{3.635961in}{2.047991in}}%
\pgfpathlineto{\pgfqpoint{3.667753in}{2.148505in}}%
\pgfpathlineto{\pgfqpoint{3.680470in}{2.190169in}}%
\pgfpathlineto{\pgfqpoint{3.680470in}{2.190169in}}%
\pgfusepath{stroke}%
\end{pgfscope}%
\begin{pgfscope}%
\pgfpathrectangle{\pgfqpoint{0.507620in}{0.352669in}}{\pgfqpoint{3.172850in}{1.925000in}}%
\pgfusepath{clip}%
\pgfsetrectcap%
\pgfsetroundjoin%
\pgfsetlinewidth{1.204500pt}%
\definecolor{currentstroke}{rgb}{0.619608,0.188235,0.172549}%
\pgfsetstrokecolor{currentstroke}%
\pgfsetdash{}{0pt}%
\pgfpathmoveto{\pgfqpoint{0.507620in}{2.190169in}}%
\pgfpathlineto{\pgfqpoint{0.533053in}{2.054886in}}%
\pgfpathlineto{\pgfqpoint{0.558487in}{1.929348in}}%
\pgfpathlineto{\pgfqpoint{0.583921in}{1.813186in}}%
\pgfpathlineto{\pgfqpoint{0.602996in}{1.731999in}}%
\pgfpathlineto{\pgfqpoint{0.622071in}{1.655730in}}%
\pgfpathlineto{\pgfqpoint{0.641146in}{1.584230in}}%
\pgfpathlineto{\pgfqpoint{0.660222in}{1.517352in}}%
\pgfpathlineto{\pgfqpoint{0.679297in}{1.454951in}}%
\pgfpathlineto{\pgfqpoint{0.698372in}{1.396883in}}%
\pgfpathlineto{\pgfqpoint{0.717447in}{1.343007in}}%
\pgfpathlineto{\pgfqpoint{0.736523in}{1.293183in}}%
\pgfpathlineto{\pgfqpoint{0.755598in}{1.247273in}}%
\pgfpathlineto{\pgfqpoint{0.774673in}{1.205143in}}%
\pgfpathlineto{\pgfqpoint{0.793748in}{1.166657in}}%
\pgfpathlineto{\pgfqpoint{0.812824in}{1.131685in}}%
\pgfpathlineto{\pgfqpoint{0.831899in}{1.100096in}}%
\pgfpathlineto{\pgfqpoint{0.850974in}{1.071762in}}%
\pgfpathlineto{\pgfqpoint{0.870049in}{1.046557in}}%
\pgfpathlineto{\pgfqpoint{0.889125in}{1.024358in}}%
\pgfpathlineto{\pgfqpoint{0.908200in}{1.005041in}}%
\pgfpathlineto{\pgfqpoint{0.927275in}{0.988486in}}%
\pgfpathlineto{\pgfqpoint{0.946350in}{0.974575in}}%
\pgfpathlineto{\pgfqpoint{0.965426in}{0.963192in}}%
\pgfpathlineto{\pgfqpoint{0.984501in}{0.954221in}}%
\pgfpathlineto{\pgfqpoint{1.003576in}{0.947551in}}%
\pgfpathlineto{\pgfqpoint{1.022651in}{0.943069in}}%
\pgfpathlineto{\pgfqpoint{1.041727in}{0.940669in}}%
\pgfpathlineto{\pgfqpoint{1.060802in}{0.940242in}}%
\pgfpathlineto{\pgfqpoint{1.079877in}{0.941684in}}%
\pgfpathlineto{\pgfqpoint{1.098952in}{0.944891in}}%
\pgfpathlineto{\pgfqpoint{1.118028in}{0.949763in}}%
\pgfpathlineto{\pgfqpoint{1.137103in}{0.956200in}}%
\pgfpathlineto{\pgfqpoint{1.156178in}{0.964105in}}%
\pgfpathlineto{\pgfqpoint{1.175253in}{0.973383in}}%
\pgfpathlineto{\pgfqpoint{1.200687in}{0.987727in}}%
\pgfpathlineto{\pgfqpoint{1.226121in}{1.004129in}}%
\pgfpathlineto{\pgfqpoint{1.251554in}{1.022379in}}%
\pgfpathlineto{\pgfqpoint{1.276988in}{1.042272in}}%
\pgfpathlineto{\pgfqpoint{1.308780in}{1.069150in}}%
\pgfpathlineto{\pgfqpoint{1.340572in}{1.097913in}}%
\pgfpathlineto{\pgfqpoint{1.378723in}{1.134414in}}%
\pgfpathlineto{\pgfqpoint{1.429590in}{1.185487in}}%
\pgfpathlineto{\pgfqpoint{1.512249in}{1.271232in}}%
\pgfpathlineto{\pgfqpoint{1.588550in}{1.349611in}}%
\pgfpathlineto{\pgfqpoint{1.639418in}{1.399642in}}%
\pgfpathlineto{\pgfqpoint{1.677568in}{1.435335in}}%
\pgfpathlineto{\pgfqpoint{1.715719in}{1.469066in}}%
\pgfpathlineto{\pgfqpoint{1.747511in}{1.495445in}}%
\pgfpathlineto{\pgfqpoint{1.779303in}{1.520070in}}%
\pgfpathlineto{\pgfqpoint{1.811095in}{1.542787in}}%
\pgfpathlineto{\pgfqpoint{1.842887in}{1.563457in}}%
\pgfpathlineto{\pgfqpoint{1.868321in}{1.578434in}}%
\pgfpathlineto{\pgfqpoint{1.893754in}{1.591964in}}%
\pgfpathlineto{\pgfqpoint{1.919188in}{1.603996in}}%
\pgfpathlineto{\pgfqpoint{1.944622in}{1.614486in}}%
\pgfpathlineto{\pgfqpoint{1.970055in}{1.623393in}}%
\pgfpathlineto{\pgfqpoint{1.995489in}{1.630687in}}%
\pgfpathlineto{\pgfqpoint{2.020923in}{1.636340in}}%
\pgfpathlineto{\pgfqpoint{2.046356in}{1.640332in}}%
\pgfpathlineto{\pgfqpoint{2.071790in}{1.642648in}}%
\pgfpathlineto{\pgfqpoint{2.097224in}{1.643281in}}%
\pgfpathlineto{\pgfqpoint{2.122657in}{1.642227in}}%
\pgfpathlineto{\pgfqpoint{2.148091in}{1.639491in}}%
\pgfpathlineto{\pgfqpoint{2.173525in}{1.635082in}}%
\pgfpathlineto{\pgfqpoint{2.198958in}{1.629017in}}%
\pgfpathlineto{\pgfqpoint{2.224392in}{1.621317in}}%
\pgfpathlineto{\pgfqpoint{2.249826in}{1.612010in}}%
\pgfpathlineto{\pgfqpoint{2.275259in}{1.601131in}}%
\pgfpathlineto{\pgfqpoint{2.300693in}{1.588720in}}%
\pgfpathlineto{\pgfqpoint{2.326127in}{1.574823in}}%
\pgfpathlineto{\pgfqpoint{2.351560in}{1.559493in}}%
\pgfpathlineto{\pgfqpoint{2.383353in}{1.538404in}}%
\pgfpathlineto{\pgfqpoint{2.415145in}{1.515293in}}%
\pgfpathlineto{\pgfqpoint{2.446937in}{1.490305in}}%
\pgfpathlineto{\pgfqpoint{2.478729in}{1.463595in}}%
\pgfpathlineto{\pgfqpoint{2.516879in}{1.429513in}}%
\pgfpathlineto{\pgfqpoint{2.555030in}{1.393527in}}%
\pgfpathlineto{\pgfqpoint{2.599539in}{1.349611in}}%
\pgfpathlineto{\pgfqpoint{2.663123in}{1.284455in}}%
\pgfpathlineto{\pgfqpoint{2.790291in}{1.153304in}}%
\pgfpathlineto{\pgfqpoint{2.834800in}{1.109865in}}%
\pgfpathlineto{\pgfqpoint{2.872951in}{1.074763in}}%
\pgfpathlineto{\pgfqpoint{2.904743in}{1.047479in}}%
\pgfpathlineto{\pgfqpoint{2.936535in}{1.022379in}}%
\pgfpathlineto{\pgfqpoint{2.961968in}{1.004129in}}%
\pgfpathlineto{\pgfqpoint{2.987402in}{0.987727in}}%
\pgfpathlineto{\pgfqpoint{3.012836in}{0.973383in}}%
\pgfpathlineto{\pgfqpoint{3.031911in}{0.964105in}}%
\pgfpathlineto{\pgfqpoint{3.050986in}{0.956200in}}%
\pgfpathlineto{\pgfqpoint{3.070062in}{0.949763in}}%
\pgfpathlineto{\pgfqpoint{3.089137in}{0.944891in}}%
\pgfpathlineto{\pgfqpoint{3.108212in}{0.941684in}}%
\pgfpathlineto{\pgfqpoint{3.127287in}{0.940242in}}%
\pgfpathlineto{\pgfqpoint{3.146363in}{0.940669in}}%
\pgfpathlineto{\pgfqpoint{3.165438in}{0.943069in}}%
\pgfpathlineto{\pgfqpoint{3.184513in}{0.947551in}}%
\pgfpathlineto{\pgfqpoint{3.203588in}{0.954221in}}%
\pgfpathlineto{\pgfqpoint{3.222664in}{0.963192in}}%
\pgfpathlineto{\pgfqpoint{3.241739in}{0.974575in}}%
\pgfpathlineto{\pgfqpoint{3.260814in}{0.988486in}}%
\pgfpathlineto{\pgfqpoint{3.279889in}{1.005041in}}%
\pgfpathlineto{\pgfqpoint{3.298965in}{1.024358in}}%
\pgfpathlineto{\pgfqpoint{3.318040in}{1.046557in}}%
\pgfpathlineto{\pgfqpoint{3.337115in}{1.071762in}}%
\pgfpathlineto{\pgfqpoint{3.356190in}{1.100096in}}%
\pgfpathlineto{\pgfqpoint{3.375266in}{1.131685in}}%
\pgfpathlineto{\pgfqpoint{3.394341in}{1.166657in}}%
\pgfpathlineto{\pgfqpoint{3.413416in}{1.205143in}}%
\pgfpathlineto{\pgfqpoint{3.432491in}{1.247273in}}%
\pgfpathlineto{\pgfqpoint{3.451567in}{1.293183in}}%
\pgfpathlineto{\pgfqpoint{3.470642in}{1.343007in}}%
\pgfpathlineto{\pgfqpoint{3.489717in}{1.396883in}}%
\pgfpathlineto{\pgfqpoint{3.508792in}{1.454951in}}%
\pgfpathlineto{\pgfqpoint{3.527868in}{1.517352in}}%
\pgfpathlineto{\pgfqpoint{3.546943in}{1.584230in}}%
\pgfpathlineto{\pgfqpoint{3.566018in}{1.655730in}}%
\pgfpathlineto{\pgfqpoint{3.585093in}{1.731999in}}%
\pgfpathlineto{\pgfqpoint{3.604169in}{1.813186in}}%
\pgfpathlineto{\pgfqpoint{3.623244in}{1.899443in}}%
\pgfpathlineto{\pgfqpoint{3.648677in}{2.022602in}}%
\pgfpathlineto{\pgfqpoint{3.674111in}{2.155414in}}%
\pgfpathlineto{\pgfqpoint{3.680470in}{2.190169in}}%
\pgfpathlineto{\pgfqpoint{3.680470in}{2.190169in}}%
\pgfusepath{stroke}%
\end{pgfscope}%
\begin{pgfscope}%
\pgfpathrectangle{\pgfqpoint{0.507620in}{0.352669in}}{\pgfqpoint{3.172850in}{1.925000in}}%
\pgfusepath{clip}%
\pgfsetrectcap%
\pgfsetroundjoin%
\pgfsetlinewidth{1.204500pt}%
\definecolor{currentstroke}{rgb}{0.431373,0.423529,0.462745}%
\pgfsetstrokecolor{currentstroke}%
\pgfsetdash{}{0pt}%
\pgfpathmoveto{\pgfqpoint{0.507620in}{0.440169in}}%
\pgfpathlineto{\pgfqpoint{0.526695in}{0.591449in}}%
\pgfpathlineto{\pgfqpoint{0.545770in}{0.730074in}}%
\pgfpathlineto{\pgfqpoint{0.564845in}{0.856659in}}%
\pgfpathlineto{\pgfqpoint{0.583921in}{0.971799in}}%
\pgfpathlineto{\pgfqpoint{0.602996in}{1.076073in}}%
\pgfpathlineto{\pgfqpoint{0.622071in}{1.170046in}}%
\pgfpathlineto{\pgfqpoint{0.641146in}{1.254265in}}%
\pgfpathlineto{\pgfqpoint{0.660222in}{1.329259in}}%
\pgfpathlineto{\pgfqpoint{0.672938in}{1.374386in}}%
\pgfpathlineto{\pgfqpoint{0.685655in}{1.415791in}}%
\pgfpathlineto{\pgfqpoint{0.698372in}{1.453620in}}%
\pgfpathlineto{\pgfqpoint{0.711089in}{1.488017in}}%
\pgfpathlineto{\pgfqpoint{0.723806in}{1.519119in}}%
\pgfpathlineto{\pgfqpoint{0.736523in}{1.547066in}}%
\pgfpathlineto{\pgfqpoint{0.749239in}{1.571989in}}%
\pgfpathlineto{\pgfqpoint{0.761956in}{1.594019in}}%
\pgfpathlineto{\pgfqpoint{0.774673in}{1.613286in}}%
\pgfpathlineto{\pgfqpoint{0.787390in}{1.629913in}}%
\pgfpathlineto{\pgfqpoint{0.800107in}{1.644024in}}%
\pgfpathlineto{\pgfqpoint{0.812824in}{1.655736in}}%
\pgfpathlineto{\pgfqpoint{0.825540in}{1.665168in}}%
\pgfpathlineto{\pgfqpoint{0.838257in}{1.672433in}}%
\pgfpathlineto{\pgfqpoint{0.850974in}{1.677641in}}%
\pgfpathlineto{\pgfqpoint{0.863691in}{1.680902in}}%
\pgfpathlineto{\pgfqpoint{0.876408in}{1.682321in}}%
\pgfpathlineto{\pgfqpoint{0.889125in}{1.682002in}}%
\pgfpathlineto{\pgfqpoint{0.901841in}{1.680044in}}%
\pgfpathlineto{\pgfqpoint{0.914558in}{1.676546in}}%
\pgfpathlineto{\pgfqpoint{0.927275in}{1.671603in}}%
\pgfpathlineto{\pgfqpoint{0.939992in}{1.665308in}}%
\pgfpathlineto{\pgfqpoint{0.952709in}{1.657751in}}%
\pgfpathlineto{\pgfqpoint{0.971784in}{1.644241in}}%
\pgfpathlineto{\pgfqpoint{0.990859in}{1.628375in}}%
\pgfpathlineto{\pgfqpoint{1.009935in}{1.610426in}}%
\pgfpathlineto{\pgfqpoint{1.029010in}{1.590656in}}%
\pgfpathlineto{\pgfqpoint{1.054443in}{1.561894in}}%
\pgfpathlineto{\pgfqpoint{1.079877in}{1.530898in}}%
\pgfpathlineto{\pgfqpoint{1.111669in}{1.489809in}}%
\pgfpathlineto{\pgfqpoint{1.156178in}{1.429582in}}%
\pgfpathlineto{\pgfqpoint{1.251554in}{1.299539in}}%
\pgfpathlineto{\pgfqpoint{1.289705in}{1.250358in}}%
\pgfpathlineto{\pgfqpoint{1.321497in}{1.211685in}}%
\pgfpathlineto{\pgfqpoint{1.353289in}{1.175599in}}%
\pgfpathlineto{\pgfqpoint{1.378723in}{1.148856in}}%
\pgfpathlineto{\pgfqpoint{1.404156in}{1.124187in}}%
\pgfpathlineto{\pgfqpoint{1.429590in}{1.101738in}}%
\pgfpathlineto{\pgfqpoint{1.455024in}{1.081629in}}%
\pgfpathlineto{\pgfqpoint{1.474099in}{1.068144in}}%
\pgfpathlineto{\pgfqpoint{1.493174in}{1.056065in}}%
\pgfpathlineto{\pgfqpoint{1.512249in}{1.045418in}}%
\pgfpathlineto{\pgfqpoint{1.531325in}{1.036223in}}%
\pgfpathlineto{\pgfqpoint{1.550400in}{1.028495in}}%
\pgfpathlineto{\pgfqpoint{1.569475in}{1.022240in}}%
\pgfpathlineto{\pgfqpoint{1.588550in}{1.017460in}}%
\pgfpathlineto{\pgfqpoint{1.607626in}{1.014152in}}%
\pgfpathlineto{\pgfqpoint{1.626701in}{1.012305in}}%
\pgfpathlineto{\pgfqpoint{1.645776in}{1.011905in}}%
\pgfpathlineto{\pgfqpoint{1.664851in}{1.012931in}}%
\pgfpathlineto{\pgfqpoint{1.683927in}{1.015360in}}%
\pgfpathlineto{\pgfqpoint{1.703002in}{1.019161in}}%
\pgfpathlineto{\pgfqpoint{1.722077in}{1.024300in}}%
\pgfpathlineto{\pgfqpoint{1.741152in}{1.030739in}}%
\pgfpathlineto{\pgfqpoint{1.760228in}{1.038436in}}%
\pgfpathlineto{\pgfqpoint{1.785661in}{1.050575in}}%
\pgfpathlineto{\pgfqpoint{1.811095in}{1.064748in}}%
\pgfpathlineto{\pgfqpoint{1.836529in}{1.080825in}}%
\pgfpathlineto{\pgfqpoint{1.861962in}{1.098667in}}%
\pgfpathlineto{\pgfqpoint{1.887396in}{1.118124in}}%
\pgfpathlineto{\pgfqpoint{1.919188in}{1.144478in}}%
\pgfpathlineto{\pgfqpoint{1.950980in}{1.172791in}}%
\pgfpathlineto{\pgfqpoint{1.989131in}{1.208876in}}%
\pgfpathlineto{\pgfqpoint{2.033640in}{1.253122in}}%
\pgfpathlineto{\pgfqpoint{2.109941in}{1.331600in}}%
\pgfpathlineto{\pgfqpoint{2.179883in}{1.402730in}}%
\pgfpathlineto{\pgfqpoint{2.224392in}{1.445749in}}%
\pgfpathlineto{\pgfqpoint{2.262543in}{1.480343in}}%
\pgfpathlineto{\pgfqpoint{2.294335in}{1.507115in}}%
\pgfpathlineto{\pgfqpoint{2.326127in}{1.531671in}}%
\pgfpathlineto{\pgfqpoint{2.351560in}{1.549512in}}%
\pgfpathlineto{\pgfqpoint{2.376994in}{1.565589in}}%
\pgfpathlineto{\pgfqpoint{2.402428in}{1.579762in}}%
\pgfpathlineto{\pgfqpoint{2.427861in}{1.591901in}}%
\pgfpathlineto{\pgfqpoint{2.446937in}{1.599598in}}%
\pgfpathlineto{\pgfqpoint{2.466012in}{1.606038in}}%
\pgfpathlineto{\pgfqpoint{2.485087in}{1.611177in}}%
\pgfpathlineto{\pgfqpoint{2.504162in}{1.614978in}}%
\pgfpathlineto{\pgfqpoint{2.523238in}{1.617406in}}%
\pgfpathlineto{\pgfqpoint{2.542313in}{1.618433in}}%
\pgfpathlineto{\pgfqpoint{2.561388in}{1.618033in}}%
\pgfpathlineto{\pgfqpoint{2.580463in}{1.616186in}}%
\pgfpathlineto{\pgfqpoint{2.599539in}{1.612877in}}%
\pgfpathlineto{\pgfqpoint{2.618614in}{1.608098in}}%
\pgfpathlineto{\pgfqpoint{2.637689in}{1.601843in}}%
\pgfpathlineto{\pgfqpoint{2.656764in}{1.594114in}}%
\pgfpathlineto{\pgfqpoint{2.675840in}{1.584920in}}%
\pgfpathlineto{\pgfqpoint{2.694915in}{1.574273in}}%
\pgfpathlineto{\pgfqpoint{2.713990in}{1.562193in}}%
\pgfpathlineto{\pgfqpoint{2.733065in}{1.548708in}}%
\pgfpathlineto{\pgfqpoint{2.758499in}{1.528600in}}%
\pgfpathlineto{\pgfqpoint{2.783933in}{1.506150in}}%
\pgfpathlineto{\pgfqpoint{2.809366in}{1.481481in}}%
\pgfpathlineto{\pgfqpoint{2.834800in}{1.454738in}}%
\pgfpathlineto{\pgfqpoint{2.860234in}{1.426090in}}%
\pgfpathlineto{\pgfqpoint{2.892026in}{1.387900in}}%
\pgfpathlineto{\pgfqpoint{2.923818in}{1.347480in}}%
\pgfpathlineto{\pgfqpoint{2.968327in}{1.288122in}}%
\pgfpathlineto{\pgfqpoint{3.095495in}{1.115609in}}%
\pgfpathlineto{\pgfqpoint{3.127287in}{1.076004in}}%
\pgfpathlineto{\pgfqpoint{3.152721in}{1.046633in}}%
\pgfpathlineto{\pgfqpoint{3.178155in}{1.019912in}}%
\pgfpathlineto{\pgfqpoint{3.197230in}{1.001963in}}%
\pgfpathlineto{\pgfqpoint{3.216305in}{0.986096in}}%
\pgfpathlineto{\pgfqpoint{3.235380in}{0.972586in}}%
\pgfpathlineto{\pgfqpoint{3.248097in}{0.965029in}}%
\pgfpathlineto{\pgfqpoint{3.260814in}{0.958734in}}%
\pgfpathlineto{\pgfqpoint{3.273531in}{0.953791in}}%
\pgfpathlineto{\pgfqpoint{3.286248in}{0.950293in}}%
\pgfpathlineto{\pgfqpoint{3.298965in}{0.948336in}}%
\pgfpathlineto{\pgfqpoint{3.311681in}{0.948016in}}%
\pgfpathlineto{\pgfqpoint{3.324398in}{0.949435in}}%
\pgfpathlineto{\pgfqpoint{3.337115in}{0.952696in}}%
\pgfpathlineto{\pgfqpoint{3.349832in}{0.957905in}}%
\pgfpathlineto{\pgfqpoint{3.362549in}{0.965170in}}%
\pgfpathlineto{\pgfqpoint{3.375266in}{0.974601in}}%
\pgfpathlineto{\pgfqpoint{3.387982in}{0.986314in}}%
\pgfpathlineto{\pgfqpoint{3.400699in}{1.000424in}}%
\pgfpathlineto{\pgfqpoint{3.413416in}{1.017052in}}%
\pgfpathlineto{\pgfqpoint{3.426133in}{1.036318in}}%
\pgfpathlineto{\pgfqpoint{3.438850in}{1.058349in}}%
\pgfpathlineto{\pgfqpoint{3.451567in}{1.083272in}}%
\pgfpathlineto{\pgfqpoint{3.464283in}{1.111218in}}%
\pgfpathlineto{\pgfqpoint{3.477000in}{1.142321in}}%
\pgfpathlineto{\pgfqpoint{3.489717in}{1.176717in}}%
\pgfpathlineto{\pgfqpoint{3.502434in}{1.214547in}}%
\pgfpathlineto{\pgfqpoint{3.515151in}{1.255952in}}%
\pgfpathlineto{\pgfqpoint{3.527868in}{1.301079in}}%
\pgfpathlineto{\pgfqpoint{3.540584in}{1.350076in}}%
\pgfpathlineto{\pgfqpoint{3.559660in}{1.431161in}}%
\pgfpathlineto{\pgfqpoint{3.578735in}{1.521822in}}%
\pgfpathlineto{\pgfqpoint{3.597810in}{1.622602in}}%
\pgfpathlineto{\pgfqpoint{3.616885in}{1.734056in}}%
\pgfpathlineto{\pgfqpoint{3.635961in}{1.856760in}}%
\pgfpathlineto{\pgfqpoint{3.655036in}{1.991305in}}%
\pgfpathlineto{\pgfqpoint{3.674111in}{2.138298in}}%
\pgfpathlineto{\pgfqpoint{3.680470in}{2.190169in}}%
\pgfpathlineto{\pgfqpoint{3.680470in}{2.190169in}}%
\pgfusepath{stroke}%
\end{pgfscope}%
\begin{pgfscope}%
\pgfsetbuttcap%
\pgfsetroundjoin%
\definecolor{currentfill}{rgb}{0.000000,0.000000,0.000000}%
\pgfsetfillcolor{currentfill}%
\pgfsetlinewidth{1.003750pt}%
\definecolor{currentstroke}{rgb}{0.000000,0.000000,0.000000}%
\pgfsetstrokecolor{currentstroke}%
\pgfsetdash{}{0pt}%
\pgfsys@defobject{currentmarker}{\pgfqpoint{-0.024306in}{-0.024306in}}{\pgfqpoint{0.024306in}{0.024306in}}{%
\pgfpathmoveto{\pgfqpoint{0.000000in}{-0.024306in}}%
\pgfpathcurveto{\pgfqpoint{0.006446in}{-0.024306in}}{\pgfqpoint{0.012629in}{-0.021745in}}{\pgfqpoint{0.017187in}{-0.017187in}}%
\pgfpathcurveto{\pgfqpoint{0.021745in}{-0.012629in}}{\pgfqpoint{0.024306in}{-0.006446in}}{\pgfqpoint{0.024306in}{0.000000in}}%
\pgfpathcurveto{\pgfqpoint{0.024306in}{0.006446in}}{\pgfqpoint{0.021745in}{0.012629in}}{\pgfqpoint{0.017187in}{0.017187in}}%
\pgfpathcurveto{\pgfqpoint{0.012629in}{0.021745in}}{\pgfqpoint{0.006446in}{0.024306in}}{\pgfqpoint{0.000000in}{0.024306in}}%
\pgfpathcurveto{\pgfqpoint{-0.006446in}{0.024306in}}{\pgfqpoint{-0.012629in}{0.021745in}}{\pgfqpoint{-0.017187in}{0.017187in}}%
\pgfpathcurveto{\pgfqpoint{-0.021745in}{0.012629in}}{\pgfqpoint{-0.024306in}{0.006446in}}{\pgfqpoint{-0.024306in}{0.000000in}}%
\pgfpathcurveto{\pgfqpoint{-0.024306in}{-0.006446in}}{\pgfqpoint{-0.021745in}{-0.012629in}}{\pgfqpoint{-0.017187in}{-0.017187in}}%
\pgfpathcurveto{\pgfqpoint{-0.012629in}{-0.021745in}}{\pgfqpoint{-0.006446in}{-0.024306in}}{\pgfqpoint{0.000000in}{-0.024306in}}%
\pgfpathlineto{\pgfqpoint{0.000000in}{-0.024306in}}%
\pgfpathclose%
\pgfusepath{stroke,fill}%
}%
\begin{pgfscope}%
\pgfsys@transformshift{0.507620in}{1.315169in}%
\pgfsys@useobject{currentmarker}{}%
\end{pgfscope}%
\begin{pgfscope}%
\pgfsys@transformshift{0.880342in}{1.315169in}%
\pgfsys@useobject{currentmarker}{}%
\end{pgfscope}%
\begin{pgfscope}%
\pgfsys@transformshift{1.641546in}{1.315169in}%
\pgfsys@useobject{currentmarker}{}%
\end{pgfscope}%
\begin{pgfscope}%
\pgfsys@transformshift{2.546543in}{1.315169in}%
\pgfsys@useobject{currentmarker}{}%
\end{pgfscope}%
\begin{pgfscope}%
\pgfsys@transformshift{3.307747in}{1.315169in}%
\pgfsys@useobject{currentmarker}{}%
\end{pgfscope}%
\begin{pgfscope}%
\pgfsys@transformshift{3.680470in}{1.315169in}%
\pgfsys@useobject{currentmarker}{}%
\end{pgfscope}%
\end{pgfscope}%
\begin{pgfscope}%
\pgfsetbuttcap%
\pgfsetroundjoin%
\pgfsetlinewidth{0.803000pt}%
\definecolor{currentstroke}{rgb}{0.431373,0.423529,0.462745}%
\pgfsetstrokecolor{currentstroke}%
\pgfsetstrokeopacity{0.600000}%
\pgfsetdash{{3.200000pt}{2.400000pt}}{0.000000pt}%
\pgfpathmoveto{\pgfqpoint{0.574286in}{0.589524in}}%
\pgfpathlineto{\pgfqpoint{0.657620in}{0.589524in}}%
\pgfpathlineto{\pgfqpoint{0.740953in}{0.589524in}}%
\pgfusepath{stroke}%
\end{pgfscope}%
\begin{pgfscope}%
\definecolor{textcolor}{rgb}{0.000000,0.000000,0.000000}%
\pgfsetstrokecolor{textcolor}%
\pgfsetfillcolor{textcolor}%
\pgftext[x=0.807620in,y=0.560357in,left,base]{\color{textcolor}{\sffamily\fontsize{6.000000}{7.200000}\selectfont\catcode`\^=\active\def^{\ifmmode\sp\else\^{}\fi}\catcode`\%=\active\def%{\%}$P_0$}}%
\end{pgfscope}%
\begin{pgfscope}%
\pgfsetbuttcap%
\pgfsetroundjoin%
\pgfsetlinewidth{0.803000pt}%
\definecolor{currentstroke}{rgb}{0.431373,0.423529,0.462745}%
\pgfsetstrokecolor{currentstroke}%
\pgfsetstrokeopacity{0.600000}%
\pgfsetdash{{1.600000pt}{1.600000pt}}{0.000000pt}%
\pgfpathmoveto{\pgfqpoint{0.574286in}{0.466002in}}%
\pgfpathlineto{\pgfqpoint{0.657620in}{0.466002in}}%
\pgfpathlineto{\pgfqpoint{0.740953in}{0.466002in}}%
\pgfusepath{stroke}%
\end{pgfscope}%
\begin{pgfscope}%
\definecolor{textcolor}{rgb}{0.000000,0.000000,0.000000}%
\pgfsetstrokecolor{textcolor}%
\pgfsetfillcolor{textcolor}%
\pgftext[x=0.807620in,y=0.436835in,left,base]{\color{textcolor}{\sffamily\fontsize{6.000000}{7.200000}\selectfont\catcode`\^=\active\def^{\ifmmode\sp\else\^{}\fi}\catcode`\%=\active\def%{\%}$P_1$}}%
\end{pgfscope}%
\begin{pgfscope}%
\pgfsetrectcap%
\pgfsetroundjoin%
\pgfsetlinewidth{1.204500pt}%
\definecolor{currentstroke}{rgb}{0.768627,0.556863,0.211765}%
\pgfsetstrokecolor{currentstroke}%
\pgfsetdash{}{0pt}%
\pgfpathmoveto{\pgfqpoint{0.996703in}{0.589524in}}%
\pgfpathlineto{\pgfqpoint{1.080036in}{0.589524in}}%
\pgfpathlineto{\pgfqpoint{1.163370in}{0.589524in}}%
\pgfusepath{stroke}%
\end{pgfscope}%
\begin{pgfscope}%
\definecolor{textcolor}{rgb}{0.000000,0.000000,0.000000}%
\pgfsetstrokecolor{textcolor}%
\pgfsetfillcolor{textcolor}%
\pgftext[x=1.230036in,y=0.560357in,left,base]{\color{textcolor}{\sffamily\fontsize{6.000000}{7.200000}\selectfont\catcode`\^=\active\def^{\ifmmode\sp\else\^{}\fi}\catcode`\%=\active\def%{\%}$P_2$}}%
\end{pgfscope}%
\begin{pgfscope}%
\pgfsetrectcap%
\pgfsetroundjoin%
\pgfsetlinewidth{1.204500pt}%
\definecolor{currentstroke}{rgb}{0.250980,0.400000,0.600000}%
\pgfsetstrokecolor{currentstroke}%
\pgfsetdash{}{0pt}%
\pgfpathmoveto{\pgfqpoint{0.996703in}{0.467044in}}%
\pgfpathlineto{\pgfqpoint{1.080036in}{0.467044in}}%
\pgfpathlineto{\pgfqpoint{1.163370in}{0.467044in}}%
\pgfusepath{stroke}%
\end{pgfscope}%
\begin{pgfscope}%
\definecolor{textcolor}{rgb}{0.000000,0.000000,0.000000}%
\pgfsetstrokecolor{textcolor}%
\pgfsetfillcolor{textcolor}%
\pgftext[x=1.230036in,y=0.437877in,left,base]{\color{textcolor}{\sffamily\fontsize{6.000000}{7.200000}\selectfont\catcode`\^=\active\def^{\ifmmode\sp\else\^{}\fi}\catcode`\%=\active\def%{\%}$P_3$}}%
\end{pgfscope}%
\begin{pgfscope}%
\pgfsetrectcap%
\pgfsetroundjoin%
\pgfsetlinewidth{1.204500pt}%
\definecolor{currentstroke}{rgb}{0.619608,0.188235,0.172549}%
\pgfsetstrokecolor{currentstroke}%
\pgfsetdash{}{0pt}%
\pgfpathmoveto{\pgfqpoint{1.415994in}{0.589524in}}%
\pgfpathlineto{\pgfqpoint{1.499328in}{0.589524in}}%
\pgfpathlineto{\pgfqpoint{1.582661in}{0.589524in}}%
\pgfusepath{stroke}%
\end{pgfscope}%
\begin{pgfscope}%
\definecolor{textcolor}{rgb}{0.000000,0.000000,0.000000}%
\pgfsetstrokecolor{textcolor}%
\pgfsetfillcolor{textcolor}%
\pgftext[x=1.649328in,y=0.560357in,left,base]{\color{textcolor}{\sffamily\fontsize{6.000000}{7.200000}\selectfont\catcode`\^=\active\def^{\ifmmode\sp\else\^{}\fi}\catcode`\%=\active\def%{\%}$P_4$}}%
\end{pgfscope}%
\begin{pgfscope}%
\pgfsetrectcap%
\pgfsetroundjoin%
\pgfsetlinewidth{1.204500pt}%
\definecolor{currentstroke}{rgb}{0.431373,0.423529,0.462745}%
\pgfsetstrokecolor{currentstroke}%
\pgfsetdash{}{0pt}%
\pgfpathmoveto{\pgfqpoint{1.415994in}{0.467044in}}%
\pgfpathlineto{\pgfqpoint{1.499328in}{0.467044in}}%
\pgfpathlineto{\pgfqpoint{1.582661in}{0.467044in}}%
\pgfusepath{stroke}%
\end{pgfscope}%
\begin{pgfscope}%
\definecolor{textcolor}{rgb}{0.000000,0.000000,0.000000}%
\pgfsetstrokecolor{textcolor}%
\pgfsetfillcolor{textcolor}%
\pgftext[x=1.649328in,y=0.437877in,left,base]{\color{textcolor}{\sffamily\fontsize{6.000000}{7.200000}\selectfont\catcode`\^=\active\def^{\ifmmode\sp\else\^{}\fi}\catcode`\%=\active\def%{\%}$P_5$}}%
\end{pgfscope}%
\begin{pgfscope}%
\pgfsetbuttcap%
\pgfsetroundjoin%
\definecolor{currentfill}{rgb}{0.000000,0.000000,0.000000}%
\pgfsetfillcolor{currentfill}%
\pgfsetlinewidth{1.003750pt}%
\definecolor{currentstroke}{rgb}{0.000000,0.000000,0.000000}%
\pgfsetstrokecolor{currentstroke}%
\pgfsetdash{}{0pt}%
\pgfsys@defobject{currentmarker}{\pgfqpoint{-0.024306in}{-0.024306in}}{\pgfqpoint{0.024306in}{0.024306in}}{%
\pgfpathmoveto{\pgfqpoint{0.000000in}{-0.024306in}}%
\pgfpathcurveto{\pgfqpoint{0.006446in}{-0.024306in}}{\pgfqpoint{0.012629in}{-0.021745in}}{\pgfqpoint{0.017187in}{-0.017187in}}%
\pgfpathcurveto{\pgfqpoint{0.021745in}{-0.012629in}}{\pgfqpoint{0.024306in}{-0.006446in}}{\pgfqpoint{0.024306in}{0.000000in}}%
\pgfpathcurveto{\pgfqpoint{0.024306in}{0.006446in}}{\pgfqpoint{0.021745in}{0.012629in}}{\pgfqpoint{0.017187in}{0.017187in}}%
\pgfpathcurveto{\pgfqpoint{0.012629in}{0.021745in}}{\pgfqpoint{0.006446in}{0.024306in}}{\pgfqpoint{0.000000in}{0.024306in}}%
\pgfpathcurveto{\pgfqpoint{-0.006446in}{0.024306in}}{\pgfqpoint{-0.012629in}{0.021745in}}{\pgfqpoint{-0.017187in}{0.017187in}}%
\pgfpathcurveto{\pgfqpoint{-0.021745in}{0.012629in}}{\pgfqpoint{-0.024306in}{0.006446in}}{\pgfqpoint{-0.024306in}{0.000000in}}%
\pgfpathcurveto{\pgfqpoint{-0.024306in}{-0.006446in}}{\pgfqpoint{-0.021745in}{-0.012629in}}{\pgfqpoint{-0.017187in}{-0.017187in}}%
\pgfpathcurveto{\pgfqpoint{-0.012629in}{-0.021745in}}{\pgfqpoint{-0.006446in}{-0.024306in}}{\pgfqpoint{0.000000in}{-0.024306in}}%
\pgfpathlineto{\pgfqpoint{0.000000in}{-0.024306in}}%
\pgfpathclose%
\pgfusepath{stroke,fill}%
}%
\begin{pgfscope}%
\pgfsys@transformshift{1.921258in}{0.589524in}%
\pgfsys@useobject{currentmarker}{}%
\end{pgfscope}%
\end{pgfscope}%
\begin{pgfscope}%
\definecolor{textcolor}{rgb}{0.000000,0.000000,0.000000}%
\pgfsetstrokecolor{textcolor}%
\pgfsetfillcolor{textcolor}%
\pgftext[x=2.071258in,y=0.560357in,left,base]{\color{textcolor}{\sffamily\fontsize{6.000000}{7.200000}\selectfont\catcode`\^=\active\def^{\ifmmode\sp\else\^{}\fi}\catcode`\%=\active\def%{\%}LGL nodes}}%
\end{pgfscope}%
\end{pgfpicture}%
\makeatother%
\endgroup%

\caption[Legendre polynomials]{The first six Legendre polynomials on
  $[-1,1]$.  Each $P_n$ has exactly $n$ zeros in the interior of the
  interval and satisfies the normalisation $P_n(1) = 1$.  Unlike the
  Chebyshev polynomials, whose local extrema all have magnitude $1$ by
  the cosine definition, the Legendre polynomials oscillate with
  non-uniform amplitude: the local maxima of $\abs{P_n}$ are largest
  near the endpoints and decay toward the centre of the interval,
  though $\abs{P_n(x)} \leq 1$ everywhere.  The dots on the $x$-axis
  mark the six Legendre--Gauss--Lobatto nodes (the extrema of $P_5$
  together with $\pm 1$), which cluster toward the endpoints in the
  same quadratic pattern as the CGL nodes of
  \cref{sec:spec-chebyshev}.}
\label{fig:legendre-polynomials}
\end{figure}

\paragraph{Orthogonality.}

The Legendre polynomials are orthogonal on $[-1,1]$ with unit weight:
%
\begin{equation}
  \int_{-1}^{1} P_m(x)\,P_n(x)\,\d x
  = \frac{2}{2n+1}\,\delta_{mn} \,.
  \label{eq:legendre-ortho}
\end{equation}
%
The integral $\int P_m\,P_n\,\d x$ vanishes unless $m = n$, with no
auxiliary weight function --- exactly the structure the DG mass matrix
entries $\int \phi_i\,\phi_j\,\d x$ require to decouple in a modal
basis.  The unit weight is the essential difference from the Chebyshev
case (\cref{eq:chebyshev-ortho}), where the inner product carries the
factor $(1 - x^2)^{-1/2}$.
\physnote{The Chebyshev weight $(1-x^2)^{-1/2}$ arises from the
$x = \cos\theta$ substitution and is natural for spectral collocation
on a single domain.  The Legendre unit weight is natural for
variational (weak-form) methods like DG, where the governing equations
are integrated against test functions without any auxiliary weight.}

\paragraph{Gaussian quadrature.}

The DG solver must evaluate integrals of polynomial functions at every
time step --- mass matrices, flux terms, source terms.  Quadrature
rules approximate these integrals as weighted sums of point
evaluations, and the central question is: how many points suffice for
exact results?  A rule with $n$ nodes and $n$ weights has $2n$ free
parameters, so it can in principle satisfy $2n$ independent
conditions.  Gaussian quadrature chooses both nodes and weights to
maximise the polynomial degree of exactness: the $n$-point Gauss rule
%
\begin{equation}
  \int_{-1}^{1} f(x)\,\d x
  \;\approx\; \sum_{i=0}^{n-1} w_i\,f(x_i)
  \label{eq:gauss-quadrature}
\end{equation}
%
is exact for all polynomials of degree at most $2n - 1$.  No
quadrature rule with $n$ nodes can do better --- the $2n$ conditions
(one for each monomial $1, x, \ldots, x^{2n-1}$) exactly exhaust the
available degrees of freedom.
\histnote{Gauss derived this result in 1814 for the specific case of
the Legendre weight, using continued fractions.  The connection to
orthogonal polynomials was established later by Jacobi and
Christoffel.}

\paragraph{Why exactness holds.}

The mechanism behind the $2n - 1$ exactness is a
division-with-remainder argument.  Any polynomial $q$ of degree at
most $2n - 1$ can be written
%
\begin{equation}
  q(x) = P_n(x)\,s(x) + r(x) \,,
  \label{eq:gauss-division}
\end{equation}
%
where $s$ and $r$ are polynomials of degree at most $n - 1$.
Integrating both sides and using orthogonality
(\cref{eq:legendre-ortho}), the first term vanishes:
$\int P_n\,s\,\d x = 0$ because $s$ is a linear combination of
$P_0, \ldots, P_{n-1}$, all orthogonal to $P_n$.  So
$\int q\,\d x = \int r\,\d x$.  If the quadrature nodes $x_i$ are
the zeros of $P_n$, then $q(x_i) = r(x_i)$ at every node (because
$P_n(x_i) = 0$), and a quadrature rule that integrates $r$ exactly
--- which has degree at most $n - 1$, well within reach of $n$ nodes
--- also integrates $q$ exactly.  The weights are uniquely determined
by requiring exactness on polynomials of degree at most $n - 1$; the
division argument then guarantees exactness extends to degree
$2n - 1$.  In effect, choosing the zeros of $P_n$ as nodes halves
the problem: the orthogonal polynomial absorbs half the integrand's
complexity, leaving the quadrature rule with only the remainder to
handle.
\warnnote{The $n$-point Gauss rule uses the zeros of $P_n$, which
lie strictly inside $(-1,1)$ --- the endpoints are not included.
For boundary-value problems and inter-element coupling in a DG
solver, one needs the endpoints in the node set.  This motivates the
Gauss--Lobatto variant below.}

\paragraph{Gauss--Lobatto nodes.}

The Legendre--Gauss--Lobatto (LGL) nodes fix the endpoints
$x_0 = -1$ and $x_N = 1$ and choose the remaining $N - 1$ interior
nodes to maximise exactness.  With two of the $N + 1$ node positions
fixed, the $N - 1$ free interior positions and $N + 1$ free weights
give $2N$ parameters, enough for exactness on polynomials of degree
at most $2N - 1$.  The interior nodes are the zeros of $P_N'$, so the
full node set satisfies
%
\begin{equation}
  (1 - x_i^2)\,P_N'(x_i) = 0 \,,
  \qquad i = 0, 1, \ldots, N \,.
  \label{eq:lgl-nodes}
\end{equation}
%
The factor $(1 - x_i^2)$ vanishes at the endpoints and $P_N'(x_i)$
vanishes at the interior nodes, so this single condition captures
both the fixed endpoints and the optimally chosen interior points.
The corresponding quadrature weights are
%
\begin{equation}
  w_i = \frac{2}{N(N+1)\,[P_N(x_i)]^2} \,,
  \qquad i = 0, 1, \ldots, N \,.
  \label{eq:lgl-weights}
\end{equation}
%
At the endpoints, $P_N(\pm 1) = (\pm 1)^N$, so
$w_0 = w_N = 2/[N(N+1)]$.  The interior weights are larger,
reflecting the wider effective spacing of the clustered interior
nodes.
\xrefnote{\Cref{sec:spec-lgl} uses these nodes and weights to
build the LGL differentiation matrix and the Vandermonde matrices
that the DG solver uses to convert between nodal and modal
representations.}

\paragraph{Computing LGL nodes.}

The interior LGL nodes --- the zeros of $P_N'$ --- cannot be written
in closed form for general $N$.  The codebase computes them by Newton
iteration, using the Bonnet recurrence
(\cref{eq:legendre-recurrence}) to evaluate $P_N$ and $P_N'$ at each
iterate.  A good initial guess comes from the Chebyshev extrema
(\cref{eq:chebyshev-extrema}), which cluster in the same quadratic
pattern near the endpoints.  For interior nodes ($\abs{x} < 1$), the
second derivative $P_N''$ needed by Newton's method is not computed by
differentiating the recurrence again; instead, the Legendre
differential equation $(1-x^2)P_N'' - 2xP_N' + N(N+1)P_N = 0$ is
rearranged to give $P_N''$ directly from $P_N$ and $P_N'$:
%
\begin{lstlisting}[language=Rust]
// Initial guess from Chebyshev extrema
let mut x = -(i as f64 * PI / N as f64).cos();
// Newton on P'_N: use Legendre ODE for P''_N
for _ in 0..MAX_ITER {
    let (p, dp) = legendre_and_deriv(N, x);
    // (1-x^2)P'' = 2xP' - N(N+1)P
    let d2p = (2.0*x*dp - N*(N+1) as f64 * p)
              / (1.0 - x*x);
    x -= dp / d2p;
}
\end{lstlisting}
%
Newton converges quadratically from the Chebyshev initial guess,
typically reaching machine precision in 5--8 iterations even at
$N = 50$.  Once the nodes are known, the weights follow from
\cref{eq:lgl-weights} using the $P_N$ values already accumulated
during the recurrence.
\implnote{The \texttt{LglBasis} struct in \codemodule{dg} stores the
LGL nodes, weights, differentiation matrix, and Vandermonde matrices.
The constructor calls \texttt{lgl\_nodes\_weights} for the nodes and
weights, then \texttt{lgl\_diff\_matrix} for the differentiation
matrix --- the same row-sum-diagonal construction as the Chebyshev
case (\cref{sec:spec-diffmat}).}

\paragraph{Quadrature precision in the DG solver.}

The $2N - 1$ exactness of the LGL rule has a direct consequence for
the DG mass matrix.  The mass matrix entry
$M_{ij} = \int \phi_i\,\phi_j\,\d x$ involves products of degree-$N$
basis polynomials, giving integrands of degree $2N$ --- one degree
beyond the exactness threshold.  The resulting quadrature error is
small but nonzero, and makes the discrete mass matrix slightly
different from the exact one.  In the Legendre modal basis, the exact
mass matrix is diagonal by orthogonality (\cref{eq:legendre-ortho}).
The LGL-quadrature approximation is also diagonal: for $m \neq n$,
the product $P_m\,P_n$ has degree $m + n \leq 2N - 1$ (since at
least one index is less than $N$), so the orthogonal zero is
reproduced exactly; the quadrature error affects only the diagonal
$P_N \cdot P_N$ entry, whose degree $2N$ exceeds the threshold.
This is why mass lumping --- replacing the full mass matrix by the
diagonal matrix of quadrature weights --- is exact in the modal
basis for all but the highest mode.  Unweighted orthogonality and
near-exact quadrature combine to give a mass matrix that is diagonal
by construction --- the payoff of choosing Legendre over Chebyshev
for the DG basis.

% sec:spec-convergence — Spectral Convergence
\section{Spectral Convergence}
\label{sec:spec-convergence}

The previous sections built the machinery of spectral methods ---
polynomial bases, differentiation matrices, quadrature rules --- but
left a central claim unsubstantiated: that spectral methods converge
exponentially for smooth functions.  Finite-difference methods gain a
fixed number of digits of accuracy each time the grid is refined by a
constant factor; spectral methods gain a fixed number of digits each
time a single node is added.  This qualitative difference --- algebraic
versus exponential convergence --- is the reason that twenty CGL nodes
can match the accuracy of hundreds of finite-difference grid points
(\cref{fig:diffmat-convergence}).  The mechanism is the rapid decay of
Chebyshev expansion coefficients, and that decay is governed entirely
by the smoothness and analyticity of the function being approximated.

\paragraph{Coefficient decay for smooth functions.}

The Chebyshev expansion coefficients $a_n$ defined in
\cref{eq:chebyshev-expansion} are Fourier cosine coefficients in the
variable $\theta = \arccos x$.  Integration by parts applied to the
integral $a_n = (2/\pi c_n)\int_0^\pi
f(\cos\theta)\cos(n\theta)\,\d\theta$ (with $c_0 = 2$, $c_n = 1$ for
$n \geq 1$) yields a factor of $1/n$ for
each derivative of $f$ that exists and is continuous.  A function
$f \in C^k[-1,1]$ with a $(k+1)$-th derivative of bounded variation
therefore has coefficients satisfying
%
\begin{equation}
  \abs{a_n} = \order{n^{-(k+1)}}
  \qquad \text{as } n \to \infty \,.
  \label{eq:coeff-decay-smooth}
\end{equation}
%
Each additional derivative of $f$ buys one more power of decay.  A
$C^\infty$ function has coefficients that decay faster than any
polynomial in $1/n$ --- superalgebraic decay --- because the
integration-by-parts argument can be applied arbitrarily many times.
\physnote{The integration-by-parts argument is the same mechanism that
makes Fourier coefficients of smooth periodic functions decay rapidly.
The Chebyshev case inherits this directly through the
$x = \cos\theta$ substitution.}

\paragraph{Geometric convergence for analytic functions.}

Functions that are merely smooth have superalgebraic coefficient decay,
but analytic functions do better: their coefficients decay
geometrically.  The key idea is that an analytic function can be
extended into the complex plane, and the Chebyshev coefficients can
be estimated using contour integrals on an elliptical contour rather
than on $[-1,1]$ itself.

Under the substitution $x = \cos\theta$, the interval
$[-1,1]$ lifts to the real axis of the $\theta$-plane, and the
substitution $z = \cos\zeta$ (with $\zeta = \theta + i\sigma$)
extends the map into the complex $x$-plane.  The image of the strip
$\abs{\operatorname{Im}\zeta} < \sigma_0$ under $x = \cos\zeta$ is
an ellipse in the complex $x$-plane with foci at $\pm 1$ and
semi-axis sum
%
\begin{equation}
  \rho = e^{\sigma_0} \,.
  \label{eq:bernstein-rho}
\end{equation}
%
This ellipse, called the \emph{Bernstein ellipse} $\mathcal{E}_\rho$,
has semi-major axis $a_\rho = (\rho + \rho^{-1})/2$ along the real
axis and semi-minor axis $b_\rho = (\rho - \rho^{-1})/2$ along the
imaginary axis.  When $\rho$ is close to $1$ the ellipse collapses
onto $[-1,1]$; as $\rho$ grows, the ellipse inflates into the complex
plane.
\histnote{The connection between analyticity in a Bernstein ellipse and
Chebyshev coefficient decay appears in Bernstein's work from the early
twentieth century.  The modern treatment follows Trefethen
\cite{Trefethen2000} and the detailed exposition in
\cite{Trefethen2019}.}%
\warnnote{The symbol $\rho$ here denotes the Bernstein ellipse
parameter, not the convergence ratio of a finite-difference stencil
(\cref{sec:fd-convergence}).  Both usages are standard in their
respective fields.}

\paragraph{The geometric decay bound.}

If $f$ is analytic and bounded inside the Bernstein ellipse
$\mathcal{E}_\rho$ for some $\rho > 1$, then its Chebyshev
coefficients satisfy
%
\begin{equation}
  \abs{a_n} \leq \frac{2M}{\rho^n} \,,
  \label{eq:coeff-decay-analytic}
\end{equation}
%
where $M = \max_{z \in \mathcal{E}_\rho}\abs{f(z)}$.
The proof uses Cauchy's integral formula: the Chebyshev coefficient
$a_n$ is a contour integral of $f(z)T_n(z)$ around $[-1,1]$, and
deforming the contour to $\mathcal{E}_\rho$ introduces the factor
$\rho^{-n}$ from the growth of $T_n$ on the ellipse
($\abs{T_n(z)} \leq \rho^n/2$ on $\mathcal{E}_\rho$ for $n \geq 1$).  The
coefficients decay like powers of $1/\rho$ --- geometric convergence,
with the rate set by the size of the largest ellipse inside which $f$
remains analytic.
\warnnote{The constant $\rho$ is determined by the nearest singularity
of $f$ in the complex plane, not by any property of $f$ on the real
interval alone.  A function that appears perfectly smooth on $[-1,1]$
can have slow spectral convergence if it has a pole or branch point
near the real axis.}

\paragraph{From coefficients to interpolation error.}

The Chebyshev interpolant $p_N(x) = \sum_{n=0}^{N}
\hat{a}_n\,T_n(x)$ reproduces the first $N+1$ coefficients of $f$
exactly (up to aliasing), so the interpolation error is controlled by
the tail of the expansion.  For an analytic function bounded in
$\mathcal{E}_\rho$, the truncation error satisfies
%
\begin{equation}
  \norm{f - p_N}_\infty
    \leq \frac{2M\,\rho^{-N}}{\rho - 1} \,,
  \label{eq:spec-interp-error}
\end{equation}
%
since $\sum_{n=N+1}^{\infty} \abs{a_n} \leq
2M\sum_{n=N+1}^{\infty}\rho^{-n} = 2M\rho^{-N}/(\rho - 1)$ and
$\abs{T_n(x)} \leq 1$ on $[-1,1]$.  The error decreases by a
factor of $\rho$ for each additional node, confirming the exponential
convergence observed in \cref{fig:diffmat-convergence}: the gold curve
drops by roughly a constant factor per node until reaching machine
precision.  The aliasing error from the discrete projection adds terms
of order $\rho^{-2N}$, which is negligible compared to the truncation
tail.
\xrefnote{The Lebesgue constant (\cref{eq:lebesgue-bound}) provides an
alternative route to the interpolation error: for CGL nodes,
$\Lambda_N \sim (2/\pi)\ln N$, so the interpolant is within a
logarithmic factor of the best polynomial approximation.  The bound
\cref{eq:spec-interp-error} is tighter for analytic functions because
it exploits the specific decay rate of the coefficients.}

\paragraph{Runge's function revisited.}

Runge's function $f(x) = 1/(1 + 25x^2)$ from
\cref{sec:spec-interpolation} provides a concrete illustration.  Its
singularities are the simple poles at $x = \pm i/5$, which lie at
distance $1/5$ from the real axis.  The largest Bernstein ellipse inside which $f$ remains analytic has
semi-axis sum $\rho$ determined by $b_\rho = 1/5$, where
$b_\rho = (\rho - \rho^{-1})/2$.  Solving gives
%
\begin{equation}
  \rho = \frac{1}{5} + \sqrt{1 + \frac{1}{25}}
  = \frac{1 + \sqrt{26}}{5}
  \approx 1.220 \,.
  \label{eq:runge-rho}
\end{equation}
%
The poles lie \emph{on} $\mathcal{E}_\rho$, so the bound
\cref{eq:coeff-decay-analytic} holds for any $\rho' < \rho$; in
practice the convergence rate approaches $\rho$ from below.
Each additional Chebyshev node reduces the interpolation error by
roughly the factor $\rho \approx 1.22$, or about $0.86$ digits per
node.
At $N = 15$, the predicted error is roughly $1.22^{-15} \approx
0.037$, which matches the $\approx 4\%$ maximum error visible in
\cref{fig:runge-phenomenon}(b).  The modest convergence rate
--- compared to the machine-precision accuracy achieved at $N = 20$ for
$f(x) = e^{\sin(\pi x)}$ in \cref{fig:diffmat-convergence} ---
reflects the poles lurking close to the real axis.  An entire
function (no finite singularities) has $\rho = \infty$, and the bound
\cref{eq:spec-interp-error} must be replaced by a supergeometric
estimate; for $e^{\sin(\pi x)}$, the coefficients decay
roughly as $\order{(2/(\pi e))^n / \sqrt{n}}$.

\paragraph{The convergence hierarchy.}

The relationship between smoothness and convergence rate is a
hierarchy with sharp boundaries:

\begin{center}
\begin{tabular}{lll}
  \toprule
  Smoothness of $f$ & Coefficient decay & Convergence rate \\
  \midrule
  $C^k$, $(k{+}1)$-th deriv.\ of BV
    & $\abs{a_n} = \order{n^{-(k+1)}}$
    & algebraic, $\order{N^{-(k+1)}}$ \\
  $C^\infty$
    & faster than any $n^{-k}$
    & superalgebraic \\
  Analytic in $\mathcal{E}_\rho$
    & $\abs{a_n} = \order{\rho^{-n}}$
    & geometric, $\order{\rho^{-N}}$ \\
  Entire
    & faster than any $\rho^{-n}$
    & supergeometric \\
  \bottomrule
\end{tabular}
\end{center}

Each step up in smoothness produces qualitatively faster convergence.
A function with a discontinuity in its $k$-th derivative converges at
the algebraic rate $\order{N^{-(k+1)}}$, indistinguishable from a
finite-difference method of order $k + 1$ --- spectral methods buy
nothing for non-smooth functions.  The exponential advantage emerges
only when $f$ is analytic, and the rate improves as the nearest complex
singularity moves further from the real axis.
\physnote{The metric functions of the Schwarzschild and Kerr spacetimes
in Kerr--Schild coordinates are rational functions of the Cartesian
coordinates --- analytic everywhere except at the physical singularity
($r = 0$).  On any domain that excludes the singularity, the spectral
solver achieves geometric convergence with $\rho$ determined by the
distance to $r = 0$ in the mapped coordinate.  This is why the XCTS
solver reaches machine precision at moderate polynomial degrees
($N \sim 20$--$30$).}

\paragraph{Convergence of spectral derivatives.}

Differentiating the Chebyshev expansion term by term introduces
algebraic factors of $n$ in the derivative's coefficients, worsening
the decay by at most one algebraic power.  For an analytic function
with geometric coefficient decay $\abs{a_n} = \order{\rho^{-n}}$,
these algebraic factors are absorbed by the geometric term:
$n\rho^{-n} \to 0$ at the same geometric rate.  The spectral
derivative therefore converges geometrically at essentially the same
rate as the function itself --- the differentiation matrix amplifies
errors by $\order{N^2}$ (\cref{eq:diffmat-norm}), but the
exponentially small interpolation error shrinks fast enough to
overwhelm this amplification.

Quantitatively, the derivative error satisfies
%
\begin{equation}
  \norm{f' - p_N'}_\infty
    \leq C\,N^2\,\rho^{-N}
  \label{eq:spec-deriv-error}
\end{equation}
%
for an analytic function in $\mathcal{E}_\rho$, where $C$ depends on
$\rho$ and $M$ but not on $N$ \cite{Trefethen2000}.  The algebraic
prefactor $N^2$ is negligible compared to the geometric decay for any
$\rho > 1$: at $N = 20$ and $\rho = 2$, the bound is roughly
$400 \times 10^{-6} = 4 \times 10^{-4}$, while the actual error is
typically much smaller because the bound is not tight.  This is the
precise sense in which spectral convergence ``overwhelms'' the
conditioning of the differentiation matrix, as claimed in
\cref{sec:spec-diffmat}.
\warnnote{The geometric convergence of spectral derivatives breaks
down at discontinuities.  If $f$ has a jump in its $k$-th derivative,
the spectral derivative converges only at rate
$\order{N^{-k}}$ --- one power slower than the function itself.
For the DG solver, each element carries its own polynomial, so
discontinuities are confined to element interfaces and handled by
the numerical flux, not by the spectral basis.}

% sec:spec-lgl — LGL Differentiation and Mass Lumping
\section{LGL Differentiation and Mass Lumping}
\label{sec:spec-lgl}

The DG solver doesn't use a spectral differentiation matrix in the
same way as the collocation solver of \cref{sec:spec-diffmat}.
Collocation methods apply the $D$ matrix directly to function values;
DG methods work in weak form, multiplying by test functions and
integrating.  The weak-form machinery requires three matrices
assembled from the LGL nodes and weights of
\cref{sec:spec-legendre}: a differentiation matrix, a mass matrix,
and a stiffness matrix.  Together with the Vandermonde matrix
(\cref{sec:spec-vandermonde}) for modal--nodal conversion, all four
live inside the \texttt{LglBasis} struct that the codebase constructs
once at startup and reuses at every time step.
\coderef{dg}{LglBasis}

\paragraph{The LGL differentiation matrix.}

The LGL differentiation matrix is constructed by the same recipe
as the CGL version (\cref{sec:spec-diffmat}): the entry $D_{ij} =
\ell_j'(x_i)$ records the derivative of the $j$-th Lagrange
interpolant evaluated at the $i$-th node.  What changes is the node
set --- LGL nodes instead of CGL --- and consequently the formula for
the off-diagonal entries.  On CGL nodes the closed-form entry uses
endpoint flags $c_i$ (\cref{eq:diffmat-cgl}); on LGL nodes the
analogous formula uses evaluations of the Legendre polynomial $P_N$:
%
\begin{equation}
  D_{ij} = \frac{P_N(x_i)}{P_N(x_j)}
    \cdot \frac{1}{x_i - x_j} \,,
  \qquad i \neq j \,.
  \label{eq:lgl-diffmat-offdiag}
\end{equation}
%
The values $P_N(x_i)$ are already available from the Newton iteration
that computed the nodes (\cref{sec:spec-legendre}), so the matrix
construction requires no additional polynomial evaluations.  The
diagonal entries follow from the same row-sum property as before
(\cref{eq:diffmat-rowsum}): setting $D_{ii} = -\sum_{j \neq i}
D_{ij}$ guarantees that the matrix differentiates constants exactly.
The two endpoint entries have closed forms that bypass the row sum:
%
\begin{equation}
  D_{00} = -\frac{N(N+1)}{4} \,, \qquad
  D_{NN} = \frac{N(N+1)}{4} \,,
  \label{eq:lgl-diffmat-endpoints}
\end{equation}
%
with all interior diagonals computed from the negative row sum.
\implnote{The codebase function \texttt{lgl\_diff\_matrix} in
\codemodule{dg} implements \cref{eq:lgl-diffmat-offdiag} with the
endpoint formulas \cref{eq:lgl-diffmat-endpoints}, then fills
interior diagonals by negative row summation --- identical in
structure to the CGL code in \codemodule{spectral}.}

\paragraph{The mass matrix.}

The DG weak form multiplies the governing PDE by a test function
$\phi_i$ and integrates over a single element.  In the Lagrange nodal
basis, the test functions are the Lagrange interpolants $\ell_i$
themselves, and the mass matrix that results from the left-hand side
has entries
%
\begin{equation}
  M_{ij} = \int_{-1}^{1} \ell_i(x)\,\ell_j(x)\,\d x \,.
  \label{eq:mass-matrix}
\end{equation}
%
This integral involves a product of two degree-$N$ polynomials ---
an integrand of degree $2N$.  The LGL quadrature rule with $N + 1$
nodes is exact for polynomials of degree at most $2N - 1$
(\cref{sec:spec-legendre}), so it misses the mark by one degree.
Approximating the integral by the LGL rule gives the
\emph{lumped mass matrix}
%
\begin{equation}
  \hat{M}_{ij}
  = \sum_{k=0}^{N} w_k\,\ell_i(x_k)\,\ell_j(x_k)
  = w_i\,\delta_{ij} \,,
  \label{eq:mass-lumped}
\end{equation}
%
where the last step uses $\ell_i(x_k) = \delta_{ik}$ (the cardinal
property of Lagrange interpolants at the collocation nodes).  The
lumped mass matrix is diagonal: its only entries are the LGL
quadrature weights $w_i$ from \cref{eq:lgl-weights}.  Inverting it
requires only $N + 1$ divisions, not a matrix solve.
\physnote{Mass lumping is exact for all but the highest mode.  The
quadrature error affects only the $(P_N, P_N)$ entry in the modal
basis, as established in \cref{sec:spec-legendre}.  For smooth
solutions, the highest mode carries exponentially small amplitude,
so the lumping error is negligible.}

\paragraph{Why mass lumping works.}

The one-degree shortfall deserves scrutiny: if the quadrature doesn't
integrate $\ell_i\,\ell_j$ exactly, why trust the result?  The
argument from \cref{sec:spec-legendre} showed that the Legendre modal
mass matrix is diagonal by orthogonality, and that the LGL quadrature
reproduces this diagonality for all entries except $(P_N, P_N)$.
In the nodal basis the same conclusion holds by a different route.
The Lagrange interpolants at LGL nodes satisfy an exact
discrete orthogonality --- the identity \cref{eq:mass-lumped}
holds because of the cardinal property, regardless of the
quadrature's polynomial exactness.  The lumped mass matrix is the
identity (up to the weight factors) by construction --- it is the
\emph{exact} integral of $\ell_i\,\ell_j$ that differs from it, and
the difference is confined to a single mode.  Mass lumping doesn't
introduce error by under-resolving a difficult integral; it
\emph{defines} the inner product that the DG scheme uses, and that
inner product is spectrally close to the exact one.
\warnnote{Mass lumping is not free: it modifies the numerical scheme
by replacing the exact $L^2$ inner product with the discrete LGL
inner product.  For $p$-refinement studies comparing different $N$
at fixed element count, the mass-lumping error decreases
exponentially with $N$ and is invisible in practice.  For very
coarse resolutions ($N = 1$ or $2$), the error can be noticeable.}

\paragraph{The stiffness matrix.}

The right-hand side of the DG weak form produces the stiffness
matrix, which encodes differentiation weighted by the test function.
Applying LGL quadrature to the integral
$\int \ell_i(x)\,\ell_j'(x)\,\d x$ --- an integrand of degree
$2N - 1$, so the quadrature is exact --- gives
%
\begin{equation}
  S_{ij} = \sum_{k=0}^{N} w_k\,\ell_i(x_k)\,\ell_j'(x_k)
  = w_i\,D_{ij} \,,
  \label{eq:stiffness-matrix}
\end{equation}
%
using $\ell_i(x_k) = \delta_{ik}$ and the differentiation matrix
definition $D_{ij} = \ell_j'(x_i)$ from \cref{eq:diffmat-def}.
The stiffness matrix is the mass-weighted differentiation matrix:
$S = \operatorname{diag}(w_0, \ldots, w_N) \cdot D$.  The
combination $\hat{M}^{-1} S = D$ --- the lumped mass inverse times
the stiffness matrix is just the differentiation matrix itself.
This identity is the algebraic payoff of mass lumping: the weak-form
DG operator reduces to a simple matrix-vector multiply by $D$,
with no matrix inversion at each time step.
\xrefnote{The DG solver applies $D$ in each coordinate direction
independently, producing the volume derivative terms.  Face
contributions from the numerical flux are added separately
(\cref{ch:dg-methods}).}

\paragraph{Assembly in the codebase.}

The \texttt{lgl\_basis} constructor in \codemodule{dg} assembles all
four components in sequence: nodes and weights from
\texttt{lgl\_nodes\_weights} (\cref{sec:spec-legendre}), the
differentiation matrix from \texttt{lgl\_diff\_matrix}
(\cref{eq:lgl-diffmat-offdiag}), and the Vandermonde matrix from
\texttt{legendre\_all} (\cref{sec:spec-vandermonde}).  The resulting
\texttt{LglBasis} struct stores:
%
\begin{lstlisting}[language=Rust]
pub struct LglBasis {
    pub order: usize,          // polynomial order p
    pub n: usize,              // number of nodes: p+1
    pub nodes: Vec<f64>,       // LGL nodes in [-1,1]
    pub weights: Vec<f64>,     // LGL quadrature weights
    pub d_matrix: Vec<f64>,    // D, n*n row-major
    pub v_matrix: Vec<f64>,    // Vandermonde V (modal->nodal)
    pub v_inv: Vec<f64>,       // V^{-1}
}
\end{lstlisting}
%
The struct is constructed once per simulation and shared across all
elements of the same polynomial order.  At each Runge--Kutta stage,
the volume derivative $D\,\mathbf{u}$ is computed by a dense
matrix-vector multiply; the mass matrix never appears explicitly
because the lumped inverse is absorbed into the operator.  For the
polynomial orders used in practice ($N = 3$--$7$), the $D$-matrix
multiply dominates the cost of each element update, and its
$\order{N^2}$ operation count per node is modest --- a single
element with $N = 5$ requires $36$ multiplies per derivative
direction.
\implnote{In three dimensions each element stores
$(N+1)^3$ unknowns per field variable.  The volume derivative
applies $D$ along each axis independently via tensor-product
structure, reducing the per-element cost from
$\order{N^6}$ (dense matrix) to $\order{N^4}$ (three
one-dimensional passes) --- the sum-factorisation technique of
\cite{Kopriva2009}.}

% sec:spec-vandermonde — Vandermonde Matrices and Modal-Nodal Representations
\section{Vandermonde Matrices and Modal-Nodal Representations}
\label{sec:spec-vandermonde}

The LGL basis stores function values at nodes — a \emph{nodal}
representation.  The Legendre polynomials $P_0, P_1, \ldots, P_N$
provide an alternative \emph{modal} representation in which a
polynomial is described by its expansion coefficients.  Converting
between the two is a matrix multiply, and the matrix that performs
the conversion appears in every spectral filtering or limiting
operation.  When the DG solver encounters steep gradients near a
black hole puncture, it needs to damp high-frequency oscillations
without corrupting the smooth part of the solution — and damping is
far more natural in the modal picture, where ``high frequency'' simply
means ``large index $k$''.
\coderef{dg}{LglBasis}

\paragraph{The Vandermonde matrix.}

A polynomial of degree $N$ can be written in the Legendre modal basis
as
%
\begin{equation}
  u(x) = \sum_{j=0}^{N} \hat{u}_j\,P_j(x) \,,
  \label{eq:modal-expansion}
\end{equation}
%
where $\hat{u}_j$ are the modal coefficients.  Evaluating this
expansion at the $i$-th LGL node $x_i$ gives
$u(x_i) = \sum_j V_{ij}\,\hat{u}_j$, where the Vandermonde matrix
has entries
%
\begin{equation}
  V_{ij} = P_j(x_i) \,, \qquad
  i, j = 0, \ldots, N \,.
  \label{eq:vandermonde}
\end{equation}
%
Since the $N + 1$ LGL nodes are distinct and $\{P_0, \ldots, P_N\}$
is a basis for the polynomials of degree $\leq N$, the Vandermonde
matrix is nonsingular --- this is the unisolvence theorem for
polynomial interpolation (\cref{sec:spec-interpolation}).  In matrix
notation, $\mathbf{u} = V\hat{\mathbf{u}}$: the Vandermonde matrix
maps modal coefficients to nodal values.  The inverse transform
$\hat{\mathbf{u}} = V^{-1}\mathbf{u}$ recovers the modal
coefficients from the nodal data.
\implnote{The codebase constructs $V$ in \texttt{lgl\_basis} by
evaluating \texttt{legendre\_all} at each LGL node, then computes
$V^{-1}$ via LU decomposition.  For the polynomial orders used in
practice ($N \leq 7$), the matrix is at most $8 \times 8$ and the
inversion is negligible.}

\paragraph{Conditioning.}

The Vandermonde matrix on LGL nodes is well-conditioned for moderate
polynomial orders because the LGL nodes are near-optimal
interpolation points for the Legendre basis: they cluster near the
endpoints in a pattern closely related to the Chebyshev distribution
that minimises the Lebesgue constant
(\cref{sec:spec-interpolation}).  For $N \leq 7$ the condition number
$\kappa(V)$ remains modest --- of order $10$ --- so the matrix
inversion introduces no significant round-off amplification.  At very
high orders ($N \gtrsim 20$), the condition number grows and the
transform should be performed via a recurrence or fast algorithm
rather than dense inversion, but the DG solver in the codebase
operates at $N = 3$--$7$ where the dense approach is both fast and
stable.
\warnnote{The Vandermonde matrix on \emph{equidistant} nodes is
notoriously ill-conditioned: $\kappa(V)$ grows exponentially with $N$.
LGL (and CGL) nodes avoid this by clustering at the endpoints, just
as they tame the Runge phenomenon for interpolation
(\cref{sec:spec-interpolation}).}

\paragraph{Modal filtering.}

The primary use of the Vandermonde transform in the codebase is
\emph{modal filtering}: selectively damping the highest-order
expansion modes to remove oscillations near discontinuities or sharp
features.  The idea is straightforward.  Transform the nodal
solution to modal space via $\hat{\mathbf{u}} = V^{-1}\mathbf{u}$,
multiply each coefficient by a filter function $\sigma(k)$ that
equals $1$ for low modes and decays to $0$ for the highest mode,
then transform back via $\mathbf{u}_{\text{filtered}} =
V\,\hat{\mathbf{u}}_{\text{filtered}}$.  In the nodal basis the
entire operation is a single matrix multiply by the filter matrix
%
\begin{equation}
  F = V \cdot \operatorname{diag}\bigl(\sigma(0), \sigma(1),
    \ldots, \sigma(N)\bigr) \cdot V^{-1} \,.
  \label{eq:filter-matrix}
\end{equation}
%
Because $V$, $V^{-1}$, and the filter weights are all fixed for a
given polynomial order, $F$ can be assembled once at startup and
reused --- filtering reduces to a single matrix-vector multiply per
element per coordinate direction.  The filter function used in the
codebase is the exponential filter of order $s$:
%
\begin{equation}
  \sigma(k) = \exp\!\Bigl(-\alpha\,\Bigl(\frac{k}{N}\Bigr)^{\!s}\Bigr) \,,
  \label{eq:exponential-filter}
\end{equation}
%
where $\alpha = -\ln(10^{-16}) = 16\ln 10 \approx 36.8$ sets the
attenuation of the highest mode to $\sigma(N) = 10^{-16} \approx
\epsilon_{\text{m}}$, and $s$ controls the steepness of the roll-off.
For $s = 1$ the filter resembles a low-pass with gradual decay; for
large $s$ it acts as a near-ideal brick-wall that leaves low modes
untouched and kills the highest modes.  The codebase uses $s = 16$,
which leaves modes $k \lesssim N/2$ essentially unchanged
($1 - \sigma < 10^{-14}$) while concentrating the attenuation in
the top few modes --- $\sigma(N) = 10^{-16}$.
\physnote{Modal filtering is the spectral-element analogue of
Kreiss--Oliger dissipation (\cref{sec:fd-dissipation}).  Both add
controlled damping to high-frequency modes to stabilise the scheme
near under-resolved features.  The exponential filter targets the
modal spectrum directly; Kreiss--Oliger dissipation targets the grid
wavenumber via finite-difference operators.}

\paragraph{Tensor-product filtering in three dimensions.}

In three dimensions the solution on a hexahedral element has
$(N+1)^3$ degrees of freedom.  A naive application of
\cref{eq:filter-matrix} would require multiplying by the full
$(N+1)^3 \times (N+1)^3$ filter matrix --- an $\order{N^6}$ operation
per element.  The tensor-product structure of the LGL basis reduces
this to three one-dimensional passes, one along each coordinate
direction.  Each pass applies $V^{-1}$, multiplies by the diagonal
filter, and applies $V$; the total cost is $\order{N^4}$ per element,
the same asymptotic scaling as the volume derivative from
\cref{sec:spec-lgl}.
\implnote{The function \texttt{apply\_modal\_filter} in
\codemodule{dg} implements the tensor-product filter using $s = 16$.
It loops over $x$, $y$, $z$ lines in sequence, applying the 1D
forward transform $V^{-1}$, the diagonal filter, and the backward
transform $V$ to each line --- the same sum-factorisation strategy
used for derivatives (\cref{sec:spec-lgl}).}

\paragraph{When to filter.}

Not every element needs filtering --- applying it uniformly wastes
work and unnecessarily diffuses smooth solutions.  The codebase
applies the exponential filter only to elements near a black hole
puncture, where the initial data or evolved fields contain steep
gradients that the polynomial basis cannot resolve.  In such
elements the highest Legendre modes carry non-negligible amplitude
and, without damping, produce Gibbs-like oscillations that can
corrupt the evolution.  The filter removes these oscillations at the
cost of spectral accuracy in the filtered region --- a deliberate
trade-off, since the solution near a puncture is not smooth and
spectral accuracy was never available there.

The Vandermonde matrix completes the LGL toolbox: nodes, weights,
differentiation matrix, and now the modal-nodal transform that
enables filtering.  Together these components form the
\texttt{LglBasis} struct that the DG solver constructs once and reuses
throughout the simulation.
\xrefnote{The DG chapter (\cref{ch:dg-methods}) discusses the
filtering strategy in the context of puncture initial data and the
interplay with the Rusanov numerical flux.}


\paragraph{Key References.}
The spectral-methods framework follows \cite{Hesthaven2007} and
\cite{Kopriva2009}; Chebyshev theory draws on \cite{Trefethen2000};
the DG application follows \cite{HesthavenWarburton2008}.
